\documentclass[11pt,a4paper,notitlepage,oneside]{amsart}
\usepackage{etoolbox}
\usepackage{microtype} %espansione font + protusione caratteri per miglior riempimento riga
\usepackage{amsmath}
\usepackage{amsthm}
\usepackage{braket}
\usepackage{pdfpages}
\usepackage[psamsfonts]{amssymb}
\usepackage[T1]{fontenc}
\usepackage[english]{babel}
\usepackage[utf8]{inputenc}
\usepackage{enumitem}
\usepackage{mathtools}
\usepackage{caption,subfig} %per figure e didascalie
\usepackage{rotating} %per ruotare figure
\usepackage{scalerel}
\usepackage{relsize} % make math symbol larger (smaller) with \mathlarger (\mathsmaller}
\usepackage{cancel} % to cancel text with \xcancel{}

\usepackage[autostyle,italian=guillemets]{csquotes}
\usepackage[style=alphabetic,sorting=nyt,firstinits=true,maxnames=5,maxalphanames=5,backend=bibtex]{biblatex}
\bibliography{biblio}
\allowdisplaybreaks %Long formulas can break in 2 different pages.

%%%for GeoGebra
\usepackage{pgf,tikz,pgfplots}
\usepackage{mathrsfs}
\usetikzlibrary{arrows}
%%%

%%% For arxiv - to avoid url of bibliography to run over the margins
%\usepackage[hyphens]{url}
\usepackage[hyphenbreaks]{breakurl}
%\setcounter{biburllcpenalty}{7000}
%\setcounter{biburlucpenalty}{8000}
%%%

\usepackage{fancyhdr}
\usepackage{tikz-cd}
\usepackage{amscd}
\usepackage[a4paper, bottom=3cm, left=3cm,right=3cm]{geometry}
\usepackage{hyperref}
\hypersetup{hidelinks} %nasconde i quadrati colorati delle referenze.

%%%for separated sections.
\usepackage{subfiles}
%%%
\pagestyle{plain}
\usepackage{my_shortcuts}



\author{Jan Hendrik Bruinier}
\address{
Fachbereich Mathematik, Technische Universität Darmstadt, Schlossgartenstraße 7, D–
64289 Darmstadt, Germany.
}
\email{bruinier@mathematik.tu-darmstadt.de}
\author{Riccardo Zuffetti}
\address{
Fachbereich Mathematik, Technische Universität Darmstadt, Schlossgartenstraße 7, D–
64289 Darmstadt, Germany.
}
\email{zuffetti@mathematik.tu-darmstadt.de}



\title{Lefschetz decompositions of Kudla--Millson theta functions}
%\date{\today}

\setcounter{tocdepth}{1} %Only main sections in the index of contents
%\setcounter{tocdepth}{2} %Also subsections in the index of contents

\begin{document}
	\maketitle
	%%%%%%%%%%%%%%%%%%%%%%%%%%%%%%%%%%%%%%%%%%%%%%%%%%%%%%%%%%%%%%%%%%%%%%%%%%%%%%%%%%%%%%%%%
	\begin{abstract}
	In the 80's Kudla and Millson introduced a theta function in two variables.
	It behaves as a Siegel modular form with respect to the first variable, and is a closed differential form on an orthogonal Shimura variety with respect to the other variable.
	We prove that the Lefschetz decomposition of the cohomology class of that theta function is also its modular decomposition in Eisenstein, Klingen and cuspidal parts.
	We also show that the image of the Kudla--Millson lift is contained in the primitive cohomology of~$X$.
	As an application, we deduce dimension formulas for the cohomology of~$X$ in low degree.
	Our results cover the cases of moduli spaces of compact hyperkähler manifolds.
%	The Kudla--Millson theta function may be considered as is a Siegel modular form with values in the space of closed forms on some orthogonal Shimura variety~$X$.
%	differential form on an orthogonal Shimura variety~$X$, with respect to the second it is as a Siegel modular form.
%	We compute the Lefschetz decomposition of its cohomology class, and show that it is also modular.
%	This means that each term of the decomposition behaves as either a Klingen Eisenstein series or a cusp form.
%	As application, we show that the image of the Kudla--Millson lift equals the primitive cohomology group of~$X$, prove dimension formulas for cohomology groups of~$X$, as well as for subspaces of special cycles in Chow groups.
	\end{abstract}
	%%%%%%%%%%%%%%%%%%%%%%%%%%%%%%%%%%%%%%%%%%%%%%%%%%%%%%%%%%%%%%%%%%%%%%%%%%%%%%%%%%%%%%%%%
	
	\tableofcontents
	
	\emergencystretch 3em %to avoid problems with arXiv compiler
	%%%%%%%%%%%%%%%%%%%%%%%%%%%%%%%%	
	\section{Introduction}\label{sec;outline}
	\subfile{sections/outline}
	%%%%%%%%%%%%%%%%%%%%%%%%%%%%%%%%
	\section{Vector-valued Siegel modular forms}\label{sec;KMtheta}
	\subfile{sections/backgroundweil}
	\subfile{sections/onPhiandKling}
	\subsection{The Petersson inner product}
	\subfile{sections/petersson_in_prod}
	\subsection{Poincaré series}\label{sec;poincser}
	\subfile{sections/poincare_series}
	\subsection{The space of almost cusp forms and its modular decomposition}\label{sec;mod_dec}
	\subfile{sections/mod_decomposition}
%	\subsection{Diagonal restrictions of Siegel Eisenstein series}\label{sec;diagonalrestr}
%	\subfile{sections/doubling}
	\subsection{Diagonal restrictions of Siegel Eisenstein series}\label{sec;diagonalrestr_vv}
	\subfile{sections/doubling_vv}
	%%%%%%%%%%%%%%%%%%%%%%%%%%%%%%%%
	\section{Orthogonal Shimura varieties and their cohomology}\label{sec;digressdiffgeo}
	\subfile{sections/digdiffgeo}
	\subfile{sections/KudlaMillsontheta}
	%%%%%%%%%%%%%%%%%%%%%%%%%%%%%%%%
%	\section{Digression on Mumford's good forms}\label{sec;digonMumfordgf}
%	\subfile{sections/mumfordgf}
	%%%%%%%%%%%%%%%%%%%%%%%%%%%%%%%%
	\section{The generating series of volumes of special cycles}\label{sec;volofspeccyinKM}
	\subfile{sections/volofspeccyinKM}
	%%%%%%%%%%%%%%%%%%%%%%%%%%%%%%%%
%	\section{VECTVAL - Volume of special cycles}\label{sec;volofspeccyinKMvv}
%	\subfile{sections/volofspeccyinKM_vv}
	%%%%%%%%%%%%%%%%%%%%%%%%%%%%%%%%
	%%%%%%%%%%%%%%%%%%%%%%%%%%%%%%%%
%	\section{VECTVAL - The case of genus 1}\label{sec;gen1vv}
%	\subfile{sections/case_codim_1_vv}
	%%%%%%%%%%%%%%%%%%%%%%%%%%%%%%%%
	\section{The modular Lefschetz decomposition of~$[\KMtheta{\genus}]$}\label{sec;gen2}
	\subfile{sections/case_codim_2}
	%%%%%%%%%%%%%%%%%%%%%%%%%%%%%%%%
%	\section{VECTVAL - The case of general genus}\label{sec;gen2_vv}
%	\subfile{sections/case_codim_2_vv}
	%%%%%%%%%%%%%%%%%%%%%%%%%%%%%%%%
	\section{Geometric applications of the Lefschetz decomposition of~$[\KMtheta{\genus}]$}\label{sec;appgen1}
	\subfile{sections/applic_decomp}
	%%%%%%%%%%%%%%%%%%%%%%%%%%%%%%%%
%	\section{Non-holomorphic theta functions}\label{sec;nonholtheta}
%	\subfile{sections/nonholtheta}
	%%%%%%%%%%%%%%%%%%%%%%%%%%%%%%%%
	
	\printbibliography
	

\end{document}

\documentclass[../Lefschetz_decomposition.tex]{subfiles}

	\begin{document}
	Orthogonal Shimura varieties provide a large family of varieties of geometric interest.
	Among them there are modular and Shimura curves, Hilbert modular surfaces, Siegel~$3$-folds, and moduli spaces of compact hyperkähler manifolds.
	
	Every orthogonal Shimura variety~$X$ comes with a large supply of algebraic cycles, given by embedded orthogonal Shimura subvarieties of smaller dimension.
	By celebrated work of Kudla--Millson, the cohomology classes of such special cycles are related to Siegel modular forms via the geometric theta correspondence, see e.g.~\cite{kudlamillson-intersection}.
	It can be realized by an integral kernel, the so-called \emph{Kudla--Millson theta function}~$\KMtheta{\genus}$ of genus~$g$.
	It is a (non-holomorphic) Siegel modular form of genus~$g$ taking values in the space of closed differential~$(g,g)$-forms on~$X$.
	Its cohomology class~$[\KMtheta{\genus}]$ is a holomorphic Siegel modular form whose Fourier coefficients are the fundamental classes of codimension~$g$ special cycles on~$X$, generalizing the generating series of intersection numbers of Hirzebruch--Zagier divisors on Hilbert modular surfaces studied by Hirzebruch and Zagier in~\cite{hirzebruch-zagier}.
%%	The geometry of an orthogonal Shimura variety~$X$ may be studied in terms of the arithmetic properties of Siegel modular forms, due to the theta correspondence.
%	Every orthogonal Shimura variety~$X$ comes with a large supply of \emph{special cycles}, given by (combinations of) immersions in~$X$ of orthogonal Shimura varieties of smaller dimension.
%	The cohomology classes of these cycles are related to Siegel modular forms via the geometric theta correspondence, as illustrated by Kudla and Millson in their celebrated work~\cite{kudlamillson-intersection}.
%	They constructed the so-called \emph{Kudla--Millson theta function}~$\KMtheta{\genus}$ of genus~$\genus$, with~$\genus$ a positive integer.
%%	This is realized by the so-called \emph{Kudla--Millson theta function}~$\KMtheta{\genus}$ of genus~$\genus$, with~$\genus$ a positive integer.
%	It is a (non-holomorphic) Siegel modular form of genus~$\genus$ with values in the space of closed differential~$(\genus,\genus)$-forms on~$X$.
%%	
%%	The construction of~$\KMtheta{\genus}$ is due to Kudla and Millson~\cite{kudlamillson-harmonic1}, \cite{kudlamillson-harmonic2}, \cite{kudlamillson-tubes}, \cite{kudlamillson-intersection}.
%	Its cohomology class~$[\KMtheta{\genus}]$ is a \emph{holomorphic} Siegel modular form whose Fourier coefficients are the fundamental classes of codimension~$\genus$ special cycles in~$X$.
%	For this reason,~$[\KMtheta{\genus}]$ may be regarded as a generalization of the generating series of intersection numbers of curves on Hilbert modular surfaces considered by Hirzebruch and Zagier~\cite{hirzebruch-zagier}.
%%	Although the Kudla--Millson theta function is nowadays a classical tool in the study of Shimura varieties, there are still several open questions regarding its properties.
%%	It is the purpose of the present article to provide a decomposition of~$[\KMtheta{\genus}]$ of both arithmetic and geometric interest, as we now clarify.
	
	Although~$X$ may be non-compact, the cohomology group~$H^{2\genus}(X)$ admits a \emph{Lefschetz decomposition} for small~$g$.
	Also the space of Siegel modular forms of large weight admits a natural decomposition, the \emph{modular decomposition}.
	It splits as a direct sum of subspaces given by cusp forms, Siegel Eisenstein series, and Klingen Eisenstein series.
%	It is known that the Siegel Eisenstein part of~$[\KMtheta{\genus}]$ may be computed by integrating the theta function against a suitable power of the Kähler form of~$X$; see~\cite{kudla-integrals} and~\cite{kudla-speccyc}.
%	The Lefschetz decomposition of~$[\KMtheta{\genus}]$, as well as its modular decomposition, are currently not available in the literature.
	
	It is the goal of the present paper to provide the Lefschetz decomposition of the cohomology class~$[\KMtheta{\genus}]$, and to show that it is at the same time the modular decomposition of~$[\KMtheta{\genus}]$ as a Siegel modular form.
	More specifically, we prove that each term of the Lefschetz decomposition of~$[\KMtheta{\genus}]$ is given by the Klingen Eisenstein series arising from the primitive part of a Kudla--Millson theta function of intermediate genus.
	We also illustrate several geometric applications on orthogonal Shimura varieties.
	
%	\subsection{The Kudla--Millson theta function in the literature}
%	Since the seminal work of Kudla and Millson, the theta function~$\KMtheta{\genus}$ has been a fundamental tool for several research investigations.
%	For instance, it has been used as theta kernel of a lift of cusp forms~\cite{bruinier-converse}, \cite{zuffetti-unfolding}, \cite{stein-converse}, \cite{kieferzuffetti}, \cite{metzlerzuffetti}, and a lift of compactly supported differential forms~\cite{kudlamillson-intersection}.
%	This is known as the \emph{Kudla--Millson lift}, and is strictly related with Borcherds' lift~\cite{bruinierfunke-ontwogeometric}, \cite{bruinierfunke}.
%	
%	Effective cones of cycles are geometric invariants of algebraic varieties, notoriously hard to study.
%	Since special cycles span subcones of the effective cones, it is an interesting problem to understand their geometric properties.
%	These are studied by means of the Kudla--Millson theta function in~\cite{bruiniermöller}, \cite{zuffetti-conescod2} and~\cite{zuffetti-equid}.
%	
%	Further instances where the Kudla--Millson theta function is of fundamental importance is the proof of the Noether--Lefschetz Conjecture~\cite{noetherlefconj} and of Kudla's Conjecture~\cite{bruinierraum}.
%	The rigid meromorphic cocycles introduced by Darmon and Vonk~\cite{darmonvonk} provide test cases of an eventual analogue of modular generating series in the $p$-adic setting.
%	This emerging~``$p$-adic Kudla's Program'' has recently attracted much interest.
%	\\\textcolor{red}{\textbf{Question:} Shall we add other citations to some other papers?}
	
	\subsection{Lefschetz and Modular decompositions}	
	To illustrate the main results of the present paper we need to introduce some notation.
	Let~$L$ be an even lattice of signature~$(n,2)$.
	To simplify the exposition, through this introduction we assume that~$L$ is \emph{unimodular}.
	This condition will be dropped in the main body of the article.
	
	Let~$X$ be a (connected) orthogonal Shimura variety arising from~$L$, i.e.\ the quotient~${X=\Gamma\backslash D}$ for some finite index subgroup~$\Gamma\subseteq \OO^+(L)$, where~$D$ is the Hermitian symmetric domain of type IV attached to~$\OO(L)$ and~$\OO^+(L)$ is the group of isometries of the lattice~$L$ with positive spinor norm.
	Then~$X$ is a quasi-projective variety of dimension~$n$, which may have finite quotient singularities.
	If~$\Gamma$ is torsion free, then~$X$ is smooth.
%	If~$X$ is not smooth, it is usual to identify its cohomology with some invariant subspace of the cohomology of a smooth finite cover of~$X$.
	We denote by~$\omega$ the Kähler form on~$X$ arising as the first Chern form of the Hodge bundle.
	
	Let~$H^{2\genus}(X)$, resp.~$\sqH^{2\genus}(X)$, be the de Rham cohomology group, resp.~$L^2$-cohomology group, of~$X$ with complex coefficients.
%	Let~$H^{\genus,\genus}(X)$, resp.~$\sqH^{\genus,\genus}(X)$, be the~$(\genus,\genus)$-term of the Hodge decomposition of the de Rham cohomology group~$H^{2\genus}(X)$, resp.\ the~$L^2$ cohomology group~$\sqH^{2\genus}(X)$, with coefficients in~$\CC$.
	If~$\genus<(n-1)/2$, then the natural homomorphism
	\[
	\sqH^{2\genus}(X)\longrightarrow H^{2\genus}(X)
	\]
	is an isomorphism preserving the Hodge structures.
	The \emph{Lefschetz decomposition} of~$\sqH^{\genus,\genus}(X)$ induces a Lefschetz decomposition on~$H^{\genus,\genus}(X)$ of the form
	\be
	H^{\genus,\genus}(X)\label{eq;lefdecintrodisp}
%	=
%	\bigoplus_{j\ge 0} \lefschetz^j \Hprim^{r-2j}(X)
	=
	\CC[\omega]^\genus
	\oplus
	\lefschetz^{\genus-1} \Hprim^{1,1}(X)
	\oplus
	\cdots
	\oplus
	\lefschetz \Hprim^{\genus-1,\genus-1}(X)
	\oplus
	\Hprim^{\genus,\genus}(X).
	\ee
	Here we denoted by~$\Hprim^{*}(X)$ the primitive cohomology groups of~$X$ and by~$\lefschetz$ the Lefschetz operator given by wedging with~$\omega$; see Section~\ref{sec;digressdiffgeo} for further information.
	
	Let~$k=1+n/2$.
	From the well-known classification of unimodular lattices, it is easy to see that~$k$ is an \emph{even} integer.
	Let~$M^k_\genus$ be the space of weight~$k$ and genus~$\genus$ Siegel modular forms.
%	If all Fourier coefficients of a modular form~$f\in M^k_\genus$ indexed over singular matrices vanish, we say that~$f$ is a cusp form.
	The subspace of cusp forms of~$M^k_\genus$ is denoted by~$S^k_\genus$.
	
	From now on, we assume that~$k>2\genus$.
%	Concrete examples of modular forms are notoriously difficult to construct.
	The first classical example of an element of~$M^k_\genus$ is the Siegel Eisenstein series~$E^k_\genus$. 
	These series have been generalized by Klingen~\cite{klingen-defKES} to the nowadays called Klingen Eisenstein series; see Section~\ref{sec;KMtheta} for details.
	They arise from genus~$r$ cusp forms with~$r\leq\genus$.
	We denote by~$E^k_{\genus,r}(f)$ the Klingen Eisenstein series of weight~$k$ and genus~$g$ arising from~$f\in S^k_r$.
	The subspace of such Klingen Eisenstein series is denoted by~$M^k_{\genus,r}$.
	We remark that in the degenerate case of~$g=r$, the series~$E^k_{\genus,\genus}(f)$ equals~$f$, and that if~$\genus=0$, then~$E^k_{\genus,0}(1)=E^k_\genus$.
	
	The space of modular forms~$M^k_\genus$ admits a decomposition
	\be\label{eq;moddecintro}
	M^k_\genus
	= \CC E^k_\genus
	\oplus
	M^k_{\genus,1}
	\oplus
	\dots
	\oplus
	M^k_{\genus,\genus-1}
	\oplus
	S^k_\genus,
	\ee
	which we call the \emph{modular decomposition} of~$M^k_\genus$.
	
	\subsection{The main result}
	The cohomology class~$[\KMtheta{\genus}]$ of the Kudla--Millson theta function may be regarded as an element of~$M^k_\genus\otimes H^{\genus,\genus}(X)$.
	It is then a natural problem to find its Lefschetz decomposition with respect to~\eqref{eq;lefdecintrodisp}, as well as its modular decomposition with respect to~\eqref{eq;moddecintro}.
	
	The main result of the present paper, stated below, answers these questions.
%	In fact, we prove that the Lefschetz decomposition of~$[\KMtheta{\genus}]$ is also the modular decomposition of that class considered as a modular form.
	For~$0\leq r \leq \genus$, we denote by~$\klin_{\genus,r}^k\colon M^k_r\to M^k_\genus$ the linear map defined with respect to the modular decomposition of~$M^k_r$ as
	\[
	\klin_{\genus,r}^k\colon\sum_{i=0}^r
	E^k_{r,i}(f_i)
	\longmapsto
	\sum_{i=0}^r
	E^k_{\genus,i}(f_i),
	\qquad
	\text{for every~$f_i\in S^k_i$.}
	\]
	\begin{thm}\label{thm;mainresintro}
	Let~$g<(n+2)/4$ be a positive integer.
	The \emph{primitive part} of~$[\KMtheta{\genus}]$ is
	\be\label{eq;primpartintro}
	[\KMtheta{\genus}]_0
	=
	[\KMtheta{\genus}]
	+
	\lefschetz \circ \klin^k_{\genus,\genus-1}\big([\KMtheta{\genus-1}]\big)
	\ee
	and lies in~$S^k_\genus\otimes\Hprim^{\genus,\genus}(X)$.
	In particular, it is a Siegel cusp form with respect to the symplectic variable.
	The \emph{Lefschetz decomposition} of~$[\KMtheta{\genus}]$ equals the \emph{modular decomposition} of~$[\KMtheta{\genus}]$, and is given by
	\be\label{eq;lefdecintro}
	[\KMtheta{\genus}]
	=
	\sum_{r=0}^\genus
	(-\lefschetz)^{\genus-r}
	\circ
	E^k_{\genus,r}\big([\KMtheta{r}]_0\big).
	\ee
	Here the summand of index~$r$ lies in~$M^k_{\genus,r}\otimes\lefschetz^{\genus-r}\Hprim^{r,r}(X)$.
	\end{thm}
	\begin{rem}
	The Klingen Eisenstein series appearing in~\eqref{eq;lefdecintro} converge only if~$k>2\genus$.
	Since~$k=1+n/2$, the latter bound is equivalent to~$g<(n+2)/4$, which is the hypothesis of Theorem~\ref{thm;mainresintro}.
	It seems reasonable that the Lefschetz decomposition of~$[\KMtheta{\genus}]$ for larger genus~$g$ may still reflect some modular property.
	 To bypass convergence issues, one could implement the method of Weissauer~\cite{weissauer-88} and consider Klingen Eisenstein series in their domain of analytic continuation.
	We intend to come back to this problem in future work.
%	This problem will be addressed in a later publication.
	\end{rem}
	
	In this paper we prove a more general version of Theorem~\ref{thm;mainresintro}, in which~$L$ is not assumed to be unimodular; see Theorem~\ref{thm;lefschetzdecgengenus}.
	To do so, it is necessary to deal with vector-valued Siegel modular forms associated to the Weil representation, which are introduced in Section~\ref{sec;KMtheta}.
	
	We now provide an idea of the proof of Theorem~\ref{thm;mainresintro}.
	The first step is to compute the modular decomposition of~$[\KMtheta{\genus}]$ by induction on~$\genus$.
%	This is deduced by using that the Fourier coefficients of~$[\KMtheta{\genus}]$ are fundamental classes of special cycles, as well as arithmetic properties of the Fourier coefficients of Klingen Eisenstein series.
	We show that the image of~$[\KMtheta{\genus}]$ under the Siegel~$\Phi$ operator equals (up to a sign) the image of~$[\KMtheta{\genus-1}]$ under the Lefschetz operator.
	This is used to prove that the right-hand side of~\eqref{eq;primpartintro} is the cuspidal part of~$[\KMtheta{\genus}]$.
	
	The technical heart of the proof is to show that each term of the modular decomposition of~$[\KMtheta{\genus}]$ lies in a different term of the Lefschetz decomposition of~$H^{\genus,\genus}(X)$.
	To achieve this goal, it is sufficient to prove that the image of the \emph{Kudla--Millson lift}
	\[
	\KMlift{\genus}\colon S^k_\genus\longrightarrow H^{\genus,\genus}(X),
	\qquad
	f\longmapsto\pet{f}{[\KMtheta{\genus}]},
	\]
	is contained in the primitive cohomology~$\Hprim^{\genus,\genus}(X)$.
	Here~$\petempty$ denotes the Petersson inner product on~$S^k_\genus$.
	Theorem~\ref{thm;mainresintro} may then be deduced from the following result of independent interest.
	\begin{thm}\label{thm;imageKMliftintro}
	Let~$\genus<(n+2)/4$.
	The image of the Kudla--Millson lift~$\KMlift{\genus}$ is a subspace of the primitive cohomology group~$\Hprim^{\genus,\genus}(X)$.
	If the codimension~$\genus$ special cycles on~$X$ arising from~$L$ span~$H^{\genus,\genus}(X)$, then the image of~$\KMlift{\genus}$ is the whole~$\Hprim^{\genus,\genus}(X)$.
	\end{thm}
	We prove Theorem~\ref{thm;imageKMliftintro} by showing that
	\be\label{eq;vanintegintro}
	\int_X \KMlift{\genus}(f)\wedge\KMtheta{\genus-1}(\tau)\wedge \omega^{n-(2\genus-1)}=0
	\ee
	for all~$f\in S^k_\genus$ and all~$\tau\in \HH_{\genus-1}$.
	By means of the doubling method, we prove that~\eqref{eq;vanintegintro} follows from the vanishing of Garrett's map
	\[
	S^k_\genus\longrightarrow M^k_{\genus-1},\qquad
	f\longmapsto\bigpet{f}{E^k_{2\genus-1}\big(\begin{smallmatrix}
	\tau & 0 \\
	0 & \ast
	\end{smallmatrix}\big)},
	\]
	where we denote by~$\ast$ the variable of integration in~$\HH_\genus$.
	
	In Section~\ref{sec;adelicspcy} of this paper we also provide criteria on~$L$ so that the codimension~$g$ special cycles arising from~$L$ span~$H^{\genus,\genus}(X)$.
	This is deduced by means of the Noether--Lefschetz Conjecture, proved in~\cite{noetherlefconj}, and a study of the irreducible components of special cycles.
	
	\subsection{Applications of geometric interest}
	Immediate consequences of our main results are the following bounds for the dimension of cohomology groups of orthogonal Shimura varieties.
%	Let~$\varphi_\genus$ be the linear map defined on each term of the modular decomposition of~$M^k_\genus$ as
%	\[
%	\varphi_\genus\colon M^k_\genus\longrightarrow H^{\genus,\genus}(X),
%	\qquad
%	E^k_{\genus,r}(f)\longmapsto (-\lefschetz)^{\genus-r}\circ\KMlift{r}(f)\quad\text{for all~$f\in S^k_r$}.
%	\]
%	\textcolor{red}{Is it necessary to construct~$\varphi_g$?}
	\begin{cor}\label{cor;ifinjthenformcohointro}
	Let~$\genus<(n+2)/4$.
	Then
%	\min \{(\dim X+1)/4,(\dim X-1)/3\}.
%	then~$\modcoho{\genus}{\lattice}$ is \emph{surjective} and
	\[
	\dim\lefschetz^{\genus-r} H_0^{r,r}(X)\geq\dim \KMlift{r}\big(S^k_{r}\big)\qquad\text{for every~$0\leq r\leq\genus$}.
	\]
	If the Kudla--Millson lift~$\KMlift{r}$ is injective for every~$r\leq \genus$, then
%	so is~$\varphi_\genus$ and
	\[
	\dim\lefschetz^{\genus-r} H_0^{r,r}(X)\geq\dim S^k_{r}
	\qquad\text{and}\qquad
	\dim H^{\genus,\genus}(X)\geq\dim M^k_\genus.
	\]
%	If moreover~$L$ splits off~$E_8^{\oplus\genus-1}\oplus U^{\oplus 2}$, then~$\modcoho{\genus}{\lattice}$ is also surjective and
%	\[
%	\dim\lefschetz^{\genus-r} H_0^{r,r}(X)=\dim S^k_{r}
%	\qquad\text{and}\qquad
%	\dim H^{\genus,\genus}(X)=\dim M^k_\genus.
%	\]
%	---------------
%	Let~$\genus<(n+2)/4$ and let~$X=\OO^+(L)\backslash D$.
%	Then~$\varphi_\genus$ is \emph{surjective} and
%	\[
%	\dim\lefschetz^{\genus-r} \Hprim^{r,r}(X)=\dim \KMlift{r}(S^k_r)\qquad\text{for every~$0\leq r\leq\genus$}.
%	\]
%	If the Kudla--Millson lift~$\KMlift{r}$ is injective for every~$r\leq \genus$, then~$\varphi_\genus$ is an \emph{isomorphism} and
%	\[
%	\dim\lefschetz^{\genus-r} \Hprim^{r,r}(X)=\dim S^k_{r}
%	\qquad\text{and}\qquad
%	\dim H^{\genus,\genus}(X)=\dim M^k_\genus.
%	\]
	\end{cor}
	
	The injectivity of~$\KMlift{\genus}$ in the case of genus~$\genus=1$ has been proved in~\cite{bruinierbook}, \cite{bruinierfunke}, \cite{zuffetti-unfolding}.
%	 and has been generalized to locally symmetric spaces of orthogonal type in~\cite{bruinierfunke}, \cite{stein-converse} and~\cite{metzlerzuffetti}.
%	
%	It is known that~$\KMlift{\genus}$ is injective if~$\genus=1$ and if~$L$ splits off two hyperbolic planes~\cite{bruinierbook}.
%	The latter condition is always satisfied for~$L$ unimodular.
%	This injectivity criterion in genus~$1$ has been generalized to a larger families of lattices in~\cite{bruinier-converse} and~\cite{zuffetti-unfolding}, as well as to locally symmetric spaces of orthogonal type in~\cite{bruinierfunke}, \cite{stein-converse} and~\cite{metzlerzuffetti}.
	Recently, Kiefer and the second author proved in~\cite{kieferzuffetti} the injectivity of~$\KMlift{\genus}$ in the case of genus~$\genus=2$ under the assumption that~$n>9$.
%	splits off two hyperbolic planes and a positive definite lattice satisfying certain local conditions.
	The papers cited above provide criteria on the injectivity of the lift that work even without assuming~$L$ to be unimodular.
	
	As an interesting example, we consider the moduli spaces~$X_d$ of degree~$2d$ quasi-polarized~K3 surfaces.
	They can be viewed as orthogonal Shimura varieties arising from even lattices of the form
	\[
	L_d=\langle 2d\rangle \oplus U\oplus U\oplus E_8\oplus E_8,
	\]
	where~$U$ and~$E_8$ are respectively the hyperbolic lattice of signature~$(1,1)$ and the~$8$th rooth lattice.
	The lattices~$L_d$ are not unimodular, but satisfy the known injectivity criteria for~$\KMlift{1}$ and~$\KMlift{2}$.
	Since the spaces~$H^{1,1}(X_d)$ and~$H^{2,2}(X_d)$ are spanned by classes of special cycles arising from~$L_d$, we may deduce the following result from Theorem~\ref{thm;imageKMliftintro} and Corollary~\ref{cor;ifinjthenformcohointro}.
	Let~$S^k_{\genus,\lattice}$ be the space of weight~$k$ and genus~$\genus$ vector-valued cusp forms with respect to the Weil representation arising from~$\lattice$.
	\begin{cor}\label{cor;cohomodspintro}
	Let~$X_d$ be the moduli space of degree~$2d$ quasi-polarized K3 surfaces.
	Then
	\[
	\dim H^2(X_d)= 1 + \dim S^k_{1,L_d}
	\qquad\text{and}\qquad
	\dim H^4(X_d)=1+\dim S^k_{1,L_d} + \dim S^k_{2,L_d}.
	\]
%	\begin{align*}
%	& \dim H^{1,1}_0(X)=\dim S^k_{1,L_d},\qquad
%	\dim H^{1,1}(X)= 1 + \dim S^k_{1,L_d},
%	\\
%	&\dim H^{2,2}_0(X)=\dim S^k_{2,L_d},\qquad
%	\dim \lefschetz H^{1,1}_0(X)=\dim S^k_{1,L_d},
%	\\
%	&\dim H^{2,2}(X)=1+\dim S^k_{1,L_d} + \dim S^k_{2,L_d}.
%	\end{align*}
	\end{cor}
	The dimension formula for~$H^2(X_d)$ above was proved in~\cite{bruinier-converse} and~\cite{noetherlefconj}.
	Our article provides a different proof, which has the advantage of working also for cohomology groups of higher degree.
	
	
%	\textcolor{red}{
%	\begin{itemize}
%	\item
%	Moduli spaces of K3s are not smooth...
%	We defined a Lefschetz decomposition only for smooth OSV.
%	Perhaps it is better to rephrase Corollary~\ref{cor;cohomodspintro} with only~$\dim H^{4}(X)$?
%	Also~$H^{2,2}(X)$ perhaps is something we could not really talk about...
%	\\
%	\textcolor{blue}{\textbf{I think we could argue as follows:}
%	Let~$\Gamma'\leq \Gamma\coloneqq\widetilde{\SO}^+(L)$ be neat, so that~$X_{\Gamma'}\coloneqq \Gamma'\backslash\domain$ is a smooth orthogonal Shimura variety, and let~$X=\Gamma\backslash\domain$ be a moduli space of quasi-polarized K3 surfaces.
%	Let
%	\[
%	\pi\colon X_{\Gamma'}\longrightarrow X
%	\]
%	be the quotient map.
%	This is a finite cover of~$X$.
%	The pullback map is injective, and one has that
%	\[
%	H^{2\genus}(X) \cong H^{2\genus}(X_{\Gamma'})^{\Gamma/\Gamma'}.
%	\]
%	If we assume that~$\genus < n/4$, then by~\cite[Theorem~$3.6$]{noetherlefconj} we deduce that~$H^{2\genus}(X_{\Gamma'})=H^{\genus,\genus}(X_{\Gamma'})$ and it is generated by special cycles.
%	\\
%	We now prove that every element of~$H^{2\genus}(X_{\Gamma'})$ is~$\Gamma$-invariant, which then implies that
%	\[
%	\dim H^{2\genus}(X)=\dim H^{2\genus}(X_{\Gamma'})
%	=
%	\dim\mathcal{M}^k_{\genus,\lattice}.
%	\]
%	We remark that
%	\[
%	H^{2\genus}(X_{\Gamma'})
%	=
%	\varphi_{\genus,\lattice}(\mathcal{M}^k_{\genus,\lattice}),
%	\]
%	where~$\varphi_{\genus,\lattice}\colon E^k_{\genus,r,\lattice}(f)=(-\lefschetz)^{\genus-r}\circ\KMlift{r}(f)$ for every~$f\in S^k_{r,L}$.
%	Since the Kähler form~$\omega$ is~$\Gamma$-invariant, and since~$\KMtheta{r}$ (and hence also its cohomology class) is invariant for every~$r\leq \genus$, it is easy to see that by our main theorem also the primitive part~$[\KMtheta{r}]_0$ is~$\Gamma$-invariant.
%	This implies that both~$\lefschetz$ and~$\KMlift{r}$ are operators whose images are~$\Gamma$-invariant.
%	This implies that the image of~$\varphi_{\genus,\lattice}$, i.e.~$H^{2\genus}(X_{\Gamma'})$, is made of classes arising from~$\Gamma$-invariant differential forms.
%	}
%	\end{itemize}
%	\textbf{To do:}\\
%	Correct dimension formulas for cohomology groups accordingly.
%	}
	
	\subsection*{Acknowledgments}
	We would like to thank Siegfried Böcherer, Soumya Das, Jens Funke, Paul Kiefer, Zhiyuan Li, Martin Möller, Riccardo Salvati-Manni, Oliver Stein and Brandon Williams for useful conversations on the topics of the present paper.
	The authors acknowledge support by Deutsche Forschungsgemeinschaft (DFG, German Research Foundation) through the Collaborative Research Centre TRR 326 \textit{Geometry and Arithmetic of Uniformized Structures}, project number 444845124.
%	This work has been partially funded by the Collaborative Research Centre TRR~$326$ ``Geometry and Arithmetic of Uniformized Structures'', project number~$444845124$.
	\end{document}
	
	
	
	
	
	
	
\documentclass[../Lefschetz_decomposition.tex]{subfiles}

	\begin{document}
	In this section we gather all the properties of vector-valued Siegel modular forms with respect to the Weil representation relevant for the purposes of the paper.
%	Along the way, we define the Siegel operator, Klingen Eisenstein and Poincaré series, and Garret's map.
	Although the following theory of vector-valued modular forms can be extended to lattices of general signature, we prefer to work in signature~$(n,2)$, since this is the only case needed in the rest of the present article.
	For the same reason, we only consider weights~$k\in\frac{1}{2}\ZZ$ such that~$2k\equiv n-2$ mod~$4$.
%	\\\textcolor{blue}{\textbf{Some remarks:}
%	\begin{itemize}
%	\item It is easy to see that if~$2k\not\equiv n-2$ mod~$2$, then the space of weight~$k$ and genus~$\genus$ Siegel modular form~$M^k_{\genus,\lattice}$ is trivial.
%	In fact, this follows from the modular transformation of Siegel modular forms and the fact that if~$R=\big(\begin{smallmatrix}
%	A & \\  & A^{-t}
%	\end{smallmatrix}\big),i\big)$ with~$A=\big(\begin{smallmatrix}
%	-1 & \\
%	 & I_{\genus-1}
%	\end{smallmatrix}\big)$, then~$\weil{\lattice}{\genus}(R^2)\mathfrak{e}_\discel = (-1)^{2-n}\mathfrak{e}_\discel$.
%	\item I do not think that if~$2k\equiv n$ mod~$4$, then everything works as for the case~$2k\equiv n-2$ mod~$4$.
%	At least, the blue check that the Klingen Eisenstein series of Definition~\ref{def;KlingEisseries} are well-defined does not work without that hypothesis.
%	\end{itemize}
%	}
	
	\subsection{The Weil representation}
	
	The Siegel upper half-space of genus~$\gengenus$ is defined as
	\[
	\HH_{\gengenus}=\{\tau=x+iy\in\CC^{\gengenus\times \gengenus} : \text{$\tau=\tau^t$ and $y>0$}\}.
	\]
	Here we denote by~$x$ and~$y$ respectively the real and imaginary part of~$\tau$.
	Moreover, we denote by~$M^t$ the transpose of any matrix~$M$.
	
	If~$z\in\CC$, we write~$e(z)=\exp(2\pi i z)$ and denote by~$z^{1/2}=\sqrt{z}$ the principal branch of the square root, so that~$\arg(\sqrt{z})\in(-\pi/2,\pi/2]$.
	For every integer~$m$ we define~$z^{m/2}=(z^{1/2})^m$.
	
	We write~$\Mp_{2\gengenus}(\RR)$ for the metaplectic cover of the symplectic group~$\Sp_{2\gengenus}(\RR)$.
	This is realized by choosing a square root of the map~$\tau\mapsto \det(C\tau+D)$ for every~$\big(\begin{smallmatrix}
	A & B\\
	C & D
	\end{smallmatrix}\big)\in\Sp_{2\gengenus}(\RR)$.
	The group law of~$\Mp_{2\gengenus}(\RR)$ is
	\be\label{eq;grouplawmetgr}
	\big(M_1,\phi_1(\tau)\big)\cdot \big(M_2,\phi_2(\tau)\big)
	=
	\big(M_1M_2 , \phi_1\big(M_2\cdot\tau\big)\phi_2(\tau)\big),
	\ee
	where the action of~$\Mp_{2\gengenus}(\RR)$ on~$\HH_\gengenus$ is the one induced by the standard action of~$\Sp_{2\gengenus}(\RR)$.
	
	The standard generators of the group~$\Mp_{2\gengenus}(\ZZ)$ are the elements of the form
	\be\label{eq;stdgenofmetcov}
	S=\left(\left(\begin{smallmatrix}
	0 & -I_\gengenus\\
	I_\gengenus & 0
	\end{smallmatrix}\right),\sqrt{\det\tau}\right)
	\qquad\text{and}\qquad
	T_B=\left(\left(\begin{smallmatrix}
	I_\gengenus & B \\
	0 & I_\gengenus
	\end{smallmatrix}\right),1\right)\quad\text{where~$B\in\Sym_\genus(\ZZ)$.}
	\ee
%	The center of~$\Mp_{2\gengenus}(\ZZ)$ contains
%	\[
%	Z\coloneqq\begin{cases}
%	\left(-I_{2\gengenus},i\right), & \text{if~$g$ is odd,}\\
%	\left(-I_{2\gengenus},-1\right), & \text{if~$g$ is even.}
%	\end{cases}
%	\]
%	Note that~$S^2=Z$.
%	\textcolor{red}{\textbf{Question:} Is the last sentence true?}\\
%	\textcolor{blue}{\textbf{Answer from Jan:} Since the function on the right-hand side should be holomorphic, it must be holomorphic also on~$iI_{2\gengenus}$.
%	Then, just evaluate on that and check the sign you got.
%	Probably we should get~$Z=(I_{2\gengenus}, i^\gengenus)$.}
%	\\\textcolor{blue}
%	{
%	\textbf{Here some computations on previous question.}
%	\\
%	We have that
%	\[
%	S^2=\Big(-I_{2\gengenus} , \sqrt{\det(-\tau^{-1})}\cdot\sqrt{\det\tau}\Big).
%	\]
%	Since we use the principal branch of the square root, if we write~$\det\tau=\rho e^{i\theta}$ in polar coordinates choosing~$\theta\in(-\pi,\pi]$, then~$\sqrt{\det\tau}=\sqrt{\rho}e^{i\theta/2}$.
%	In fact~$\arg(e^{i\theta/2})\in(-\pi/2,\pi/2]$.
%	\\
%	The point is that in general~$\sqrt{\det(-\tau^{-1})}\neq\sqrt{\det\tau}^{-1}$.
%	We may check this by writing~$\det\tau=\rho e^{i(\theta + 2k\pi)}$, where~$\theta$ as above and~$k\in \ZZ$, and computing
%	\[
%	\sqrt{\det\tau}^{-1} = \rho^{-1/2}e^{-i\theta/2}
%	\]
%	and
%	\bas
%	\sqrt{\det(-\tau^{-1})}
%	=
%	\sqrt{\det(-I_{\gengenus})\det(\tau)^{-1}}
%	=
%	\sqrt{(-1)^\gengenus \rho^{-1}e^{-i(\theta + 2k\pi)}}.
%	\eas
%	Since~$-1=e^{i\pi}$, we deduce that
%	\bas
%	\sqrt{\det(-\tau^{-1})}
%	=
%	\sqrt{e^{i\gengenus\pi} \rho^{-1}e^{-i(\theta + 2k\pi)}}
%	=
%	\rho^{-1/2}
%	\sqrt{e^{-i(\theta -\gengenus\pi + 2k\pi)}}.
%	\eas
%	It is then clear that the two square roots of~$\det(-\tau^{-1})$ are
%	\bas
%	\rho^{-1/2}
%	e^{-i(\theta/2 -\gengenus\pi/2 + k\pi)}
%	=
%	\begin{cases}
%	\pm \rho^{-1/2} e^{-i\theta/2}, & \text{if~$g$ is even}\\
%	\pm i \rho^{-1/2} e^{-i\theta/2}, & \text{if~$g$ is odd}
%	\end{cases}
%	=
%	\begin{cases}
%	\pm \sqrt{\det\tau}^{-1}, & \text{if~$g$ is even,}\\
%	\pm i \sqrt{\det\tau}^{-1}, & \text{if~$g$ is odd}
%	\end{cases}
%	\eas
%	But how to choose the one that gives the principal branch?
%	Perhaps by working case by case on~$\theta$?
%	}\\
	For future purposes we introduce here also the elements of~$\Mp_{2\gengenus}(\ZZ)$ given by
	\be\label{eq;defRgenmetapl}
	R_A=\left(\left(\begin{smallmatrix}
	A & 0\\
	0 & A^{-t}
	\end{smallmatrix}\right), \det A^{1/2}\right)\qquad\text{for every~$A\in\GL_\gengenus(\ZZ)$,}
	\ee
	where~$A^{-t}=(A^{-1})^t = (A^t)^{-1}$.
	
	Let~$(\lattice,q)$ be an even indefinite lattice of signature~$(n,2)$, and let~$L'$ be its dual lattice.
	We denote by~$(\cdot{,}\cdot)$ the bilinear form associated to~$q$ such that~$q(\lambda)=(\lambda,\lambda)/2$ for every~$\lambda\in\lattice$.
	If~$\lambda\in \lattice^\genus$, we denote by~$q(\lambda)$ the~$\genus\times\genus$ moment matrix of the tuple~$\lambda$.
	The discriminant group of~$\lattice$ is the quotient~$D_\lattice=\lattice'/\lattice$.
	The form~$q$ induces a~$\QQ/\ZZ$-valued moment form on~$D_\lattice^\genus$ which we still denote by~$q$.
	
	We endow the group algebra~$\CC[D_\lattice]$ with the Hermitian scalar product~$\langle\cdot{,}\cdot\rangle$ defined as
	\be\label{eq;stinprodgroupalg}
	\Big\langle
	\sum_{\discel\in D_\lattice}\lambda_\discel \mathfrak{e}_\discel
	,
	\sum_{\discel\in D_\lattice}\mu_\discel \mathfrak{e}_\discel
	\Big\rangle
	=
	\sum_{\discel\in D_\lattice}\lambda_\discel\overline{\mu_\discel},
	\ee
	where~$(\mathfrak{e}_\discel)_\discel$ is the standard basis of~$\CC[D_\lattice]$.
	We denote with the same symbol also the extension of that scalar product to~$\CC[D_\lattice^\gengenus]$ which makes different copies of~$D_\lattice$ orthogonal to each others. 
	
	\begin{defi}
	The \emph{Weil representation} of genus~$\gengenus$ associated to the lattice~$\lattice$ is the unitary representation
	\[
	\weil{\lattice}{\gengenus}\colon \Mp_{2\gengenus}(\ZZ) \longrightarrow \Aut(\CC[D_\lattice^\gengenus])
	\]
	defined by
	\begin{align*}
	\weil{\lattice}{\gengenus}(T_B)\mathfrak{e}_\discel &=
	e\big(
	\trace (q(\discel)B)
	\big)\mathfrak{e}_\discel\qquad\text{for every~$B\in\Sym_\gengenus(\ZZ)$,}
	\\
	\weil{\lattice}{\gengenus}(S)\mathfrak{e}_\discel &=
	\frac{\sqrt{i}^{\gengenus(2-n)}}{|D_\lattice|^{\gengenus/2}}
	\sum_{\disceltwo\in D_\lattice^\gengenus} e\big(-\tr ( \disceltwo,\discel)\big)\mathfrak{e}_\disceltwo,
	\\
	\weil{\lattice}{\gengenus}(R_A)\mathfrak{e}_\discel
	&=
	(\det A)^{(2-n)/2}\mathfrak{e}_{\discel A^{-1}}\qquad
	\text{for every~$A\in\GL_\gengenus(\ZZ)$},
	\end{align*}
	see e.g.~\cite{borcherds-grassmannians} and~\cite{zh;phd}.
	\end{defi}
	
	We remark that
	\be\label{eq;howto-phi}
	\weil{\lattice}{\gengenus}\big((M,-\phi)\big)\mathfrak{e}_\discel
	=
	(-1)^{2-n}\weil{\lattice}{\gengenus}\big( (M,\phi)\big)\mathfrak{e}_\discel
	\qquad\text{for every~$(M,\phi)\in\Mp_{2\gengenus}(\ZZ)$.}
	\ee
	This follows from the fact that~$(M,-\phi)=(M,\phi)\cdot R_A^2$, where~$R_A$ is chosen as in~\eqref{eq;defRgenmetapl} with~$A=\big(
	\begin{smallmatrix}
	-1 & \\
	 & I_{\gengenus-1}
	\end{smallmatrix}
	\big)$.
	
	Weil representations of different genera associated to the same lattice~$L$ are strictly related.
	To clarify this, we need to introduce another piece of notation.
	We identify~$\CC[D_\lattice^r]\otimes\CC[D_\lattice^{\gengenus-r}]$ with~$\CC[D_\lattice^\gengenus]$ under the isomorphism
	\be\label{eq;identgroupalg}
	\CC[D_\lattice^r]\otimes\CC[D_\lattice^{\gengenus-r}]\longrightarrow\CC[D_\lattice^\gengenus],\qquad \mathfrak{e}_\discel\otimes\mathfrak{e}_\disceltwo\longmapsto\mathfrak{e}_{(\discel,\disceltwo)}.
	\ee
	Let~$\iota$ be the inclusion of metaplectic groups given as
	\ba\label{eq;inclprodsimpl}
	\iota\colon\Mp_{2r}(\ZZ)\times\Mp_{2(\gengenus - r)}(\ZZ)&\longrightarrow\Mp_{2\gengenus}(\ZZ),
	\\
	\Big(\big(\big(\begin{smallmatrix}
	A & B\\
	C & D
	\end{smallmatrix}
	\big),\phi\big),\big(\big(
	\begin{smallmatrix}
	A' & B'\\
	C' & D'
	\end{smallmatrix}\big),\phi'\big)\Big)
	&\longmapsto
	\left(\left(\begin{smallmatrix}
	A & & B & \\
	 & A' & & B' \\
	C & & D & \\
	 & C' & & D'
	\end{smallmatrix}\right),\tilde{\phi}\right),
	\ea
	where the function~$\tilde{\phi}(\tilde{\tau})$, with~$\tilde{\tau}\in\HH_\gengenus$, is~$\pm\det\big(\big(\begin{smallmatrix}
	C & \\
	 & C'
	\end{smallmatrix}\big)\tilde{\tau} + \big(\begin{smallmatrix}
	D & \\
	 & D'
	\end{smallmatrix}\big)\big)^{1/2}$ with sign chosen such that~$\tilde{\phi}\big(\begin{smallmatrix}
	\tau & 0 \\
	0 & \tau'
	\end{smallmatrix}\big)= \phi(\tau)\cdot\phi'(\tau')$.
%	\\\textcolor{blue}{\textbf{Note:} If~$R,R'$ are the standard block-diagonal generators as above, then~$\iota(R,R')$ is in general not a standard generator.
%	Take the case where~$R=\big(\big(\begin{smallmatrix}
%	-1 & 0\\
%	0 & -1
%	\end{smallmatrix}\big),i\big)=R'$.}\\
	
	Weil representations of different genera attached to the same lattice behave well with respect to the inclusion~$\iota$.
	This is illustrated in the following result, whose proof is left to the reader and may be deduced with the same idea of~\cite[Lemma~$3.2$]{kieferzuffetti}.
	\begin{lemma}\label{lemma;proofweilreponsubmeta}
	Let~$0\le r\le g$, and let~$\iota$ be the inclusion as in~\eqref{eq;inclprodsimpl}.
	If~$\gamma\in\Mp_{2r}(\ZZ)$ and~$\gamma'\in\Mp_{2(\gengenus-r)}(\ZZ)$, then
	\be\label{eq;inlcfmet2}
	\weil{\lattice}{\gengenus}\big(\iota(\gamma,\gamma')\big)(\mathfrak{e}_\discel\otimes\mathfrak{e}_\disceltwo)
	=
	\big(\weil{\lattice}{r}(\gamma)\mathfrak{e}_\discel\big)\otimes
	\big(\weil{\lattice}{\gengenus-r}(\gamma')\mathfrak{e}_\disceltwo\big)
	\ee
	for every~$\discel\in D_\lattice^r$ and~$\disceltwo\in D_\lattice^{\gengenus-r}$.
	\end{lemma}
%	\begin{proof}
%	\textcolor{blue}{By~\eqref{eq;howto-phi} it is enough to work with~$\gamma=(M,\phi)$ and~$\gamma'=(M',\phi')$ chosen such that~$\phi$ and~$\phi'$ are principal square roots.
%	Hence, it is enough to check~\eqref{eq;inlcfmet2} replacing~$\gamma$ and~$\gamma'$ with the standard generators of respectively~$\Mp_{2r}(\ZZ)$ and~$\Mp_{2(\gengenus-r)}(\ZZ)$.}
%	
%	\textcolor{blue}{If~$\gamma=T$ and~$\gamma'=T'$, where~$T$ and~$T'$ arise from respectively~$B\in\Sym_r(\ZZ)$ and~$B'\in\Sym_{\gengenus-r}(\ZZ)$ as in~\eqref{eq;stdgenofmetcov}, then~$\tilde{T}\coloneqq\iota(T,T')$ is as in~\eqref{eq;stdgenofmetcov} as well, and arises from the symmetric matrix~$\tilde{B}=\big(\begin{smallmatrix}
%	B & 0\\
%	0 & B'
%	\end{smallmatrix}\big)$.
%	We may then compute that
%	\bas
%	\weil{\lattice}{\gengenus}\big(\tilde{T}\big)(\mathfrak{e}_\discel\otimes\mathfrak{e}_\disceltwo)
%	&=
%	e\Big(
%	\tr \big(q(\discel,\disceltwo)\big(\begin{smallmatrix}
%	B & \\
%	 & B'
%	\end{smallmatrix}\big)
%	\big)\Big) \mathfrak{e}_\discel\otimes\mathfrak{e}_\disceltwo
%	=
%	e\Big(\tr \Big(
%	\begin{smallmatrix}
%	q(\discel)B & \ast \\
%	\ast & q(\disceltwo)B'
%	\end{smallmatrix}
%	\Big)\Big)\mathfrak{e}_\discel\otimes\mathfrak{e}_\disceltwo
%	\\
%	&=
%	\Big(e\big(\tr(q(\discel)B)\big)\mathfrak{e}_\discel\Big) \otimes
%	\Big(e\big(\tr(q(\disceltwo)B')\big)\mathfrak{e}_\disceltwo\Big)
%	\\
%	&=
%	\big(\weil{\lattice}{r}(T)\mathfrak{e}_\discel\big)\otimes\big(\weil{\lattice}{\gengenus-r}(T')\mathfrak{e}_\disceltwo\big).
%	\eas}
%	
%	\textcolor{blue}{If~$\gamma=R$ and~$\gamma'=R'$, where~$R$ and~$R'$ are as in~\eqref{eq;defRgenmetapl} arising from respectively~$A\in\GL_r(\ZZ)$ and~$A'\in\GL_{\gengenus-r}(\ZZ)$, then~$\tilde{R}\coloneqq\iota(R,R')$ is as in~\eqref{eq;defRgenmetapl} as well, and arises from~$\tilde{A}=\big(\begin{smallmatrix}
%	A & \\
%	 & A'
%	\end{smallmatrix}\big)$.
%	In this case we have that
%	\bas
%	\weil{\lattice}{\gengenus}(\tilde R)(\mathfrak{e}_\alpha\otimes\mathfrak{e}_\beta)
%	&=
%	(\det \tilde A)^{(2-n)/2}\mathfrak{e}_{(\alpha,\beta)\tilde A^{-1}}
%	=
%	(\det A)^{(2-n)/2}(\det A')^{(2-n)/2}\mathfrak{e}_{\alpha A^{-1}}\otimes\mathfrak{e}_{\beta A'^{-1}}
%	\\
%	&=
%	\big(\weil{\lattice}{r}(R)\mathfrak{e}_\alpha\big)\otimes
%	\big(\weil{\lattice}{\gengenus - r}(R')\mathfrak{e}_\beta\big).
%	\eas}
%	
%	\textcolor{blue}{Since~$\iota(\gamma,\gamma')=\iota(\gamma,I_{\gengenus-r})\cdot\iota(I_r,\gamma')$, the previous check implies that~\eqref{eq;inlcfmet2} is satisfied also if~$(\gamma,\gamma')$ is made of standard generators which are not of the same kind, e.g.~$\gamma=T$ and~$\gamma'=R'$.
%	In fact, to conclude the proof it is enough to check~\eqref{eq;inlcfmet2} with~$\gamma=S$ and~$\gamma'=I_{\gengenus-r}$.
%	Since~$\iota(S,I_{\gengenus-r})$ is not a standard generator in~$\Mp_{2\gengenus}(\ZZ)$, the check of~\eqref{eq;inlcfmet2} is not immediate as in the previous cases.
%	We proceed as in the proof of~\cite[Lemma~$3.2$]{kieferzuffetti}.}
%	
%	\textcolor{blue}{To avoid confusion, we denote by~$S_r$ the standard generator of~$\Mp_{2r}(\ZZ)$ arising from $\big(\begin{smallmatrix}
%	 & -I_r \\
%	 I_r &
%	\end{smallmatrix}\big)$.
%	Let~$T_r(-I_r)\in\Mp_{2r}(\ZZ)$ be as in~\eqref{eq;stdgenofmetcov} with~$B=-I_r$.
%	We may rewrite
%	\be\label{eq;lowtriang}
%	S_r=T_r(-I_r)\cdot\underbrace{\Big(\big(\begin{smallmatrix}
%	I_r & 0\\
%	I_r & I_r
%	\end{smallmatrix}\big) , \sqrt{\det(\tau + I_r)}\Big)}_{M}{}\cdot T_r(-I_r).
%	\ee
%	We compute the action of~$\iota(M,I_{\gengenus-r})\in\Mp_{2\gengenus}(\ZZ)$, where~$M$ is the lower-triangular metaplectic element appearing in~\eqref{eq;lowtriang}.
%	We denote by~$\tau$ also the top-left~$r\times r$ block of~$\tilde\tau\in\HH_\gengenus$, and rewrite
%	\bas
%	\iota(M,I_{\gengenus-r})
%	&=
%	\left(
%	\left(
%	\begin{smallmatrix}
%	I_r & & & \\
%	& I_{\gengenus-r} & & \\
%	I_r & & I_r & \\
%	& & & I_{\gengenus-r}
%	\end{smallmatrix}
%	\right),\sqrt{\det(\tau + I_r)}
%	\right)
%	=
%	S_{\gengenus}^{-1} \cdot T_{\gengenus}(-I_r) \cdot S_{\gengenus},
%	\eas
%	where~$S_{\gengenus}$ is the standard generator of~$\Mp_{2\gengenus}(\ZZ)$ that arises from~$\big(\begin{smallmatrix}
%	 & -I_{\gengenus} \\
%	I_{\gengenus} & 
%	\end{smallmatrix}\big)$ as in~\eqref{eq;stdgenofmetcov},
%	\bas
%	S_{\gengenus}^{-1}\coloneqq\left(\left(\begin{smallmatrix}
%	& I_\gengenus\\
%	-I_\gengenus &
%	\end{smallmatrix}\right),
%	\det (-\tilde\tau)^{1/2}\right)
%	\qquad\text{and}\qquad
%	T_{\gengenus}(-I_r)\coloneqq\left(\left(\begin{smallmatrix}
%	I_r & & -I_r & \\
%	& I_{\gengenus-r} & & \\
%	& & I_r & \\
%	& & & I_{\gengenus-r}
%	\end{smallmatrix}\right),1\right).
%	\eas
%	We may then compute
%	\bas
%	\weil{\lattice}{\gengenus}\big(\iota(M,I_{\gengenus-r})\big)\mathfrak{e}_\discel
%	&=
%	\weil{\lattice}{\gengenus}(S_\gengenus^{-1})
%	\weil{\lattice}{\gengenus}\big(T_\gengenus(-I_r)\big)
%	\weil{\lattice}{\gengenus}(S_\gengenus)\mathfrak{e}_\discel
%	\\
%	&=
%	\weil{\lattice}{\gengenus}(S_\gengenus^{-1})\frac{i^{\gengenus(2-n)/2}}{|D_\lattice|^{g/2}}\sum_{\disceltwo\in D_L^\gengenus} e\big(-\tr(\disceltwo,\discel)\big)
%	\weil{\lattice}{\gengenus}\big(T_\gengenus(-I_r)\big)\mathfrak{e}_\disceltwo
%	\\
%	&=
%	\frac{i^{\gengenus(2-n)/2}}{|D_\lattice|^{\gengenus/2}}\sum_{\disceltwo\in D_\lattice^\gengenus} e\Big(
%	-\tr(\disceltwo,\discel)-\tr \big(q(\disceltwo_r)\big)
%	\Big)
%	\weil{\lattice}{\gengenus}(S_\gengenus^{-1})\mathfrak{e}_\disceltwo,
%	\eas
%	where we denote by~$\disceltwo_r$ the tuple given by the first~$r$ entries of~$\disceltwo\in D_\lattice^\gengenus$.
%	Since
%	\[
%	\weil{\lattice}{\gengenus}(S_\gengenus^{-1})\mathfrak{e}_\disceltwo
%	=
%	(-1)^{\gengenus(n-2)/2}\weil{\lattice}{\gengenus}(S_\gengenus)\mathfrak{e}_{-\disceltwo},
%	\]
%	we may further compute
%	\bas
%	\weil{\lattice}{\gengenus}\big(\iota(M,I_{\gengenus-r})\big)\mathfrak{e}_\discel
%	&=
%	\frac{(-1)^{\gengenus(n-2)/2} i^{\gengenus(2-n)/2}}{|D_\lattice|^{\gengenus/2}}
%	\sum_{\disceltwo\in D_\lattice^\gengenus}
%	e\Big(
%	\tr (\disceltwo,\discel) - \tr\big(q(\disceltwo_r)\big)
%	\Big)
%	\weil{\lattice}{\gengenus}(S_\gengenus)\mathfrak{e}_\disceltwo
%	\\
%	&=
%	\frac{1}{|D_\lattice|^{\gengenus}}
%	\sum_{\disceltwo\in D_\lattice^\gengenus}
%	e\Big(
%	\tr (\disceltwo,\discel) - \tr\big(q(\disceltwo_r)\big)
%	\Big)
%	\sum_{\discelthree\in D_\lattice^\gengenus}e\big(-\tr(\disceltwo,\discelthree)\big)\mathfrak{e}_\discelthree
%	\\
%	&=
%	\frac{1}{|D_\lattice|^\gengenus}\Big[
%	\sum_{\disceltwo_r,\discelthree_r\in D_\lattice^r}
%	e\Big(
%	\tr(\disceltwo_r,\discel_r-\discelthree_r) - \tr\,q(\disceltwo_r)
%	\Big)
%	\mathfrak{e}_{\discelthree_r}\Big]
%	\\
%	&\quad\otimes\Big[
%	\sum_{\disceltwo_{\gengenus-r},\discelthree_{\gengenus-r}\in D_\lattice^{\gengenus-r}}
%	e\Big(
%	\tr (\disceltwo_{\gengenus-r},\discel_{\gengenus-r} - \discelthree_{\gengenus-r})
%	\Big)
%	\mathfrak{e}_{\discelthree_{\gengenus-r}}
%	\Big].
%	\eas
%	Since by a standard argument on multiplicative characters
%	\[
%	\sum_{\disceltwo_{\gengenus-r}\in D_\lattice^{\gengenus-r}}
%	e\Big(
%	\tr (\disceltwo_{\gengenus-r},\discel_{\gengenus-r}-\discelthree_{\gengenus-r})
%	\Big)
%	=
%	\begin{cases}
%	|D_\lattice|^{\gengenus - r}, & \text{if~$\discelthree_{\gengenus-r}=\discel_{\gengenus-r}$,}
%	\\
%	0, & \text{otherwise,}
%	\end{cases}
%	\]
%	we deduce that
%	\ba\label{eq;stopaftersomecomp}
%	\weil{\lattice}{\gengenus}\big(\iota(M,I_{\gengenus-r})\big)\mathfrak{e}_\discel
%	=
%	\frac{1}{|D_\lattice|^r}
%	\sum_{\discelthree\in D_\lattice^r}\Big(
%	\sum_{\disceltwo \in D_\lattice^r}
%	e\big(
%	-\tr (q(\disceltwo)) + \tr (\disceltwo,\discel_r - \discelthree)
%	\big)
%	\Big)
%	\mathfrak{e}_\discelthree\otimes\mathfrak{e}_{\discel_{\gengenus-r}}.
%	\ea
%	Recall from~\cite[Corollary~$4.2$]{borcherds-grassmannians} that~$\sum_{\disceltwo\in D_\lattice} e\big(q(\disceltwo)\big) = i^{(n-2)/2}|D_\lattice|^{1/2}$, which may be rewritten over~$D_\lattice^r$ as
%	\[
%	\sum_{\disceltwo\in D_\lattice^r}e\big(\tr(q(\disceltwo))\big) = i^{r(n-2)/2}|D_\lattice|^{r/2}.
%	\]
%	This implies that
%	\bas
%	e\big(-\tr\, q(\discel_r - \discelthree)\big)i^{r(n-2)/2} |D_\lattice|^{r/2}
%	&=
%	e\big(-\tr\, q(\discel_r - \discelthree)\big)\sum_{\disceltwo\in D_\lattice^r}e\Big(\tr\, q\big(\disceltwo - (\discel_r - \discelthree)\big)\Big)
%	\\
%	&=
%	\sum_{\disceltwo\in D_\lattice^r}
%	e\Big(
%	\tr\, q(\disceltwo) - \tr(\disceltwo,\discel_r - \discelthree)
%	\Big),
%	\eas
%	which we conjugate to deduce that
%	\[
%	\sum_{\disceltwo\in D_\lattice^r}
%	e\Big(
%	-\tr\, q(\disceltwo) + \tr (\disceltwo,\discel_r-\discelthree)
%	\Big)
%	=
%	i^{r(2-n)/2} |D_\lattice|^{r/2}
%	e\big(\tr\,q(\discel_r-\discelthree)\big).
%	\]
%	We then deduce that~\eqref{eq;stopaftersomecomp} may be further simplified as
%	\ba\label{eq;finaleqactofiotaMI}
%	\weil{\lattice}{\gengenus}\big(\iota(M,I_{\gengenus-r})\big)\mathfrak{e}_\discel
%	&=
%	\frac{i^{r(2-n)/2}}{|D_\lattice|^{r/2}}\sum_{\discelthree\in D_\lattice^r}e\Big(
%	\tr\, q(\discel_r - \discelthree)
%	\Big)\mathfrak{e}_\discelthree\otimes\mathfrak{e}_{\discel_{\gengenus-r}}.
%	\ea}
%	
%	\textcolor{blue}{We are now ready to compute the action of~$\iota(S_r,I_{\gengenus-r})$ on~$\CC[D_\lattice^\gengenus]$ under the Weil representation~$\weil{\lattice}{\gengenus}$.
%	By~\eqref{eq;lowtriang} and~\eqref{eq;finaleqactofiotaMI} we deduce that
%	\begin{align*}
%	\weil{\lattice}{\gengenus}&\big(\iota(S_r,I_{\gengenus-r})\big)\mathfrak{e}_\discel
%	\\
%	&=
%	\weil{\lattice}{\gengenus}\big(T_\gengenus(-I_r)\big)
%	\weil{\lattice}{\gengenus}\big(\iota(M,I_{\gengenus-r})\big)
%	\weil{\lattice}{\gengenus}\big(T_\gengenus(-I_r)\big)
%	\mathfrak{e}_\discel
%	\\
%	&=
%	e\big(-\tr\, q(\discel_r)\big)
%	\weil{\lattice}{\gengenus}\big(T_\gengenus(-I_r)\big)
%	\weil{\lattice}{\gengenus}\big(\iota(M,I_{\gengenus-r})\big)
%	\mathfrak{e}_\discel
%	\\
%	&=
%	e\big(-\tr\, q(\discel_r)\big)
%	\frac{i^{r(2-n)/2}}{|D_\lattice|^{r/2}}
%	\sum_{\discelthree\in D_\lattice^r}
%	e\Big(
%	\tr\, q(\discel_r-\discelthree)
%	\Big)
%	\weil{\lattice}{\gengenus}\big(T_\gengenus(-I_r)\big)
%	\mathfrak{e}_\discelthree\otimes\mathfrak{e}_{\discel_{\gengenus-r}}
%	\\
%	&=
%	e\big(-\tr\, q(\discel_r)\big)
%	\frac{i^{r(2-n)/2}}{|D_\lattice|^{r/2}}
%	\sum_{\discelthree\in D_\lattice^r}
%	e\Big(
%	\tr\, q(\discel_r-\discelthree)
%	-\tr\, q(\discelthree)
%	\Big)
%	\mathfrak{e}_\discelthree\otimes\mathfrak{e}_{\discel_{\gengenus-r}}
%	\\
%	&=
%	\frac{i^{r(2-n)/2}}{|D_\lattice|^{r/2}}
%	\sum_{\discelthree\in D_\lattice^r}
%	e\big(
%	-\tr(\discelthree,\discel_r)
%	\big)
%	\mathfrak{e}_\discelthree\otimes\mathfrak{e}_{\discel_{\gengenus-r}}
%	\\
%	&=
%	\Big(\weil{\lattice}{r}(S_r)\mathfrak{e}_{\discel_r}\Big)
%	\otimes\mathfrak{e}_{\discel_{\gengenus-r}}.
%	\qedhere
%	\end{align*}}
%	\end{proof}
	
	\subsection{Vector-valued modular forms}\label{sec;backgrvvmodforms}
	Let~$k\in\frac{1}{2}\ZZ$ be such that~$2k\equiv n-2$ mod~$4$.
	A Siegel modular form of weight~$k$ and genus~$\gengenus$ with respect to the Weil representation~$\weil{\lattice}{\gengenus}$ is a holomorphic function~$f\colon\HH_\gengenus\to\CC[D_\lattice^\gengenus]$ such that
	\[
	f(\gamma\cdot\tau)=\phi(\tau)^{2k} \weil{\lattice}{\gengenus}(\gamma) f(\tau)
	\qquad
	\text{for every~$\gamma=\big(M , \phi(\tau)\big)\in\Mp_{2\gengenus}(\ZZ)$},
	\]
	with the usual requirement that~$f$ is holomorphic at the cusp~$\infty$ for the case of genus~$1$.
	We write~$f=\sum_{\discel\in D_\lattice^\gengenus}f_\discel\mathfrak{e}_\discel$, denoting by~$f_\discel$ the~$\discel$-component of~$f$.
	Such components are scalar-valued Siegel modular forms of level~$N$, where~$N$ is the level of~$\lattice$, i.e.\ the smallest positive integer such that~$N q(\lambda)\in\ZZ$ for every~$\lambda\in \lattice'$.
	
	Let~$\halfint{\gengenus}$
	%, (resp~$\halfint{\gengenus}^+$),
	be the set of symmetric half-integral~$\gengenus\times\gengenus$-matrices.
	%positive semi-definite (resp.\ positive definite)
	The invariance of~$f$ with respect to the translations on~$\HH_\gengenus$ induced by~$\Mp_{2\gengenus}(\ZZ)$ implies that~$f$ admits a Fourier expansion of the form
	\be\label{eq;Fexpfirstappvvc}
	f(\tau)=\sum_{\discel\in D_\lattice^\gengenus}
	\sum_{\substack{T\in q(\discel) + \halfint{\gengenus} \\ T\ge 0}}
	c_T(f_\discel) q^T \mathfrak{e}_\discel,
	\ee
	where~$q^T=e \big(\tr (T\tau) \big)$.
	If~$\genus=1$, then the fact that the series~\eqref{eq;Fexpfirstappvvc} is indexed over non-negative rational numbers is implied by the request that elliptic modular forms are holomorphic at~$\infty$.
	For larger genus, one can show that the Fourier expansion is indexed only over positive semidefinite matrices by the Koecher principle for Siegel modular forms with level.
%	The Fourier coefficients may be computed as
%	\[
%	c_T(f_\discel)=\int_{\text{$x$ mod $1$}}f_\discel(\tau) e\big(-\tr(T\tau)\big)\,dx
%	=
%	\int_{\text{$x$ mod $1$}}\big\langle
%	f(\tau),e\big(\tr(T\overline{\tau})\big)\mathfrak{e}_\discel
%	\big\rangle\,dx.
%	\]
	
	We denote by
	\[
	c_T(f)=\sum_{\discel\in D_\lattice^\gengenus} c_T(f_\discel) \mathfrak{e}_\discel
	\]
	the vector in~$\CC[D_\lattice^\gengenus]$ gathering all Fourier coefficients of index~$T$ of~$f$.
	The Fourier coefficients of Siegel modular forms with respect to~$\weil{\lattice}{\gengenus}$ satisfy symmetric relations with respect to the action~$T\mapsto A^t T A$ induced by matrices~$A\in\GL_\gengenus(\ZZ)$.
	More precisely, we have
	\be\label{eq;symmofcoef}
	c_T(f) = \phi^{2k} \weil{\lattice}{\gengenus}\big(\big(\begin{smallmatrix}
	A & \\
	& A^{-t}
	\end{smallmatrix}\big),\phi\big) c_{A^t T A}(f)\qquad \text{for all~$A\in\GL_{\gengenus}(\ZZ),$}
	\ee
	where~$\phi$ is any square-root of~$\det A$.
%	\\
%	\textcolor{blue}{
%	---------------------\\
%	\textbf{Here a check of the previous formula.}\\
%	If~$T\not\in q(\discel) + \halfint{\gengenus}$ or~$T<0$, then~$c_T(f_\discel)=0$.
%	We may then rewrite the Fourier expansion of~$f$ as
%	\[
%	f(\tau)=\sum_{T\in\halfint{\gengenus}}
%	\Big(\sum_{\discel\in D_\lattice^\gengenus}
%	c_T(f_\discel)\mathfrak{e}_\discel
%	\Big) q^T
%	=
%	\sum_{T\in\halfint{\gengenus}}
%	c_T(f)
%	q^T.
%	\]
%	If~$\gamma=\big(\big(\begin{smallmatrix}
%	A & \\
%	& A^{-t}
%	\end{smallmatrix}\big),\phi\big)$, then we may rewrite~$f(\gamma\cdot\tau)$ in two different ways, namely as
%	\bas
%	f(\gamma\cdot\tau)
%	&=
%	f(A\tau A^{t})
%	=
%	\sum_{T\in\halfint{\gengenus}}
%	c_T(f)
%	q^{A^{t}TA}
%	=
%	\sum_{T\in\halfint{\gengenus}}
%	c_{A^{-t}TA}(f) q^T,
%	\\
%	f(\gamma\cdot\tau)
%	&=
%	\sum_{T\in\halfint{\gengenus}}
%	\phi^{2k} \weil{\lattice}{\gengenus}(\gamma) f(\tau)
%	=
%	\sum_{T\in\halfint{\gengenus}}
%	\phi^{2k}\weil{\lattice}{\gengenus}(\gamma) c_T(f) q^T.
%	\eas
%	Here we used that the action~$\halfint{\gengenus}\to\halfint{\gengenus}$, $T\mapsto A^tTA$ is bijective for every~$A\in\GL_{\gengenus}(\ZZ)$.
%	If we read the previous two rewriting of~$f(\gamma\cdot\tau)$ replacing~$T$ by~$A^{t}TA$, we can read the relation we want to prove.
%	In fact, two Fourier expansions of the same function must have Fourier coefficients that are pairwise equal.
%	\\---------------------
%	}
	It is easy to see that~\eqref{eq;symmofcoef} implies that
	\be\label{eq;symmofcoef2}
	c_T(f_\discel)=c_{A^t T A}(f_{\discel A})
	\qquad
	\text{for all~$\discel\in D_\lattice^\genus$}.
	\ee
%	\textcolor{blue}{
%	---------------------\\
%	\textbf{Here a check of the previous formula.}\\
%	By choosing~$\phi$ to be the positive square root and replacing~$f$ with~$\sum_{\discel\in D_\lattice^\genus}f_\discel\mathfrak{e}_\discel$, we deduce from~\eqref{eq;symmofcoef} that
%	\bas
%	\sum_{\discel\in D_\lattice^\genus} c_T(f_\discel) \mathfrak{e}_\discel
%	&=
%	\sum_{\discel\in D_\lattice^\genus}
%	c_{A^t T A}(f_\discel)\cdot \det A^k\cdot
%	\weil{\lattice}{\genus}\big(\big(\begin{smallmatrix}
%	A & \\
%	 & A^{-t}
%	 \end{smallmatrix}\big), \det A^{1/2}\big)
%	 \mathfrak{e}_\discel
%	 \\
%	&=
%	\sum_{\discel\in D_\lattice^\genus}
%	c_{A^t T A}(f_\discel)
%	\mathfrak{e}_{\discel A^{-1}}
%	=
%	\sum_{\discel\in D_\lattice^\genus}
%	c_{A^t T A}(f_{\discel A})
%	\mathfrak{e}_\discel,
%	\eas
%	from which we see~\eqref{eq;symmofcoef2}.
%	\\---------------------
%	}
	
	If all Fourier coefficients~$c_T(f)$ of~$f$ with respect to singular matrices~$T$ vanish, then~$f$ is said to be a \emph{Siegel cusp form}.	
	We denote by~$M^k_{\gengenus,\lattice}$ and~$S^k_{\gengenus,\lattice}$ respectively the space of Siegel modular forms and cusp forms of weight~$k$ and genus~$\gengenus$ with respect to~$\weil{\lattice}{\gengenus}$, with the convention that~$S^k_{0,\lattice}=M^
	k_{0,\lattice}=\CC$ as in~\cite{klingen}.	
%	\textcolor{blue}{\textbf{Note:} We agreed to put the text below on Siegel Eisenstein series in blue, and add a paragraph in the section of Klingen Eisenstein series.
%	There, we define Siegel Eisenstein series as special cases of Klingen Eisenstein series.}
%	
%	\textcolor{blue}{The first examples of modular forms are the Siegel Eisenstein series, which we now recall.
%	We denote by~$\sieg{\gengenus}$, resp.~$\metsieg{\gengenus}$, the Siegel parabolic subgroup of~$\Sp_{2\gengenus}(\ZZ)$, resp.~the preimage of~$\sieg{\gengenus}$ in~$\Mp_{2\gengenus}(\ZZ)$.
%	Explicitly, the Siegel parabolic is
%	\[
%	\sieg{\gengenus}
%	=
%	\left\{
%	\left(\begin{smallmatrix}
%	A & B A^{-t}\\
%	0 & A^{-t}
%	\end{smallmatrix}\right)
%	:
%	A\in\GL_\gengenus(\ZZ)\quad\text{and}\quad B\in\Sym_\gengenus(\ZZ)
%	\right\},
%	\]
%	and may be characterized by the subgroup of~$\Sp_{2\gengenus}(\ZZ)$ consisting of matrices where the lower-left~$\gengenus\times\gengenus$-block vanishes.}
%%	
%%	Note that if~$\mathfrak{e}_0$ is the basis element of~$\CC[D_\lattice^\gengenus]$ arising from~$0\in D_\lattice^\gengenus$, then
%%	\[
%%	\det(C\tau+D)^k\weil{\lattice}{\gengenus}(\gamma)\mathfrak{e}_0=\mathfrak{e}_0\qquad\text{for every~$\gamma=\big(\begin{smallmatrix}
%%	A & B\\
%%	C & D
%%	\end{smallmatrix}\big)\in\metsieg{\gengenus}$}.
%%	\]
%	\textcolor{blue}{
%	\begin{defi}\label{def:SeigEisvv}
%	Let~$k\in\frac{1}{2}\ZZ$ be such that~$2k\equiv n-2$ mod~$4$ and~$k>g+1$.
%	The genus~$\gengenus$ and weight~$k$ \emph{Siegel Eisenstein series}~$\eiszero{k}{\gengenus}{\lattice}$ with respect to~$\weil{\lattice}{\gengenus}$ (and arising from the trivial isotropic vector~$0\in D_\lattice^\gengenus$) is
%	\be\label{eq;defvvEis}
%	\eiszero{k}{\gengenus}{\lattice}(\tau)
%	=
%	\sum_{\gamma\in\metsieg{\gengenus}\backslash\Mp_{2\gengenus}(\ZZ)}
%	%\det(C\tau + D)^{-k}
%	\phi(\tau)^{-2k}
%	\weil{\lattice}{\gengenus}(\gamma)^{-1}\mathfrak{e}_0,
%	\ee
%	where we write~$\gamma=(M,\phi)\in\Mp_{2\gengenus}(\ZZ)$.
%%	$\gamma=\big(\left(\begin{smallmatrix}
%%	A & B\\
%%	C & D
%%	\end{smallmatrix}\right), \det (C\tau+D)^{1/2}\big)$.
%	\end{defi}
%	The convergence of the series will follow easily from the more general case of Klingen Eisenstein series treated in Section~\ref{sec;KEseries}.
%	The Eisenstein series~$\eiszero{k}{\gengenus}{\lattice}$ is normalized such that its constant Fourier coefficient is~$c_0(\eiszero{k}{\gengenus}{\lattice},\discel)=\delta_{0,\discel}$, for every~$\discel\in D_L^\gengenus$.}
%	
%	\textcolor{blue}{\textbf{Check that the Eisenstein series~$\eiszero{k}{\gengenus}{\lattice}$ is well-defined:}\\
%	Let~${\gamma_1=(M_1,\pm\phi_1)\in\metsieg{\gengenus}}$ for some~$M_1=\big(\begin{smallmatrix}
%	A & BA^{-t}\\
%	0 & A^{-t}
%	\end{smallmatrix}\big)$, where~$\phi_1(\tau)=\det A^{1/2}$.
%	Let~~$F(\gamma)\coloneqq \phi(\tau)^{-2k}\weil{\lattice}{\gengenus}(\gamma)^{-1}\mathfrak{e}_0$ for every~$\gamma=(M,\phi)\in\Mp_{2\gengenus}(\ZZ)$.
%	Since~$\gamma_1\cdot\gamma=\big(M_1M , \pm\phi_1(M\cdot\tau)\phi(\tau)\big)$, we deduce that
%	\bas
%	F(\gamma_1\cdot\gamma)
%	&=
%	\big(\pm\phi_1(M\cdot\tau)\big)^{-2k}\phi(\tau)^{-2k}\weil{\lattice}{\gengenus}(\gamma)^{-1}\weil{\lattice}{\gengenus}(\gamma_1)^{-1}\mathfrak{e}_0
%	\\
%	&=
%	(\pm 1)^{-2k}\det A ^{-k}\phi(\tau)^{-2k} (\pm 1)^{n-2}\det A^{(n-2)/2}\weil{\lattice}{\gengenus}(\gamma)^{-1}\mathfrak{e}_0
%	\\
%	&=
%	\phi(\tau)^{-2k}\weil{\lattice}{\gengenus}(\gamma)^{-1}\mathfrak{e}_0
%	=
%	F(\gamma),
%	\eas
%	where in the first to the last equality we used that~$2k\equiv n-2$ mod~$4$.}
%	
%	\textcolor{blue}{
%	\begin{lemma}
%	The constant Fourier coefficient of~$\eiszero{k}{\gengenus}{\lattice}$ is
%	\[
%	c_0(\eiszero{k}{\gengenus}{\lattice},\discel)=\begin{cases}
%	1, & \text{if $\discel=0\in D_\lattice^\gengenus$,}\\
%	0, & \text{otherwise}.
%	\end{cases}
%	\]
%	\end{lemma}
%	\begin{proof}
%	Let~$\lambda\in\RR_{>0}$.
%	We rewrite the Fourier expansion of~$\eiszero{k}{\gengenus}{\lattice}$ at~$\tau=i\lambda I_\gengenus\in\HH_\gengenus$ as
%	\bas
%	\eiszero{k}{\gengenus}{\lattice}(i\lambda I_\gengenus)
%	=
%	\sum_{\substack{\discel\in D_\lattice^\gengenus \\ q(\discel)=0}} c_0(\eiszero{k}{\gengenus}{\lattice},\discel)\mathfrak{e}_\discel
%	+
%	\sum_{\discel\in D_\lattice^\gengenus}
%	\sum_{\substack{T\in q(\discel) + \halfint{\gengenus} \\ T\ge 0,\, T\neq 0}}
%	c_T(\eiszero{k}{\gengenus}{\lattice},\discel) \exp\big(
%	-2\pi \lambda \trace(T)
%	\big) \mathfrak{e}_\discel.
%	\eas
%	Since every non-zero positive-semidefinite matrix~$T$ has positive trace, we deduce that
%	\ba\label{eq;limFexpinspecvalEis}
%	\lim_{\lambda\to +\infty}\eiszero{k}{\gengenus}{\lattice}(i\lambda I_\gengenus)
%%	=
%%	\sum_{\discel\in D_\lattice^\gengenus} c_0(\eiszero{k}{\gengenus}{\lattice},\discel)\mathfrak{e}_\discel
%	=
%	\sum_{\substack{\discel\in D_\lattice^\gengenus \\ q(\discel)=0}} c_0(\eiszero{k}{\gengenus}{\lattice},\discel)\mathfrak{e}_\discel.
%	\ea
%	We now proceed with an analogous result replacing the Fourier expansion of~$\eiszero{k}{\gengenus}{\lattice}$ with its defining series.
%	Since every element of~$\metsieg{\gengenus}\backslash\Mp_{2\gengenus}(\ZZ)$ can be represented by a matrix of the form~$\big(\begin{smallmatrix}
%	\ast & \ast
%	\\
%	C & D
%	\end{smallmatrix}\big)$ for unique~$C$ and~$D$ up to multiplying both by~$-I_\gengenus$, we rewrite
%	\ba\label{eq;inproofreplfdofeisvv}
%	\eiszero{k}{\gengenus}{\lattice}(i\lambda I_\gengenus)
%	&=
%	\weil{\lattice}{\gengenus}(I_{2\gengenus})^{-1}\mathfrak{e}_0
%	+
%	\sum_{\substack{\gamma\in\metsieg{\gengenus}\backslash\Mp_{2\gengenus}(\ZZ) \\ \gamma\neq I_{2\gengenus}}}
%	\lambda^{-k}\det\Big(Ci + \frac{D}{\lambda}\Big)^{-k} \weil{\lattice}{\gengenus}(\gamma)^{-1}\mathfrak{e}_0.
%	\ea
%	If~$\gamma\neq I_{2\gengenus}$ in~$\metsieg{\gengenus}\backslash\Mp_{2\gengenus}(\ZZ)$, then the~$\gengenus\times\gengenus$-blocks decomposition of~$\gamma$ is such that~$C\neq 0$.
%	This implies that if~$\lambda\to+\infty$, then the second summand of the right-hand side of~\eqref{eq;inproofreplfdofeisvv} tends to~$0$.
%	We then deduce that
%	\bas
%	\lim_{\lambda\to+\infty}\eiszero{k}{\gengenus}{\lattice}(i\lambda I_\gengenus)
%	=
%	\mathfrak{e}_0.
%	\eas
%	The computation of the constant term of the Fourier expansion of~$\eiszero{k}{\gengenus}{\lattice}$ may be then deduce by comparing the previous formula with~\eqref{eq;limFexpinspecvalEis}.
%	\end{proof}
%	The space of Siegel Eisenstein series with respect to~$\weil{\lattice}{\gengenus}$ can be of dimension greater than~$1$; see e.g.~\cite[Section~$1.2.3$]{bruinierbook} for the case of genus~$1$.
%	For the purposes of this article, we do not need to introduce other Eisenstein series than~\eqref{eq;defvvEis}, i.e.\ the one arising from the zero of~$D_\lattice^\gengenus$.}
%	
%	\subsection{Some bounds for cusp forms}\label{sec;someboundscusp}
%	\textcolor{blue}{In this section we generalize the well-known bounds for cusp forms given e.g.\ in~\cite[Section~$5$]{klingen} to vector-valued cusp forms with respect to the Weil representation.
%	The only purpose of this section is to provide the necessary background to prove Lemma~\ref{lemma;bound2blafromKling}.}
%	\\
%	
%	\textcolor{blue}{The imaginary parts of points~$\tau=x+iy\in \HH_\gengenus$ form a convex cone
%	\be\label{eq;convconeimpart}
%	\cone{\gengenus}\coloneqq\{y : y>0\}
%	\ee
%	in~$\RR^{\gengenus(\gengenus+1)/2}$ with vertex at the origin, consisting of the~$\gengenus\times\gengenus$ positive definite symmetric matrices with coefficients in~$\RR$.
%	The group~$\GL_\gengenus(\ZZ)$ acts on~$\cone{\gengenus}$ as
%	\[
%	\GL_\gengenus(\ZZ)\times\cone{\gengenus}\longrightarrow\cone{\gengenus},
%	\qquad
%	(A,y)\longmapsto A y A^t.
%	\]
%	This is induced by the action of the unimodular matrices~$\big(\begin{smallmatrix}
%	A & 0 \\
%	0 & A^{-t}
%	\end{smallmatrix}\big)\in\Sp_{2\gengenus}(\ZZ)$ on~$\HH_\gengenus$.}
%	
%	\textcolor{blue}{A fundamental set in~$\cone{\gengenus}$ with respect to the action of~$\GL_\gengenus(\ZZ)$ is given by the Minkowski reduced domain, which we now recall; see~\cite[Section~$2$]{klingen} for details.
%	\begin{defi}\label{def;minkowskireddom}
%	The \emph{Minkowski reduced domain} (of genus~$\genus$) is the set
%	\[
%	\mink{\gengenus}=\{ y \in \cone{\gengenus} : \text{$y$ satisfies~(i) and~(ii)} \},
%	\]
%	where
%	\begin{enumerate}[label=(\roman*)]
%	\item $\intvect y\intvect^t\ge y_{k,k}$ for all integral vectors~$\intvect\in\ZZ^{\gengenus}$ such that~$\gcd(\intvect_j,\dots,v_\genus)=1$ for all~$1\le j \le \gengenus$,
%	\item $y_{j,j+1}\ge 0$ for all~$1\le j \le \gengenus-1$.
%	\end{enumerate}
%	The points of~$\cone{\gengenus}$ that lie in~$\mink{\gengenus}$ are said to be \emph{reduced in the sense of Minkowski}.
%	\end{defi}
%	Recall that the norm on~$\CC[D_\lattice^\gengenus]$ is the one induced by the standard inner product~\eqref{eq;stinprodgroupalg}.
%	The following result can be easily proved following the wording of~\cite[p.~$56$, Proposition~$3$]{klingen}.
%	\begin{prop}\label{prop;bound1blafromKling}
%	Let~$f\in S^k_{\gengenus,L}$, and let~$c>0$.
%	There exist positive numbers~$c_1$ and~$c_2$ such that
%	\[
%	\lvert f(\tau) \rvert \le c_1 e^{-c_2\tr y}
%	\]
%	for all~$\tau\in\HH_\gengenus$ such that~$y$ is reduced in the sense of Minkowski and~$y\ge c I_\gengenus$.
%	\end{prop}
%	\begin{proof}
%	By replacing~$f$ with its Fourier expansion, we may rewrite~$\lvert f(\tau)\rvert$ as
%	\bas
%	\lvert f(\tau)\rvert
%	\le
%	\sum_{\discel\in D_\lattice^\gengenus}
%	\sum_{\substack{T\in q(\discel) + \halfint{\gengenus} \\ T> 0}}
%	\lvert c_T(f,\discel)\rvert
%	e^{-2\pi\tr (Ty)}.
%	\eas
%	Since~$y\ge c I_\gengenus$, then~$\tr (Ty)\ge c\tr T$.
%	We may deduce another bound of~$\tr (Ty)$ as follows.
%	By~\cite[p.~$20$, Lemma~$2$]{klingen} we may assume that the reduced matrix~$y$ is diagonal.
%	Let~$N$ be the smallest positive integer such that~$Nq(\discel)\in\ZZ$ for every~$\discel\in D_\lattice$.
%	The diagonal elements of~$T$ are positive rational numbers such that, when rewritten in minimal terms, have a denominator bounded by~$N$.
%	This implies that~$\tr (Ty)\ge c'\tr y$ for some positive constant~$c'$ depending only on~$\gengenus$ and~$L$.
%	Combining the previous estimates, we deduce that
%	\[
%	\lvert f(\tau)\rvert
%	\le
%	e^{-\pi c'\tr y}
%	\sum_{\discel\in D_\lattice^\gengenus}
%	\sum_{\substack{T\in q(\discel) + \halfint{\gengenus} \\ T> 0}}
%	\lvert c_T(f,\discel)\rvert
%	e^{-\pi c \tr T}
%	=
%	c_1 e^{-c_2\tr y},
%	\]
%	where~$c_1=\sum_\discel\sum_T\lvert c_T(f,\discel)\rvert e^{-\pi c \tr T}$ and~$c_2=\pi c'$.
%	We remark that the series defining~$c_1$ converges since the Fourier expansion of~$f$ converges absolutely.
%	\end{proof}}
	
	\subsection{Relations with sublattices}\label{sec;relwithsublat}
	Let~$M\subset \lattice$ be a sublattice of finite index.
	There are two natural maps that relate the spaces~$M^k_{\genus,\lattice}$ and~$M^k_{\genus,M}$, as illustrated in the following result.
	This may be considered as a generalization to higher genus of~\cite[Section~$4$]{scheit} and~\cite[Lemma~$3.1$]{bruinier-yang09}.
	The natural surjective map~$\lattice'/M\to\lattice'/\lattice$, $\mu\mapsto\overline{\mu}$, induces a map~$(\lattice'/M)^\genus\to(\lattice'/\lattice)^\genus$ for every~$\genus$.
	\begin{lemma}\label{lemma;genofBYandS}
	There are two natural maps
	\[
	\res_{\genus,\lattice/M}\colon M^k_{\genus,\lattice}\longrightarrow M^k_{\genus,M},
	\qquad
	f\longmapsto f_M
	\]
	and
	\[
	\tr_{\genus,\lattice/M}\colon M^k_{\genus,M}\longrightarrow M^k_{\genus,\lattice},
	\qquad
	h\longmapsto h^\lattice
	\]
	such that~$\langle f, h^\lattice\rangle=\langle f_M,h\rangle$ for every~$f\in M^k_{\genus,\lattice}$ and~$h\in M^k_{\genus,M}$.
	They are given as
	\[
	\big(f_M\big)_\mu
	=
	\begin{cases}
	f_{\overline{\mu}}, & \text{if~$\mu\in (\lattice'/M)^\genus$,}
	\\
	0, & \text{otherwise,}
	\end{cases}
	\qquad
	\text{for every~$f\in M^k_{\genus,\lattice}$ and~$\mu\in D_M^\genus$}
	\]
	and
	\[
	\big(h^\lattice\big)_{\overline{\mu}}
	=
	\sum_{\discel\in (\lattice/M)^\genus}
	h_{\discel + \mu}
	\qquad
	\text{for every~$h\in M^k_{\genus,M}$,}
	\]
	where~$\mu$ is any fixed preimage of~$\overline{\mu}$ in~$(\lattice'/M)^\genus$.
	\end{lemma}
	\begin{proof}
	It is analogous to~\cite[Theorem~$4.1$]{scheit} and~\cite[Lemma~$3.1$]{bruinier-yang09}.
%	
%	\textcolor{blue}{\textbf{Check that if~$\boldsymbol{f\in M^k_{\genus,\lattice}}$, then~$\boldsymbol{f_M\in M^k_{\genus,M}}$:}
%	\\
%	It is enough to prove that~$f_M$ transforms correctly under the generators of~$\Mp_{2\genus}(\ZZ)$.
%	The check for~$T_B$ with~$B\in\Sym_\genus(\ZZ)$ is easy and is omitted.
%	We check the transformation with respect to~$S$ following the wording of~\cite[Proposition~$6.9$]{scheit-habil}.
%	\\
%	Let~$H\coloneqq (\lattice/M)^\genus$.
%	Note that~$H\subset D_M^\genus$ and~$H\subset H^\perp$ in~$D_M^\genus$.
%	Moreover
%	\be\label{eq;inbluechecknhab0}
%	H^\perp = (\lattice'/M)^\genus,
%	\qquad
%	H^\perp/H\cong D_\lattice^\genus
%	\qquad
%	\text{and}
%	\qquad
%	\lvert D_M\rvert^\genus
%	=
%	\lvert H\rvert^2
%	\cdot
%	\lvert D_L\rvert^\genus.
%	\ee
%	\underline{\emph{Case of~$\mu\in H^\perp=(L'/M)^\genus$:}}
%	\\
%	Since~$f\in M^k_{\genus,\lattice}$, we know that
%	\[
%	f_{\overline{\mu}}(S\cdot\tau)
%	=
%	\det\tau^k
%	\frac{\sqrt{i}^{\genus(n-2)}}{\lvert D_\lattice\rvert^{\genus/2}}
%	\sum_{\overline{\beta}\in D_L^\genus}
%	f_{\overline{\beta}}(\tau)
%	\cdot
%	e\big(-\tr(\overline{\mu},\overline{\beta})\big).
%	\]
%	Note that
%	\ba\label{eq;inbluechecknhab}
%	\sum_{\beta\in D_M^\genus}
%	e\big(-\tr(\mu,\beta)\big)
%	(f_M)_\beta
%	&=
%	\sum_{\beta\in(L'/M)^\genus}
%	e\big(-\tr(\mu,\beta)\big)
%	(f_M)_\beta
%	\\
%	&=
%	\sum_{\beta\in H^\perp/H}
%	\sum_{\alpha\in H}
%	e\big(-\tr(\mu,\beta + \alpha)\big)
%	(f_M)_{\beta + \alpha}
%	\\
%	&=
%	\sum_{\beta\in H^\perp/H}
%	\sum_{\alpha\in H}
%	e\big(-\tr(\mu,\beta)\big)
%	f_\beta
%	\\
%	&=
%	\lvert H\rvert
%	\sum_{\beta\in D_\lattice^\genus}
%	e\big(-\tr(\mu,\beta)\big)
%	f_\beta.
%	\ea
%	Since~$\mu\in H^\perp$, we have
%	\ba\label{eq;inblnhabthird}
%	(f_M)_\mu (S\cdot\tau)
%	&=
%	f_{\overline{\mu}}(S\cdot\tau)
%	=
%	\det \tau^k
%	\frac{\sqrt{i}^{\genus(n-2)}}{\lvert D_L\rvert^{\genus/2}}
%	\sum_{\overline{\beta}\in D_\lattice^\genus}
%	f_{\overline{\beta}}(\tau)
%	\cdot
%	e\big(-\tr(\overline{\mu},\overline{\beta})\big)
%	\\
%	&=
%	\det \tau^k
%	\frac{\sqrt{i}^{\genus(n-2)}}{\underbrace{\lvert D_L\rvert^{\genus/2}\cdot \lvert H\rvert}_{=\lvert D_M\rvert^{\genus/2}}}
%	\sum_{\beta\in D_M^\genus}
%	e\big(-\tr(\mu,\beta)\big)
%	(f_M)_\beta(\tau),
%	\ea
%	where in the last passage we used~\eqref{eq;inbluechecknhab} and~\eqref{eq;inbluechecknhab0}.
%	\\
%	\underline{\emph{Case of~$\mu\not\in H^\perp=(L'/M)^\genus$:}}
%	\\
%	We have~$(f_M)_\mu=0$, hence~$(f_M)_\mu(S\cdot\tau)=0$ for all~$\tau\in\HH_\genus$.
%	Moreover, the sum appearing on the last term in~\eqref{eq;inblnhabthird} may be computed as
%	\bas
%	\sum_{\beta\in D_M^\genus}
%	e\big(-\tr(\mu,\beta)\big)
%	(f_M)_\beta (\tau)
%	&=
%	\sum_{\beta\in H^\perp}
%	e\big(-\tr(\mu,\beta)\big)
%	(f_M)_\beta (\tau)
%	\\
%	&=
%	\sum_{\overline{\beta}\in H^\perp/H}
%	\sum_{\alpha\in H}
%	e\big(-\tr(\mu,\overline{\beta}+\alpha)\big)
%	(f_M)_{\overline{\beta}+\alpha}(\tau)
%	\\
%	&=
%	\sum_{\overline{\beta}\in H^\perp/H}
%	e\big(-\tr(\mu,\overline{\beta})\big)
%	f_{\overline{\beta}}(\tau)
%	\underbrace{\sum_{\alpha\in H}
%	e\big(-\tr(\mu,\alpha)\big)}_{=0}
%	\\
%	&=
%	0,
%	\eas
%	where in the last step we deduce the vanishing from the vanishing of a sum over~$H$.
%	The latter is due to the fact that the sum over~$H$ is a sum over a non-trivial character of~$H$.
%	}
%	\\
%	
%	\textcolor{blue}{\textbf{Check that if~$\boldsymbol{h\in M^k_{\genus,M}}$, then~$\boldsymbol{h^L\in M^k_{\genus,\lattice}}$:}
%	\\
%	It is enough to prove that~$h^\lattice$ transforms correctly under the generators of~$\Mp_{2\genus}(\ZZ)$.
%	The check for~$T_B$ with~$B\in\Sym_\genus(\ZZ)$ is easy and is omitted.
%	We check the transformation with respect to~$S$ following the same wording (and notation!) of the previous check.
%	\\
%	Since~$h\in M^k_{\genus,M}$, we know that
%	\[
%	h_{\beta}(S\cdot\tau)
%	=
%	\det\tau^k
%	\frac{\sqrt{i}^{\genus(n-2)}}{\lvert D_M\rvert^{\genus/2}}
%	\sum_{\alpha\in D_M^\genus}
%	h_{\alpha}(\tau)
%	\cdot
%	e\big(-\tr(\beta,\alpha)\big)
%	\qquad
%	\text{for all~$\beta\in D_M^\genus$.}
%	\]
%	Let~$\overline{\mu}\in D_\lattice^\genus$, and let~$\mu$ be any fixed preimage of~$\overline{\mu}$ in~$(\lattice'/M)^\genus$.
%	Then
%	\bas
%	(h^\lattice)_{\overline{\mu}}(S\cdot\tau)
%	&=
%	\sum_{\alpha\in (L/M)^\genus} h_{\alpha + \mu}(S\cdot\tau)
%	\\
%	&=
%	\sum_{\alpha\in (L/M)^\genus}
%	\det\tau^k
%	\frac{\sqrt{i}^{\genus(n-2)}}{\lvert D_M\rvert^{\genus/2}}
%	\sum_{\beta\in D_M^\genus}
%	h_\beta(\tau)
%	\cdot
%	e\big(-\tr(\alpha + \mu , \beta)\big)
%	\\
%	&=
%	\det\tau^k
%	\frac{\sqrt{i}^{\genus(n-2)}}{\lvert D_M\rvert^{\genus/2}}
%	\bigg[
%	\sum_{\alpha\in (L/M)^\genus}
%	\sum_{\beta\in H^\perp}
%	h_\beta(\tau)
%	\underbrace{e\big(-\tr(\alpha,\beta)\big)}_{=1}
%	e\big(-\tr(\mu,\beta)\big)
%	\\
%	&\quad
%	+
%	\sum_{\beta\not\in H^\perp}
%	h_\beta(\tau)
%	e\big(-\tr(\mu,\beta)\big)
%	\underbrace{
%	\sum_{\alpha\in(L/M)^\genus}
%	e\big(-\tr(\alpha,\beta)\big)
%	}_{=0}
%	\bigg]
%	\\
%	&=
%	\det\tau^k
%	\frac{\sqrt{i}^{\genus(n-2)}}{\underbrace{\lvert D_M\rvert^{\genus/2} / \lvert H \lvert}_{=\lvert D_L\rvert^{\genus/2}}}
%	\underbrace{\sum_{\beta\in H^\perp}
%	h_\beta(\tau)\cdot e\big(-\tr(\mu,\beta)\big)}_{\eqqcolon\clubsuit}.
%	\eas
%	For every~$\beta\in D_\lattice^\genus$, we denote by~$\mu_\beta$ any fixed preimage of~$\beta$ under the natural map~$H^\perp=(\lattice'/M)^\genus\to D_L^\genus$.
%	Then, we may compute the term~$\clubsuit$ above as
%	\bas
%	\clubsuit
%	&=
%	\sum_{\beta\in \underbrace{H^\perp/H}_{\cong D_L^\genus}}
%	\sum_{\alpha\in H}
%	h_{\alpha + \mu_\beta}(\tau)
%	\cdot
%	e\big(-\tr(\mu,\beta+\alpha)\big)
%	\\
%	&=
%	\sum_{\beta\in D_\lattice^\genus}
%	\sum_{\alpha\in H}
%	h_{\alpha + \mu_\beta}(\tau)
%	\cdot
%	e\big(-\tr(\mu,\beta)\big)
%	\cdot
%	\underbrace{e\big(-\tr(\mu,\alpha)\big)}_{=1}
%	\\
%	&=
%	\sum_{\beta\in D_\lattice^\genus}
%	e\big(-\tr(\mu,\beta)\big)
%	\underbrace{
%	\sum_{\alpha\in H}
%	h_{\alpha + \mu_\beta}(\tau)
%	}_{=(h^L)_{\beta}(\tau)}.
%	\eas
%	This concludes the proof.
%	}
%	\\
%	
%	\textcolor{blue}{\textbf{Check that~$\boldsymbol{\langle f, h^L\rangle = \langle f_M,h\rangle}$ for every~$\boldsymbol{f\in M^k_{\genus,\lattice}}$ and~$\boldsymbol{h\in M^k_{\genus,M}}$:}
%	\\
%	We want to show that
%	\be\label{eq;inblTYJBinpr}
%	\langle f, h^L\rangle = \langle f_M,h\rangle,
%	\ee
%	where we denote by~$\langle\cdot{,}\cdot\rangle$ the standard inner product of~$\CC[D_L^\genus]$ as well as the one of~$\CC[D_M^\genus]$.
%	For every~$\alpha\in D_\lattice^\genus$ we denote by~$\mu_\alpha$ a fixed element of the preimage of~$\alpha$ under~$H^\perp\to D_L^\genus$, so that~$\alpha=\overline{\mu_\alpha}$.
%	Since
%	\bas
%	\langle f,h^L\rangle
%	=
%	\sum_{\alpha\in D_\lattice^\genus}
%	f_\alpha\cdot \overline{(h^L)_\alpha}
%	=
%	\sum_{\alpha\in D_\lattice^\genus}
%	f_\alpha
%	\sum_{\beta\in H}
%	\overline{h_{\beta+\mu_\alpha}}
%	\eas
%	and
%	\bas
%	\langle f_M,h\rangle
%	=
%	\sum_{\alpha\in D_M^\genus}
%	(f_M)_\alpha\cdot\overline{h_\alpha}
%	=
%	\sum_{\alpha\in H^\perp}
%	f_{\overline{\alpha}}\cdot\overline{h_\alpha}
%	=
%	\sum_{\alpha\in \underbrace{H^\perp/H}_{\cong D_L^\genus}}
%	\sum_{\beta\in H}
%	f_\alpha\cdot\overline{h_{\beta+\mu_\alpha}},
%	\eas
%	we deduce~\eqref{eq;inblTYJBinpr}.
%	}
	\end{proof}
	
	\subsection{Klingen Eisenstein series}\label{sec;KEseries}
	In this section we introduce vector-valued Klingen Eisenstein series with respect to the Weil representation.
	In the case of unimodular lattices, these series are the classical scalar-valued Klingen Eisenstein series of~\cite[Chapter~II, Section~$5$]{klingen}.
	
	Let~$r\le\gengenus$ be a non-negative integer.
	We denote by~$\klipar{\gengenus}{r}$ and~$\klimet{\gengenus}{r}$ respectively the Klingen parabolic subgroup
	\[
	\klipar{\gengenus}{r}\coloneqq\left\{
	\gamma=\left(\begin{smallmatrix}
	A & B\\
	C & D
	\end{smallmatrix}\right)\in\Sp_{2\gengenus}(\ZZ) :
	A=\left(\begin{smallmatrix}
	\ast & 0\\
	\ast & \ast
	\end{smallmatrix}\right),
	C=\left(\begin{smallmatrix}
	\ast & 0\\
	0 & 0
	\end{smallmatrix}\right),
	D=\left(\begin{smallmatrix}
	\ast & \ast\\
	0 & \ast
	\end{smallmatrix}\right)
	\right\}
	\]
	and its preimage in~$\Mp_{2\gengenus}(\ZZ)$, where the top-left corners of the block decompositions of~$A$, $C$ and~$D$ above are~$r\times r$ matrices.
	
	The following well-known result provides generators of the Klingen parabolic subgroup.
	\begin{lemma}\label{lemma;genklingenpar}
	The Klingen parabolic subgroup~$\klipar{\gengenus}{r}$ is generated by the following matrices.
	\begin{enumerate}[label=(\roman*)]
	\item $\left(\begin{smallmatrix}
	A & & B & \\
	 & U & & \\
	C & & D & \\
	& & & U^{-t}
	\end{smallmatrix}\right)$, where~$\left(\begin{smallmatrix}
	A & B \\
	C & D
	\end{smallmatrix}\right)\in \Sp_{2r}(\ZZ)$ and~$U\in\GL_{\gengenus-r}(\ZZ)$.
	\label{item;genklingenpar1}
	\item $\left(\begin{smallmatrix}
	I_r & & & N \\
	M^t & I_{\gengenus-r} & N^t & B \\
	& & I_r & -M \\
	& & & I_{\gengenus-r}
	\end{smallmatrix}\right)$, where~$N,M\in\ZZ^{r\times(\gengenus-r)}$ and~$B\in\ZZ^{(\gengenus-r)\times(\gengenus-r)}$ satisfy the relation~$N^t M + B = M^t N + B^t$.
	\label{item;genklingenpar2}
	\end{enumerate}
	\end{lemma}
	\begin{proof}
	See e.g.~\cite[Proposition~$4.10$~(i)]{namikawa}.
	\end{proof}
	
	We consider the projections
	\ba
	\ast&\colon\HH_\gengenus \longrightarrow\HH_r,
	\quad
	&&
	\tau=\big(\begin{smallmatrix}
	\tau^*  & \ast\\
	\ast & \ast
	\end{smallmatrix}\big)\longmapsto \tau^*,
	\\
	\ast&\colon \klipar{\gengenus}{r}\longrightarrow \Sp_{2r}(\ZZ),\quad
	&& M=\big(\begin{smallmatrix}
	A & B \\
	C & D
	\end{smallmatrix}\big)
	\longmapsto
	M^*=\big(\begin{smallmatrix}
	A_1 & B_1\\
	C_1 & D_1
	\end{smallmatrix}\big),
	\\
	\ast & \colon\klimet{\gengenus}{r}\longrightarrow \Mp_{2r}(\ZZ),\quad
	&&
	\big(M,\phi(\tau)\big)\longmapsto \big(M^*,\phi^*(\tau^*)\big),
	\ea
	where~$\phi^*(\tau^*)$ is the square root function chosen such that~$\phi(\tau)
%	\left(\left(\begin{smallmatrix}
%	\tau^* & 0 \\
%	0 & \tau'
%	\end{smallmatrix}\right)\right)
	={\det D_4} ^{1/2} \phi^*(\tau^*)$.
%	\[
%	\text{if~$\phi(\tau)=\pm \det(C\tau+D)^{1/2}$,}\qquad\text{then~$\phi^*(\tau^*)\coloneqq\pm \Big((\det D_4) \det(C_1\tau^* + D_1)\Big)^{1/2}$}.
%	\]
	Here $A_1,\dots,D_1$ are the upper-left~$r\times r$ blocks of~$A,\dots,D$, and~$D_4$ is the lower-right~$(\gengenus-r)\times(\gengenus-r)$ block of~$D$.
	We remark that if~$M\in\klipar{\gengenus}{r}$, then
	\bas
	\phi(\tau)^2
	&=
%	\det(C\tau + D) 
%	=
	\det\left(
	\left(\begin{smallmatrix}
	C_1 & 0 \\
	0 & 0
	\end{smallmatrix}\right)\left(\begin{smallmatrix}
	\tau^* & \ast \\
	\ast & \ast
	\end{smallmatrix}\right) + \big(\begin{smallmatrix}
	D_1 & D_2 \\
	0 & D_4
	\end{smallmatrix}\big)
	\right)
	= \det D_4 \cdot \det(C_1\tau^* + D_1)
%	\\
%	&
	= \det D_4 \cdot \phi^*(\tau^*)^2
	,
	\eas
	hence~$\phi^*(\tau^*)$ is defined to be equal to~$\phi(\tau)$ up to a square root of~$\det D_4$, which we choose to be the principal one.
%	This implies that
%	\be\label{eq;iotaand*forKlingPar}
%	\iota\left(\gamma*,\gamma')=\gamma
%	\qquad\text{for every~$\gamma=\left(\begin{smallmatrix}
%	A & & B & \\
%	  & U & & \\
%	C & & D & \\
%	  & &   & U^{-t}
%	\end{smallmatrix}\right)\in\klipar{\gengenus}{r}$.}
%	\ee
	Note also that~$(\gamma\cdot \tau)^* = \gamma^*\cdot\tau^*$ for every~$\gamma\in \klipar{\gengenus}{r}$ and~$\tau\in\HH_\gengenus$.
%	\begin{rem}\label{rem;on*andiota}
%	If~$\gamma_1=\left(
%	\left(\begin{smallmatrix}
%	A & B \\
%	C & D
%	\end{smallmatrix}\right)
%	,\phi_1(\tau_1)\right)$ and~$\gamma_2=\left(
%	\left(\begin{smallmatrix}
%	U & \\
%	  & U^{-t}
%	\end{smallmatrix}\right)
%	, \det U ^{1/2}\right),$
%	then
%	\be\label{eq;remforjan*}
%	\gamma\coloneqq\iota(\gamma_1,\gamma_2)\qquad \text{is such that}\qquad\gamma^*
%	=
%	\gamma_1.
%%	\left(\left(\begin{smallmatrix}
%%	A & B \\
%%	C & D
%%	\end{smallmatrix}\right)
%%	,\phi_1(\tau_1)\right).
%	\ee
%	Note that~$\gamma$ is one of the standard generators of the Klingen parabolic~$\klipar{\gengenus}{r}$; see Lemma~\ref{lemma;genklingenpar}.
%	\end{rem}
%	\textcolor{blue}{\textbf{IMPORTANT REMARK regarding the definition of~$\ast \colon\klimet{\gengenus}{r}\to \Mp_{2r}(\ZZ)$.}
%	\\
%	A priory one could define~$\phi^*(\tau^*)$ in the following way:
%	\[
%	\phi(\tau)=\pm\det(C\tau+D)^{1/2} \longmapsto
%	\phi^*(\tau^*)=\pm\det(C_1\tau^* + D_1)^{1/2}.
%	\]
%	In this case there is no twist of~$\det D_4$, but I think this leads to the same definition of~$\phi^*$ given above.}
	If it is unclear on which space~$\mathfrak{e}_0$ lies, we write~$\zeroe{r}$ to remark that it lies in~$\CC[D_\lattice^r]$.
	\begin{defi}\label{def;KlingEisseries}
	Let~$\gengenus,r\in\ZZ_{\ge0}$ be such that~$r\le\gengenus$.
%	, and let~$2k\equiv n-2$ mod~$4$.
	The \emph{Klingen Eisenstein series} of genus~$\gengenus$ arising from the cusp form~$f\in S^k_{r,\lattice}$ is defined as
	\[
	\klinser{k}{\gengenus}{r}{\lattice}(\tau,f)
	=
	\sum_{\gamma\in\klimet{\gengenus}{r}\backslash \Mp_{2\gengenus}(\ZZ)}
	\phi(\tau)^{-2k}
	\weil{\lattice}{\gengenus}(\gamma)^{-1}
	\Big(f(\gamma^*\cdot\tau^*)
	\otimes \zeroe{\gengenus-r}\Big).
	\]
	The \emph{Siegel Eisenstein series} of genus~$g$ and weight~$k$ with respect to~$\weil{\lattice}{\genus}$
%	(and the zero of~$D_\lattice^\genus$)
	is defined as~$\eiszero{k}{\gengenus}{\lattice}(\tau)\coloneqq\klinser{k}{\gengenus}{0}{\lattice}(\tau,1)$.
	\end{defi}
	It is easy to see by using Lemma~\ref{lemma;genklingenpar} that the series~$\klinser{k}{\gengenus}{r}{\lattice}(f)$ is well-defined.
%	As for the scalar-valued case, the Siegel Eisenstein series~$\eiszero{k}{\gengenus}{\lattice}$ may be constructed as a Klingen Eisenstein series arising from the constant~$f=1$.
	A standard argument shows that~$\eiszero{k}{\gengenus}{\lattice}$ is normalized such that its Fourier coefficients of index~$T=0$ are
	\[
	c_0\big((\eiszero{k}{\gengenus}{\lattice})_\discel\big)
	=\begin{cases}
	1, & \text{if $\discel=0\in D_\lattice^\gengenus$,}\\
	0, & \text{otherwise}.
	\end{cases}
	\]
	
	The space of Siegel Eisenstein series with respect to~$\weil{\lattice}{\gengenus}$ can be of dimension greater than~$1$; see e.g.~\cite[Section~$1.2.3$]{bruinierbook} for the case of genus~$1$.
	For the purposes of the present article, we do not need to introduce other Eisenstein series than~$\eiszero{k}{\gengenus}{\lattice}$, i.e.\ the one arising from the zero of~$D_\lattice^\gengenus$.
%	\textcolor{blue}{\textbf{Here the check that the Klingen Eisenstein series are well defined.}\\
%	We show here that~$\klinser{k}{\gengenus}{r}{\lattice}(f)$ is well-defined, and postpone convergence properties to Proposition~\ref{prop;convpropKEseries} below.
%	Let
%	\[
%	F(\gamma)\coloneqq\phi(\tau)^{-2k}
%	\weil{\lattice}{\gengenus}(\gamma)^{-1}
%	\Big(f(\gamma^*\cdot\tau^*)
%	\otimes \zeroe{\gengenus-r}\Big),\qquad\text{$\gamma=(M,\phi)\in\Mp_{2\gengenus}(\ZZ)$.}
%	\]
%	We want to show that
%	\[
%	F(\gamma_1\cdot\gamma)=F(\gamma)\qquad
%	\text{for every~$\gamma_1=(M_1,\phi_1)\in\klimet{\gengenus}{r}$.}
%	\]
%	By~\eqref{eq;howto-phi} we may assume without loss of generality that~$\phi_1$ is the principal branch of the square root, i.e.\ it is given as~$\phi_1(\tau)=\det(C\tau+D)^{1/2}$.
%	Since~$\gamma_1\cdot\gamma=\big(M_1M , \phi_1(M\cdot\tau)\phi(\tau)\big)$, we deduce that
%	\bas
%	F(\gamma_1\cdot\gamma)
%	&=
%	\phi_1(\gamma\cdot\tau)^{-2k}\phi(\tau)^{-2k}\weil{\lattice}{\gengenus}(\gamma)^{-1}\weil{\lattice}{\gengenus}(\gamma_1)^{-1}
%	\Big(f(\gamma_1^*\cdot\gamma^*\cdot\tau^*)\otimes\zeroe{\gengenus-r}\Big)
%	\\
%	&=
%	\phi_1(\gamma\cdot\tau)^{-2k}\phi(\tau)^{-2k}\weil{\lattice}{\gengenus}(\gamma)^{-1}\weil{\lattice}{\gengenus}(\gamma_1)^{-1}
%	\Big(
%	\phi_1^*(\gamma^*\cdot\tau^*)^{2k}\weil{\lattice}{r}(\gamma_1^*)f(\gamma^*\tau^*)\otimes\zeroe{\gengenus-r}
%	\Big)
%	\\
%	&=
%	\phi_1(\gamma\cdot\tau)^{-2k}\phi_1^*(\gamma^*\cdot\tau^*)^{2k}\phi(\tau)^{-2k}\weil{\lattice}{\gengenus}(\gamma)^{-1}\weil{\lattice}{\gengenus}(\gamma_1)^{-1}
%	\Big(
%	\weil{\lattice}{r}(\gamma_1^*)f(\gamma^*\tau^*)\otimes\zeroe{\gengenus-r}
%	\Big).
%	\eas
%	Since~$\phi_1(\tau)=\det D_4^{1/2}\phi_1^*(\tau^*)$, we deduce that
%%	We may compute
%%	\ba\label{eq;checkklingeisphiphi^*}
%%	\phi_1(\tau)
%%	&=
%%	\det\Big(
%%	\big(\begin{smallmatrix}
%%	C_1 & 0\\
%%	0 & 0
%%	\end{smallmatrix}\big)
%%	\tau + \big(\begin{smallmatrix}
%%	D_1 & D_2\\
%%	0 & D_4
%%	\end{smallmatrix}\big)
%%	\Big)^{1/2}
%%	=
%%	\det\big(
%%	\begin{smallmatrix}
%%	C_1\tau^* + D_1 & \ast \\
%%	0 & D_4
%%	\end{smallmatrix}
%%	\big)^{1/2}
%%	\\
%%	&=
%%	\Big((\det D_4)\det(C_1\tau^* + D_1)\Big)^{1/2}
%%	=
%%	\textcolor{blue}{\det D_4^{-1/2}}\phi_1^*(\tau^*),
%%	\ea
%%	\textcolor{blue}{(The factor~$\det D_4^{-1/2}$ is a test.
%%	Let's see whether it works...
%%	It means that we define~$\phi_1^*(\tau^*)=\det D_4^{1/2}\phi_1(\tau)$.)}\\
%%	hence we deduce that
%	\[
%	F(\gamma_1\cdot\gamma)
%	=
%	{\det D_4}^{-k}\phi(\tau)^{-2k}\weil{\lattice}{\gengenus}(\gamma)^{-1}\weil{\lattice}{\gengenus}(\gamma_1)^{-1}
%	\Big(
%	\weil{\lattice}{r}(\gamma_1^*)f(\gamma^*\tau^*)\otimes\zeroe{\gengenus-r}
%	\Big).
%	\]
%	To prove that~$F(\gamma_1\cdot\gamma)=F(\gamma)$ it is then enough to show that
%	\be\label{eq;todocheckweldefklingeis}
%	\weil{\lattice}{\gengenus}(\gamma_1)
%	(\mathfrak{e}_\discel \otimes\zeroe{\gengenus-r})
%	=
%	{\det D_4} ^{-k}\cdot
%	\big(
%	\weil{\lattice}{r}(\gamma_1^*)
%	\mathfrak{e}_\discel
%	\big)\otimes\zeroe{\gengenus-r}
%	\ee
%	for every~$\discel\in D_\lattice^{r}$.
%	By means of Lemma~\ref{lemma;proofweilreponsubmeta}, this can be easily checked for the standard generators of the (inclusion in~$\Mp_{2\gengenus}(\ZZ)$ of the) Klingen parabolic subgroup~$\klipar{\gengenus}{r}$ provided by Lemma~\ref{lemma;genklingenpar}.}
%	
%	\textcolor{blue}{We check~\eqref{eq;todocheckweldefklingeis} choosing~$\gamma_1$ to be a standard generator as in Lemma~\ref{lemma;genklingenpar}.	
%	We begin with the case of
%	\[
%	\gamma_1=\left(\left(\begin{smallmatrix}
%	A & & B & \\
%	 & U & & \\
%	 C & & D & \\
%	 & & & U^{-t}
%	\end{smallmatrix}\right),\phi_1\right)
%	=
%	\iota\big(\gamma, R(U)\big),
%	\]
%	where~$\iota$ is the inclusion~\eqref{eq;inclprodsimpl},~$\gamma=\big(\big(\begin{smallmatrix}
%	A & B \\
%	C & D
%	\end{smallmatrix}\big), \phi\big)$, and~$R(U)$ is the generator of~$\Mp_{2(\gengenus-r)}(\ZZ)$ with~$U$ and~$U^{-t}$ as diagonal blocks.}
%	
%	\textcolor{blue}{By Lemma~\ref{lemma;proofweilreponsubmeta} we have
%	\bas
%	\weil{\lattice}{\gengenus}\big(\iota(\gamma,R(U))\big)\big(\mathfrak{e}_\discel\otimes\zeroe{\gengenus-r}\big)
%	&=
%	\big(\weil{\lattice}{r}(\gamma)\mathfrak{e}_\discel\big)\otimes
%	\big(
%	\weil{\lattice}{\gengenus-r}(R(U))\zeroe{\gengenus-r}
%	\big)
%	\\
%	&=\det U^{(2-n)/2} \big(\weil{\lattice}{r}(\gamma)\mathfrak{e}_\discel\big)\otimes \zeroe{\gengenus-r}.
%	\eas
%	Since~$2k\equiv n-2$ mod~$4$, we deduce that~$\det U^{(2-n)/2} = \det U ^{-k}$. (\textbf{If~$\boldsymbol{k}$ satisfies only the weaker condition~$\boldsymbol{2k\equiv n-2}$ mod~2, then I do not think this is true anymore.})
%	As remarked in~\eqref{eq;remforjan*}, we have that~$(\iota(\gamma,R(U))^*=\gamma$.
%	These imply~\eqref{eq;todocheckweldefklingeis}, since~$D_4$ in this case equals~$U^{-t}$ and hence~$\det D_4 = \det U ^{-1}$.}
%	
%	\textcolor{blue}{We conclude with the second kind of generators given by Lemma~\ref{lemma;genklingenpar}, namely we choose~$\gamma_1=(M_1,\phi_1)$, where
%	\[
%	M_1=
%	\left(\begin{smallmatrix}
%	I_r & & & N \\
%	M^t & I_{\gengenus-r} & N^t & B \\
%	& & I_r & -M \\
%	& & & I_{\gengenus-r}
%	\end{smallmatrix}\right)
%	\qquad \text{and} \qquad
%	\phi_1(\tau)=1,
%	\]
%	with~$N,M\in\ZZ^{r\times(\gengenus-r)}$ and~$B\in\ZZ^{(\gengenus-r)\times(\gengenus-r)}$ satisfying the relation~$N^t M + B = M^t N + B^t$.
%	Let~$\tilde B\coloneqq B - M^t N$.
%	Because of the relation satisfied by~$B$, the matrix~$\tilde B$ is symmetric.
%	We rewrite~$\gamma_1$ in terms of the standard generators of~$\Mp_{2\gengenus}(\ZZ)$ as follows.
%	Note that
%	\bas
%	\underbrace{\left(\begin{smallmatrix}
%	I_r & & & \\
%	M^t & I_{\gengenus-r} & & \\
%	& & I_r & -M \\
%	& & & I_{\gengenus-r}
%	\end{smallmatrix}\right)}_{
%	M_2
%	}
%	\underbrace{\left(\begin{smallmatrix}
%	I_r & & & N \\
%	& I_{\gengenus-r} & N^t & \tilde B \\
%	& & I_r & \\
%	& & & I_{\gengenus-r}
%	\end{smallmatrix}\right)}_{
%	M_3
%	}
%	=
%	\left(\begin{smallmatrix}
%	I_r & & & N\\
%	M^t & I_{\gengenus-r} & N^t & M^t N + \tilde B \\
%	 & & I_r & -M \\
%	 & & & I_{\gengenus-r}
%	\end{smallmatrix}\right)
%	=M_1.
%	\eas
%	}
%	
%	\textcolor{blue}{Let~$\gamma_j=(M_j,\phi_j)$, for~$j=2,3$, be the standard generators of~$M_{2\gengenus}(\ZZ)$ arising from~$M_j$.
%	These are such that~$\gamma_1=\gamma_2\cdot\gamma_3$.
%	We may then compute
%	\bas
%	\weil{\lattice}{\gengenus}(\gamma_1)
%	(\mathfrak{e}_\discel \otimes\zeroe{\gengenus-r})
%	&=
%	\weil{\lattice}{\gengenus}(\gamma_2)\weil{\lattice}{\gengenus}(\gamma_3)(\mathfrak{e}_\discel\otimes\zeroe{\gengenus-r})
%	\\
%	&=
%	\underbrace{e\Big( \trace \big(\begin{smallmatrix}
%	q(\discel) & 0 \\
%	0 & 0
%	\end{smallmatrix}\big)
%	\big(\begin{smallmatrix}
%	0 & N \\
%	N^t & \tilde B
%	\end{smallmatrix}\big)\Big)}_{=1}\weil{\lattice}{\gengenus}(\gamma_2)(\mathfrak{e}_\discel\otimes\zeroe{\gengenus-r})
%	\\
%	&=
%	\mathfrak{e}_{(\discel,0,\dots,0)\big(
%	\begin{smallmatrix}
%	I_r & 0 \\
%	-M^t & I_{\gengenus-r}
%	\end{smallmatrix}\big)}
%	=\mathfrak{e}_\discel \otimes \zeroe{\gengenus-r}.
%	\eas
%	Since in this case~$\det D_4=1$ and~$\gamma_1^*=(I_r,1)$, we deduce~\eqref{eq;todocheckweldefklingeis}.}
%	\\
	
	Scalar-valued Eisenstein series have excellent convergence properties.
	These are satisfied also by the vector-valued Klingen Eisenstein series defined above, as illustrated in the following generalization of~\cite[p.~$67$, Theorem~$1$]{klingen}.
%	The proof relies on the bounds provided in Section~\ref{sec;someboundscusp}.
%	Following the wording of~\cite{klingen}, 
	For~$c>0$, we define the \emph{vertical strip of height~$c$} as
	\[
	V_\gengenus(c)
	\coloneqq
	\{
	\tau\in\HH_\gengenus : \text{$\tr(x^2)\le c^{-1}$ and~$y\ge cI_\gengenus$}
	\}.
	\]
	\begin{prop}\label{prop;convpropKEseries}
	Let~$f\in S^k_{r,\lattice}$.
	If~$k>\gengenus+r+1$, then the Klingen Eisenstein series~$\klinser{k}{\gengenus}{r}{\lattice}(f)$ converges absolutely and uniformly on any vertical strip of positive height, and is a modular form of genus~$\gengenus$ and weight~$k$ with respect to~$\weil{\lattice}{\gengenus}$.
	\end{prop}
	
	To prove Proposition~\ref{prop;convpropKEseries} we need the following auxiliary result, which may be easily shown as in~\cite[p.~$57$, Proposition~$4$]{klingen}.
	\begin{lemma}\label{lemma;bound2blafromKling}
	Let~$f\in S^k_{\gengenus,\lattice}$.
	There exist positive numbers~$a_1$ and~$a_2$ such that
	\be\label{eq;bound2blafromKling}
	\det y^{k/2} \lvert f(\tau)\rvert \le a_1 e^{-a_2 \det y^{1/\gengenus}}
	\qquad
	\text{for every~$\tau\in\HH_\gengenus$.}
	\ee
	\end{lemma}
%	\textcolor{blue}{\begin{proof}[Proof of Lemma~\ref{lemma;bound2blafromKling}]
%	It is easy to see that~$\det y^{k/2} \lvert f(\tau)\rvert$ is~$\Sp_{2\gengenus}(\ZZ)$-invariant function, where we consider~$\Sp_{2\gengenus}(\ZZ)$ to be included in~$\Mp_{2\gengenus}(\ZZ)$ in the standard way.
%	To bound the left-hand side of~\eqref{eq;bound2blafromKling}, we may then choose~$\gamma\in\Sp_{2\gengenus}(\ZZ)$ such that~$\tilde\tau=\tilde x + i \tilde y = \gamma\cdot\tau$ lies in the Siegel fundamental domain of~$\HH_\gengenus$.
%	This implies that~$\tilde y$ is reduced in the sense of Minkowski and, by~\cite[p.~$30$, Lemma~$2$]{klingen}, that~$\tilde y\ge c I_\gengenus$ for some positive~$c$.
%	We deduce from Proposition~\ref{prop;bound1blafromKling} that there exist two positive numbers~$c_1$ and~$c_2$ such that
%	\be\label{eq;boundinprooffscf}
%	\det y^{k/2} \lvert f(\tau)\rvert \le
%	c_1 \det \tilde y^{k/2} e^{-c_2\tr \tilde y}.
%	\ee
%	As remarked in~\cite[p.~$58$, Proof of Proposition~$4$]{klingen}, we have
%	\[
%	\det y\le\det\tilde y
%	\qquad\text{and}\qquad
%	\det\tilde y ^{1/\gengenus}\le (\tr \tilde y )/ \gengenus.
%	\]
%	These imply that we may further bound~\eqref{eq;boundinprooffscf} as
%	\[
%	\det y^{k/2} \lvert f(\tau)\rvert
%	\le 
%	c_1 \det \tilde y^{k/2} e^{-c_2 \gengenus \det \tilde y ^{1/\gengenus}}
%	\le 
%	a_1 e^{-a_2 \det\tilde y ^{1/\gengenus}}
%	\le 
%	a_1 e^{-a_2 \det y^{1/\gengenus}}.
%	\qedhere
%	\]
%	\end{proof}}
	\begin{proof}[Proof of Proposition~\ref{prop;convpropKEseries}]
	Under the assumption of local uniform convergence, the series gives a holomorphic function.
	It is easy to see that it transforms as a modular form.
%	 by using the group law~\eqref{eq;grouplawmetgr} of the metaplectic group.
	
	We prove the absolute and uniform convergence.
	Since the Weil representation~$\weil{\lattice}{\gengenus}$ is unitary with respect to the standard inner product of~$\CC[D_\lattice^\gengenus]$, we deduce that
	\[
	\lvert
	\klinser{k}{\gengenus}{r}{\lattice}(f)
	\rvert
	\le 
	\sum_{\gamma\in\klimet{\gengenus}{r}\backslash \Mp_{2\gengenus}(\ZZ)}
	\lvert
	\phi(\tau)
	\rvert^{-2k}
	\lvert f(\gamma^*\cdot\tau^*)\rvert.
	\]
	By Lemma~\ref{lemma;bound2blafromKling} the norm of~$f(\tau^*)$ is bounded (up to a positive factor depending on~$f$) by~$\det \Im(\tau^*)^{-k/2}$.
	We deduce that the function
	\ba
	\majo{k}{\gengenus}{r}(\tau)
	&\coloneqq 
	\sum_{\gamma\in\klimet{\gengenus}{r}\backslash \Mp_{2\gengenus}(\ZZ)}
	\det \Im(\gamma^*\cdot\tau^*)^{-k/2}
	\lvert
	\phi(\tau)
	\rvert^{-2k}
%	\\
%	&=
%	\sum_{\gamma\in\klipar{\gengenus}{r}\backslash \Sp_{2\gengenus}(\ZZ)}
%	\det \Im(\gamma^*\cdot\tau^*)^{-k/2}
%	\lvert
%	\det (C\tau+D)
%	\rvert^{-k}
	,
	\ea
	if convergent, is a majorant of~$\lvert
	\klinser{k}{\gengenus}{r}{\lattice}(f)
	\rvert$.
	Note that~$\majo{k}{\gengenus}{r}$ is the same as~\cite[p.~$63$, ($11$)]{klingen}, the only difference is that we allow the weight~$k$ to be half-integral instead of even.
	This difference does not play any role, and the proof can be completed as in~\cite[p.~$67$, Theorem~$1$]{klingen}.
%	\textcolor{blue}{\textbf{For the sake of completeness, we quickly recall it.}\\
%	By~\cite[p.~$63$, Lemma~$2$]{klingen}, for every positive~$c$ there exists a positive constant~$a=a(\gengenus,c)$ such that
%	\[
%	\det \Im(\gamma^*\cdot\tau^*)
%	\cdot
%	\lvert
%	\det(C\tau+D)
%	\rvert^2
%	\ge 
%	a\det \Im\big(\gamma^*\cdot (iI_\gengenus)^*\big)
%	\cdot
%	\lvert
%	\det(CiI_\gengenus + D)
%	\rvert^2
%	\]
%	for all~$\tau\in V_\gengenus(c)$ and~$\gamma\in\Sp_{2\gengenus}(\ZZ)$.
%	This implies that if~$\majo{k}{\gengenus}{r}(iI_\gengenus)$ converges, then~$\majo{k}{\gengenus}{r}(\tau)$ and hence~$\lvert
%	\klinser{k}{\gengenus}{r}{\lattice}(f)
%	\rvert$ can be majorized by~$\majo{k}{\gengenus}{r}(iI_\gengenus)$ uniformly on any vertical strip of positive height.
%	The rest of this proof is devoted to show that~$\majo{k}{\gengenus}{r}(iI_\gengenus)$ converges.}
%	\textcolor{blue}{Let~$\eta_1$, $\eta_2$,~$\eta_3$ be symmetric matrices such that
%	\[
%	\Im(\gamma\cdot \tau)
%	=
%	\left(\begin{smallmatrix}
%	I_r & 0 \\
%	\eta_3 & I_{\gengenus-r}
%	\end{smallmatrix}\right)
%	\left(\begin{smallmatrix}
%	\eta_1 & 0 \\
%	0 & \eta_2
%	\end{smallmatrix}\right)
%	\left(\begin{smallmatrix}
%	I_r & \eta_3 \\
%	0 & I_{\gengenus-r}
%	\end{smallmatrix}\right)
%	=
%	\left(
%	\begin{smallmatrix}
%	\eta_1 & \eta_1\eta_3 \\
%	\eta_3\eta_1 & \eta_3\eta_1\eta_3 + \eta_2
%	\end{smallmatrix}\right).
%	\]
%	Here~$\eta_1\in\RR^{r\times r}$ and~$\eta_2\in\RR^{(\gengenus-r)\times(\gengenus-r)}$.
%	To emphasize that these arise from~$\tau$, we also denote them by~$\eta_j(\tau)$.
%	Using the well-known that if~$A$ is an invertible matrix, then~$\det\big(\begin{smallmatrix}
%	A & B \\ C & D
%	\end{smallmatrix}\big)
%	=\det(A)\cdot\det(D-CA^{-1}B)$, It is easy to see that
%	\[
%	\majo{k}{\gengenus}{r}(i I_\gengenus)
%	=
%	\sum_{\gamma\in\klipar{\gengenus}{r}\backslash \Sp_{2\gengenus}(\ZZ)}
%	\big(\det \eta_2(iI_\gengenus)\big)^{k/2}.
%	\]
%	The bound given by~\cite[p.~$64$, Lemma~$3$]{klingen} does not depend on the choice of~$k$.
%	Hence, we may apply it and deduce that for every compact subset~$K$ of positive volume contained in the Siegel fundamental domain, there exists a constant~$\beta_1$ such that
%	\[
%	\big(\det \eta_2(iI_\gengenus)\big)^{k/2}
%	\le 
%	\beta_1\big( \det \eta_2(\tau)\big)^{k/2}
%	\qquad
%	\text{for all~$\tau\in K$ and~$\gamma\in\Sp_{2\gengenus}(\RR)$.}
%	\]
%	We integrate over~$K$ both sides of the previous inequality and divide by~$\Vol(K)$, obtaining
%	\bas
%	\big(\det\eta_2(iI_\gengenus)\big)^{k/2}
%	\le 
%	\frac{\beta_1}{\Vol(K)}\int_K\big(\det\eta_2(\tau)\big)^{k/2}\,dv_\gengenus
%	=
%	\beta_2\int_{\gamma(K)} \big(\det y_2\big)^{k/2}\, dv_\gengenus.
%	\eas
%	Here~$y_2$ is constructed as for~$\eta_2$ but arising from~$\gamma= I_{2\gengenus}$, $dv_\gengenus=dx\,dy/\det y^{\gengenus+1}$ and~$\beta_2\coloneqq\beta_1/\Vol(K)$.
%	This implies that we may bound
%	\bas
%	\majo{k}{\gengenus}{r}(i I_\gengenus)
%	&\le 
%	\beta_2
%	\sum_{\gamma\in\klipar{\gengenus}{r}\backslash \Sp_{2\gengenus}(\ZZ)}
%	\int_{\gamma(K)} \big(\det y_2\big)^{k/2}\, dv_\gengenus
%%	\\
%%	&
%	=
%	\beta_2\underbrace{\int_{\bigcup_{\gamma\in\klipar{\gengenus}{r}\backslash\Sp_{2\gengenus}(\ZZ)} \gamma(K)}
%	(\det y_2) ^{k/2} dv_\gengenus}_{(\star)}.
%	\eas
%	Hence, it is enough to show that the integral~$(\star)$ converges for all~$k>\gengenus+r+1$.
%	This is illustrated in~\cite[p.~$66$, Paragraphs after~($14$)]{klingen}, and does not depend on whether~$k$ is integral or not.}
	\end{proof}
		
	\subsection{The Siegel operator}\label{sec;Siegelphi}
	Let~$0\le r\le\gengenus$.	
%	Recall that we may identify~$\CC[D_\lattice^{r}]\otimes\CC[D_\lattice^{\gengenus-r}]$ with~$\CC[D_\lattice^\gengenus]$ under the isomorphism~$\iota$; see~\eqref{eq;inclprodsimpl}.
%	\be\label{eq;identgroupalgg-11}
%	\CC[D_\lattice^{\gengenus-1}]\otimes\CC[D_\lattice]\longrightarrow\CC[D_\lattice^\gengenus],\qquad \mathfrak{e}_\discel\otimes\mathfrak{e}_\disceltwo\longmapsto\mathfrak{e}_{(\discel,\disceltwo)}.
%	\ee
	\begin{defi}
	Let~$\disceltwo\in D_\lattice^{\gengenus-r}$.
	The \emph{Siegel~$\Phi^{\gengenus-r}_\disceltwo$ operator} is defined as
	\be\label{eq;vvSiegelphi}
	\Phi^{\gengenus-r}_\disceltwo f(\tau)
	=
	\sum_{\discel\in D_\lattice^{r}}
	\lim_{\lambda\to\infty}
	f_{(\discel,\disceltwo)}\left(\left(\begin{smallmatrix}
	\tau & 0\\
	0 & i\lambda I_{\gengenus-r}
	\end{smallmatrix}\right)\right)\mathfrak{e}_\discel,\qquad \tau\in\HH_{r},
	\ee
	for every~$f=\sum_{\discelthree\in D_\lattice^\gengenus}f_\discelthree \mathfrak{e}_\discelthree\in M^k_{\gengenus,\lattice}$.
	If~$\gengenus-r=1$, then we simply denote the operator as~$\Phi_\disceltwo$.
	\end{defi}
	As for the scalar-valued case, the limit defining the Siegel operator converges since the Fourier expansion of~$f$ converges uniformly on vertical strips of positive heights.
	
	Note that if~$q(\disceltwo)\neq 0$, then~$\Phi^{\gengenus-r}_\disceltwo$ maps every Siegel modular form to~$0$.
	Moreover, if~$\disceltwo=(\disceltwo_j)_j$ with~$\beta_j\in D_\lattice$, then~$\Phi_\beta^{\gengenus-r}$ may be written as the composition of Siegel operators~$\Phi^{\gengenus-r}_\beta=\Phi_{\beta_{\gengenus-r}}\circ\dots\circ\Phi_{\beta_1}$.
%	\textcolor{blue}{
%	\textbf{Here a quick check of the last assertion.}
%	\\
%	For simplicity, we work out the case of~$\beta=(\beta_1,\beta_2)$.
%	We write
%	\[
%	\Phi_{\beta_2}f(\tilde\tau)=
%	\sum_{\tilde\alpha\in D_\lattice^{\gengenus-1}}\lim_{\tilde\lambda\to\infty}
%	f_{(\tilde\alpha,\beta_2)}\big(\big(\begin{smallmatrix}
%	\tilde\tau & 0\\
%	0 & i\tilde\lambda
%	\end{smallmatrix}\big)\big)\mathfrak{e}_{\tilde\alpha},\qquad
%	\tilde\tau\in\HH_{\gengenus-1}.
%	\]
%	Since the Fourier expansion of~$f$ converges uniformly on vertical strips of positive heights in~$\HH_{\gengenus}$, we may compute
%	\bas
%	\Phi_{\beta_1}\circ\Phi_{\beta_2} f(\tau)
%	&=
%	\Phi_{\beta_1}\Big(
%	\sum_{\tilde\alpha\in D_\lattice^{\gengenus-1}}\lim_{\tilde\lambda\to\infty}
%	f_{(\tilde\alpha,\beta_2)}\big(\big(\begin{smallmatrix}
%	\bullet & 0\\
%	0 & i\tilde\lambda
%	\end{smallmatrix}\big)\big)\mathfrak{e}_{\tilde\alpha}
%	\Big)(\tau)
%	\\
%	&=
%	\sum_{\discel\in D_\lattice^{\gengenus-2}}
%	\lim_{\lambda\to\infty}
%	\Big(
%	\sum_{\tilde\alpha\in D_\lattice^{\gengenus-1}}\lim_{\tilde\lambda\to\infty}
%	f_{(\tilde\alpha,\beta_2)}\big(\big(\begin{smallmatrix}
%	\bullet & 0\\
%	0 & i\tilde\lambda
%	\end{smallmatrix}\big)\big)\mathfrak{e}_{\tilde\alpha}
%	\Big)_{(\alpha,\beta_1)}
%	\big(\begin{smallmatrix}
%	\tau & 0\\
%	0 & i\lambda
%	\end{smallmatrix}\big)\mathfrak{e}_\discel
%	\\
%	&=
%	\sum_{\alpha\in D_\lattice^{\gengenus-2}}
%	\lim_{\substack{\lambda\to\infty \\ \tilde\lambda\to\infty}}
%	f_{(\alpha,\beta_1,\beta_2)}
%	\Big(\begin{smallmatrix}
%	\tau & 0 & 0\\
%	0 & i\lambda & 0\\
%	0 & 0 & i\tilde\lambda
%	\end{smallmatrix}\Big)\mathfrak{e}_\discel
%	\\
%	&=
%	\sum_{\alpha\in D_\lattice^{\gengenus-2}}
%	\lim_{\lambda\to\infty}
%	f_{(\alpha,\beta_1,\beta_2)}
%	\Big(\begin{smallmatrix}
%	\tau & 0 \\
%	0 & i\tilde\lambda I_2
%	\end{smallmatrix}\Big)\mathfrak{e}_\discel
%	\\
%	&=
%	\Phi_\beta^2 f (\tau)
%	\eas
%	for every~$\tau\in\HH_{\gengenus-2}$.
%	}\\
%	With a proof analogous to~\cite[Section~$5$, Proposition~$1$]{klingen} one can deduce the following result.
	\begin{lemma}\label{lemma;phivviswelldef}
	Let~$\disceltwo\in D_\lattice^{\gengenus-r}$.
	If~$f\in M^k_{\gengenus,\lattice}$, then~$\Phi^{\gengenus-r}_\disceltwo f\in M^k_{r,L}$.
%	The function~$\Phi_\disceltwo f$ is a Siegel modular form of weight~$k$ and genus~$\gengenus-1$ with respect to the Weil representation~$\weil{\lattice}{\gengenus-1}$.
	\end{lemma}
	\begin{proof}
	It is analogous to~\cite[Section~$5$, Proposition~$1$]{klingen}.
	\end{proof}
%	\begin{proof}
%	\textcolor{blue}{Let~$\gamma^*=\left(\left(\begin{smallmatrix}
%	A^* & B^*\\
%	C^* & D^*
%	\end{smallmatrix}\right), \det(C^*\tau^* + D^*)^{1/2}\right)\in\Mp_{2r}(\ZZ)$, where we denote by~$\tau^*$ the variable in~$\HH_{r}$.
%	Let
%	\[
%	\tau=\left(\begin{smallmatrix}
%	\tau^* & 0\\
%	0 & i\lambda I_{\gengenus-r}
%	\end{smallmatrix}\right)\in\HH_\gengenus,\qquad\text{for some~$\lambda\in\RR_{>0}$,}
%	\]
%	and let~$\gamma=\left(\left(\begin{smallmatrix}
%	A & B\\
%	C & D
%	\end{smallmatrix}\right),\det(C\tau+D)^{1/2}\right)\in\Mp_{2\gengenus}(\ZZ)$, where
%	\[
%	A=\left(\begin{smallmatrix}
%	A^* & 0\\
%	0 & I_{\gengenus-r}
%	\end{smallmatrix}\right),\qquad
%	B=\left(\begin{smallmatrix}
%	B^* & 0\\
%	0 & 0
%	\end{smallmatrix}\right),
%	\qquad
%	C=\left(\begin{smallmatrix}
%	C^* & 0\\
%	0 & 0
%	\end{smallmatrix}\right),
%	\qquad
%	D=\left(\begin{smallmatrix}
%	D^* & 0\\
%	0 & I_{\gengenus-r}
%	\end{smallmatrix}\right).
%	\]
%	This means that~$\gamma=\iota\big(\gamma^*, (I_{2(\gengenus-r)},1)\big)$.
%	In fact, it is easy to check that
%	\[
%	A\tau+B = \left(
%	\left(\begin{smallmatrix}
%	A^*\tau^* + B^* & 0\\
%	0 & i\lambda I_{\gengenus-r}
%	\end{smallmatrix}\right)
%	\right)
%	\qquad\text{and}\qquad
%	C\tau + D = \left(
%	\left(\begin{smallmatrix}
%	C^*\tau^* + D^* & 0\\
%	0 & I_{\gengenus-r}
%	\end{smallmatrix}\right)
%	\right),
%	\]
%	from which we deduce that
%	\[
%	\det(C\tau+D)=\det(C^*\tau^* + D^*)\qquad\text{and}\qquad
%	\gamma\cdot\tau = \left(
%	\left(\begin{smallmatrix}
%	\gamma^*\cdot\tau^* & 0\\
%	0 & i\lambda I_{\gengenus-r}
%	\end{smallmatrix}\right)
%	\right).
%	\]
%	Since~$f$ is a modular form of weight~$k$ with respect to~$\weil{\lattice}{\gengenus}$, we obtain that for the previous choices of~$\tau$ and~$\gamma$ we have
%	\ba\label{eq;inproofofweldefvvphi}
%	f\left(
%	\left(\begin{smallmatrix}
%	\gamma^*\tau^* & 0\\
%	0 & i\lambda I_{\gengenus-r}
%	\end{smallmatrix}\right)
%	\right)
%	&=
%	f(\gamma\cdot\tau)
%	=
%	\det(C\tau+D)^k\weil{\lattice}{\gengenus}(\gamma) f (\tau)
%	\\
%	&=
%	\det(C^*\tau^* + D^*)^k \weil{\lattice}{\gengenus}(\gamma) f \left(
%	\left(\begin{smallmatrix}
%	\tau^* & 0\\
%	0 & i\lambda I_{\gengenus-r}
%	\end{smallmatrix}\right)
%	\right).
%	\ea
%	This implies that
%	\bas
%	\Phi_\disceltwo^{\gengenus-r} f(\gamma^*\cdot\tau^*)
%	&=
%	\sum_{\discel\in D_\lattice^{r}}\lim_{\lambda\to\infty}f_{(\discel,\disceltwo)}\left(
%	\left(\begin{smallmatrix}
%	\gamma^*\cdot\tau^* & 0\\
%	0 & i\lambda I_{\gengenus-r}
%	\end{smallmatrix}\right)
%	\right)\mathfrak{e}_\discel
%	\\
%	&=
%	\det(C^*\tau^* + D^*)^k
%	\sum_{\discel\in D_\lattice^{r}}
%	\lim_{\lambda\to\infty}
%	\left(
%	\weil{\lattice}{\gengenus}(\gamma) f
%	\right)_{(\discel,\disceltwo)}\left(
%	\left(\begin{smallmatrix}
%	\tau^* & 0\\
%	0 & i\lambda I_{\gengenus-r}
%	\end{smallmatrix}\right)
%	\right)\mathfrak{e}_\discel.
%	\eas
%	To prove that~$\Phi_\disceltwo^{\gengenus-r} f$ behaves as a modular form of weight~$k$ with respect to~$\weil{\lattice}{r}$, it is then enough to show that~$\weil{\lattice}{\gengenus}(\gamma)\mathfrak{e}_{(\discel,\disceltwo)}
%	=
%	\big(
%	\weil{\lattice}{r}(\gamma^*)\mathfrak{e}_\discel
%	\big)\otimes\mathfrak{e}_\disceltwo$ for every~$\discel\in D_\lattice^{r}$.
%	This follows from Lemma~\ref{lemma;proofweilreponsubmeta}.}
%	\be\label{eq;wwwinwdofphi}
%	\bigg(
%	\weil{\lattice}{\gengenus}(\gamma)\Big(
%	\sum_{\disceltwo\in D_\lattice^{\gengenus}}f_\disceltwo\mathfrak{e}_\disceltwo
%	\Big)
%	\bigg)_{(\discel,\disceltwo)}
%	=
%	\bigg(\weil{\lattice}{\gengenus-1}(\gamma^*)\Big(
%	\sum_{\disceltwo\in D_\lattice^{\gengenus-1}} f_{(\disceltwo,\disceltwo)}\mathfrak{e}_\disceltwo
%	\Big)\bigg)_{\discel}
%%	\sum_{\discel\in D_\lattice^{\gengenus-1}}\bigg(
%%	\weil{\lattice}{\gengenus}(\gamma)\Big(
%%	\sum_{\disceltwo\in D_\lattice^{\gengenus}}f_\disceltwo\mathfrak{e}_\disceltwo
%%	\Big)
%%	\bigg)_{(\discel,\disceltwo)}\mathfrak{e}_\discel
%%	=
%%	\bigg(\weil{\lattice}{\gengenus-1}(\gamma^*)\Big(
%%	\sum_{\disceltwo\in D_\lattice^{\gengenus-1}} f_{(\disceltwo,\disceltwo)}\mathfrak{e}_\disceltwo
%%	\Big)\bigg).
%	\ee
%	for every~$\delta\in D_\lattice^{\gengenus-1}$.
%	In fact, it is enough to prove~\eqref{eq;wwwinwdofphi} for~$\gamma^*$ to be a standard generator of~$\Mp_{2(\gengenus-1)}(\ZZ)$.
%	
%	We begin with the first kind of standard generators~$\gamma^*=T^*=\big(\big(\begin{smallmatrix}
%	I_{\gengenus-1} & B^*\\
%	0 & I_{\gengenus-1}
%	\end{smallmatrix}\big),1\big)$, for some symmetric integral~$B^*$.
%	In this case we have~$\gamma=\big(\big(
%	\begin{smallmatrix}
%	I_{\gengenus} & B\\
%	0 & I_{\gengenus-1}
%	\end{smallmatrix}
%	\big),1
%	\big)$, with~$B$ as above.
%	It is easy to compute the left-hand and right-hand side of~\eqref{eq;wwwinwdofphi}, and see that they are both equal to~$f_{\discel,\mu} e\big(\tr (q(\discel)B^*)\big)$.
%	
%	We now conclude with~$\gamma^*=S^*=\big(\big(\begin{smallmatrix}
%	0 & -I_{\gengenus-1}\\
%	I_{\gengenus-1} & 0
%	\end{smallmatrix}\big),\det\tau^{1/2}\big)$.
%	In this case we have
%	\[
%	\gamma=\left(\left(\begin{smallmatrix}
%	0 & 0 & -I_{\gengenus-1} & 0\\
%	0 & 1 & 0 & 0\\
%	I_{\gengenus-1} & 0 & 0 & 0\\
%	0 & 0 & 0 & 1
%	\end{smallmatrix}
%	\right),\dots\right).
%	\]
%	\textcolor{red}{\textbf{Problem:}
%	How to compute~$\weil{\lattice}{\gengenus}(\gamma)$? That~$\gamma$ is not a standard generator of~$\Mp_{2\gengenus}(\ZZ)$.
%	}\\
%	\textcolor{blue}{
%	From Jan: check Wei Zhang thesis (email from Jan). Weil rep should be compatible with respect to block matrices embeddings!
%	}
%	\end{proof}
	
	\end{document}
\documentclass[../Lefschetz_decomposition.tex]{subfiles}

	\begin{document}
	\begin{thm}\label{thm;Phi0onkling}
	Let~$f\in S^k_{r,\lattice}$.
	If~$k>\gengenus+r+1$, then~$\Phi_0^{\gengenus-r} \klinser{k}{\gengenus}{r}{\lattice} (f) =f$.
Moreover, if~$\disceltwo\in D_\lattice$ is non-zero, then~$\Phi_\disceltwo \klinser{k}{\gengenus}{r}{\lattice} (f) =0$.
	\end{thm}
	\begin{proof}
%	We rewrite~$\big(\Phi_\disceltwo^{\gengenus-r}\klinser{k}{\gengenus}{r}{\lattice}(f)\big)(\tau)$ as
%	\bas
%%	\big(\Phi_\disceltwo^{\gengenus-r}
%%	&
%%	\klinser{k}{\gengenus}{r}{\lattice} (f)\big)(\tau)
%%	=
%%	\sum_{\discel\in D_\lattice^{r}}
%%	\lim_{\lambda\to\infty}
%%	\Big(\klinser{k}{\gengenus}{r}{\lattice} \left( \big(\begin{smallmatrix}
%%	\tau & 0 \\
%%	0 & i\lambda I_{\gengenus-r}
%%	\end{smallmatrix}\big),f\right)\Big)_{(\discel,\disceltwo)}
%%	\mathfrak{e}_\discel
%%	\\
%%	&=
%	\sum_{\discel\in D_\lattice^{r}}
%	\lim_{\lambda\to\infty}
%	\sum_{\gamma\in\klimet{\gengenus}{r}\backslash \Mp_{2\gengenus}(\ZZ)}
%	\phi\left(\big(\begin{smallmatrix}
%	\tau & 0 \\
%	0 & i\lambda I_{\gengenus-r}
%	\end{smallmatrix}\big)\right)^{-2k}
%	\Big(
%	\weil{\lattice}{\gengenus}(\gamma)^{-1}
%	\big(f(\gamma^*\cdot\tau)
%	\otimes \zeroe{\gengenus-r}\big)
%	\Big)_{(\discel,\disceltwo)}
%	\mathfrak{e}_\discel.
%	\eas
	By Proposition~\ref{prop;convpropKEseries} the Klingen Eisenstein series converges uniformly on any vertical strip of positive height.
	We may then rewrite~$\big(\Phi_\disceltwo^{\gengenus-r} \klinser{k}{\gengenus}{r}{\lattice} (f)\big)(\tau)$ as
	\be\label{eq;underunifconv}
	\sum_{\discel\in D_L^{r}}\sum_{\gamma\in\klimet{\gengenus}{r}\backslash \Mp_{2\gengenus}(\ZZ)}
	\lim_{\lambda\to\infty}
	\phi\left(\big(\begin{smallmatrix}
	\tau & 0 \\
	0 & i\lambda I_{\gengenus-r}
	\end{smallmatrix}\big)\right)^{-2k}
	\Big(
	\weil{\lattice}{\gengenus}(\gamma)^{-1}
	\big(f(\gamma^*\cdot\tau)
	\otimes \zeroe{\gengenus-r}\big)
	\Big)_{(\discel,\disceltwo)}
	\mathfrak{e}_\discel.
	\ee
	We extract from~\eqref{eq;underunifconv} the terms associated to~$\gamma\in\klimet{\gengenus}{r}$.
	Recall that if~$\gamma=(M,\phi)$ with~$M\in\klipar{\gengenus}{r}$, then~$\phi(\tilde\tau)
	={\det D_4} ^{1/2} \phi^*(\tilde\tau^*)$ for every~$\tilde\tau\in\HH_\gengenus$, where~$D_4$ is the lower-right~$(\gengenus-r)\times (\gengenus-r)$ block of~$M$.
	This implies that the sum of the terms associated to~$\gamma\in\klimet{\gengenus}{r}$ extracted from~\eqref{eq;underunifconv} is
	\begin{align*}
	\sum_{\discel\in D_L^{r}}
	&
	\lim_{\lambda\to\infty}
	\phi\left(\big(\begin{smallmatrix}
	\tau & 0 \\
	0 & i\lambda I_{\gengenus-r}
	\end{smallmatrix}\big)\right)^{-2k}
	\Big(
	\weil{\lattice}{\gengenus}(\gamma)^{-1}
	\big(f(\gamma^*\cdot\tau)
	\otimes \zeroe{\gengenus-r}\big)
	\Big)_{(\discel,\disceltwo)}
	\mathfrak{e}_\discel
	\\
	&=
	\sum_{\discel\in D_L^{r}}
	\det D_4^{-k} \phi^*(\tau)^{-2k}
	\Big(
	\weil{\lattice}{\gengenus}(\gamma)^{-1}
	\phi^*(\tau)^{2k}\weil{r}{\lattice}(\gamma^*)f(\tau)
	\otimes \zeroe{\gengenus-r}\Big)_{(\discel,\disceltwo)}
	\mathfrak{e}_\discel
	\\
	&=
	\sum_{\discel\in D_L^{r}}
	\Big(
	f(\tau)
	\otimes \zeroe{\gengenus-r}\Big)_{(\discel,\disceltwo)}
	\mathfrak{e}_\discel
	=
	\begin{cases}
	f(\tau), & \text{if~$\disceltwo=0$,}\\
	0, & \text{otherwise,}
	\end{cases}
	\end{align*}
	where to deduce the first to the last equality above we used that
	\bes%\label{eq;todocheckweldefklingeis}
	\weil{\lattice}{\gengenus}(\gamma)
	(\mathfrak{e}_\discel \otimes\zeroe{\gengenus-r})
	=
	{\det D_4} ^{-k}\cdot
	\big(
	\weil{\lattice}{r}(\gamma^*)
	\mathfrak{e}_\discel
	\big)\otimes\zeroe{\gengenus-r}
	\qquad
	\text{for every~$\discel\in D_\lattice^{r}$, $\gamma\in\klimet{\gengenus}{r}$.}
	\ees
	As one can easily show following the idea of~\cite[p.~$68$, Proof of Proposition~$5$]{klingen}, the terms in~\eqref{eq;underunifconv} arising from~$\gamma\not\in\klimet{\gengenus}{r}$ are zero.
%	\textcolor{blue}{We now show that the terms in~\eqref{eq;underunifconv} arising from~$\gamma\not\in\klimet{\gengenus}{r}$ are zero.
%	To do so, we note that by Lemma~\ref{lemma;bound2blafromKling} their norm can be majorized by
%	\begin{align*}
%	&\sum_{\substack{\gamma\in\klimet{\gengenus}{r}\backslash\Mp_{2\gengenus}(\ZZ)
%	\\
%	\gamma\not\in\klimet{\gengenus}{r}}}
%	\lvert\phi\left(\big(\begin{smallmatrix}
%	\tau & 0 \\
%	0 & i\lambda I_{\gengenus-r}
%	\end{smallmatrix}\big)\right)\rvert^{-2k}
%	\lvert f(\gamma^*\cdot\tau)\rvert
%	\\
%	&\leq
%	\sum_{\substack{\gamma\in\klipar{\gengenus}{r}\backslash \Sp_{2\gengenus}(\ZZ)
%	\\
%	\gamma\not\in\klipar{\gengenus}{r}}}
%	\det \Im(\gamma^*\cdot\tau)^{-k/2}
%	\lvert
%	\det \big(C\big(\begin{smallmatrix}
%	\tau & 0 \\
%	0 & i\lambda I_{\gengenus-r}
%	\end{smallmatrix}\big)+D\big)
%	\rvert^{-k}.
%	\end{align*}
%	Hence, it is enough to show that if~$\gamma\not\in\klimet{\gengenus}{r}$, then
%	\[
%	\lim_{\lambda\to\infty}
%	\det\Im(\gamma^*\cdot\tau)^{-1}
%	\big\lvert
%	\big(C\big(\begin{smallmatrix}
%	\tau & 0 \\
%	0 & i\lambda I_{\gengenus-r}
%	\end{smallmatrix}\big)+D\big)
%	\big\rvert^{-2}
%	=
%	0
%	\qquad\text{for every~$\tau\in\HH_r$.}
%	\]
%	This statement does not depend on the choice of~$k$ nor of~$L$, and has been proved in e.g.~\cite[p.~$65$, Proof of Proposition~$5$]{klingen}.}
	\end{proof}
	\begin{lemma}\label{lemma;FexpofimPhiop}
	Let~$f\in M^k_{\gengenus,\lattice}$.
	If~$\disceltwo\in D_\lattice^{\gengenus-r}$ is isotropic, then
	\[
	\Phi_\disceltwo^{\gengenus-r} f(\tau)=
	\sum_{\discel\in D_\lattice^{r}}
	\sum_{\substack{T\in q(\discel) + \halfint{r} \\ T\ge 0}} 
	c_{\big(\begin{smallmatrix}
	T & 0\\
	0 & 0
	\end{smallmatrix}\big)}(f_{(\discel,\disceltwo)}) \cdot e\big(
	\trace(T\tau)
	\big)\mathfrak{e}_\discel,\qquad\text{$\tau\in\HH_{\gengenus-1}$.}
	\]
	In particular, a Siegel modular form in~$M^k_{\gengenus,\lattice}$ is a cusp form if and only if it lies in~$\ker\Phi_\disceltwo$ for every~$\disceltwo\in D_{\lattice}$.
	\end{lemma}
	\begin{proof}
	As in the scalar-valued case, the limit~$\lim_{\lambda\to\infty}
	f_{(\discel,\disceltwo)}\left(\left(\begin{smallmatrix}
	\tau & 0\\
	0 & i\lambda I_{\gengenus-r}
	\end{smallmatrix}\right)\right)$ extracts the part of the Fourier expansion of~$f_{(\discel,\disceltwo)}$ associated to singular matrices.
	This can be checked as follows.
	
	Without loss of generality we consider only the case of~$r=\gengenus-1$.
	Let~$\tau\in\HH_{\gengenus-1}$.
	The sequence of points~$\left(\begin{smallmatrix}
	\tau & 0\\
	0 & i\lambda
	\end{smallmatrix}\right)$ with~$\lambda$ large is contained in a vertical strip of positive height in~$\HH_{\gengenus}$, on which the Fourier series of~$f$ converges uniformly.
	We may then compute the limit term by term as
	\ba\label{eq;inproofeasyahuu}
	\Phi_\beta f(\tau)
%	&=
%	\sum_{\discel\in D_\lattice^{\gengenus-1}}
%	\lim_{\lambda\to\infty} f_{(\discel,\beta)}\left(\left(\begin{smallmatrix}
%	\tau & 0\\
%	0 & i\lambda
%	\end{smallmatrix}\right)\right)\mathfrak{e}_\discel
%	\\
	&=
	\sum_{\discel\in D_\lattice^{\gengenus-1}}
	\sum_{\substack{T\in q(\discel,\beta)+\halfint{\gengenus} \\ T\ge 0}}c_T(f_{(\discel,\beta)})
	e\big(\tr(T^*\tau)\big)
	\underbrace{\lim_{\lambda\to\infty} e\big(i\lambda T_*\big)}_{=\delta_{T_*,0}}\mathfrak{e}_\discel
%	\\
%	&=
%	\sum_{\discel\in D_\lattice^{\gengenus-1}}
%	\sum_{\substack{T\in q(\discel,\beta)+\halfint{\gengenus} \\ T\ge 0,\, T_*=0}}c_T(f_{(\discel,\beta)})
%	e\big(\tr(T^*\tau)\big)
%	\mathfrak{e}_\discel,
	\ea
	where we denoted by~$T^*$ and~$T_*$ respectively the top-left~$(\gengenus-1)\times(\gengenus-1)$-block and the bottom-right entry of~$T$.
	It is easy to see that if~$T\ge0$ and~$T_*=0$, then~$T=\left(\begin{smallmatrix}
	T^* & 0\\
	0 & 0
	\end{smallmatrix}\right)$.
	This and~\eqref{eq;inproofeasyahuu} give the Fourier expansion of~$\Phi_\disceltwo f$.
%	\textcolor{blue}{
%	\textbf{Here a quick check of the last assertion.}
%	\\
%	Suppose that the symmetric matrix~$T$ is such that~$T\ge0$ and of the form
%	\[
%	T=\left(\begin{smallmatrix}
%	T^* & b\\
%	b^t & 0
%	\end{smallmatrix}\right)
%	\]
%	for some vector~$b\in\QQ^{\gengenus-1}$.
%	Let~$b_j$ be the~$j$-th entry of~$b$.
%	Assume that~$b_j\neq 0$.
%	We denote by~$\sgn(b_j)$ its sign.
%	Let~$c\in\QQ_{>0}$ be such that
%	\[
%	c T_{j,j} - 2\sgn(b_j) b_j < 0.
%	\]
%	Let~$v\in\QQ^\gengenus\setminus\{0\}$ be the vector with~$j$-th entry~$c$,~$\gengenus$-th entry~$-\sgn(b_j)$, and zero everywhere else.
%	Since~$T$ is positive semidefinite, we must have
%	\[
%	v^t\cdot T\cdot v \ge 0.
%	\]
%	However, we may compute
%	\bas
%	v^t\cdot T\cdot v
%	&=
%	(0,\dots,0,c,0,\dots,0,-\sgn(b_j) )
%	\begin{pmatrix}
%	T^* & b\\
%	b^t & 0
%	\end{pmatrix}\cdot v
%	\\
%	&=
%	(cT_{j,1}-\sgn(b_j)b_1,\dots,cT_{j,\gengenus-1}-\sgn(b_j)b_{\gengenus-1},cb_j)\cdot v
%	\\
%	&=
%	c^2T_{j,j}-2c\sgn(b_j)b_j
%	\\
%	&=
%	c\big(	\underbrace{cT_{j,j}-2\sgn(b_j)b_j}_{<0}
%	\big)<0.
%	\eas
%	This is in contrast with the fact that~$T$ is positive semidefinite, hence it must be~$b_j=0$ for every~$j$.
%	}\\
%	This implies that~\eqref{eq;inproofeasyahuu} is the Fourier expansion of~$\Phi_\disceltwo f$.
	Clearly, it vanishes for all~$\beta$ if and only if~$f$ is a cusp form.
	\end{proof}
	
	
	\end{document}
\documentclass[../Lefschetz_decomposition.tex]{subfiles}

	\begin{document}
	We quickly review the Petersson inner product and its properties for vector-valued modular forms with respect to the Weil representation.
	\begin{defi}
	Let~$f_1,f_2\in M^k_{\genus,\lattice}$, at least one of which is a cusp form.
	The \emph{Petersson inner product} of~$f_1$ and~$f_2$ is
	\[
	\pet{f_1}{f_2}\coloneqq\int_{\Sp_{2\genus}(\ZZ)\backslash\HH_\genus}
	\langle f_1(\tau), f_2(\tau)\rangle
	\det y^k
	\frac{dx\, dy}{\det y^{\genus+1}}.
	\]
	\end{defi}
	The Petersson inner product is well-defined by Lemma~\ref{lemma;bound2blafromKling}, and makes~$S^k_{\genus,\lattice}$ a Hermitian space.
%	\\
%	
%	\textcolor{blue}{\textbf{Check that the integral above converges:}\\
%	As explained in~\cite[Part~$3$, p.~$191$, End of proof of Theorem~$1$]{1-2-3} we may majorize
%	\ba\label{eq;inbluecfjt}
%	\lvert f_2(\tau)\rvert
%	\leq
%	\sum_{\discel\in D_\lattice^\genus}
%	\sum_{\substack{T\in\halfint{\genus} \\ T\ge 0}}
%	\lvert c_T(f_\discel)\rvert \cdot \lvert e^{-2\pi\tr Ty} \rvert,
%	\ea
%	from which we see that~$\lvert f_2(\tau)\rvert$ is bounded on vertical strips of positive height in~$\HH_\genus$, in which we then have uniform convergence.
%	Suppose w.l.o.g.\ that~$f_1$ is a cusp form.
%	By Cauchy--Schwartz and Lemma~\ref{lemma;bound2blafromKling} we have
%	\[
%	\det y^k \lvert \langle f_1(\tau),f_2(\tau)\rangle \rvert
%	\leq
%	\det y^k \lvert f_1(\tau)\rvert \cdot\lvert f_2(\tau)\rvert
%	\leq
%	a_1 \det y^{k/2}\cdot e^{-a_2 \det y^{1/\genus}}\cdot\lvert f_2(\tau)\rvert
%	\]
%	for some positive constants~$a_1$ and~$a_2$.
%	We deduce from~\eqref{eq;inbluecfjt} that on vertical strips of positive heights we may uniformly bound~$\det y^k \lvert \langle f_1(\tau),f_2(\tau)\rangle \rvert$, up to a constant, with
%	\[
%	\det y^{k/2}\cdot e^{-a_2 \det y^{1/\genus}},
%	\]
%	which tends to zero if~$\det y\to\infty$.
%	\\
%	We have just proved that~$\det y^k \lvert \langle f_1(\tau),f_2(\tau)\rangle \rvert$ is bounded on every vertical strip of positive height.
%	Since the Siegel fundamental domain is contained in a vertical strip of positive height by~\cite[p.~$30$, Lemma~$2$]{klingen}, we deduce that the integral given by~$\pet{f_1}{f_2}$ converges.
%	}
	
	As in the classical setting, Klingen Eisenstein series are orthogonal to cusp forms with respect to~$\petempty$.
	This may be proved as in~\cite[Proof of Proposition~$8$, Part~(ii)]{klingen}.
	\begin{prop}\label{prop;ortcusp&kling}
	Let~$0\leq r <\genus$ and let~$f\in S^k_{r,\lattice}$.
	Then~$\klinser{k}{\gengenus}{r}{\lattice}(f)$ is orthogonal to~$S^k_{\genus,\lattice}$ with respect to the Petersson inner product.
	\end{prop}
%	\begin{proof}
%	\textcolor{blue}{Let~$h\in S^k_{\genus,\lattice}$.
%	We denote by~$\mathcal{F}_\genus$ the Siegel fundamental domain and compute
%	\bas
%	&\pet{\klinser{k}{\gengenus}{r}{\lattice}(f)}{h}
%	\\
%	&=
%	\int_{\mathcal{F}_\genus}
%	\det y^k
%	\left\langle\klinser{k}{\gengenus}{r}{\lattice}(\tau,f) , h(\tau)
%	\right\rangle
%	\frac{dx\,dy}{\det y^{\genus+1}}
%	\\
%	&=
%	\int_{\mathcal{F}_\genus}
%	\det y^k
%	\sum_{\gamma\in\klimet{\gengenus}{r}\backslash \Mp_{2\gengenus}(\ZZ)}
%	\phi(\tau)^{-2k}
%	\Big\langle
%	\weil{\lattice}{\gengenus}(\gamma)^{-1}
%	\Big(f(\gamma^*\cdot\tau^*)
%	\otimes \zeroe{\gengenus-r}\Big)
%	, h(\tau)
%	\Big\rangle
%	\frac{dx\,dy}{\det y^{\genus+1}}\\
%	&=
%	\int_{\mathcal{F}_\genus}
%	\det y^k
%	\sum_{\gamma\in\klimet{\gengenus}{r}\backslash \Mp_{2\gengenus}(\ZZ)}
%	\lvert\phi(\tau)\rvert^{-4k}
%	\Big\langle
%	\Big(f(\gamma^*\cdot\tau^*)
%	\otimes \zeroe{\gengenus-r}\Big)
%	, h(\gamma\cdot\tau)
%	\Big\rangle
%	\frac{dx\,dy}{\det y^{\genus+1}}
%	\\
%	&=
%	\int_{\mathcal{F}_\genus}
%	\det \Im (\gamma\cdot\tau)^k
%%	\cdot\lvert\det (C\tau+D)\rvert^{2k}
%	\sum_{\gamma\in\klimet{\gengenus}{r}\backslash \Mp_{2\gengenus}(\ZZ)}
%%	\lvert\phi(\tau)\rvert^{-4k}
%	\Big\langle
%	\Big(f(\gamma^*\cdot\tau^*)
%	\otimes \zeroe{\gengenus-r}\Big)
%	, h(\gamma\cdot\tau)
%	\Big\rangle
%	\frac{dx\,dy}{\det y^{\genus+1}}.
%	\eas
%	We recall that Klingen Eisenstein series converge uniformly on vertical strips of positive height by Proposition~\ref{prop;convpropKEseries}, that the Siegel fundamental domain~$\mathcal{F}_\genus$ is contained in a vertical strip of positive height by~\cite[p.~$30$, Lemma~$2$]{klingen}, and that~$\det y^k \lvert h(\tau) \rvert$ is bounded by Lemma~\ref{lemma;bound2blafromKling}, we are allowed to switch summation with integration and then unfold the integral as
%	\bas
%	\pet{\klinser{k}{\gengenus}{r}{\lattice}(f)}{h}
%	&=
%	\sum_{\gamma\in\klimet{\gengenus}{r}\backslash \Mp_{2\gengenus}(\ZZ)}
%	\int_{\gamma(\mathcal{F}_\genus)}
%	\det \Im (\tau)^k
%	\Big\langle
%	\Big(f(\tau^*)
%	\otimes \zeroe{\gengenus-r}\Big)
%	, h(\tau)
%	\Big\rangle
%	\frac{dx\,dy}{\det y^{\genus+1}}
%	\\
%	&=
%	2^{-\varepsilon_r}\int_{\mathcal{Q}}
%	\det \Im (\tau)^k
%	\Big\langle
%	\Big(f(\tau^*)
%	\otimes \zeroe{\gengenus-r}\Big)
%	, h(\tau)
%	\Big\rangle
%	\frac{dx\,dy}{\det y^{\genus+1}},
%	\eas
%	where~$\mathcal{Q}$ is the standard fundamental domain of~$\klipar{\gengenus}{r}$ in~$\HH_\genus$ and~$\varepsilon=\begin{cases}
%	0, & r=0,
%	\\ 
%	1, & r=1,\dots,\genus-1.
%	\end{cases}$
%	As for the scalar-valued case, the factor~$2^{-\varepsilon_r}$ is caused by the fact that~$\mathcal{Q}$ maps to itself by the matrix~$\gamma=\big(\begin{smallmatrix}
%	A & \\
%	 & A^{-t}
%	\end{smallmatrix}\big)\in\klipar{\gengenus}{r}$ with~$A=\big(\begin{smallmatrix}
%	I_r & \\
%	 & \pm I_{\genus-r}
%	\end{smallmatrix}\big)$.
%	Such a matrix does not fix~$\tau\in\HH_\genus$ in general, in fact
%	\[
%	\big(\begin{smallmatrix}
%	A & \\
%	 & A^{-t}
%	\end{smallmatrix}\big)
%	\colon \mathcal{Q}\longrightarrow\mathcal{Q},
%	\qquad
%	\tau=\big(\begin{smallmatrix}
%	\tau_1 & \tau_2\\
%	\tau_2^t & \tau_3
%	\end{smallmatrix}\big)\longmapsto
%	\big(\begin{smallmatrix}
%	\tau_1 & -\tau_2\\
%	-\tau_2^t & \tau_3
%	\end{smallmatrix}\big).
%	\]
%	(In particular, $\mathcal{Q}$ is not really a fundamental domain.)}
%	
%	\textcolor{blue}{We recall from~\cite[p.~$66$, $(15)$]{klingen} that
%	\[
%	\mathcal{Q}
%	=
%	\{
%	\tau\in\HH_\genus : \text{$\tau^*\in\mathcal{F}_r$, $x$ and~$y_3$ are reduced mod~$1$,~$y_2\in\mink{\genus-r}$}
%	\},
%	\]
%	where~$y=\big(\begin{smallmatrix}
%	y_1 & 0 \\
%	0 & y_2
%	\end{smallmatrix}\big)\big[\big(
%	\begin{smallmatrix}
%	I_\genus & y_3 \\
%	0 & I_{\genus}
%	\end{smallmatrix}
%	\big)\big]$.
%	We now apply Fubini--Tonelli to the integral over~$\mathcal{Q}$.
%	The integration over the lower-right~$(\genus-r)\times(\genus-r)$ block~$x_4$ of~$x$ leads to
%	\bes
%	\int_{\text{$x_4$ mod~$1$}}
%	\langle \zeroe{\gengenus-r} , h(\tau) \rangle dx_4.
%	\ees
%	This integral vanishes since~$h$ is a cusp form, as one can see by replacing~$h$ with its Fourier expansion.}
%	\end{proof}
	\end{document}
	
	
	
	
	
	
	
	
	
	
	
	
	
	
	
	
	
	
	
	
	
	
	
	
	
	
	
	
	
	
	
	
	
	
	
	
	
	
	
	
	
	
	
	
	
\documentclass[../Lefschetz_decomposition.tex]{subfiles}

	\begin{document}
	In this section we provide vector-valued generalizations of Poincaré series associated to positive definite matrices.
	Further information regarding the scalar-valued case can be found in~\cite[Section~$75$]{klingen}.

%	OLD VERSION - NON WELL-DEFINED!
%	We denote by~$\trans{\gengenus}$ and~$\mettrans{\gengenus}$ respectively the subgroup of translation
%	\[
%	\trans{\gengenus}
%	\coloneqq
%	\left\{
%	\left(\begin{smallmatrix}
%	\pm I_\gengenus & B \\
%	0 & \pm I_\gengenus
%	\end{smallmatrix}\right)
%	: B\in\Sym_{\gengenus}(\ZZ)
%	\right\}
%	\]
%	in~$\Sp_{2\gengenus}(\ZZ)$ and its metaplectic cover.
%	\begin{defi}
%	Let~$\discel\in D_\lattice^\gengenus$ and let~$T\in q(\discel)+\halfint{\gengenus}$ be a \emph{positive definite} matrix.
%	The Poincaré series~$\poincv{\gengenus}{\discel}{T}{\lattice}$ of index~$(\discel,T)$ is defined as
%	\[
%	\poincv{\gengenus}{\discel}{T}{\lattice}(\tau)
%	=
%	\sum_{\gamma\in\mettrans{\gengenus}\backslash\Mp_{2\gengenus}(\ZZ)} e\big(\tr(T\gamma\cdot\tau)\big) \cdot \phi(\tau)^{-2k} \weil{\lattice}{\gengenus}(\gamma)^{-1}\mathfrak{e}_\discel,
%	\qquad\tau\in\HH_\gengenus.
%	\]
%	\end{defi}
%\textcolor{blue}{\textbf{Here a quick check that is formally~$\poincv{\gengenus}{\discel}{T}{\lattice}$ well-defined.
%	We could drop it before publishing.}\\
%	Let~${\gamma_1=(M_1,\pm\phi_1)\in\metsieg{\gengenus}}$ for some~$M_1=\big(\begin{smallmatrix}
%	\pm I_\gengenus & B\\
%	0 & \pm I_\gengenus
%	\end{smallmatrix}\big)$, where~$\phi_1(\tau)=\sqrt{\det (\pm I_\gengenus)}$.
%	Let
%	\[
%	F(\gamma)\coloneqq e\big(\tr T\gamma\cdot\tau\big)\cdot\phi(\tau)^{-2k}\weil{\lattice}{\gengenus}(\gamma)^{-1}\mathfrak{e}_\discel\qquad \text{for every~$\gamma=(M,\phi)\in\Mp_{2\gengenus}(\ZZ)$}.
%	\]
%	Since~$\gamma_1\cdot\gamma=\big(M_1M , \pm\phi_1(M\cdot\tau)\phi(\tau)\big)$, we deduce that
%	\bas
%	F&(\gamma_1\cdot\gamma)
%	\\
%	&=
%	e\big(
%	\tr T \gamma\cdot\tau
%	\big)\cdot e(\pm\tr TB) \cdot
%	\big(\pm\phi_1(M\cdot\tau)\big)^{-2k}\phi(\tau)^{-2k}\weil{\lattice}{\gengenus}(\gamma)^{-1}\weil{\lattice}{\gengenus}(\gamma_1)^{-1}\mathfrak{e}_\discel
%	\\
%	&=
%	e(\tr T\gamma\cdot\tau)\cdot e(\pm\tr TB)\cdot (\pm 1)^{-2k}\det A ^{-k}\phi(\tau)^{-2k}
%	(\pm1)^{2-n}\weil{\lattice}{\gengenus}(\gamma)^{-1}
%	\\
%	&\quad\times
%	\weil{\lattice}{\gengenus}\left(\left(\begin{smallmatrix}
%	\pm I_\gengenus & \\
%	 & \pm I_\gengenus
%	\end{smallmatrix}\right), \sqrt{\det(\pm I_\gengenus)}\right)^{-1}
%	\weil{\lattice}{\gengenus}\left(\left(\begin{smallmatrix}
%	I_\gengenus & \pm B\\
%	 & I_\gengenus
%	\end{smallmatrix}\right), 1 \right)^{-1}
%	\mathfrak{e}_\discel
%	\\
%	&=
%	e(\tr T\gamma\cdot\tau)\cdot e(\pm\tr TB)\cdot (\pm 1)^{-2k}\det A ^{-k}\phi(\tau)^{-2k}
%	(\pm1)^{2-n}\weil{\lattice}{\gengenus}(\gamma)^{-1}
%	\\
%	&\quad\times
%	\eas
%	where in the first to the last equality we used that~$k=1+n/2$.}
%	
%	\textcolor{red}{To do: Show that these Poincaré series are well defined. Show they are mod forms. Show they are cusp forms. Prove they converge uniformly.}
%	
%	\begin{lemma}
%	Let~$\discel\in D_\lattice^\gengenus$ and let~$T\in\halfint{\gengenus}-q(\discel)$ be \emph{positive definite}.
%	Let~$f\in S^k_{\gengenus,\lattice}$.
%	The Petersson inner product of~$f$ and the Poincaré series~$\poincv{\gengenus}{\discel}{T}{\lattice}$ is
%	\bas
%	\pet{f}{\poincv{\gengenus}{\discel}{T}{\lattice}}= \constant_{k,\gengenus}\cdot \det T^{(\gengenus+1)/2 - k} \cdot c_T(f_\discel),
%	\eas
%	where~$\constant_{k,\gengenus}$ is the factor depending only on~$k$ and~$\gengenus$ given by
%	\[
%	\constant_{k,\gengenus}
%	=
%	\pi^{\gengenus(\gengenus-1)/4}(4\pi)^{\gengenus(\gengenus+1)/2 - \gengenus k}
%	\prod_{\ell=1}^\gengenus \Gamma\Big(k - \frac{\gengenus + \ell}{2}\Big).
%	\]
%	\textcolor{red}{Clarify hypothesis for convergence of Poincaré series.}
%	\end{lemma}
%	\begin{proof}
%	Since~$f(\tau)=\phi(\tau)^{-2k}\weil{\lattice}{\gengenus}(\gamma)^{-1} f(\gamma\cdot\tau)$ for every~$\gamma=(M,\phi)\in\Mp_{2\gengenus}(\ZZ)$, we may compute 
%	\bas
%	\pet{f}{\poincv{\gengenus}{\discel}{T}{\lattice}}
%	&=
%	\int_{\Sp_{2\gengenus}\backslash\HH_\gengenus}
%	\det y^k
%\langle f(\tau),\poincv{\gengenus}{\discel}{T}{\lattice}(\tau) \rangle\,\frac{dx\,dy}{\det y^{\gengenus+1}}
%\\
%&=
%	\int_{\Sp_{2\gengenus}(\ZZ)\backslash\HH_\gengenus}
%	\det y^k\Big\langle
%	\phi(\tau)^{-2k}\weil{\lattice}{\gengenus}(\gamma)^{-1} f(\gamma\cdot\tau)
%	,
%	\\
%	&\quad
%	\sum_{\gamma\in\mettrans{\gengenus}\backslash\Mp_{2\gengenus}(\ZZ)}e\big(\tr T\gamma\cdot\tau\big)
%	\phi(\tau)^{-2k}\weil{\lattice}{\gengenus}(\gamma)^{-1}\mathfrak{e}_\discel
%	\Big\rangle
%	\,\frac{dx\,dy}{\det y^{\gengenus+1}}
%	\eas
%	Since the Hermitian inner product~$\langle\cdot{,}\cdot\rangle$ on~$\CC[D_\lattice^\gengenus]$ is unitary with respect to the Weil representation, and since
%	\[
%	\det y = \det\big(\Im(\gamma\cdot\tau)\big)\cdot \lvert\det(C\tau+D)\rvert^2,
%	\]
%	we may proceed with the computation as
%	\bas
%	\pet{f}{\poincv{\gengenus}{\discel}{T}{\lattice}}
%	&=
%	\int_{\Sp_{2\gengenus}(\ZZ)\backslash\HH_\gengenus}
%	\sum_{\gamma\in\mettrans{\gengenus}\backslash\Mp_{2\gengenus}(\ZZ)}
%	\det\big(\Im(\gamma\cdot\tau)\big)^k
%	\frac{\lvert \det(C\tau+D)\rvert^{2k}}{\lvert\phi(\tau)\rvert^{4k}}
%	\overline{e\big(\tr T\gamma\cdot\tau\big)}
%	\\
%	&\quad\times
%	\langle f(\gamma\cdot\tau),\mathfrak{e}_\discel\rangle
%	\,\frac{dx\,dy}{\det y^{\gengenus+1}}.
%	\eas
%	Since the convergence of the sum over~$\gamma\in\mettrans{\gengenus}\backslash\Mp_{2\gengenus}(\ZZ)$ is uniform, we may move it outside the integral.
%	Let~$\mathcal{F}_\gengenus$ be the Siegel fundamental domain in~$\HH_\gengenus$ with respect to~$\Sp_{2\gengenus}(\ZZ)$.
%	Since~$|\phi(\tau)|^{4k}=|\det(C\tau+D)|^{2k}$, and since
%	\[
%	\langle f(\tau),\mathfrak{e}_\discel\rangle
%	=
%	\sum_{\disceltwo\in D_\lattice^\gengenus}f_\disceltwo(\tau)\langle\mathfrak{e}_\disceltwo,\mathfrak{e}_\discel\rangle
%	=
%	f_\discel(\tau)\qquad\text{for every~$\tau\in\HH_g$},
%	\]
%	we may apply the unfolding trick and deduce that
%	\bas
%	\pet{f}{\poincv{\gengenus}{\discel}{T}{\lattice}}
%	&=
%	\sum_{\gamma\in\mettrans{\gengenus}\backslash\Mp_{2\gengenus}(\ZZ)}
%	\int_{\mathcal{F}_\gengenus}\det\big(\Im(\gamma\cdot\tau)\big)^k \cdot
%	\overline{e\big(\tr T\gamma\cdot\tau\big)} \cdot
%	f_\discel(\gamma\cdot\tau)
%	\,\frac{dx\,dy}{\det y^{\gengenus+1}}
%	\\
%	&=
%	\sum_{\gamma\in\mettrans{\gengenus}\backslash\Mp_{2\gengenus}(\ZZ)}
%	\int_{\gamma(\mathcal{F}_\gengenus)}
%	\det y^k\cdot e(-\tr T\overline\tau)\cdot f_\discel(\tau)
%	\,\frac{dx\,dy}{\det y^{\gengenus+1}}
%	\\
%	&=
%	\int_{\trans{\gengenus}\backslash\HH_\gengenus}
%	\det y^k\cdot e(-\tr T\overline\tau)\cdot f_\discel(\tau)
%	\,\frac{dx\,dy}{\det y^{\gengenus+1}}.
%	\eas
%	Let~$\cone{\gengenus}$ be the convex cone of the imaginary parts of matrices in~$\HH_{\gengenus}$; see~\eqref{eq;convconeimpart}.
%	A fundamental domain in~$\HH_\gengenus$ with respect to~$\trans{\gengenus}$ is given by
%	\[
%	\left\{
%	\tau=x+iy\in\HH_\gengenus
%	:
%	\text{$x$ is reduced modulo~$1$ and~$y\in\cone{\gengenus}$}
%	\right\}.
%	\]
%	We may then replace~$f_\discel$ with its Fourier expansion and compute
%	\ba\label{eq;TheLastEqgenJanpoinpet}
%	\pet{f}{\poincv{\gengenus}{\discel}{T}{\lattice}}
%	&=
%	\int_{y\in\cone{\gengenus}}
%	\det y^{k-\gengenus-1}
%	\int_{\text{$x$ mod $1$}}
%	e\big(-\tr T \overline\tau\big)
%	f_\discel(\tau)\,dx\,dy
%	\\
%	&=
%	\int_{y\in\cone{\gengenus}}
%	\det y^{k-\gengenus-1}
%	\exp(-2\pi\tr Ty)
%	\sum_{\substack{M\in q(\discel)+\halfint{\gengenus}\\ M>0}}
%	c_M(f_\discel)
%	\\
%	&\quad\times\exp(-2\pi\tr My)\underbrace{\int_{\text{$x$ mod $1$}}e\big(\tr(M-T)x\big)\,dx}_{=\delta_{M,T}}\,dy
%	\\
%	&=
%	c_T(f_\discel)\int_{y\in\cone{\gengenus}}
%	\det y^{k-\gengenus-1}
%	\exp(-4\pi\tr Ty)
%	\,dy.
%	\ea
%	As remarked in~\cite[p.~$90$,~(iv)]{klingen}, as well as in the last equality of the proof of~\cite[Proposition~$1.5$]{bruinierbook} for the case of genus~$g=1$, the integral over~$\trans{\gengenus}$ appearing on the right-hand side of the last equality in~\eqref{eq;TheLastEqgenJanpoinpet} equals~$\constant_{k,g}\det T^{(\gengenus+1)/2-k}$.
%	\end{proof}
	
	We denote by~$\trans{\gengenus}$ and~$\mettrans{\gengenus}$ respectively the subgroup
	\be\label{eq;translinsympl}
	\trans{\gengenus}
	\coloneqq
	\left\{
	\left(\begin{smallmatrix}
	I_\gengenus & B \\
	0 & I_\gengenus
	\end{smallmatrix}\right)
	: B\in\Sym_{\gengenus}(\ZZ)
	\right\}
	\ee
	of translations in~$\Sp_{2\gengenus}(\ZZ)$ and its metaplectic cover.
%	Recall that the signature of~$\lattice$ is~$(n,2)$ and the weight~$k\in\frac{1}{2}\ZZ$ is chosen such that~$2k\equiv n-2$ mod~$4$.
	\begin{defi}
	Let~$\discel\in D_\lattice^\gengenus$ and let~$T\in q(\discel)+\halfint{\gengenus}$ be a \emph{positive definite} matrix.
	Assume that~$k>2\genus$.
	The weight~$k$ Poincaré series~$\poincv{\gengenus}{\discel}{T}{\lattice}$ is defined as
	\[
	\poincv{\gengenus}{\discel}{T}{\lattice}(\tau)
	=
	\frac{1}{2}\sum_{\gamma\in\mettrans{\gengenus}\backslash\Mp_{2\gengenus}(\ZZ)} e\big(\tr(T\gamma\cdot\tau)\big) \cdot \phi(\tau)^{-2k} \weil{\lattice}{\gengenus}(\gamma)^{-1}\mathfrak{e}_\discel,
	\qquad\tau\in\HH_\gengenus.
	\]
	\end{defi}
	The Weil representation~$\weil{\lattice}{\genus}$ is trivial on the principal subgroup of level~$N$ of~$\Sp_\genus(\ZZ)$, where~$N$ is the level of the lattice~$L$.
	Hence, the vector-valued Poincare series~$\poincv{\gengenus}{\discel}{T}{\lattice}$ is a finite combination of scalar valued Poincare series for that principal congruence subgroup.
	The absolute and local uniform convergence of~$\poincv{\gengenus}{\discel}{T}{\lattice}$ is therefore implied by the convergence of scalar-valued Poincaré series; see e.g.~\cite{christian}.
	With an analogous argument, we may deduce that~$\poincv{\gengenus}{\discel}{T}{\lattice}$ is a cusp form.
	We also remark that~$\poincv{\gengenus}{\discel}{T}{\lattice}=\poincv{\gengenus}{-\discel}{T}{\lattice}$.
%	\textcolor{blue}{\textbf{Quick check that~$\poincv{\gengenus}{\discel}{T}{\lattice}$ is formally well-defined:}\\
%	Let~${\gamma_1=(M_1,\pm1)\in\mettrans{\gengenus}}$ for some~$M_1=\big(\begin{smallmatrix}
%	I_\gengenus & B\\
%	0 & I_\gengenus
%	\end{smallmatrix}\big)$.
%	Let
%	\[
%	F(\gamma)\coloneqq e\big(\tr T\gamma\cdot\tau\big)\cdot\phi(\tau)^{-2k}\weil{\lattice}{\gengenus}(\gamma)^{-1}\mathfrak{e}_\discel,\qquad \gamma=(M,\phi)\in\Mp_{2\gengenus}(\ZZ).
%	\]
%	Since~$\gamma_1\cdot\gamma=\big(M_1M , \pm\phi(\tau)\big)$, we deduce that
%	\bas
%	F&(\gamma_1\cdot\gamma)
%	\\
%	&=
%	e\big(
%	\tr T \gamma\cdot\tau
%	\big)\cdot e(\tr TB) \cdot
%	(\pm1)^{-2k}\phi(\tau)^{-2k}\weil{\lattice}{\gengenus}(\gamma)^{-1}\weil{\lattice}{\gengenus}(\gamma_1)^{-1}\mathfrak{e}_\discel
%	\\
%	&=
%	e(\tr T\gamma\cdot\tau)\cdot e(\tr TB)\cdot (\pm 1)^{-2k}\phi(\tau)^{-2k}
%	(\pm1)^{n-2}\weil{\lattice}{\gengenus}(\gamma)^{-1}
%%	\\
%%	&\quad\times
%	\weil{\lattice}{\gengenus}\left(\left(\begin{smallmatrix}
%	I_\gengenus & B\\
%	 & I_\gengenus
%	\end{smallmatrix}\right), 1 \right)^{-1}
%	\mathfrak{e}_\discel
%	\\
%	&=
%	e(\tr T\gamma\cdot\tau)\cdot e(\tr TB)\cdot \phi(\tau)^{-2k}
%	e(-\tr q(\discel)B)\cdot\weil{\lattice}{\gengenus}(\gamma)^{-1}\mathfrak{e}_\discel,
%	\eas
%	where in the first to the last equality we used that~$2k\equiv n-2$ mod~$4$.
%	Since~$e(\tr TB)\cdot e(-\tr q(\discel)B) = 1$, we deduce that~$F(\gamma_1\cdot\gamma)=F(\gamma)$ for every~$\gamma\in\Mp_{2\gengenus}(\ZZ)$.}
%	
%	\textcolor{blue}{\textbf{Check that~$\poincv{\gengenus}{\discel}{T}{\lattice}$ behaves as a vector-valued modular form with respect to~$\weil{\lattice}{\gengenus}$:} Easy, same as in genus~$1$.}
%	
%	\textcolor{blue}{\textbf{Check that~$\poincv{\gengenus}{\discel}{T}{\lattice}\in S^k_{\genus,\lattice}$:} The fact that~$\poincv{\gengenus}{\discel}{T}{\lattice}$ is holomorphic follows from the local uniform convergence, since the summands of the series are holomorphic functions.
%	The fact that the Poincaré series behaves as we want w.r.t.~$\weil{\lattice}{\genus}$ has been already checked above, hence~$\poincv{\gengenus}{\discel}{T}{\lattice}\in M^k_{\genus,\lattice}$.
%	To show that they are cusp forms: it follows from the fact that their scalar-valued analogs with level are known to be cusp forms~\cite{christian}.}
%	
%	\textcolor{blue}{\textbf{Check that~$\poincv{\gengenus}{\discel}{T}{\lattice}=\poincv{\gengenus}{-\discel}{T}{\lattice}$:} Let~$\trans{\genus}'$ be the subgroup generated by~$\trans{\genus}$ and~$-I_{2\genus}$.
%	We denote by~$\mettrans{\genus}'$ its metaplectic cover.
%	It is easy to check that
%	\[
%	\poincv{\gengenus}{\discel}{T}{\lattice}(\tau)
%	=
%	\frac{1}{2}
%	\sum_{\gamma\in\mettrans{\genus}'\backslash\Mp_{2\genus}(\ZZ)}
%	e(\tr T\gamma\cdot\tau)
%	\phi(\tau)^{-2k}
%	\weil{\lattice}{\genus}(\gamma)^{-1}
%	\big(
%	\mathfrak{e}_\discel
%	+
%	\mathfrak{e}_{-\discel}
%	\big).
%	\]
%	}
	\begin{lemma}\label{lemma;petSiegpoinandcusp}
	Let~$\discel\in D_\lattice^\gengenus$, let~$T\in q(\discel) + \halfint{\gengenus}$ be \emph{positive definite}, and let~$f\in S^k_{\gengenus,\lattice}$.
	If~$k>2\genus$, then
	\bas
	\bigpet{f}{\poincv{\gengenus}{\discel}{T}{\lattice}}= \constant_{k,\gengenus}\cdot \det T^{(\gengenus+1)/2 - k} \cdot c_T(f_\discel),
	\eas
	where~$\constant_{k,\gengenus}$ is the factor depending only on~$k$ and~$\gengenus$ given by
	\[
	\constant_{k,\gengenus}
	=
	\pi^{\gengenus(\gengenus-1)/4}(4\pi)^{\gengenus(\gengenus+1)/2 - \gengenus k}
	\prod_{\ell=1}^\gengenus \Gamma\Big(k - \frac{\gengenus + \ell}{2}\Big).
	\]
	\end{lemma}
	\begin{proof}
	Since the idea is as for the classical case, we only give a brief overview of the proof.
	We may rewrite~$\pet{f}{\poincv{\gengenus}{\discel}{T}{\lattice}}$ as
	\bas
	\frac{1}{2}
	\int_{\Sp_{2\gengenus}(\ZZ)\backslash\HH_\gengenus}
	\sum_{\gamma\in\mettrans{\gengenus}\backslash\Mp_{2\gengenus}(\ZZ)}
	\det\big(\Im(\gamma\cdot\tau)\big)^k
	\frac{\lvert \det(C\tau+D)\rvert^{2k}}{\lvert\phi(\tau)\rvert^{4k}}
	\overline{e(\tr T\gamma\cdot\tau)}
	\langle f(\gamma\cdot\tau),\mathfrak{e}_\discel\rangle
	\,\frac{dx\,dy}{\det y^{\gengenus+1}}.
	\eas
	By Lemma~\ref{lemma;bound2blafromKling}
%	 the convergence of the sum over~$\gamma\in\mettrans{\gengenus}\backslash\Mp_{2\gengenus}(\ZZ)$ is uniform, hence
	 we may move the sum outside the integral and apply the unfolding trick to deduce that
	\bas
	\pet{f}{\poincv{\gengenus}{\discel}{T}{\lattice}}
	=
	\int_{\trans{\gengenus}\backslash\HH_\gengenus}
	\det y^k\cdot e(-\tr T\overline\tau)\cdot f_\discel(\tau)
	\,\frac{dx\,dy}{\det y^{\gengenus+1}}.
	\eas
	Let~$\cone{\gengenus}$ be the convex cone of the imaginary parts of matrices in~$\HH_{\gengenus}$; see~\cite[Section~$2$, p.~$10$]{klingen}.
	A fundamental domain in~$\HH_\gengenus$ with respect to~$\trans{\gengenus}$ is given by
	\[
	\left\{
	\tau=x+iy\in\HH_\gengenus
	:
	\text{$x$ is reduced modulo~$1$ and~$y\in\cone{\gengenus}$}
	\right\}.
	\]
	It is easy to replace~$f_\discel$ with its Fourier expansion and compute
	\ba\label{eq;TheLastEqgenJanpoinpet}
	\pet{f}{\poincv{\gengenus}{\discel}{T}{\lattice}}
%	&=
%	\int_{y\in\cone{\gengenus}}
%	\det y^{k-\gengenus-1}
%	\int_{\text{$x$ mod $1$}}
%	e\big(-\tr T \overline\tau\big)
%	f_\discel(\tau)\,dx\,dy
%	\\
%	&=
%	\int_{y\in\cone{\gengenus}}
%	\det y^{k-\gengenus-1}
%	\exp(-2\pi\tr Ty)
%	\sum_{\substack{M\in q(\discel)+\halfint{\gengenus}\\ M>0}}
%	c_M(f_\discel)
%	\\
%	&\quad\times\exp(-2\pi\tr My)\underbrace{\int_{\text{$x$ mod $1$}}e\big(\tr(M-T)x\big)\,dx}_{=\delta_{M,T}}\,dy
%	\\
	&=
	c_T(f_\discel)\int_{y\in\cone{\gengenus}}
	\det y^{k-\gengenus-1}
	\exp(-4\pi\tr Ty)
	\,dy.
	\ea
	As remarked in~\cite[p.~$90$,~(iv)]{klingen}, the integral over~$\cone{\gengenus}$ appearing on the right-hand side of~\eqref{eq;TheLastEqgenJanpoinpet} equals~$\constant_{k,g}\det T^{(\gengenus+1)/2-k}$.
%	\textcolor{blue}{\textbf{Here a more detailed proof:}\\
%	Since~$f(\tau)=\phi(\tau)^{-2k}\weil{\lattice}{\gengenus}(\gamma)^{-1} f(\gamma\cdot\tau)$ for every~$\gamma=(M,\phi)\in\Mp_{2\gengenus}(\ZZ)$, we may compute 
%	\bas
%	\pet{f}{\poincv{\gengenus}{\discel}{T}{\lattice}}
%	&=
%	\int_{\Sp_{2\gengenus}\backslash\HH_\gengenus}
%	\det y^k
%\langle f(\tau),\poincv{\gengenus}{\discel}{T}{\lattice}(\tau) \rangle\,\frac{dx\,dy}{\det y^{\gengenus+1}}
%\\
%&=\frac{1}{2}
%	\int_{\Sp_{2\gengenus}(\ZZ)\backslash\HH_\gengenus}
%	\det y^k\Big\langle
%	\phi(\tau)^{-2k}\weil{\lattice}{\gengenus}(\gamma)^{-1} f(\gamma\cdot\tau)
%	,
%	\\
%	&\quad
%	\sum_{\gamma\in\mettrans{\gengenus}\backslash\Mp_{2\gengenus}(\ZZ)}e\big(\tr T\gamma\cdot\tau\big)
%	\phi(\tau)^{-2k}\weil{\lattice}{\gengenus}(\gamma)^{-1}\mathfrak{e}_\discel
%	\Big\rangle
%	\,\frac{dx\,dy}{\det y^{\gengenus+1}}
%	\eas
%	Since the Hermitian inner product~$\langle\cdot{,}\cdot\rangle$ on~$\CC[D_\lattice^\gengenus]$ is unitary with respect to the Weil representation, and since
%	\[
%	\det y = \det\big(\Im(\gamma\cdot\tau)\big)\cdot \lvert\det(C\tau+D)\rvert^2,
%	\]
%	we may proceed with the computation as
%	\bas
%	\pet{f}{\poincv{\gengenus}{\discel}{T}{\lattice}}
%	&=\frac{1}{2}
%	\int_{\Sp_{2\gengenus}(\ZZ)\backslash\HH_\gengenus}
%	\sum_{\gamma\in\mettrans{\gengenus}\backslash\Mp_{2\gengenus}(\ZZ)}
%	\det\big(\Im(\gamma\cdot\tau)\big)^k
%	\frac{\lvert \det(C\tau+D)\rvert^{2k}}{\lvert\phi(\tau)\rvert^{4k}}
%	\overline{e\big(\tr T\gamma\cdot\tau\big)}
%	\\
%	&\quad\times
%	\langle f(\gamma\cdot\tau),\mathfrak{e}_\discel\rangle
%	\,\frac{dx\,dy}{\det y^{\gengenus+1}}.
%	\eas
%	By Lemma~\ref{lemma;bound2blafromKling} the convergence of the sum over~$\gamma\in\mettrans{\gengenus}\backslash\Mp_{2\gengenus}(\ZZ)$ is uniform, hence we may move the sum outside the integral.
%	Let~$\mathcal{F}_\gengenus$ be the Siegel fundamental domain in~$\HH_\gengenus$ with respect to~$\Sp_{2\gengenus}(\ZZ)$.
%	Since~$|\phi(\tau)|^{4k}=|\det(C\tau+D)|^{2k}$, and since
%	\[
%	\langle f(\tau),\mathfrak{e}_\discel\rangle
%	=
%	\sum_{\disceltwo\in D_\lattice^\gengenus}f_\disceltwo(\tau)\langle\mathfrak{e}_\disceltwo,\mathfrak{e}_\discel\rangle
%	=
%	f_\discel(\tau)\qquad\text{for every~$\tau\in\HH_g$},
%	\]
%	we may apply the unfolding trick and deduce that
%	\bas
%	\pet{f}{\poincv{\gengenus}{\discel}{T}{\lattice}}
%	&=
%	\frac{1}{2}
%	\sum_{\gamma\in\mettrans{\gengenus}\backslash\Mp_{2\gengenus}(\ZZ)}
%	\int_{\mathcal{F}_\gengenus}\det\big(\Im(\gamma\cdot\tau)\big)^k \cdot
%	\overline{e\big(\tr T\gamma\cdot\tau\big)} \cdot
%	f_\discel(\gamma\cdot\tau)
%	\,\frac{dx\,dy}{\det y^{\gengenus+1}}
%	\\
%	&=
%	\frac{1}{2}
%	\sum_{\gamma\in\mettrans{\gengenus}\backslash\Mp_{2\gengenus}(\ZZ)}
%	\int_{\gamma(\mathcal{F}_\gengenus)}
%	\det y^k\cdot e(-\tr T\overline\tau)\cdot f_\discel(\tau)
%	\,\frac{dx\,dy}{\det y^{\gengenus+1}}
%	\\
%	&=
%	\int_{\trans{\gengenus}\backslash\HH_\gengenus}
%	\det y^k\cdot e(-\tr T\overline\tau)\cdot f_\discel(\tau)
%	\,\frac{dx\,dy}{\det y^{\gengenus+1}},
%	\eas
%	where in the last passage we use that~$I_{2\gengenus}$ and~$-I_{2\gengenus}$ give different classes in~$\mettrans{\gengenus}\backslash\Mp_{2\gengenus}(\ZZ)$, explaining why we dropped the factor of~$1/2$.
%	Let~$\cone{\gengenus}$ be the convex cone of the imaginary parts of matrices in~$\HH_{\gengenus}$; see~\eqref{eq;convconeimpart}.
%	A fundamental domain in~$\HH_\gengenus$ with respect to~$\trans{\gengenus}$ is given by
%	\[
%	\left\{
%	\tau=x+iy\in\HH_\gengenus
%	:
%	\text{$x$ is reduced modulo~$1$ and~$y\in\cone{\gengenus}$}
%	\right\}.
%	\]
%	We may then replace~$f_\discel$ with its Fourier expansion and compute
%	\ba\label{eq;TheLastEqgenJanpoinpet}
%	\pet{f}{\poincv{\gengenus}{\discel}{T}{\lattice}}
%	&=
%	\int_{y\in\cone{\gengenus}}
%	\det y^{k-\gengenus-1}
%	\int_{\text{$x$ mod $1$}}
%	e\big(-\tr T \overline\tau\big)
%	f_\discel(\tau)\,dx\,dy
%	\\
%	&=
%	\int_{y\in\cone{\gengenus}}
%	\det y^{k-\gengenus-1}
%	\exp(-2\pi\tr Ty)
%	\sum_{\substack{M\in q(\discel)+\halfint{\gengenus}\\ M>0}}
%	c_M(f_\discel)
%	\\
%	&\quad\times\exp(-2\pi\tr My)\underbrace{\int_{\text{$x$ mod $1$}}e\big(\tr(M-T)x\big)\,dx}_{=\delta_{M,T}}\,dy
%	\\
%	&=
%	c_T(f_\discel)\int_{y\in\cone{\gengenus}}
%	\det y^{k-\gengenus-1}
%	\exp(-4\pi\tr Ty)
%	\,dy.
%	\ea
%	As remarked in~\cite[p.~$90$,~(iv)]{klingen}, as well as in the last equality of the proof of~\cite[Proposition~$1.5$]{bruinierbook} for the case of genus~$g=1$, the integral over~$\cone{\gengenus}$ appearing on the right-hand side of the last equality in~\eqref{eq;TheLastEqgenJanpoinpet} equals~$\constant_{k,g}\det T^{(\gengenus+1)/2-k}$.}
	\end{proof}
	With the usual method, one can deduce from Lemma~\ref{lemma;petSiegpoinandcusp} the following result.
	\begin{lemma}\label{eq;spanpoincseries}
	Let~$k>2\genus$.
	The Poincaré series of weight~$k$ and genus~$\genus$ span the space of cusp forms~$S^k_{\genus,\lattice}$.
	\end{lemma}
%	\begin{proof}\textcolor{blue}{
%	Let~$\mathcal{S}$ be the space spanned by Poincaré series.
%	The space~$S^k_{\genus,\lattice}$, endowed with the Petersson inner product, is a Hermitian space.
%	In particular, it splits as~$S^k_{\genus,\lattice}=\mathcal{S}\oplus \mathcal{S}^\perp$.
%	\\
%	Let~$f\in\mathcal{S}^\perp$.
%	By Lemma~\ref{lemma;petSiegpoinandcusp} we deduce that~$c_T(f_\discel)=0$ for every~$\discel\in D_\lattice^\genus$ and every~$T\in\Sym_\genus(\QQ)$ positive definite.
%	This implies that~$f=0$, hence~$S^k_{\genus,\lattice}=\mathcal{S}^\perp$.}
%	\end{proof}
	\end{document}
	
	
	
	
	
	
	
	
	
	
	
	
	
	
	
	
	
	
	
	
	
	
	
	
	
	
	
	
	
	
	
	
	
	
	
	
	
	
	
	
\documentclass[../Lefschetz_decomposition.tex]{subfiles}

	\begin{document}
	Let~$k>2\gengenus$, and let~$\klinser{k}{\gengenus}{r}{\lattice}\colon S^k_{r,\lattice}\to M^k_{\gengenus,\lattice}$ be the linear operator lifting vector-valued Siegel cusp forms of genus~$r$ and weight~$k$ to the associated Klingen Eisenstein series of genus~$\gengenus$.
%	Recall that~$S^k_{0,L}=M^k_{0,L}=\CC$, so that we may construct the weight~$k$ normalized Siegel Eisenstein series of genus~$\gengenus$ (and index~$0\in D_\lattice^\genus$) as~$E^k_{\gengenus,\lattice}=\klinser{k}{\gengenus}{0}{\lattice}(1)$.
	We define%~$\vklisub{k}{0}{0}{\lattice}\coloneqq M^k_{0,L}$ and
	\[
	\vklisub{k}{\gengenus}{r}{\lattice}
	\coloneqq
	\klinser{k}{\gengenus}{r}{\lattice}(S^k_{r,\lattice})
%	\begin{dcases}
%	S^k_{\gengenus,\lattice}, & \text{if~$r=\gengenus$,}\\
%	\{f\in N^k_{\gengenus,\lattice} : \text{$\Phi_0(f)\in \vklisub{k}{\gengenus-1}{r}{\lattice}$ for every~$\disceltwo\in D_\lattice^\genus$}\}, & \text{if~$0\le r <\gengenus$.}
%	\end{dcases}
\qquad\text{for every~$0\leq r\leq \genus$.}
	\]
	If~$L$ is unimodular, then these are the subspaces of Klingen Eisenstein series introduced in~\cite[Section~$5$, Definition~$4$]{klingen}.
	We remark that~$\vklisub{k}{\gengenus}{0}{\lattice}=\CC \eiszero{k}{\genus}{L}$ and~$\vklisub{k}{\gengenus}{\genus}{\lattice}=S^k_{\genus,\lattice}$.
	
	By Theorem~\ref{thm;Phi0onkling} the restriction of the Siegel operator~$\Phi_0^{\genus-r}\colon \vklisub{k}{\gengenus}{r}{\lattice}\to S^k_{r,\lattice}$ is an isomorphism whose inverse is~$\klinser{k}{\gengenus}{r}{\lattice}$.
	The subspace~$\vklisub{k}{\gengenus}{r}{\lattice}$ is orthogonal to~$S^k_{\genus,\lattice}$ for every~$0\leq r < \genus$ by Proposition~\ref{prop;ortcusp&kling}.
	\begin{defi}
	The space of \emph{almost cusp forms}~$\almcusp{k}{\genus}{\lattice}$ is the subspace of~$M^k_{\genus,\lattice}$ defined as
	\[
	\almcusp{k}{\genus}{\lattice}
	=
	\bigoplus_{r=0}^\genus
	\vklisub{k}{\gengenus}{r}{\lattice}.
	\]
	The \emph{modular decomposition of~$f\in\almcusp{k}{\genus}{\lattice}$} is the rewriting of~$f$ with respect to the direct sum above .
	\end{defi}
	
	If~$g=1$, then the space of almost cusp forms is the same as its homonym in~\cite[Definition~$2.7$]{peterson} and~\cite[p.~$500$,~$(2)$]{bruiniermöller}.
	Since the subspace of Eisenstein series in~$M^k_{1,L}$ is in general of dimension greater than~$\dim \vklisub{k}{1}{0}{\lattice}=1$, the space~$\almcusp{k}{1}{\lattice}$ is usually smaller than~$M^k_{1,L}$.
	The same holds in higher genus.
	
	As we will see later, the cohomology class~$[\KMtheta{\genus}]$ of the Kudla--Millson theta function is an almost cusp form.
	It is the purpose of the present article to describe the modular decomposition of~$[\KMtheta{\genus}]$.
	This is given in Section~\ref{sec;gen2}.
	
	We define the map~$\vklin{k}{\genus}{r}{\lattice}\colon \almcusp{k}{r}{\lattice}\to \almcusp{k}{\genus}{\lattice}$ as the inverse of the restriction of~$\Phi^{\genus-r}_0$ to~$\vklisub{k}{\gengenus}{0}{\lattice}\oplus\cdots\oplus \vklisub{k}{\gengenus}{r}{\lattice}$.
	Concretely, it is defined on modular decompositions as
	\be\label{eq;defofKlinblamap}
	\vklin{k}{\genus}{r}{\lattice}\colon \sum_{j=0}^r
	\klinser{k}{r}{j}{\lattice}(f_j)
	\longmapsto
	\sum_{j=0}^{r}
	\klinser{k}{\gengenus}{j}{\lattice}(f_j),
	\qquad\text{$f_j\in S^k_{j,L}$.}
	\ee
%	We remark that if~$\disceltwo\neq 0$, then~$\Phi_\disceltwo \circ \vklin{k}{\genus}{r}{\lattice} = 0$ by Theorem~\ref{thm;Phi0onkling}.
%	\begin{ex}
%	If~$a\in \CC$ and~$f\in S^k_{1,L}$, then~$\vklin{k}{2}{1}{L}(a\cdot \eiszero{k}{1}{L} + f) = a\cdot \eiszero{k}{2}{L} + \klinser{k}{2}{1}{L}(f)$.
%	\end{ex}
	\end{document}
	
	
	
	
	
	
	
	
	
	
	
	
	
	
	
	
	
	
	
	
	
	
	
\documentclass[../Lefschetz_decomposition.tex]{subfiles}

	\begin{document}
	Böcherer proved in~\cite[Satz~$12$]{doubling} that certain diagonal restrictions of (essential subseries of) Siegel Eisenstein series can be written in terms of Klingen Eisenstein series arising from Poincaré series.
	The goal of this section is to extend Böcherer's result to the case of vector-valued Siegel modular forms with respect to the Weil representation.
	This is then used to show that the vanishing of Garrett's map stated in~\cite[Satz~$13$~$(1)$]{doubling} holds also in the vector-valued setting.
	This is a key ingredient to prove the modular decomposition of the Kudla--Millson theta function.
	\\
	
%	The Siegel Eisenstein series of genus~$g$ and weight~$k$ with respect to~$\weil{\lattice}{\gengenus}$ may be written as
%	\[
%	\eiszero{k}{\gengenus}{\lattice}(\tau)
%	=
%	\sum_{\gamma=\left(\begin{smallmatrix}
%	\ast & \ast \\
%	C & D
%	\end{smallmatrix}\right)\in\sieg{\gengenus}\backslash\Sp_{2\gengenus}(\ZZ)}
%	\det(C\tau + D)^{-k}
%	\weil{\lattice}{\gengenus}(\gamma)^{-1}\mathfrak{e}_0,
%	\]
%	where we include~$\Sp_{2\genus}(\ZZ)$ in~$\Mp_{2\genus}(\ZZ)$ by mapping~$\gamma$ to~$(\gamma,\phi)$, with~$\phi(\tau)=\sqrt{\det(C\tau+D)}$.
%	A full set of representatives of~$\sieg{\genus}\backslash\Sp_{2\genus}(\ZZ)$ is given by the set of pairs of symmetric matrices~$\{C,D\}$ which are non-associated and without common divisors.
	
	Let~$r\leq\genus$ be a non-negative integer.
%	We denote by~$T^(r)$ any~$r\times r$ matrix.
	We denote the \emph{Fourier--Jacobi expansion of cogenus~$r$ of~$\eiszero{k}{\gengenus}{\lattice}$} as
	\be\label{eq;FJexpsvcvv}
	\eiszero{k}{\gengenus}{\lattice}(\tau)
	=
	\sum_{\discel\in D_\lattice^r}
	\sum_{\substack{T\in q(\discel) + \halfint{r} \\ T\ge 0}}
	\mathfrak{e}_\discel
	\otimes
	\jaccoefvv{\genus}{r}{\discel}(\tau_4,\tau_2,T)\cdot e(\tr T\tau_1),
	\ee
	where~$\tau=\big(\begin{smallmatrix}
	\tau_1 & \tau_2\\
	\tau_2^t & \tau_4
	\end{smallmatrix}\big)$ is such that~$\tau_1\in\HH_r$ and~$\tau_4\in\HH_{\genus-r}$.
	The Fourier--Jacobi coefficients
	\[
	\jaccoefvv{\genus}{r}{\discel}\colon \HH_{\genus-r}\times\CC^{r\times(\genus-r)}\to\CC[D_\lattice^{\genus-r}]
	\]
	are Jacobi forms of weight~$k$ with respect to a Weil representation associated to the metaplectic Jacobi group of cogenus~$r$; see e.g.~\cite[Section~$3.2$]{kieferzuffetti} for the case of genus~$2$.
	\begin{defi}\label{def;Essentialpart}
	The \emph{essential part} of the Fourier--Jacobi expansion~\eqref{eq;FJexpsvcvv} of~$\eiszero{k}{\gengenus}{\lattice}$ is
	\[
	\espart{k}{\genus}{r}(\tau)
	=
	\sum_{\discel\in D_\lattice^r}
	\sum_{\substack{T\in q(\discel) + \halfint{r} \\ T>0}}
	\mathfrak{e}_\discel
	\otimes
	\jaccoefvv{\genus}{r}{\discel}(\tau_4,\tau_2,T)\cdot e(\tr T\tau_1).
	\]
	\end{defi}
	If~$M$ is a~$m\times n$ matrix with coefficients in~$\QQ$, in short~$M\in\QQ^{(m, n)}$, then we write the block decomposition of~$M$, where the top-left block is a~$\nu\times\nu$ matrix, as
	\[
	M=\Big(\begin{smallmatrix}
	M_1^{(\nu)} & M_2^{(\nu,n-\nu)}\\
	M_3^{(m-\nu,\nu)} & M_4^{(m-\nu,n-\nu)}
	\end{smallmatrix}\Big)\qquad\text{for every~$0\leq \nu \leq \min\{m,n\}$.}
	\]
%	\textcolor{blue}{
%	We define the following piece of notation for every~$0\leq \nu \leq\genus$.
%	\bas
%	ST(\genus)&\coloneqq\{(C,D)\in\ZZ^{(\genus,2\genus)} : \text{$C$ and $D$ are symmetric and~$\gcd(C,D)=1$}\},
%	\\
%	ST(\genus,\nu)&\coloneqq\{(C,D)\in ST(n) : \rk (C) = \nu \}
%	\\
%	ST(\genus)_\nu &\coloneqq\{(C,D)\in ST(\genus) : C_3^{(\genus-\nu,\nu)}=0\},
%	\\
%	ST(\genus)_\nu^* &\coloneqq\{(C,D)\in ST(\genus)_\nu : \rk C_1^{(\nu)} = \nu\}.
%	\eas
%	The group~$\GL(\genus,\ZZ)$ acts on~$ST(\genus)$ and~$ST(\genus,\nu)$ as~$U\colon (C,D)\mapsto(UC,UD)$.
%	An analogous action is given by
%	\[
%	\GL(\genus,\ZZ)_\nu\coloneqq
%	\{
%	U\in\GL(\genus,\ZZ) : U_3^{(\genus-\nu,\nu)}=0
%	\}
%	\]
%	on~$ST(\genus)_\nu$ and~$ST(\genus)_\nu^*$.}
%	
%	\textcolor{blue}{To study the properties of~$\espart{k}{\genus}{r}$ we first consider
%	\[
%	\Psi_{\genus,\nu}^k(\tau)
%	\coloneqq
%	\sum_{(C,D)\in\GL(\genus,\ZZ)\backslash ST(\genus,\nu)}
%	\phi_{(C,D)}(\tau)^{-2k}
%	\weil{\lattice}{\gengenus}(\gamma_{(C,D)})^{-1}\mathfrak{e}_0,
%	\]
%	where we denote by~$\gamma_{(C,D)}=\big(\big(\begin{smallmatrix}
%	\ast & \ast \\
%	C & D
%	\end{smallmatrix}\big),\phi_{(C,D)}\big)$ any of the fibers of a representative~$\big(\begin{smallmatrix}
%	\ast & \ast \\
%	C & D
%	\end{smallmatrix}\big)\in\Sp_{2\genus}(\ZZ)$ under the metaplectic covering map.}
%	
%	\textcolor{blue}{It is easy to see by using the properties of~$\weil{\lattice}{\gengenus}$ on matrices~$\big(\begin{smallmatrix}
%	U & \\
%	& U^{-t}
%	\end{smallmatrix}\big)$, with~$U\in\GL(\genus,\ZZ)$, that~$\Psi_{\genus,\nu}^k$ is well-defined.
%	Note that if we replace~$\gamma_{(C,D)}$ by~$\gamma'_{(C,D)}\coloneqq\big(\big(\begin{smallmatrix}
%	\ast & \ast \\
%	C & D
%	\end{smallmatrix}\big),-\phi_{(C,D)}\big)$, then by~\eqref{eq;howto-phi} we have
%	\bas
%	(-\phi_{(C,D)}(\tau))^{-2k}
%	\weil{\lattice}{\genus}(\gamma'_{(C,D)})^{-1}\mathfrak{e}_0
%	&=
%	\underbrace{(-1)^{-2k}(-1)^{n-2}}_{=1}
%	\phi_{(C,D)}(\tau)^{-2k}
%	\weil{\lattice}{\genus}(\gamma_{(C,D)})^{-1}\mathfrak{e}_0
%	\\
%	&=
%	\phi_{(C,D)}(\tau)^{-2k}
%	\weil{\lattice}{\genus}(\gamma_{(C,D)})^{-1}\mathfrak{e}_0.
%	\eas
%	Hence, it does not matter if we choose~$\gamma_{(C,D)}$ to be the standard inclusion of~$\big(\begin{smallmatrix}
%	\ast & \ast \\
%	C & D
%	\end{smallmatrix}\big)$ in~$\Mp_{2\genus}(\ZZ)$ or not.
%	}
%	\textcolor{blue}{\begin{lemma}\label{lemma;onPsivv}
%	Let~$\coeis{k}{\genus}{T}{\discel}$ denote the~$T$-th Fourier coefficient of the~$\discel$ component of~$\eiszero{k}{\gengenus}{\lattice}$.
%	The function~$\Psi_{\genus,\nu}^k$ is a sub-series of~$\eiszero{k}{\gengenus}{\lattice}$, whose Fourier expansion is
%	\be\label{eq;Psibövv}
%	\Psi_{\genus,\nu}^k(\tau)
%	=\sum_{\discel\in D_\lattice^\genus}
%	\sum_{\substack{T\in q(\discel)+\halfint{\genus} \\ \rk T = \nu \\ T\geq 0}} \coeis{k}{\genus}{T}{\discel} q^T \mathfrak{e}_\discel.
%	\ee
%	\end{lemma}}
%	\begin{proof}
%\textcolor{blue}{	We rewrite~$\eiszero{k}{\gengenus}{\lattice}$ as
%	\[
%	\eiszero{k}{\gengenus}{\lattice}(\tau)
%	=
%	\sum_{(C,D)\in\GL(\genus,\ZZ)\backslash ST(\genus)}
%	\phi_{(C,D)}(\tau)^{-2k}
%	\weil{\lattice}{\gengenus}(\gamma_{(C,D)})^{-1}\mathfrak{e}_0,
%	\]
%	where~$\gamma_{(C,D)}=\left(\left(\begin{smallmatrix}
%	\ast & \ast \\
%	C & D
%	\end{smallmatrix}\right),\phi_{(C,D)}\right)$ is any element mapping to~$\left(\begin{smallmatrix}
%	\ast & \ast \\
%	C & D
%	\end{smallmatrix}\right)\in \Sp_{2\gengenus}(\ZZ)$ under the metaplectic covering map.
%	Since~$\rk(UC)=\rk C$ for every~$U\in\GL(\genus,\ZZ)$, we may split~$\GL(\genus,\ZZ)\backslash ST(\genus)$ as the disjoint union of the cosets~$\GL(\genus,\ZZ)\backslash ST(\genus,\nu)$ with~$\nu=0,\dots,\genus$.
%	From this we see that~$\Psi_{\genus,\nu}^k$ is a sub-series of~$E^k_{\genus,\lattice}$.}
%	
%	\textcolor{blue}{By~\cite[Lemma~$1$]{doubling} a complete set of representatives of~$\GL(\genus,\ZZ)\backslash ST(\genus,\nu)$ is given by
%	\be\label{eq;setofrepvv}
%	\left\{
%	\left(
%	\big(\begin{smallmatrix}
%	C_1 & 0 \\
%	0 & 0_{\genus-\nu}
%	\end{smallmatrix}\big) W^t
%	,
%	\big(\begin{smallmatrix}
%	D_1 & 0 \\
%	0 & I_{\genus-\nu}
%	\end{smallmatrix}\big) W^{-1}
%	\right)
%	:
%	\substack{(C_1,D_1)\in\GL(\nu,\ZZ)\backslash ST(\nu,\nu) \\ W\in\GL(\genus,\ZZ)/\GL(\genus,\ZZ)_\nu}
%	\right\}.
%	\ee
%	Let
%	\[
%	\gamma_{W,(C_1,D_1)}
%	=
%	\left(\left(\begin{smallmatrix}
%	\ast & \ast \\
%	\big(\begin{smallmatrix}
%	C_1 & 0 \\
%	0 & 0_{\genus-\nu}
%	\end{smallmatrix}\big) W^t & \big(\begin{smallmatrix}
%	D_1 & 0 \\
%	0 & I_{\genus-\nu}
%	\end{smallmatrix}\big) W^{-1}
%	\end{smallmatrix}\right),
%	\phi_{W,(C_1,D_1)}
%	\right)
%	\]
%	be an element mapping under the metaplectic covering map to a representative of the quotient~$\sieg{\genus}\backslash\Sp_{2\genus}(\ZZ)$ arising as in~\eqref{eq;setofrepvv}.
%	This can be written as~${\iota\big(\gamma_{(C_1,D_1)},I_{\genus-\nu}\big)\cdot\gamma_W}$, where~$\iota$ is the inclusion defined in~\eqref{eq;inclprodsimpl}, for some~$\gamma_{(C_1,D_1)}=\big(\big(\begin{smallmatrix}
%	\ast& \ast \\
%	C_1 & D_1
%	\end{smallmatrix}\big),\phi_{(C_1,D_1)}\big)\in\Mp_{2\nu}(\ZZ)$ and some~$\gamma_W
%	=
%	\big(\big(\begin{smallmatrix}
%	W^t & \\
%	 & W^{-1}
%	\end{smallmatrix}\big),\phi_W\big)$ in~$\Mp_{2\genus}(\ZZ)$.}
%	
%	\textcolor{blue}{By Lemma~\ref{lemma;proofweilreponsubmeta} we may compute
%	\ba\label{eq;vcefexpvv}
%	\Psi_{\genus,\nu}^k(\tau)
%	&=
%	\sum_{W\in\GL(\genus,\ZZ)/\GL(\genus,\ZZ)_\nu}
%	\sum_{(C_1,D_1)\in\GL(\nu,\ZZ)\backslash ST(\nu,\nu)}
%	\phi_{W,(C_1,D_1)}(\tau)^{-2k}
%	\weil{\lattice}{\gengenus}(\gamma_{W,(C_1,D_1)})^{-1}\mathfrak{e}_0
%	\\
%	&=
%	\sum_{W\in\GL(\genus,\ZZ)/\GL(\genus,\ZZ)_\nu}
%	\phi_W(\tau)^{-2k}
%	\weil{\lattice}{\genus}(\gamma_W)^{-1}
%	\\
%	&\quad\times
%	\Big(\sum_{(C_1,D_1)\in\GL(\nu,\ZZ)\backslash ST(\nu,\nu)}
%	\phi_{(C_1,D_1)}\big( (W^t\tau W)^*\big)^{-2k}
%	\weil{\lattice}{\nu}(\gamma_{(C_1,C_2)})^{-1}\zeroe{\nu}\Big)\otimes\zeroe{\genus-\nu}
%	\\
%	&=
%	\sum_{\discel\in D_\lattice^\nu}
%	\sum_{\substack{T\in q(\discel) + \halfint{\nu} \\ T>0}}
%	\sum_{W\in\GL(\genus,\ZZ)/\GL(\genus,\ZZ)_\nu}
%	\coeis{k}{\nu}{T}{\discel}
%	\cdot
%	e\left(\tr
%	\big(\begin{smallmatrix}
%	T & 0\\
%	0 & 0_{\genus-\nu}
%	\end{smallmatrix}\big)\cdot W^t \tau W
%	\right)
%	\\
%	&\quad\times
%	\det W^{k}\weil{\lattice}{\genus}(\gamma_W)^{-1}
%	\mathfrak{e}_\discel\otimes\zeroe{\genus-\nu},
%	\ea
%	where in the last passage we used that if~$\tau^*\in\HH_\nu$, then
%	\ba\label{eq;wwwtpFKsolved}
%%	\sum_{(C_1,D_1)\in\GL(\nu,\ZZ)\backslash ST(\nu,\nu)}
%%	\phi_{(C_1,D_1)}(\tau^*)^{-2k}
%%	\weil{\lattice}{\nu}(\gamma_{(C_1,C_2)})^{-1}\zeroe{\nu}
%%	\\
%%	=
%	\Psi_{\nu,\nu}^k(\tau^*)
%	=
%	\sum_{\discel\in D_\lattice^{\nu}}\sum_{\substack{T\in q(\discel)+\halfint{\nu} \\ T>0}} a^k_\nu(T,\discel) e(\tr T\tau^*)\mathfrak{e}_\discel.
%	\ea
%	We prove~\eqref{eq;wwwtpFKsolved} following the wording of~\cite[Kapitel~IV, Bemerkung~$7.4$]{freitag}.
%	The invariance of~$\Psi_{\nu,\nu}^k$ with respect to the translations on~$\HH_\nu$ induced by~$\Mp_{2\nu}(\ZZ)$ implies that~$\Psi_{\nu,\nu}^k$ admits a Fourier expansion of the form
%	\[
%	\Psi_{\nu,\nu}^k(\tau^*)
%	=
%	\sum_{\discel\in D_\lattice^\nu}
%	\sum_{T\in q(\discel) + \halfint{\nu}}
%	c_T((\Psi_{\nu,\nu}^k)_\discel) q^T \mathfrak{e}_\discel.
%	\]
%	By the Koecher Principle and the fact that~$\Psi_{\nu,\nu}^k$ maps to~$0$ under the Siegel operators, if a Fourier coefficient of~$\Psi_{\nu,\nu}^k$ is non-zero, then it is indexed over a positive-definite matrix~$T$.
%	\\
%	\textbf{Remark on the Koecher Principle:}
%	$\Psi_{\nu,\nu}^k$ is invariant w.r.t.~$\GL_\nu(\ZZ)$, where the latter is embedded in~$\Sp_{2\nu}(\ZZ)$ in the usual way.
%	This invariance follows in the same way as proving that~$E^k_{\nu,L}$ is well-defined.
%	In fact, the embedding of~$\GL_\nu(\ZZ)$ in~$\Sp_\nu(\ZZ)$ lands in the Siegel parabolic~$\sieg{\nu}$, and~$\Psi_{\nu,\nu}^k$ can be rewritten as a sum over~$\sieg{\nu}\backslash\Sp_{2\nu}(\ZZ)$.
%	Thanks to the invariance w.r.t.~$\GL_\nu(\ZZ)$, one can deduce the invariance property of the Fourier coefficients of~$\Psi_{\nu,\nu}^k$, like~\eqref{eq;symmofcoef}.
%	The absolute convergence of~$\Psi_{\nu,\nu}^k$ implies a bound on Fourier coefficients, which together with the invariance property of the Fourier coefficients implies the vanishing of the ones indexed over non-positive-semidefinite matrices.
%	This standard argument can be found e.g.~in~\cite[Part~$3$, p.~$190$, Proof of Theorem~$1$]{1-2-3}.
%	\\
%	Consider now the differential operator
%	\[
%	\lvert \partial \rvert = \det\Big(e_{m,n}\frac{\partial}{\partial \tau^*_{m,n}}\Big),
%	\qquad 1\leq m,n\leq \nu,\qquad
%	\text{where}\quad e_{m,n}=\begin{cases}
%	1 & \text{if~$m=n$,}
%	\\
%	\frac{1}{2} & \text{otherwise.}
%	\end{cases}
%	\]
%	It has the property that~$\lvert\partial\rvert e^{2\pi i \tr (T\tau^*)}=(\det T) e^{2\pi i \tr(T\tau^*)}$ for every symmetric real~$\nu\times\nu$ matrix~$T$.
%	To prove~\eqref{eq;wwwtpFKsolved} it is then enough to show that
%	\[
%	\lvert \partial \rvert \big(E^k_{\nu,L} - \Psi_{\nu,\nu}^k\big)=0,
%	\]
%	where~$\lvert \partial \rvert$ is applied componentwise.
%	That follows from the fact that
%	\[
%	\lvert \partial\rvert\big(
%	\det(C_1\tau^* + D_1)^{-k}
%	\big) = 0\qquad\text{if~$\det C_1=0$.}
%	\]
%	}
%	
%	\textcolor{blue}{Since the Siegel operator~$\Phi^{\genus-\nu}_0$ applied to~$\eiszero{k}{\gengenus}{\lattice}$ gives~$\eiszero{k}{\nu}{\lattice}$, we deduce that
%	\[
%	\coeis{k}{\nu}{T}{\discel}
%	=
%	\Bigcoeis{k}{\genus}{\big(\begin{smallmatrix}
%	T & 0\\
%	0 & 0
%	\end{smallmatrix}\big)}{(\discel,0,\dots,0)}
%	\qquad\text{for every~$\discel\in D_\lattice^\nu$ and~$T\in q(\discel)+\halfint{\nu}$.}
%	\]
%	Furthermore, since~$2k\equiv n-2$ mod~$4$ we have
%	\be\label{eq;mtacifvv}
%	\phi_W(\tau)^{-2k}\weil{\lattice}{\genus}(\gamma_W)^{-1}
%	\mathfrak{e}_\discel\otimes\zeroe{\genus-\nu}
%	=
%	\underbrace{(\pm 1)^{-2k+n-2}}_{=1}
%	\underbrace{\det W^{-k+(n-2)/2}}_{=1}
%	\mathfrak{e}_{(\discel,0,\dots, 0)W^t},
%	\ee
%	\\\textcolor{blue}{\textbf{Note:} Here it is not enough that~$2k\equiv n-2$ mod~$2$ only.}\\
%	and by~\eqref{eq;symmofcoef2} we also have
%	\[
%	\Bigcoeis{k}{\genus}{\big(\begin{smallmatrix}
%	T & 0\\
%	0 & 0
%	\end{smallmatrix}\big)}{(\discel,0,\dots,0)}
%	=
%	\Bigcoeis{k}{\genus}{W\big(\begin{smallmatrix}
%	T & 0\\
%	0 & 0
%	\end{smallmatrix}\big)W^t}{(\discel,0,\dots,0)W^t},
%	\]
%	so that we may rewrite~\eqref{eq;vcefexpvv} as
%	\bas
%	\Psi_{\genus,\nu}^k(\tau)
%	&=
%	\sum_{\discel\in D_\lattice^\nu}
%	\sum_{\substack{T\in q(\discel) + \halfint{\nu} \\ T>0}}
%	\sum_{W\in\GL(\genus,\ZZ)/\GL(\genus,\ZZ)_\nu}
%	\Bigcoeis{k}{\genus}{W\big(\begin{smallmatrix}
%	T & 0\\
%	0 & 0
%	\end{smallmatrix}\big)W^t}{(\discel,0,\dots,0)W^t}
%	\\
%	&\quad \times
%	e\left(\tr
%	W\big(\begin{smallmatrix}
%	T & 0\\
%	0 & 0_{\genus-\nu}
%	\end{smallmatrix}\big)W^t \tau 
%	\right)
%	\cdot \mathfrak{e}_{(\discel,0,\dots 0)W^t},
%	\eas}
%	
%	\textcolor{blue}{The matrices~$W\big(\begin{smallmatrix}
%	T & 0\\
%	0 & 0
%	\end{smallmatrix}\big)W^t$, with~$W\in\GL(\genus,\ZZ)/\GL(\genus,\ZZ)_\nu$ and~$T\in\halfint{\nu}$ positive definite, span the set of positive semidefinite matrices in~$\halfint{\genus}$ of rank~$\nu$.
%	More generally, the matrices~$W\big(\begin{smallmatrix}
%	T & 0\\
%	0 & 0
%	\end{smallmatrix}\big)W^t$, with~$W\in\GL(\genus,\ZZ)/\GL(\genus,\ZZ)_\nu$, $T\in q(\discel)+\halfint{\nu}$ positive definite and~$\discel\in D_\lattice^\nu$, span the set of positive semidefinite matrices in~$q(\disceltwo)+\halfint{\genus}$, with~$\disceltwo\in D_\lattice^\genus$, of rank~$\nu$.}
%	
%	\textcolor{blue}{From~\eqref{eq;mtacifvv} we deduce that the elements of the form~$\phi_W(\tau)^{-2k}\weil{\lattice}{\genus}(\gamma_W)^{-1}
%	\mathfrak{e}_\discel\otimes\zeroe{\genus-\nu}$ generate the set of all~$\mathfrak{e}_\beta$, with~$\disceltwo\in D_\lattice^\genus$, admitting matrices~$T \in q(\disceltwo)+\halfint{\genus}$ of rank~$\nu$.
%%	\\\textcolor{blue}{Note: For fixed~$\nu$, we can not expect to find some positive semidefinite~$T\in q(\beta) + \halfint{\genus}$ of rank~$\nu$ for every~$\beta\in D_\lattice^\genus$.
%%	Think for instance to the case of~$\nu=0$, genus~$\genus=2$ and a lattice~$L$ admitting some~$\beta\in D_\lattice^2$ such that~$q(\beta)=\big(\begin{smallmatrix}
%%	1/3 & \\
%%	 & 1/3
%%	\end{smallmatrix}\big)$.
%%	\\
%%	Alternatively, think to the constant terms of the Eisenstein series of genus~$1$.
%%	If they are non zero, then they are associated to isotropic elements of~$D_\lattice$.}\\
%	With these properties we may deduce~\eqref{eq;Psibövv} from~\eqref{eq;vcefexpvv}.}
%%%%%%%%%%%%%%%%%%%%%% OLD VERSION WITH POSITIVE CHOICE IN METAPLECTIC FIBER %%%%%%%%%%%%	
%	Note that this can be written as~$\iota\big(\gamma_{(C_1,D_1)},I_{\genus-\nu}\big)\cdot\gamma_W$, where~$\iota$ is the inclusion defined in~\eqref{eq;inclprodsimpl}, $\gamma_{(C_1,D_1)}=\big(\begin{smallmatrix}
%	\ast& \ast \\
%	C_1 & D_1
%	\end{smallmatrix}\big)$ is in~$\Sp_{2\nu}(\ZZ)$, and
%	\[
%	\gamma_W
%	=
%	\begin{cases}
%	\big(\big(\begin{smallmatrix}
%	W^t & \\
%	 & W^{-1}
%	\end{smallmatrix}\big),1\big),
%	& \text{if~$\det W = 1$}
%	\\
%	\big(\big(\begin{smallmatrix}
%	W^t & \\
%	 & W^{-1}
%	\end{smallmatrix}\big),-i\big),
%	& \text{if~$\det W = -1$}
%	\end{cases}
%	\qquad\text{is in~$\Mp_{2\genus}(\ZZ)$}.
%	\]
%	The choice of~$\gamma_W$ is done in such a way that~$\iota\big(\gamma_{(C_1,D_1)},I_{\genus-\nu}\big)\cdot\gamma_W$ lies in the standard inclusion of~$\Sp_{2\genus}(\ZZ)$ in~$\Mp_{2\genus}(\ZZ)$.
%	
%	\textcolor{red}{IMPORTANT REMARK - This must be said.}
%	
%	\textcolor{blue}{\textbf{Explanation of why we need to choose~$\gamma_W$ in that way:}
%	Take for instance~$\gamma_{(C_1,D_1)}=\big(\begin{smallmatrix}
%	A & \\
%	 & A^{-t}
%	\end{smallmatrix}\big)$ with~$\det A=-1$.
%	We consider it as an element of~$\Mp_{2\nu}(\ZZ)$ under the natural inclusion of~$\Sp_{2\nu}(\ZZ)$, which we denote with the same symbol, namely~$\gamma_{(C_1,D_1)}=\big(\big(\begin{smallmatrix}
%	A & \\
%	 & A^{-t}
%	\end{smallmatrix}\big),i\big)$.
%	From the definition of~$\iota$ we have that
%	\[
%	\iota\big(\gamma_{(C_1,D_1)},I_{\genus-\nu}\big)
%	=
%	\left(\left(\begin{smallmatrix}
%	A & & & \\
%	  & I_{\genus-\nu} & & \\
%	  & & A^{-t} & \\
%	  & & & I_{\genus-\nu}
%	\end{smallmatrix}\right), i\right).
%	\]
%	We want~$\gamma_{W,(C_1,D_1)}$ to be an element in~$\Mp_{2\genus}(\ZZ)$ which arises from the inclusion of~$\Sp_{2\genus}(\ZZ)$, namely with the choice of the square root to be~$+\sqrt{...}$.
%	If~$\det W=-1$, thanks to the definition of~$\gamma_W$ we gave above we may use the group law of~$\Mp_{2\genus}(\ZZ)$ to deduce that
%	\[
%	\iota\big(\gamma_{(C_1,D_1)},I_{\genus-\nu}\big)\cdot \gamma_W
%	=
%	\left(\left(\begin{smallmatrix}
%	\ast & \ast\\
%	0
%	&
%	\big(\begin{smallmatrix}
%	A^{-t} & \\
%	 & I_{\genus-\nu}
%	\end{smallmatrix}\big)W^{-1}
%	\end{smallmatrix}\right),\underbrace{-i^2}_{=1}\right).
%	\]
%	If we choose~$\gamma_W$ in a different way, the highlighted square-root above would be~$-1$, which is not what we want: None of the elements of~$\Mp_{2\genus}(\ZZ)$ coming from~$\Sp_{2\genus}(\ZZ)$ can have~$-1$ as square root function.
%	\\
%	\textbf{Note:} This behavior comes from the fact that the standard inclusion~$\Sp_{2\genus}(\ZZ)\hookrightarrow\Mp_{2\genus}(\ZZ)$ is NOT a homomorphism of groups in general!
%	Hence, products of matrices in~$\Sp_{2\genus}(\ZZ)$ do not necessarily map to the products of their standard inclusions in~$\Mp_{2\genus}(\ZZ)$.
%	}
%
%	By Lemma~\ref{lemma;proofweilreponsubmeta} we may compute
%	\ba\label{eq;vcefexpvv}
%	\Psi_{\genus,\nu}^k(\tau)
%	&=
%	\sum_{W\in\GL(\genus,\ZZ)/\GL(\genus,\ZZ)_\nu}
%	\sum_{(C_1,D_1)\in\GL(\nu,\ZZ)\backslash ST(\nu,\nu)}
%	\det\left(
%	\big(\begin{smallmatrix}
%	C_1 & 0 \\
%	0 & 0_{\genus-\nu}
%	\end{smallmatrix}\big) W^t \tau + \big(\begin{smallmatrix}
%	D_1 & 0 \\
%	0 & I_{\genus-\nu}
%	\end{smallmatrix}\big) W^{-1}
%	\right)^{-k}
%	\\
%	&\quad\times
%	\weil{\lattice}{\gengenus}(\gamma_{W,(C_1,D_1)})^{-1}\mathfrak{e}_0
%	\\
%	&=
%	\sum_{W\in\GL(\genus,\ZZ)/\GL(\genus,\ZZ)_\nu}
%	\det W^{k}\weil{\lattice}{\genus}(\gamma_W)^{-1}
%	\\
%	&\quad\times
%	\Big(\sum_{(C_1,D_1)\in\GL(\nu,\ZZ)\backslash ST(\nu,\nu)}
%	\det\left(
%	\big(\begin{smallmatrix}
%	C_1 & 0 \\
%	0 & 0_{\genus-\nu}
%	\end{smallmatrix}\big) W^t \tau W + \big(\begin{smallmatrix}
%	D_1 & 0 \\
%	0 & I_{\genus-\nu}
%	\end{smallmatrix}\big)
%	\right)^{-k}
%	\\
%	&\quad\times
%	\weil{\lattice}{\nu}(\gamma_{(C_1,C_2)})^{-1}\zeroe{\nu}\Big)\otimes\zeroe{\genus-\nu}
%	\\
%	&=
%	\sum_{\discel\in D_\lattice^\nu}
%	\sum_{\substack{T\in q(\discel) + \halfint{\nu} \\ T>0}}
%	\sum_{W\in\GL(\genus,\ZZ)/\GL(\genus,\ZZ)_\nu}
%	\coeis{k}{\nu}{T}{\discel}
%	\cdot
%	e\left(\tr
%	\big(\begin{smallmatrix}
%	T & 0\\
%	0 & 0_{\genus-\nu}
%	\end{smallmatrix}\big)\cdot W^t \tau W
%	\right)
%	\\
%	&\quad\times
%	\det W^{k}\weil{\lattice}{\genus}(\gamma_W)^{-1}
%	\mathfrak{e}_\discel\otimes\zeroe{\genus-\nu},
%	\ea
%	where in the last passage we used that if~$\tau^*\in\HH_\nu$, then
%	\bas
%	{}&{}\sum_{(C_1,D_1)\in\GL(\nu,\ZZ)\backslash ST(\nu,\nu)}
%	\det
%	( C_1 \tau^* + D_1 )^{-k}
%	\weil{\lattice}{\nu}(\gamma_{(C_1,C_2)})^{-1}\zeroe{\nu}
%	\\
%	&=
%	\sum_{\discel\in D_\lattice^{\nu}}\sum_{\substack{T\in q(\discel)+\halfint{\nu} \\ T>0}} a^k_\nu(T,\discel) e(\tr T\tau^*)\mathfrak{e}_\discel.
%	\eas
%	\textcolor{red}{Why is this true? Perhaps like in proof that~$E^k_\nu$ is normalized.}
%	
%	Since the Siegel operator~$\Phi^{\genus-\nu}_0$ applied to~$\eiszero{k}{\gengenus}{\lattice}$ gives~$\eiszero{k}{\nu}{\lattice}$, we deduce that
%	\[
%	\coeis{k}{\nu}{T}{\discel}
%	=
%	\Bigcoeis{k}{\genus}{\big(\begin{smallmatrix}
%	T & 0\\
%	0 & 0
%	\end{smallmatrix}\big)}{(\discel,0,\dots,0)}
%	\qquad\text{for every~$\discel\in D_\lattice^\nu$ and~$T\in q(\discel)+\halfint{\nu}$.}
%	\]
%	Furthermore, we have
%	\be\label{eq;mtacifvv}
%	\det W^{k}\weil{\lattice}{\genus}(\gamma_W)^{-1}
%	\mathfrak{e}_\discel\otimes\zeroe{\genus-\nu}
%	=
%	\underbrace{\det W^{2(n-1)}}_{=1} \mathfrak{e}_{(\discel,0,\dots, 0)W^t},
%	\ee
%	and by~\eqref{eq;symmofcoef2} we also have
%	\[
%	\Bigcoeis{k}{\genus}{\big(\begin{smallmatrix}
%	T & 0\\
%	0 & 0
%	\end{smallmatrix}\big)}{(\discel,0,\dots,0)}
%	=
%	\Bigcoeis{k}{\genus}{W\big(\begin{smallmatrix}
%	T & 0\\
%	0 & 0
%	\end{smallmatrix}\big)W^t}{(\discel,0,\dots,0)W^t},
%	\]
%	so that we may rewrite~\eqref{eq;vcefexpvv} as
%	\bas
%	\Psi_{\genus,\nu}^k(\tau)
%	&=
%	\sum_{\discel\in D_\lattice^\nu}
%	\sum_{\substack{T\in q(\discel) + \halfint{\nu} \\ T>0}}
%	\sum_{W\in\GL(\genus,\ZZ)/\GL(\genus,\ZZ)_\nu}
%	\Bigcoeis{k}{\genus}{W\big(\begin{smallmatrix}
%	T & 0\\
%	0 & 0
%	\end{smallmatrix}\big)W^t}{(\discel,0,\dots,0)W^t}
%	\\
%	&\quad \times
%	e\left(\tr
%	W\big(\begin{smallmatrix}
%	T & 0\\
%	0 & 0_{\genus-\nu}
%	\end{smallmatrix}\big)W^t \tau 
%	\right)
%	\cdot \mathfrak{e}_{(\discel,0,\dots 0)W^t},
%	\eas
%	
%	The matrices~$W\big(\begin{smallmatrix}
%	T & 0\\
%	0 & 0
%	\end{smallmatrix}\big)W^t$, with~$W\in\GL(\genus,\ZZ)/\GL(\genus,\ZZ)_\nu$ and~$T\in\halfint{\nu}$ positive definite, span the set of positive semidefinite matrices in~$\halfint{\genus}$ of rank~$\nu$.
%	More generally, the matrices~$W\big(\begin{smallmatrix}
%	T & 0\\
%	0 & 0
%	\end{smallmatrix}\big)W^t$, with~$W\in\GL(\genus,\ZZ)/\GL(\genus,\ZZ)_\nu$, $T\in q(\discel)+\halfint{\nu}$ positive definite and~$\discel\in D_\lattice^\nu$, span the set of positive semidefinite matrices in~$q(\disceltwo)+\halfint{\genus}$, with~$\disceltwo\in D_\lattice^\genus$, of rank~$\nu$.
%	
%	From~\eqref{eq;mtacifvv} we deduce that the elements of the form~$\det W^{k}\weil{\lattice}{\genus}(\gamma_W)^{-1}
%	\mathfrak{e}_\discel\otimes\zeroe{\genus-\nu}$ generate the set of all~$\mathfrak{e}_\beta$, with~$\disceltwo\in D_\lattice^\genus$, admitting matrices~$T \in q(\disceltwo)+\halfint{\genus}$ of rank~$\nu$.
%%	\\\textcolor{blue}{Note: For fixed~$\nu$, we can not expect to find some positive semidefinite~$T\in q(\beta) + \halfint{\genus}$ of rank~$\nu$ for every~$\beta\in D_\lattice^\genus$.
%%	Think for instance to the case of~$\nu=0$, genus~$\genus=2$ and a lattice~$L$ admitting some~$\beta\in D_\lattice^2$ such that~$q(\beta)=\big(\begin{smallmatrix}
%%	1/3 & \\
%%	 & 1/3
%%	\end{smallmatrix}\big)$.
%%	\\
%%	Alternatively, think to the constant terms of the Eisenstein series of genus~$1$.
%%	If they are non zero, then they are associated to isotropic elements of~$D_\lattice$.}\\
%	With these properties we may deduce~\eqref{eq;Psibövv} from~\eqref{eq;vcefexpvv}.
%	\end{proof}
%	\textcolor{blue}{\begin{lemma}\label{lemma;frombbvv}
%	For every~$0\leq r\leq\genus$ we have
%	\[
%	\espart{k}{\genus}{r}(\tau)
%	=
%	\sum_{\substack{(C,D)\in\GL(\genus,\ZZ)_r\backslash ST(\genus)_r^*}}
%	\phi_{(C,D)}(\tau)^{-2k}
%	\weil{\lattice}{\genus}\big(\gamma_{(C,D)}\big)^{-1}\mathfrak{e}_0,
%	\]
%	where~$\gamma_{(C,D)}=\left(\left(\begin{smallmatrix}
%	\ast & \ast \\
%	C & D
%	\end{smallmatrix}\right),\phi_{(C,D)}\right)$ is any element mapping to~$\left(\begin{smallmatrix}
%	\ast & \ast \\
%	C & D
%	\end{smallmatrix}\right)\in \Sp_{2\gengenus}(\ZZ)$ under the metaplectic covering map.
%	\\\textcolor{blue}{\textbf{Note:} The reason why the RHS above does not depend on the choice of the branch of the square root~$\phi_{(C,D)}$ is the same as for~$E^k_{\genus,\lattice}$.}
%	\end{lemma}}
%	\begin{proof}
%	\textcolor{blue}{We follow the wording of~\cite[Lemma~$3$]{doubling}.
%	For every~$0\le\nu\le\genus$ we study the possible contribution of the Fourier expansion of~$\Psi^k_{\genus,\nu}$ to~$\espart{k}{\genus}{r}$.
%	By Lemma~\ref{lemma;onPsivv} we may consider only the cases where~$\nu\ge r$.
%	Without loss of generality, we choose the representatives of~$\GL(\genus,\ZZ)/\GL(\genus,\ZZ)_\nu$ to be of the form
%	\[
%	W=\left(\begin{smallmatrix}
%	w_1^{(\nu)} & \ast \\
%	\ast & \ast
%	\end{smallmatrix}\right)
%	\qquad
%	\text{with}
%	\qquad
%	w_1^{(\nu)}=\left(\begin{smallmatrix}
%	w_{1,1}^{(r)} & 0^{(r,\nu-r)} \\
%	\ast & \ast
%	\end{smallmatrix}\right).
%	\]
%	Recall from~\eqref{eq;vcefexpvv} that the Fourier expansion of~$\Psi^k_{\genus,\nu}$ can be rewritten as
%	\[
%	\Psi^k_{\genus,\nu}(\tau)
%	=
%	\sum_{\discel\in D_\lattice^\genus}
%	\sum_{\substack{T\in q(\discel) + \halfint{\nu} \\ T>0}}
%	\sum_{W\in\GL(\genus,\ZZ)/\GL(\genus,\ZZ)_\nu}
%	a^k_\nu(T,\discel)
%	\cdot e\left(
%	W\big(\begin{smallmatrix}
%	T & 0\\
%	0 & 0_{\genus-\nu}
%	\end{smallmatrix}\big) W^t\cdot \tau
%	\right)\mathfrak{e}_{(\discel,0,\dots,0)W^t}.
%	\]
%	The contributions of this expansion in the one of~$\espart{k}{\genus}{r}$ are the ones where~$W$ is such that~$w^{(r)}_{1,1}$ has maximal rank.
%	Without loss of generality we may also assume that the representative in~\eqref{eq;setofrepvv} are such that~$(C_1,D_1)\in\GL(\nu,\ZZ)\backslash ST(\nu,\nu)$ are chosen to lie in~$\GL(\nu,\ZZ)_r\backslash ST(\nu,\nu)\cap ST(\nu)_r$, so that
%	\[
%	\left(
%	\big(\begin{smallmatrix}
%	C_1 & 0 \\
%	0 & 0_{\genus-\nu}
%	\end{smallmatrix}\big) W^t
%	,
%	\big(\begin{smallmatrix}
%	D_1 & 0 \\
%	0 & I_{\genus-\nu}
%	\end{smallmatrix}\big) W^{-1}
%	\right)
%	\]
%	lies in~$ST(\genus)_r^*$ exactly when~$w_{1,1}^{(r)}$ is of maximal rank.
%	Since two elements of~$ST(\genus)_r^*$ are associated if and only if they are~$\GL(\genus,\ZZ)_r$-associated, we deduce the assertion.}
%	\end{proof}
%	
%	\textcolor{blue}{We want to rewrite~$\espart{k}{\genus}{r}$ isolating the dependence from~$\tau_4$.
%	To do so, we need to introduce the following notation.
%	We denote by~$\embGL$ the embedding of~$\GL(\genus,\RR)$ in~$\Sp_{2\genus}(\RR)$ given by~$\embGL\colon U\mapsto \big(\begin{smallmatrix}
%	U^{-t} & \\
%	 & U
%	\end{smallmatrix}\big)$.
%	We also say that two matrices are \emph{associated} if they differ from a left factor of~$\GL(\genus,\ZZ)$.}
%	
%	\textcolor{blue}{We define the following subgroups of~$\Sp_{2\genus}(\ZZ)$, following the wording of~\cite[p.~$25$]{doubling}.
%	\bas
%	\mathfrak{C}_{\genus,r}&\coloneqq\Big\{
%	M\in\Sp_{2\genus}(\ZZ) : C=\Big(\begin{smallmatrix}
%	C_1^{(r)} & 0 \\
%	0 & 0
%	\end{smallmatrix}\Big),\quad D=\Big(
%	\begin{smallmatrix}
%	D_1^{(r)} & \ast \\
%	0 & \ast
%	\end{smallmatrix}
%	\Big)
%	\Big\},
%	\\
%	\mathfrak{D}_{\genus,r}&\coloneqq\Big\{
%	M\in\mathfrak{C}_{\genus,r} : D=\Big(
%	\begin{smallmatrix}
%	D_1^{(r)} & \ast \\
%	0 & I_{\genus-r}
%	\end{smallmatrix}
%	\Big)
%	\Big\},
%	\\
%	\mathfrak{E}_{\genus,r}&\coloneqq
%	\left\{
%	M\in\Sp_{2\genus}(\ZZ) : M=\left(\begin{smallmatrix}
%	A^{(r)} & & B^{(r)} & \\
%	 & I_{\genus-r} & & \\
%	 C^{(r)} & & D^{(r)} & \\
%	  & & & I_{\genus-r}
%	\end{smallmatrix}\right)
%	\right\},
%	\\
%	\mathfrak{F}_{\genus,r}
%	&\coloneqq
%	\left\{
%	M\in\Sp_{2\genus}(\ZZ) : M=\left(\begin{smallmatrix}
%	I_{\genus-r} & & & \\
%	 & A^{(r)} & & B^{(r)} \\
%	 & & I_{\genus-r} & \\	 
%	 & C^{(r)} & & D^{(r)}
%	\end{smallmatrix}\right)
%	\right\}.
%	\eas
%	\textcolor{blue}{\textbf{Note:} The group~$\mathfrak{C}_{\genus,r}$ is just the Klingen parabolic subgroup.
%	In fact, it is used by Böcherer~\cite[Definition on p.~$32$]{doubling} to define Klingen Eisenstein series, and it can be checked that if~$M\in\mathfrak{C}_{\genus,r}$, then~$A=\big(\begin{smallmatrix}
%	A_1 & 0 \\
%	A_3 & A_4
%	\end{smallmatrix}\big)$, where~$A_1\in\ZZ^{(r,r)}$.}
%	\\
%	For every subgroup~$\Delta$ of~$\Sp_{2\genus}(\RR)$, we define~$\Delta_\infty\coloneqq\{M\in\Delta : C=0_\genus\}$.
%	We also set
%	\bas
%	\Gamma^{(\genus)}
%	&\coloneqq\Sp_{2\genus}(\ZZ),
%	\\
%	\Gamma^{(\genus,r)}
%	&\coloneqq
%	\left\{
%	M\in\Gamma^{(\genus)} : C = \big(\begin{smallmatrix}
%	C_1^{(r)} & \ast \\
%	0 & \ast
%	\end{smallmatrix}\big)
%	\right\},
%	\\
%	\Gamma^{(\genus,r)}_*
%	&\coloneqq
%	\{
%	M\in\Gamma^{(\genus,r)} : \rk(C_1^{(r)})=r
%	\},
%	\\
%	\mathfrak{E}_{\genus,r}^*
%	&\coloneqq
%	\mathfrak{E}_{\genus,r}\cap\Gamma^{(\genus,r)}_*,
%	\\
%	\GL_r(\genus,\ZZ)^*
%	&
%	\coloneqq\{
%	U\in\GL(\genus,\ZZ) : \rk \big( (U^{-1})^{(r)}_1\big)=r
%	\}.
%	\eas}
%	\textcolor{blue}{\begin{lemma}\label{lemma;form(11)bövv}
%	For every~$0\leq r\leq\genus$ we have
%	\bas
%	\espart{k}{\genus}{r}(\tau)
%	&=
%	\sum_{N\in (\mathfrak{E}_{\genus,r})_\infty\backslash\mathfrak{E}_{\genus,r}^*}
%	\sum_{U\in\GL(\genus,\ZZ)_r\backslash\GL_r(\genus,\ZZ)^*}
%	\sum_{F\in(\mathfrak{F}_{\genus,\genus-r})_\infty\backslash\mathfrak{F}_{\genus,\genus-r}}
%	\phi_{N\embGL(U)F}(\tau)^{-2k}
%	\\
%	&\quad\times
%	\weil{\lattice}{\genus}\big(\gamma_{N\embGL(U)F}\big)^{-1}\mathfrak{e}_0
%	,
%	\eas
%	where~$\gamma_{N\embGL(U)F}=\big(N\embGL(U)F,\phi_{N\embGL(U)F}(\tau)\big)\in\Mp_{2\genus}(\ZZ)$ is any element mapping to~$N\embGL(U)F$ under the metaplectic cover.
%	\\\textcolor{blue}{\textbf{Note:} The reason why the RHS above does not depend on the choice of the branch of the square root~$\phi_{(C,D)}$ is the same as for~$E^k_{\genus,\lattice}$.}
%	\end{lemma}
%	\begin{proof}
%	If~$M=\big(\begin{smallmatrix}
%	A & B\\
%	C & D
%	\end{smallmatrix}\big)$ runs in a complete system of representatives for~$(\mathfrak{C}_{\genus,r})_\infty\backslash\Gamma^{(\genus,r)}_*$, then the couples~$(C,D)$ give a complete system of representatives of~$\GL(\genus,\ZZ)_r\backslash ST(\genus)_r^*$.
%	We then deduce from Lemma~\ref{lemma;frombbvv} that
%	\[
%	\espart{k}{\genus}{r}(\tau)
%	=
%	\sum_{M\in(\mathfrak{C}_{\genus,r})_\infty\backslash\Gamma^{(\genus,r)}_*}
%	\phi_M(\tau)^{-2k}
%	\weil{\lattice}{\genus}(\gamma_M)^{-1}\mathfrak{e}_0,
%	\]
%	where~$\gamma_M=(M,\phi_M)\in\Mp_{2\genus}(\ZZ)$, and~$\phi_M$ does not necessarily arise from the principal branch of the square root.
%	The assertion we want to prove follows immediately from the study of complete sets of representatives of~$(\mathfrak{C}_{\genus,r})_\infty\backslash\Gamma^{(\genus,r)}_*$ given in~\cite[Proposition~$5$]{doubling}.
%	\end{proof}
%	For simplicity we define~$A[B]\coloneqq B^t A B$.}
%	\textcolor{blue}{\begin{lemma}\label{lemma;fbiamttle}
%	Let~$k>\genus+1$.
%	For every~$0\leq r\leq\genus$ we have
%	\bas
%	\espart{k}{\genus}{r}(\tau)
%	&=
%	\sum_{\discel\in D_\lattice^r}
%	\sum_{\substack{T\in q(\discel) + \halfint{r} \\ T>0}}
%	\sum_{\substack{U\in\GL(\genus,\ZZ)_r\backslash\GL_r(\genus,\ZZ)^* \\ U^{-1}=\big(
%	\begin{smallmatrix}
%	w_1 & \ast \\
%	w_3 & \ast
%	\end{smallmatrix}
%	\big)}}
%	\sum_{M\in \Gamma^{(\genus-r)}_\infty\backslash\Gamma^{(\genus-r)}}
%	a^k_r(T,\discel)
%	\phi_M(\tau_4)^{-2k}
%	\\
%	&\quad\times 
%	e\left(
%	\tr T
%	\big(\iota(I_{2r},M)\cdot \tau\big)
%	\big[\big(\begin{smallmatrix}
%	w_1
%	\\
%	w_3
%	\end{smallmatrix}\big)\big]
%	\right)
%	\mathfrak{e}_{\discel w_1^t}
%	\otimes
%	\big(\weil{\lattice}{\genus-r}(\gamma_M)^{-1}\mathfrak{e}_{\discel w_3^t}\big),
%	\eas
%	where we denote by~$w_1$ the top-left~$r\times r$ block of~$U^{-1}$, and where~$\gamma_M\coloneqq(M,\phi_M)$ is any element in~$\Mp_{2(\genus-r)}(\ZZ)$ mapping to~$M$ under the metaplectic cover.
%	\end{lemma}}
%	\begin{proof}
%	\textcolor{blue}{Let~$\gamma_N\coloneqq(N,\phi_N)$, $\gamma_{\embGL(U)}\coloneqq(\embGL(U),\phi_{\embGL(U)})$ and~$\gamma_F\coloneqq(F,\phi_F)$ be elements respectively in the fiber of~$N$, $\embGL(U)$ and~$F$ under the metaplectic covering map, and let~$\gamma_{N\embGL(U)F}\coloneqq\gamma_N\cdot\gamma_{\embGL(U)}\cdot \gamma_F$.
%	If we write~$\gamma_{N\embGL(U)F}=\big(N\embGL(U)F,\phi_{N\embGL(U)F}\big)$, then
%	\bas
%	\phi_{N\embGL(U)F}(\tau)
%%	&=
%%	\phi_{N\embGL(U)}(F\cdot\tau)\cdot \phi_F(\tau)
%	=
%	\phi_N(\embGL(U)F\cdot\tau)
%	\cdot
%	\underbrace{\phi_{\embGL(U)}(F\cdot \tau)}_{=\phi_{\embGL(U)}(\tau)}
%	{}\cdot
%	\phi_F(\tau)
%	\eas
%	for all~$\tau\in\HH_\genus$.
%%	
%%	We denote by~$\phi_{M\,\hat{}}\colon\HH_\genus\to\CC$ the branch of the square-root chosen such that
%%	\[
%%	\phi_{M\,\hat{}}(\tau)
%%	=
%%	\phi_M(\tau_4^{(\genus-r)})
%%	\qquad\text{for all~$M\in\Gamma^{(\genus-r)}$},
%%	\]
%%	where~$\hat{}\colon \Gamma^{(\genus-r)}\to\mathfrak{F}_{\genus,\genus-r}$ is the natural isomorphism.
%	By Lemma~\ref{lemma;form(11)bövv} we may rewrite~$\espart{k}{\genus}{r}$ as
%	\bas
%	\espart{k}{\genus}{r}(\tau)
%	&=
%	\sum_{N\in (\mathfrak{E}_{\genus,r})_\infty\backslash\mathfrak{E}_{\genus,r}^*}
%	\sum_{U\in\GL(\genus,\ZZ)_r\backslash\GL_r(\genus,\ZZ)^*}
%	\sum_{F\in(\mathfrak{F}_{\genus,\genus-r})_\infty\backslash\mathfrak{F}_{\genus,\genus-r}}
%	\phi_{N\embGL(U)F}(\tau)^{-2k}
%	\\
%	&\quad\times
%	\weil{\lattice}{\genus}(\gamma_{N\embGL(U)F})^{-1}\mathfrak{e}_0
%	\\
%	&=
%	\sum_{U\in\GL(\genus,\ZZ)_r\backslash\GL_r(\genus,\ZZ)^*}
%	\sum_{F\in(\mathfrak{F}_{\genus,\genus-r})_\infty\backslash\mathfrak{F}_{\genus,\genus-r}}
%	\phi_{\embGL(U)}(\tau)^{-2k}
%	\phi_F(\tau)^{-2k}
%	\weil{\lattice}{\genus}(\gamma_F)^{-1}
%	\weil{\lattice}{\genus}(\gamma_{\embGL(U)})^{-1}
%	\\
%	&\quad\times
%	\sum_{N\in (\mathfrak{E}_{\genus,r})_\infty\backslash\mathfrak{E}_{\genus,r}^*}
%	\phi_N(\embGL(U)F\cdot\tau)^{-2k}
%	\weil{\lattice}{\genus}(\gamma_N)^{-1}
%	\mathfrak{e}_0
%	.
%	\eas
%		In~\cite[Section~$3$, p.~$27$]{doubling} the natural isomorphism~$\hat{}\colon \Gamma^{(\genus-r)}\to\mathfrak{F}_{\genus,\genus-r}$ is introduced.
%	Such a map induces an isomorphism on the metaplectic coverings, which is~$\iota(I_{2r},\cdot)$, with~$\iota$ the inclusion as in~\eqref{eq;inclprodsimpl}.
%	By Lemma~\ref{lemma;proofweilreponsubmeta} we have that
%	\bas
%	\sum_{N\in(\mathfrak{E}_{\genus,r})_\infty\backslash \mathfrak{E}_{\genus,r}^*}
%	\phi_N(\tau)^{-2k}
%	\weil{\lattice}{\genus}(\gamma_N)^{-1}
%	\mathfrak{e}_0
%	=
%	\Big(\underbrace{\sum_{R\in\Gamma_\infty^{(r)}\backslash\Gamma^{(r,r)}_*}
%	\phi_R(\tau_1^{(r)})^{-2k}
%	\weil{\lattice}{r}(\gamma_R)^{-1}
%	\zeroe{r}}_{\eqqcolon \mathfrak{G}_{\genus,r}^k(\tau)}
%	\Big)\otimes \zeroe{\genus-r}.
%	\eas
%	Here we used that the~$\gamma_R=(R,\phi_r)\in\Mp_{2r}(\ZZ)$ map bijectively to the elements~$\iota(\gamma_R,I_{2(\genus-r)})$, where the latter are of the form~$\gamma_N$ with~$N\in(\mathfrak{E}_{\genus,r})_\infty\backslash \mathfrak{E}_{\genus,r}^*$.
%	We may then use again Lemma~\ref{lemma;proofweilreponsubmeta} to compute
%	\ba\label{eq;tbaFexpofmathfrakGvv2}
%	\espart{k}{\genus}{r}(\tau)
%	&=
%	\sum_{U\in\GL(\genus,\ZZ)_r\backslash\GL_r(\genus,\ZZ)^*}
%	\sum_{M\in \Gamma^{(\genus-r)}_\infty\backslash\Gamma^{(\genus-r)}}
%	\phi_{\embGL(U)}(\tau)^{-2k}
%	\phi_M(\tau_4)^{-2k}
%	\\
%	&\quad\times
%	\weil{\lattice}{\genus}\big(
%	\iota(I_{2r},\gamma_M)
%	\big)^{-1}
%	\weil{\lattice}{\genus}(\gamma_{\embGL(U)})^{-1}
%	\mathfrak{G}^k_{\genus,r}\big(\embGL(U)\iota(I_{2r},M)\cdot\tau\big)
%	\otimes
%	\zeroe{\genus-r}.
%	\ea
%	Since~$\Gamma_*^{(r,r)}=ST(r,r)$ and~$\Gamma_\infty^{(r)}$ is the Siegel parabolic, we deduce that~$\mathfrak{G}^k_{\genus,r}(\tau)=\Psi_{r,r}^k(\tau_1^{(r)})$, hence by Lemma~\ref{lemma;onPsivv} we have
%	\[
%	\mathfrak{G}_{\genus,r}^k(\tau)
%	=
%	\sum_{\discel\in D_\lattice^r}
%	\sum_{\substack{T\in q(\discel)+\halfint{r} \\ T>0}} \coeis{k}{r}{T}{\discel}
%	\cdot
%	e(\tr T \tau_1^{(r)}) \mathfrak{e}_\discel.
%	\]
%	We may then substitute the Fourier expansion of~$\mathfrak{G}_{\genus,r}^k$ in~\eqref{eq;tbaFexpofmathfrakGvv2} and, assuming there is no convergence issue, deduce that
%	\bas
%	\espart{k}{\genus}{r}(\tau)
%	&=
%	\sum_{\discel\in D_\lattice^r}
%	\sum_{\substack{T\in q(\discel) + \halfint{r} \\ T>0}}
%	\sum_{U\in\GL(\genus,\ZZ)_r\backslash\GL_r(\genus,\ZZ)^*}
%	\sum_{M\in \Gamma^{(\genus-r)}_\infty\backslash\Gamma^{(\genus-r)}}
%	a^k_r(T,\discel)
%	\phi_{\embGL(U)}(\tau)^{-2k}
%	\phi_M(\tau_4)^{-2k} 
%	\\
%	&\quad\times
%	e\left(
%	\tr T \big(
%	\embGL(U) \iota(I_{2r},M) \cdot \tau
%	\big)_1^{(r)}
%	\right)
%	\weil{\lattice}{\genus}\big(
%	\iota(I_{2r},\gamma_M)
%	\big)^{-1}
%	\weil{\lattice}{\genus}(\gamma_{\embGL(U)})^{-1}
%	\mathfrak{e}_\discel
%	\otimes
%	\zeroe{\genus-r}
%	\\
%	&=
%	\sum_{\discel\in D_\lattice^r}
%	\sum_{\substack{T\in q(\discel) + \halfint{r} \\ T>0}}
%	\sum_{\substack{U\in\GL(\genus,\ZZ)_r\backslash\GL_r(\genus,\ZZ)^* \\ U^{-1}=\big(
%	\begin{smallmatrix}
%	w_1^{(r)} & \ast \\
%	w_3 & \ast
%	\end{smallmatrix}
%	\big)}}
%	\sum_{M\in \Gamma^{(\genus-r)}_\infty\backslash\Gamma^{(\genus-r)}}
%	a^k_r(T,\discel)
%	\phi_{\embGL(U)}(\tau)^{-2k}
%	\phi_M(\tau_4)^{-2k}
%	\\
%	&\quad\times 
%	e\left(
%	\tr T
%	\big(\iota(I_{2r},M)\cdot \tau\big)
%	\left[\left(\begin{smallmatrix}
%	w_1^{(r)}
%	\\
%	w_3
%	\end{smallmatrix}\right)\right]
%	\right)
%	\weil{\lattice}{\genus}\big(
%	\iota(I_{2r},\gamma_M)
%	\big)^{-1}
%	\weil{\lattice}{\genus}(\gamma_{\embGL(U)})^{-1}
%	\mathfrak{e}_\discel
%	\otimes
%	\zeroe{\genus-r},
%	\eas
%	where the last step can be checked easily with a direct computation.
%	We now compute~$\phi_{\embGL(U)}(\tau)^{-2k}
%	\weil{\lattice}{\genus}(\gamma_{\embGL(U)})^{-1}
%	\mathfrak{e}_\discel\otimes\zeroe{\genus-r}$ explicitly.
%	Since~$2k\equiv n-2$ mod~$4$, without loss of generality we may assume that~$\gamma_{\embGL(U)}$ is given by choosing the principal branch of the square root.
%	Hence~$\phi_{\embGL(U)}(\tau)=\det U^{1/2}$ and
%	\begin{align*}
%	\weil{\lattice}{\genus}(\gamma_{\embGL(U)})^{-1}
%	\mathfrak{e}_\discel\otimes\zeroe{\genus-r}
%	&=
%	\weil{\lattice}{\genus}\big(\big(\begin{smallmatrix}
%	U^{-t} & \\
%	 & U
%	\end{smallmatrix}\big),\det U^{1/2}\big)^{-1}
%	\mathfrak{e}_\discel\otimes\zeroe{\genus-r}
%	\\
%	&=
%	\det U^{(n-2)/2}\mathfrak{e}_{(\discel,0,\dots,0) U^{-t}}.
%	\end{align*}
%	If we write~$U^{-1} =\left(\begin{smallmatrix}
%	w_1 & \ast \\
%	w_3 & \ast
%	\end{smallmatrix}\right)$ with~$w_1\in\ZZ^{(r,r)}$, then
%	\[
%	(\discel,0,\dots,0)U^{-t}
%	=
%	(\discel w_1^t , \discel w_3^t),
%	\]
%	from which we deduce that
%	\bas
%	\phi_{\embGL(U)}(\tau)^{-2k}
%	\weil{\lattice}{\genus}(\gamma_{\embGL(U)})^{-1}
%	\mathfrak{e}_\discel\otimes\zeroe{\genus-r}
%	&=
%	\underbrace{\det U^{-k}
%	\det U ^{(n-2)/2}}_{=1}
%	\mathfrak{e}_{\discel w_1^t}
%	\otimes\mathfrak{e}_{\discel w_3^t},
%	\eas
%	where again we used that~$2k\equiv n-2$ mod~$4$.
%	\\\textcolor{blue}{\textbf{Note:} Here we really needed that~$2k\equiv n-2$ mod~$4$, and not just~$2k\equiv n-2$ mod~$2$.}}
%	
%	\textcolor{blue}{The previous series converges absolutely and uniformly on vertical strips of positive heights, as one can check in the same way as in~\cite[p.~$28$]{doubling}.
%	The idea is to first reduce the problem to the convergence of the series at block diagonal elements in~$\HH_\genus$ of the form~$\big(\begin{smallmatrix}
%	\tau_1 & 0\\
%	0 & \tau_4
%	\end{smallmatrix}\big)$, with~$\tau_1\in \HH_r$.
%	This can be done using the fact that on every compact subset~$K$ of~$\HH_\genus$ there exists a positive constant~$c(K)$ such that
%	\[
%	\Im(M\cdot\tau)\geq c(K)\cdot\Im\big(M\cdot\big(\begin{smallmatrix}
%	\tau_1 & 0\\
%	0 & \tau_4
%	\end{smallmatrix}\big)\big)
%	\qquad
%	\text{for every~$\tau=\big(\begin{smallmatrix}
%	\tau_1 & \tau_2\\
%	\tau_2^t & \tau_4
%	\end{smallmatrix}\big)\in K$.}
%	\]
%	Then one can use the good convergence properties of Siegel Eisenstein series to study the norm of~$\mathfrak{G}_{\genus,r}^k(\tau)$, whose value actually depend on the top-left~$r\times r$ block of~$\tau$, from which one deduces the convergence of the series above.}
%	\end{proof}
%	
%	\textcolor{blue}{We say that a matrix~$A\in\ZZ^{(\genus,r)}$ is \emph{primitive} if~$A$ can be completed to an element of~$\GL(\genus,\ZZ)$.
%	Let
%	\[
%	M^\genus_r(\ZZ)^*
%	=
%	\{
%	W\in\ZZ^{(r,r)} : \text{$\det W\neq 0$ and there is~$V\in \GL(\genus,\ZZ)$ with~$V_1^{(r)}=W$}
%	\}.
%	\]
%	\begin{prop}\label{prop;frombösatz7vv}
%	If~$k>\genus+1$ and~$0\leq r\leq\genus$, then
%	\bas
%	\espart{k}{\genus}{r}(\tau)
%	&=
%	\sum_{\discel\in D_\lattice^r}
%	\sum_{\substack{T\in q(\discel) + \halfint{r} \\ T>0}}
%	\sum_{w_1\in M^\genus_r(\ZZ)^*/\GL(r,\ZZ)}
%	\sum_{\substack{w_3\in\ZZ^{\genus-r,r} \\ \left(\begin{smallmatrix}
%	w_1 \\ w_3
%	\end{smallmatrix}\right)\text{ primitive}}}
%	\sum_{M\in \Gamma^{(\genus-r)}_\infty\backslash\Gamma^{(\genus-r)}}
%	a^k_r(T,\discel)
%	\phi_M(\tau_4)^{-2k}
%	\\
%	&\quad\times 
%	e\left(
%	\tr T
%	\big(\iota(I_{2r},M)\cdot \tau\big)
%	\big[\big(\begin{smallmatrix}
%	w_1
%	\\
%	w_3
%	\end{smallmatrix}\big)\big]
%	\right)
%	\mathfrak{e}_{\discel w_1^t}
%	\otimes
%	\big(\weil{\lattice}{\genus-r}(\gamma_M)^{-1}
%	\mathfrak{e}_{\discel w_3^t}\big),
%	\eas
%	where the notation is the same as in Lemma~\ref{lemma;fbiamttle}.
%	\end{prop}
%	\begin{proof}
%	This is an immediate consequence of Lemma~\ref{lemma;fbiamttle} and~\cite[Lemma~$6$]{doubling}.
%	In fact, in the latter it is given a system of representatives of~$\GL(\genus,\ZZ)_r\backslash\GL_r(\genus,\ZZ)^*$ in terms of representatives of~$M^\genus_r(\ZZ)^*/\GL(r,\ZZ)$ and~$\ZZ^{\genus-r,r}$.
%	\end{proof}}
%	
%	
%	
%	
%	
%	
%	
%	
%	
%	
%	
%	
%	
%	
%	
%	
%	
%	
%	We now construct certain generalizations of the Poincaré series introduced in Section~\ref{sec;poincser}.
	Let~$\widetilde{\mathfrak{A}}_{\genus,r}$ be the metaplectic cover of
	\begin{align*}
	\mathfrak{A}_{\genus,r}
	&=
	\left\{
	M=\big(
	\begin{smallmatrix}
	A & B \\
	C & D
	\end{smallmatrix}
	\big)\in\Sp_{2\genus}(\ZZ) : \text{$C=0$, $D=\big(\begin{smallmatrix}
	I_r & \ast \\
	0 & \ast
	\end{smallmatrix}\big)$, $A_1^{(r)}=I_r$}
	\right\}.
%	\\
%	\textcolor{blue}{\mathfrak{A}_{\genus,r}'}
%	&=
%	\textcolor{blue}{\left\{
%	M=\big(
%	\begin{smallmatrix}
%	A & B \\
%	C & D
%	\end{smallmatrix}
%	\big)\in\mathfrak{C}_{\genus,r} : \text{$C_1^{(r)}=0$, $A_1^{(r)}=D_1^{(r)}=\pm I_r$}
%	\right\}.}
	\end{align*}
%	\\\textcolor{blue}{\textbf{Remark:} The group~$\mathfrak{A}_{\genus,r}'$ is not the one we use to construct Poincaré series, it is too large.
%	Instead, we use~$\mathfrak{A}_{\genus,r}$.
%	However, we still introduce~$\mathfrak{A}_{\genus,r}'$ since it is the one used by Böcherer, and he studies certain representatives of quotients under~$\mathfrak{A}_{\genus,r}'$.
%	We will need them later.}
	\begin{defi}
	Let~$\discel\in D_\lattice^r$, and let~$t\in q(\discel) + \halfint{r}$ be a positive definite matrix.
	We define~$T\coloneqq\big(\begin{smallmatrix}
	t & 0\\
	0 & 0
	\end{smallmatrix}\big)\in q(\discel,0,\dots,0) + \halfint{\genus}$.
	For every~$k>\genus+r+1$, the \emph{generalized Poincaré series of exponential type}~$\mathfrak{g}_{\genus,r}^k(T,\discel)$ is
	\[
	\mathfrak{g}_{\genus,r}^k(\tau,T,\discel)
	\coloneqq
	\frac{1}{2}
	\sum_{\gamma\in\widetilde{\mathfrak{A}}_{\genus,r}\backslash\Mp_{2\genus}(\ZZ)}
	e(\tr T \gamma\cdot\tau)
	\phi(\tau)^{-2k}
	\weil{\lattice}{\genus}(\gamma)^{-1}
	\mathfrak{e}_\discel
	\otimes
	\zeroe{\genus-r},
	\qquad
	\tau\in\HH_\genus.
	\]
	\end{defi}
%	\textcolor{blue}{\textbf{Check that it is well-defined:} We want to show that if~$T=\big(\begin{smallmatrix}
%	t & 0\\
%	0 & 0
%	\end{smallmatrix}\big)$ and~$\gamma\in\widetilde{\mathfrak{A}}_{\genus,r}$, then
%	\be\label{eq;welldefgenPoincBö}
%	e(\tr T \gamma\cdot\tau)\cdot
%	\phi(\tau)^{-2k}
%	\weil{\lattice}{\genus}(\gamma)^{-1}
%	\mathfrak{e}_\discel
%	\otimes
%	\zeroe{\genus-r}
%	=
%	e(\tr T \tau)
%	\mathfrak{e}_\discel
%	\otimes
%	\zeroe{\genus-r}.
%	\ee
%	Decomposing the matrices in the usual way, so that the top-left block is a~$r\times r$ matrix, it is easy to see that
%	\bas
%	\gamma\cdot\tau
%	&=
%	\left(
%	\left(\begin{smallmatrix}
%	I_r & 0\\
%	A_3 & A_4
%	\end{smallmatrix}
%	\right)
%	\left(\begin{smallmatrix}
%	\tau_1 & \tau_2\\
%	\tau_2^t & \tau_4
%	\end{smallmatrix}
%	\right) + B\right)
%	\left(\begin{smallmatrix}
%	I_r & D_2\\
%	0 & D_4
%	\end{smallmatrix}
%	\right)^{-1}
%	=
%	\left(
%	\left(\begin{smallmatrix}
%	\tau_1 & \ast\\
%	\ast & \ast
%	\end{smallmatrix}
%	\right)
%	+
%	\left(\begin{smallmatrix}
%	B_1 & B_2\\
%	B_3 & B_4
%	\end{smallmatrix}
%	\right)\right)
%	\left(\begin{smallmatrix}
%	I_r & \ast\\
%	0 & \ast
%	\end{smallmatrix}
%	\right)
%	\\
%	&=
%	\left(\begin{smallmatrix}
%	\tau_1 + B_1 & \ast\\
%	\ast & \ast
%	\end{smallmatrix}
%	\right)
%	\left(\begin{smallmatrix}
%	I_r & \ast\\
%	0 & \ast
%	\end{smallmatrix}
%	\right)
%	=
%	\left(\begin{smallmatrix}
%	\tau_1 + B_1 & \ast\\
%	\ast & \ast
%	\end{smallmatrix}
%	\right),
%	\eas
%	hence
%	\[
%	e(\tr T\gamma\cdot \tau)
%	=
%	e\Big(\tr
%	\big(\begin{smallmatrix}
%	t & 0\\
%	0 & 0
%	\end{smallmatrix}\big)
%	\big(\begin{smallmatrix}
%	\tau_1 + B_1 & \ast\\
%	\ast & \ast
%	\end{smallmatrix}\big)
%	\Big)
%	=
%	e(\tr t\tau_1)\cdot\underbrace{e(\tr t B_1)}_{=e(\tr q(\discel) B_1)}
%	=
%	e(\tr T\tau)
%	\cdot
%	e(\tr q(\discel) B_1).
%	\]
%	To prove~\eqref{eq;welldefgenPoincBö} it is then enough to show that
%	\be\label{eq;welldefgenPoincBö2}
%	\phi(\tau)^{-2k}
%	\weil{\lattice}{\genus}(\gamma)^{-1}
%	\mathfrak{e}_\discel
%	\otimes
%	\zeroe{\genus-r}
%	=
%	e(-\tr q(\discel) B_1)
%	\mathfrak{e}_\discel
%	\otimes
%	\zeroe{\genus-r}.
%	\ee
%	The idea here is the same as for the check that Klingen Eisenstein series are well-defined.
%	Namely, we check~\eqref{eq;welldefgenPoincBö2} on each of the standard generators of~$\widetilde{\mathfrak{A}}_{\genus,r}$ arising as in Lemma~\ref{lemma;genklingenpar}, hence we only choose the generators given in that lemma that lie in~$\widetilde{\mathfrak{A}}_{\genus,r}$.}
%	
%	\textcolor{blue}{
%	If~$\gamma=(M,\phi_M)$ is as in Lemma~\ref{lemma;genklingenpar}~\ref{item;genklingenpar1}, then it must be such that
%	\[
%	M=\left(\begin{smallmatrix}
%	I_r & & B_1  &\\
%	 & U & & \\
%	 & & I_r    & \\
%	 &  & & U^{-t}
%	\end{smallmatrix}\right)
%	=
%	\iota\Big(\underbrace{\big(\begin{smallmatrix}
%	I_r & B_1\\
%	  & I_r
%	\end{smallmatrix}\big)}_{\eqqcolon M'},\big(\begin{smallmatrix}
%	U & \\
%	 & U^{-t}
%	\end{smallmatrix}\big)\Big).
%	\]
%	Let~$\phi_{M'}$ and~$\phi_U$ be the square roots of respectively~$\det M'$ and~$\det U$ chosen such that~$\phi_M\big(\begin{smallmatrix}
%	\tau_1 & \\
%	 & \tau_4
%	\end{smallmatrix}\big)=\phi_{M'}(\tau_1)\phi_U(\tau_4)$.
%	Let~$\gamma'\coloneqq(M',\phi_{M'})$ and~$\gamma_U\coloneqq\big(\big(
%	\begin{smallmatrix}
%	U & \\
%	 & U^{-t}
%	\end{smallmatrix}
%	\big),\phi_U\big)$.
%	Then~$\gamma=\iota(\gamma',\gamma_U)$.
%	\\
%	Without loss of generality we may assume that~$\phi_U=\det U^{1/2}$.
%	By Lemma~\ref{lemma;proofweilreponsubmeta} and~\eqref{eq;howto-phi} we may compute
%	\bas
%	\phi(\tau)^{-2k} & \weil{\lattice}{\genus}(\gamma)^{-1}
%	\mathfrak{e}_\discel\otimes\zeroe{\genus-r}
%	\\
%	&=
%	\phi(\tau)^{-2k}
%	\big(\weil{\lattice}{r}(\gamma')^{-1}\mathfrak{e}_\discel\big)
%	\otimes
%	\big(
%	\weil{\lattice}{\genus-r}(\gamma_U)^{-1}
%	\zeroe{\genus-r}
%	\big)
%	\\
%	&=
%	(\pm 1)^{-2k}
%	\det U^{-k}
%	\Big(\weil{\lattice}{r}\big(\big(
%	\begin{smallmatrix}
%	I_r & B_1 \\
%	 & I_r
%	\end{smallmatrix}\big),\pm 1\big)^{-1}\mathfrak{e}_\discel\Big)
%	\otimes
%	\Big(
%	\weil{\lattice}{\genus-r}\big(\big(
%	\begin{smallmatrix}
%	U & \\
%	 & U^{-1}
%	\end{smallmatrix}
%	\big),\det U^{1/2}\big)^{-1}
%	\zeroe{\genus-r}
%	\Big)
%	\\
%	&=
%	\underbrace{(\pm 1)^{-2k+n-2}}_{=1}
%	\underbrace{\det U^{-k+(n-2)/2}}_{=1}
%	\Big(\weil{\lattice}{r}\big(\big(
%	\begin{smallmatrix}
%	I_r & B_1 \\
%	 & I_r
%	\end{smallmatrix}\big),1\big)^{-1}\mathfrak{e}_\discel\Big)\otimes\zeroe{\genus-r}
%	\\
%	&=
%	e(-\tr q(\discel)B_1)\mathfrak{e}_\discel\otimes \zeroe{\genus-r}.
%	\eas
%	This proves~\eqref{eq;welldefgenPoincBö2}.
%	Note that we used that~$2k\equiv n-2$ mod~$4$, and not just~$2k\equiv n-2$ mod~$2$.
%	}
	%
%	\textcolor{blue}{If~$\gamma$ is as in Lemma~\ref{lemma;genklingenpar}~\ref{item;genklingenpar2}, then~$B_1=0$, hence~\eqref{eq;welldefgenPoincBö2} simplifies to
%	\[
%	\phi(\tau)^{-2k}
%	\weil{\lattice}{\genus}(\gamma)^{-1}
%	\mathfrak{e}_\discel
%	\otimes
%	\zeroe{\genus-r}
%	=
%	\mathfrak{e}_\discel
%	\otimes
%	\zeroe{\genus-r}.
%	\]
%	This can be checked exactly in the same way as for the analogous check on Klingen Eisenstein series.
%	}
	
	If~$r=\genus$, then~$\mathfrak{g}_{\genus,\genus}^k(T,\discel)$ equals the Poincaré series~$\poincv{\gengenus}{\discel}{T}{\lattice}$ introduced in Section~\ref{sec;poincser}.
	We also remark that the generalized Poincaré series~$\mathfrak{g}_{\genus,r}^k(T,\discel)$ is actually the Klingen Eisenstein series arising from the Poincaré series~$\poincv{r}{\discel}{t}{\lattice}$, namely
	\[
	\mathfrak{g}_{\genus,r}^k(\tau,T,\discel)
	=
	\klinser{k}{\gengenus}{r}{\lattice}(\tau,\poincv{r}{\discel}{t}{\lattice}).
	\]
%	\textcolor{blue}{\textbf{Check of the previous equality.}
%	\bas
%	\klinser{k}{\gengenus}{r}{\lattice}(\tau,\poincv{\gengenus}{\discel}{t}{\lattice})
%	&=
%	\sum_{\gamma\in\klimet{\gengenus}{r}\backslash \Mp_{2\gengenus}(\ZZ)}
%	\phi(\tau)^{-2k}
%	\weil{\lattice}{\gengenus}(\gamma)^{-1}
%	\Big(\poincv{r}{\discel}{t}{\lattice}(\gamma^*\cdot\tau^*)
%	\otimes \zeroe{\gengenus-r}\Big)
%	\\
%	&=
%	\frac{1}{2}
%	\sum_{\gamma\in\klimet{\gengenus}{r}\backslash \Mp_{2\gengenus}(\ZZ)}
%	\sum_{\gamma'\in\mettrans{r}\backslash\Mp_{2r}(\ZZ)}
%	\phi(\tau)^{-2k}
%	\phi'(\gamma^*\cdot\tau^*)^{-2k}
%	e\big(\tr(t\gamma'\gamma^*\cdot\tau^*)\big)
%	\\
%	&\quad\times
%	\weil{\lattice}{\gengenus}(\gamma)^{-1}
%	\Big(\weil{\lattice}{r}(\gamma')^{-1}\mathfrak{e}_\discel
%	\otimes \zeroe{\gengenus-r}\Big),
%	\eas
%	where we write~$\gamma=(M,\phi)$ and~$\gamma'=(M',\phi')$.
%	Let~$\iota$ be the inclusion~\eqref{eq;inclprodsimpl}.
%	As usual, we identify~$I_{2(\genus-r)}$ with its standard inclusion in~$\Mp_{2(\genus-r)}(\ZZ)$.
%	It is easy to see that if we write
%	\[
%	\iota(\gamma', I_{2(\genus-r)})
%	=
%	\big(
%	\iota(M',I_{2(\genus-r)}),\tilde{\phi}
%	\big),
%	\qquad
%	\text{then~$\tilde\phi(\tau)=\phi'(\tau^*)$.}
%	\]
%	Furthermore, since~$T\coloneqq\big(\begin{smallmatrix}
%	t & 0\\
%	0 & 0
%	\end{smallmatrix}\big)$, we have
%	\[
%	\tr(t\gamma'\gamma^*\cdot\tau^*)
%	=
%	\tr\big(t\gamma'\cdot(\gamma\cdot\tau)^*\big)
%	=
%	\tr\big(
%	T\iota(\gamma',I_{2(\genus-r)})\gamma\cdot\tau
%	\big).
%	\]
%	By Lemma~\ref{lemma;proofweilreponsubmeta} we may then rewrite
%	\bas
%	\klinser{k}{\gengenus}{r}{\lattice}(\tau,\poincv{\gengenus}{\discel}{t}{\lattice})
%	&=
%	\frac{1}{2}
%	\sum_{\gamma\in\klimet{\gengenus}{r}\backslash \Mp_{2\gengenus}(\ZZ)}
%	\sum_{\gamma'\in\mettrans{r}\backslash\Mp_{2r}(\ZZ)}
%	\tilde\phi(\gamma\cdot\tau)^{-2k}
%	\phi(\tau)^{-2k}
%	e\Big(\tr\big(
%	T\iota(\gamma',I_{2(\genus-r)})\gamma\cdot\tau
%	\big)\Big)
%	\\
%	&\quad\times
%	\weil{\lattice}{\gengenus}\big(\iota(\gamma',I_{2(\genus-r)})\gamma\big)^{-1}
%	\mathfrak{e}_\discel
%	\otimes \zeroe{\gengenus-r}.
%	\eas
%	We remark that
%	\[
%	\iota(\gamma',I_{2(\genus-r)})\cdot\gamma
%	=
%	\big(
%	\iota(M',I_{2(\genus-r)}),\tilde\phi(\tau)
%	\big)
%	\cdot
%	\big(
%	\iota(M',I_{2(\genus-r)})\cdot M , \tilde\phi(\gamma\cdot\tau)\cdot\phi(\tau)
%	\big),
%	\]
%	that the inclusion~$\iota(\cdot,I_{2(\genus-r)})\colon\Mp_{2r}(\ZZ)\to\Mp_{2\genus}(\ZZ)$ lands in~$\klimet{\gengenus}{r}$, and that it maps~$\mettrans{r}$ inside~$\widetilde{\mathfrak{A}}_{\genus,r}$.
%	\\
%	As remarked e.g.\ in~\cite[Section~$7$, p.~$93$]{klingen}, if~$\gamma'$ runs in a complete set of representatives of~$\trans{r}\backslash\Sp_{2r}(\ZZ)$, then~$\iota(\gamma',I_{2(\genus-r)})$ runs in a complete set of representatives of~$\mathfrak{A}_{\genus,r}\backslash\klipar{\genus}{r}$.
%	Since if~$\gamma$ runs in a complete set of representatives of~$\klimet{\gengenus}{r}\backslash \Mp_{2\gengenus}(\ZZ)$ and~$\gamma'$ runs in a complete set of representatives of~$\mettrans{r}\backslash\Mp_{2r}(\ZZ)$, then~$\iota(\gamma',I_{2(\genus-r)})\cdot\gamma$ provide a complete set of representatives of~$\widetilde{\mathfrak{A}}_{\genus,r}\backslash\Mp_{2\genus}(\ZZ)$, we deduce that
%	\bas
%	\klinser{k}{\gengenus}{r}{\lattice}(\tau,\poincv{\gengenus}{\discel}{t}{\lattice})
%	&=
%	\frac{1}{2}
%	\sum_{\gamma\in\widetilde{\mathfrak{A}}_{\genus,r}\backslash\Mp_{2\genus}(\ZZ)}
%	\phi(\tau)^{-2k}
%	e\big(
%	\tr T \gamma\cdot\tau
%	\big)
%	\weil{\lattice}{\genus}(\gamma)^{-1}
%	\mathfrak{e}_\discel\otimes\zeroe{\genus-r}
%	=
%	\mathfrak{g}_{\genus,r}^k(\tau,T,\discel).
%	\eas
%	}
	Let
	\bas
	\ZZ^{(\genus-r,r)}_s
	&\coloneqq
	\{
	w_3\in\ZZ^{(\genus-r,r)}
	:
	\rk(w_3)=s
	\},
	\\
	\ZZ^{(\genus-r,r)}_{s,0}
	&\coloneqq
	\left\{
	w_3 \in\ZZ^{(\genus-r,r)}_s
	:
	w_3=\left(\begin{smallmatrix}
	(w_3)_1^{(s)} & \ast
	\\
	0 & 0_{\genus-r-s,r-s}
	\end{smallmatrix}\right)
	\right\},
	\\
	\GL(\genus,\ZZ)_\nu
	&\coloneqq
	\{
	U\in\GL(\genus,\ZZ) : U_3^{(\genus-\nu,\nu)}=0
	\},
%	\\
%	\textcolor{blue}{\GL(\genus,\ZZ)_\nu^1}\,
%	&\textcolor{blue}{\coloneqq
%	\{
%	U\in\GL(\genus,\ZZ)_\nu : U_1^{(\nu)}=I_\nu
%	\},}
	\\
	M^\genus_r(\ZZ)^*
	&\coloneqq
	\{
	W\in\ZZ^{(r,r)} : \text{$\det W\neq 0$ and there is~$V\in \GL(\genus,\ZZ)$ with~$V_1^{(r)}=W$}
	\}.
	\eas
%	\textcolor{blue}{\textbf{Important: }The definition of~$\ZZ^{(\genus-r,r)}_{s,0}$ above is NOT the one appearing in~\cite{doubling}.
%	In fact, I think that Böcherer made a typo there, since he defined~$\ZZ^{(\genus-r,r)}_{s,0}$ as
%	\[
%	\left\{
%	w_3 \in\ZZ^{(\genus-r,r)}_s
%	:
%	w_3=\left(\begin{smallmatrix}
%	(w_3)_1^{(\genus-r)} & \ast
%	\\
%	0 & 0_{\genus-r-s}
%	\end{smallmatrix}\right)
%	\right\},
%	\]
%	which does not make sense, since~$w_3$ would be a square matrix in~$\ZZ^{(2(\genus-r),s)}$ and not in~$\ZZ^{(\genus-r,r)}$.}
	
	We are now ready to state the vector-valued generalization of Böcherer's result on diagonal restrictions of principal parts of Siegel Eisenstein series, which can be proved following the same wording of~\cite[Satz~$12$]{doubling}.
	\begin{thm}[Böcherer]\label{thm;genBöcherervv}
	If~$k>\genus+1$, $\tau_1\in\HH_r$ and~$\tau_4\in\HH_{\genus-r}$, then
	\bas
	\espart{k}{\genus}{r}
	\big(\begin{smallmatrix}
	\tau_1 & 0\\
	0 & \tau_4
	\end{smallmatrix}\big)
	&=
	\sum_{\discel\in D_\lattice^r}
	\sum_{\substack{T\in q(\discel) + \halfint{r} \\ T>0}}
	\sum_{w_1\in M^\genus_r(\ZZ)^*/\GL(r,\ZZ)}
	a_r^k(T,\discel)
	\cdot e\big(
	\tr T[w_1^t]\tau_1
	\big)
	\\
	&\quad\times
	\sum_{s=0}^{\min\{r,\genus-r\}}
	\varepsilon(s)
	\sum_{\substack{w_3\in\GL(\genus-r,\ZZ)_s\backslash\ZZ^{(\genus-r,r)}_{s,0} \\ \left(\begin{smallmatrix}
	w_1 \\ w_3
	\end{smallmatrix}\right)\text{ primitive}}}
	\mathfrak{e}_{\discel w_1^t}
	\otimes
	\mathfrak{g}_{\genus-r,s}^k\big(\tau_4,T[w_3^t],\discel(w_3^t)_{(s)}\big),
	\eas
	where~$(w_3^t)_{(s)}$ denotes the matrix given by the first~$s$ columns of~$w_3^t$ and
	\[
	\varepsilon(s)=\begin{cases}
	1, & \text{if $s=0$,}\\
	2, & \text{if $s>0$.}
	\end{cases}\]
	\end{thm}
%	\textcolor{blue}{\textbf{Note:} Our~$\varepsilon(s)$ differs from the one of Böcherer.
%	In fact, I think he made a mistake.
%	I highlight where the factor~$\varepsilon(s)$ pops out inside the proof below.}
%	\begin{proof}
%	\textcolor{blue}{If~$w_3$ runs in~$\GL(\genus-r,\ZZ)_s\backslash\ZZ^{(\genus-r,r)}_{s,0}$ and if~$v$ runs in~$\GL(\genus-r,\ZZ)/\GL(\genus-r,\ZZ)^1_s$, then~$vw_3$ gives the whole~$\ZZ^{(\genus-r,r)}_s$.
%	The same works also under the additional condition that~$\big(\begin{smallmatrix}
%	w_1 \\ w_3
%	\end{smallmatrix}\big)$ is primitive.
%	Since~$\ZZ^{(\genus-r,r)}=\bigcup_{s=0}^{\min\{\genus-r,r\}} \ZZ^{(\genus-r,r)}_s$, we deduce by Proposition~\ref{prop;frombösatz7vv} that
%	\begin{align*}
%	\espart{k}{\genus}{r}(\tau)
%	&=
%%	\sum_{\discel\in D_\lattice^r}
%%	\sum_{\substack{T\in q(\discel) + \halfint{r} \\ T>0}}
%%	\sum_{w_1\in M^\genus_r(\ZZ)^*/\GL(r,\ZZ)}
%%	\sum_{\substack{w_3\in\ZZ^{\genus-r,r} \\ \left(\begin{smallmatrix}
%%	w_1 \\ w_3
%%	\end{smallmatrix}\right)\text{ primitive}}}
%%	\sum_{M\in \Gamma^{(\genus-r)}_\infty\backslash\Gamma^{(\genus-r)}}
%%	a^k_r(T,\discel)
%%	\phi_M(\tau_4)^{-2k}
%%	\\
%%	&\quad\times 
%%	e\left(
%%	\tr T
%%	\big(\iota(I_{2r},M)\cdot \tau\big)
%%	\big[\big(\begin{smallmatrix}
%%	w_1
%%	\\
%%	w_3
%%	\end{smallmatrix}\big)\big]
%%	\right)
%%	\mathfrak{e}_{\discel w_1^t}
%%	\otimes
%%	\big(\weil{\lattice}{\genus-r}(\gamma_M)^{-1}\mathfrak{e}_{\discel w_3^t}\big)
%%	\\
%%	&=
%	\sum_{\discel\in D_\lattice^r}
%	\sum_{\substack{T\in q(\discel) + \halfint{r} \\ T>0}}
%	\sum_{w_1\in M^\genus_r(\ZZ)^*/\GL(r,\ZZ)}
%	\sum_{s=0}^{\min\{r,\genus-r\}}
%	\sum_{\substack{w_3\in\ZZ^{(\genus-r,r)}_s \\ \left(\begin{smallmatrix}
%	w_1 \\ w_3
%	\end{smallmatrix}\right)\text{ primitive}}}
%	\sum_{M\in \Gamma^{(\genus-r)}_\infty\backslash\Gamma^{(\genus-r)}}
%	a^k_r(T,\discel)
%	\\
%	&\quad\times 
%	\phi_M(\tau_4)^{-2k}
%	e\left(
%	\tr T
%	\big(\iota(I_{2r},M)\cdot \tau\big)
%	\big[\big(\begin{smallmatrix}
%	w_1
%	\\
%	w_3
%	\end{smallmatrix}\big)\big]
%	\right)
%	\mathfrak{e}_{\discel w_1^t}
%	\otimes
%	\big(\weil{\lattice}{\genus-r}(\gamma_M)^{-1}
%	\mathfrak{e}_{\discel w_3^t}\big)
%	\\
%	&=
%	\sum_{\discel\in D_\lattice^r}
%	\sum_{\substack{T\in q(\discel) + \halfint{r} \\ T>0}}
%	\sum_{w_1\in M^\genus_r(\ZZ)^*/\GL(r,\ZZ)}
%	\sum_{s=0}^{\min\{r,\genus-r\}}
%	\sum_{\substack{w_3\in\GL(\genus-r,\ZZ)_s\backslash\ZZ^{(\genus-r,r)}_{s,0} \\ \left(\begin{smallmatrix}
%	w_1 \\ w_3
%	\end{smallmatrix}\right)\text{ primitive}}}
%	\\
%	&\quad\times
%	\sum_{v\in\GL(\genus-r,\ZZ)/\GL(\genus-r,\ZZ)_s^1}
%	\sum_{M\in \Gamma^{(\genus-r)}_\infty\backslash\Gamma^{(\genus-r)}}
%	a^k_r(T,\discel)
%	\phi_M(\tau_4)^{-2k}
%	\\
%	&\quad\times 
%	e\left(
%	\tr T
%	\big(\iota(I_{2r},M)\cdot \tau\big)
%	\big[\big(\begin{smallmatrix}
%	w_1
%	\\
%	v w_3
%	\end{smallmatrix}\big)\big]
%	\right)
%	\mathfrak{e}_{\discel w_1^t}
%	\otimes
%	\big(\weil{\lattice}{\genus-r}(\gamma_M)^{-1}
%	\mathfrak{e}_{\discel w_3^tv^t}\big)
%	\end{align*}
%	We remark that if~$M=\big(\begin{smallmatrix}
%	A & B\\
%	C & D
%	\end{smallmatrix}\big)\in\Gamma^{(\genus-r)}$,
%	then
%	\bas
%	\iota(I_{2r},M)\cdot\left(
%	\begin{smallmatrix}
%	\tau_1^{(r)} &  \\
%	 & \tau_4^{(\genus-r)}
%	\end{smallmatrix}\right)
%	&=
%	\left(
%	\begin{smallmatrix}
%	\tau_1^{(r)}  &   \\
%	 &     (A\tau_4^{(\genus-r)} + B)\cdot(C\tau_4^{(\genus-r)}+D)^{-1}
%	\end{smallmatrix}\right)
%	=
%	\left(
%	\begin{smallmatrix}
%	\tau_1^{(r)}  &   \\
%	 &     M\cdot\tau_4^{(\genus-r)}
%	\end{smallmatrix}\right),
%	\eas
%	and, moreover, that
%	\begin{align*}
%	\tr\left(
%	T \left(\begin{smallmatrix}
%	\tau_1^{(r)} & \\ & M\cdot\tau_4^{(\genus-r)}
%	\end{smallmatrix}\right)
%	\left[\left(
%	\begin{smallmatrix}
%	w_1 \\ vw_3
%	\end{smallmatrix}
%	\right)\right]
%	\right)
%	&=
%	\tr\big( T\tau_1^{(r)}[w_1]\big)
%	+
%	\tr\big(T M\cdot\tau_4^{(\genus-r)}[vw_3]\big)
%	\\
%	&=
%	\tr\big( T[w_1^t]\tau_1^{(r)}\big)
%	+
%	\tr\big(T M\cdot\tau_4^{(\genus-r)}[vw_3]\big).
%	\end{align*}
%	If~$\tau_1\in\HH_r$ and~$\tau_4\in\HH_{\genus-r}$, we may then compute
%	\begin{align*}
%	\espart{k}{\genus}{r}
%	\big(\begin{smallmatrix}
%	\tau_1 & \\
%	 & \tau_4
%	\end{smallmatrix}\big)
%	&=
%	\sum_{\discel\in D_\lattice^r}
%	\sum_{\substack{T\in q(\discel) + \halfint{r} \\ T>0}}
%	\sum_{w_1\in M^\genus_r(\ZZ)^*/\GL(r,\ZZ)}
%	a_r^k(T,\discel)
%	\cdot e\big(
%	\tr T[w_1^t]\tau_1
%	\big)
%	\sum_{s=0}^{\min\{r,\genus-r\}}
%	\\
%	&\quad\times
%	\sum_{\substack{w_3\in\GL(\genus-r,\ZZ)_s\backslash\ZZ^{(\genus-r,r)}_{s,0} \\ \left(\begin{smallmatrix}
%	w_1 \\ w_3
%	\end{smallmatrix}\right)\text{ primitive}}}
%	\sum_{v\in\GL(\genus-r,\ZZ)/\GL(\genus-r,\ZZ)_s^1}
%	\sum_{M\in\Gamma_\infty^{(\genus-r)}\backslash\Gamma^{(\genus-r)}}
%	\phi_M(\tau_4)^{-2k}
%	\\
%	&\quad\times
%	e\Big(
%	\tr TM\cdot\tau_4 [vw_3]
%	\Big)
%	\mathfrak{e}_{\discel w_1^t}
%	\otimes
%	\big(\weil{\lattice}{\genus-r}(\gamma_M)^{-1}
%	\mathfrak{e}_{\discel w_3^tv^t}\big).
%	\end{align*}
%	Let~$\phi_{\embGL(v^{-1})M}=\det v^{1/2} \phi_M$, so that
%	\[
%	\underbrace{\big( \embGL(v^{-1}),\det v^{1/2} \big)}_{\eqqcolon\gamma_{\embGL(v^{-1})}}
%	{}\cdot(M,\phi_M)
%	=
%	\underbrace{\big(
%	\embGL(v^{-1})M , \phi_{\embGL(v^{-1})M}
%	\big)}_{\eqqcolon \gamma_{\embGL(v^{-1})M}}.
%	\]
%	We remark that~$e\big(
%	\tr T (M\cdot\tau_4)[vw_3]
%	\big)
%	=
%	e\big(
%	\tr T [w_3^t] (\embGL(v^{-1}) M )\cdot\tau_4)
%	\big)$, and also
%	\bas
%	\phi_M(\tau_4)^{-2k} \mathfrak{e}_{\discel w_3^tv^t}
%	&=
%	\phi_M(\tau_4)^{-2k}
%	\weil{\lattice}{\genus-r}\big(\gamma_{\embGL(v^{-1})}\big)^{-1}
%	\weil{\lattice}{\genus-r}\big(\gamma_{\embGL(v^{-1})}\big)
%	\mathfrak{e}_{\discel w_3^tv^t}
%	\\
%	&=
%	\phi_{\embGL(v^{-1})M}(\tau_4)^{-2k}
%	\underbrace{\det v^{k + (2-n)/2}}_{=1}
%	\weil{\lattice}{\genus-r}\big(\gamma_{\embGL(v^{-1})}\big)^{-1}
%	\mathfrak{e}_{\discel w_3^t},
%	\eas
%	where we used that~$2k\equiv n-2$ mod~$4$.
%	\\\textcolor{blue}{\textbf{Note:} Here it is another instance where it is not enough to use the congruence mod~$2$ only.}\\
%	We then deduce that
%	\begin{align*}
%	\espart{k}{\genus}{r}
%	\big(\begin{smallmatrix}
%	\tau_1 & \\
%	 & \tau_4
%	\end{smallmatrix}\big)
%	&=
%	\sum_{\discel\in D_\lattice^r}
%	\sum_{\substack{T\in q(\discel) + \halfint{r} \\ T>0}}
%	\sum_{w_1\in M^\genus_r(\ZZ)^*/\GL(r,\ZZ)}
%	a_r^k(T,\discel)
%	\cdot e \big(
%	\tr T[w_1^t]\tau_1
%	\big)
%	\sum_{s=0}^{\min\{r,\genus-r\}}
%	\\
%	&\quad\times
%	\sum_{\substack{w_3\in\GL(\genus-r,\ZZ)_s\backslash\ZZ^{(\genus-r,r)}_{s,0} \\ \left(\begin{smallmatrix}
%	w_1 \\ w_3
%	\end{smallmatrix}\right)\text{ primitive}}}
%	\sum_{v\in\GL(\genus-r,\ZZ)/\GL(\genus-r,\ZZ)_s^1}
%	\sum_{M\in\Gamma_\infty^{(\genus-r)}\backslash\Gamma^{(\genus-r)}}
%	\phi_{\embGL(v^{-1})M}(\tau_4)^{-2k}
%	\\
%	&\quad\times
%	e\big(
%	\tr T [w_3^t] (\embGL(v^{-1}) M )\cdot\tau_4)
%	\big)
%	\mathfrak{e}_{\discel w_1^t}
%	\otimes
%	\big(\weil{\lattice}{\genus-r}(\underbrace{\gamma_{\embGL(v^{-1})}\gamma_M}_{=\gamma_{\embGL(v^{-1})M}})^{-1}
%	\mathfrak{e}_{\discel w_3^t}\big).
%	\end{align*}}
%	
%	\textcolor{blue}{If~$v$ runs in a system of representatives for~$\GL(\genus-r,\ZZ)\backslash\GL(\genus-r,\ZZ)^1_s$ and~$M$ runs in a system of representatives for~$\Gamma^{(\genus-r)}_\infty\backslash\Gamma^{(\genus-r)}$, then~$\embGL(v^{-1}) M$ gives system of representatives for~$\mathfrak{A}_{\genus-r,s}'\backslash\Sp_{2(\genus-r)}(\ZZ)$ if~$s=0$, and twice a set of representatives if~$s>0$.
%	\\\textcolor{blue}{\textbf{Caveat regarding the previous sentence:} The set~$\mathfrak{A}_{\genus-r,s}$ in these notes differs from its homonym in~\cite{doubling}.
%	In fact, we request the top-left block of~$A$ and~$D$ to be~$I_r$, not~$\pm I_r$.
%	This means that our set~$\mathfrak{A}_{\genus-r,s}$ is "a half" of the one of Böcherer.
%	In fact, it does not matter so much: The correction factor of~$1/2$ we need to use is exactly the one we used in our definition of generalized Poincaré series.}\\
%	This means that
%	\begin{align*}
%	\espart{k}{\genus}{r}
%	\big(\begin{smallmatrix}
%	\tau_1 & \\
%	 & \tau_4
%	\end{smallmatrix}\big)
%	&=
%	\sum_{\discel\in D_\lattice^r}
%	\sum_{\substack{T\in q(\discel) + \halfint{r} \\ T>0}}
%	\sum_{w_1\in M^\genus_r(\ZZ)^*/\GL(r,\ZZ)}
%	a_r^k(T,\discel)
%	\cdot e\big(
%	\tr T[w_1^t]\tau_1
%	\big)
%	\\
%	&\quad\times
%	\sum_{s=0}^{\min\{r,\genus-r\}}
%	\sum_{\substack{w_3\in\GL(\genus-r,\ZZ)_s\backslash\ZZ^{(\genus-r,r)}_{s,0} \\ \left(\begin{smallmatrix}
%	w_1 \\ w_3
%	\end{smallmatrix}\right)\text{ primitive}}}
%	\varepsilon(s)
%	\mathfrak{e}_{\discel w_1^t}
%	\\
%	&\quad\otimes
%	\underbrace{\frac{1}{2}
%	\sum_{\gamma\in\widetilde{\mathfrak{A}}_{\genus-r,s}\backslash \Mp_{2(\genus-r)}(\ZZ)}
%	e\big(
%	\tr T [w_3^t] \gamma\cdot\tau_4
%	\big)
%	\phi(\tau_4)^{-2k}
%	\weil{\lattice}{\genus-r}(\gamma)^{-1}
%	\mathfrak{e}_{\discel w_3^t}}_{(\clubsuit)},
%	\end{align*}
%	where the factor of~$1/2$ arises from the fact that the map~$\widetilde{\mathfrak{A}}_{\genus-r,s}\backslash \Mp_{2(\genus-r)}\to\widetilde{\mathfrak{A}}_{\genus-r,s}'\backslash \Mp_{2(\genus-r)}$ has fibers of cardinality~$2$.
%	\\\textcolor{blue}{\textbf{Note: }The previous displayed formula is where~$\varepsilon(s)$ pops out.}\\
%	If~$w_3\in\ZZ^{(\genus-r,r)}_{s,0}$, then~$T[w_3^t]=\big(\begin{smallmatrix}
%	\ast & 0\\
%	 0 & 0_{\genus-r-s}
%	\end{smallmatrix}\big)$ and~$\discel w_3^t = \discel (w_3^t)_{(s)}\otimes\zeroe{\genus-r-s}$, where~$(w_3^t)_{(s)}$ denotes the matrix given by the first~$s$ columns of~$w_3^t$.
%	This implies that the sum denoted by~$(\clubsuit)$ above equals~$\mathfrak{g}_{\genus-r,s}^k\big(\tau_4,T[w_3^t],\discel(w_3^t)_{(s)}\big)$.}
%	\end{proof}
%	
%	\textcolor{blue}{Recall from~\eqref{eq;FJexpsvcvv} that we denote the Fourier--Jacobi expansion of~$\eiszero{k}{\gengenus}{\lattice}$ of cogenus~$r$ by
%	\bes
%	\eiszero{k}{\gengenus}{\lattice}(\tau)
%	=
%	\sum_{\discel\in D_\lattice^r}
%	\sum_{\substack{T\in q(\discel) + \halfint{r} \\ T\ge 0}}
%	\mathfrak{e}_\discel
%	\otimes
%	\jaccoefvv{\genus}{r}{\discel}(\tau_4,\tau_2,T)\cdot e(\tr T\tau_1),
%	\ees
%	where~$\tau_1\in\HH_r$ and~$\tau_4\in\HH_{\genus-r}$.}
	
	The proof of the following auxiliary lemma is left to the reader.
	\begin{lemma}\label{lemma;propFJcoefSEvv}
	Let~$k>\genus+1$,
%	This is same condition to make $E^k_g$ convergent; see \cite[p.~$68$, Proposition~$5$]{klingen}.
	and let~$T_r\in q(\discel) + \halfint{r}$ be a positive semidefinite matrix, where~$\discel\in D_\lattice^r$.
	Assume that there exists~$U\in \GL(\genus,\ZZ)$ such that~$U^t T_r U = \big(\begin{smallmatrix}
	0 & 0\\
	0 & T_{r-1}
	\end{smallmatrix}\big)$ for some~$T_{r-1}\in q(\discel') + \halfint{r-1}$, where~$\discel U=(\discel'',\discel')$ with~$\discel'\in D_\lattice^{r-1}$ and~$\discel''$ is isotropic.
	Then
	\be\label{eq;symmetryFjaccoefdiffgenvv}
	\jaccoefvv{\genus}{r}{\discel}(\tau_4,0,T_{r})
	=
	\begin{cases}
	0 & \text{if~$\discel''\neq 0$,}
	\\
	\jaccoefvv{\genus-1}{r-1}{\discel'}(\tau_4,0,T_{r-1}) & \text{otherwise},
	\end{cases}
	\ee
	for all~$\tau_4\in\HH_{\genus-r}$.
	Moreover, we have~$\jaccoefvv{\genus}{r}{0}(\tau_4,0,0)
	=
	\eiszero{k}{\genus-r}{\lattice}(\tau_4)$.
	\end{lemma}
%	\begin{proof}
%	\textcolor{blue}{For simplicity we define~$s\coloneqq\genus-r$.
%	The Fourier--Jacobi coefficients~$\jaccoefvv{\genus}{r}{\discel}$ encode parts of the Fourier expansion of~$\eiszero{k}{\genus}{\lattice}$, more precisely we have
%	\[
%	\jaccoefvv{\genus}{r}{\discel}(\tau_4,0,T_r)
%	=
%	\sum_{\disceltwo\in D_\lattice^s}
%	\sum_{\substack{T_s\in q(\disceltwo) + \halfint{s} \\ T_s\geq0}}
%	\sum_{\substack{T=\big(\begin{smallmatrix}
%	T_r & \tilde{T}^t \\
%	\tilde{T} & T_s
%	\end{smallmatrix}\big)\in q(\discel,\disceltwo) + \halfint{\genus} \\ T\geq 0}}
%	a^k_\genus\big(T,(\discel,\disceltwo)\big)\cdot e(
%	\tr T_s\tau_4 
%	)
%	\mathfrak{e}_\disceltwo
%	.
%	\]
%	Since~$U^t T_r U = \big(\begin{smallmatrix}
%	0 & 0\\
%	0 & T_{r-1}
%	\end{smallmatrix}\big)$, we have that every positive semidefinite matrix~$T=\big(\begin{smallmatrix}
%	T_r & \ast\\
%	\ast & T_s
%	\end{smallmatrix}\big)\in q(\discel,\disceltwo) + \halfint{\genus}$, for some~$\disceltwo\in D_\lattice^{\genus-r}$, is equivalent to a matrix of the form~$\Big(\begin{smallmatrix}
%	0 & 0 & 0 \\
%	0 & T_{r-1} & \ast \\
%	0 & \ast & T_s
%	\end{smallmatrix}\Big)$.
%	By~\eqref{eq;symmofcoef2}, the Fourier coefficients of~$\eiszero{k}{\genus}{\lattice}$ associated to equivalent matrices are equal.
%	Hence we may assume without loss of generality that~$T_r=\big(\begin{smallmatrix}
%	0 & 0\\
%	0 & T_{r-1}
%	\end{smallmatrix}\big)$ and~$\alpha=(\alpha'',\alpha')$ for some~$\alpha'\in D_\lattice^{r-1}$ and some isotropic~$\alpha''$.
%	By Theorem~\ref{thm;Phi0onkling}, if~$\varepsilon\in D_\lattice$, then
%	\be\label{eq;fromPhi0onkling}
%	\Phi_\varepsilon(\eiszero{k}{\genus}{\lattice})
%	=
%	\begin{cases}
%	\eiszero{k}{\genus-1}{\lattice} & \text{if~$\varepsilon=0$,}\\
%	0 & \text{otherwise.}
%	\end{cases}
%	\ee
%	If~$\alpha''=0$, then we may compute
%	\ba\label{eq;comPhiandJacvv}
%	\eiszero{k}{\genus-1}{\lattice}\big(\begin{smallmatrix}
%	\tau_1' & 0\\
%	0 & \tau_4
%	\end{smallmatrix}\big)
%	&=
%	\sum_{\discel'\in D_\lattice^{\genus-1}}
%	\lim_{t\to\infty}
%	\big(\eiszero{k}{\genus}{\lattice}\big)_{(0,\discel')}\Big(\begin{smallmatrix}
%	it & 0 & 0\\
%	0 & \tau_1' & 0\\
%	0 & 0 & \tau_4
%	\end{smallmatrix}\Big)
%	\mathfrak{e}_{\discel'}
%	\\
%	&=
%	\sum_{\discel'\in D_\lattice^{\genus-1}}
%	\sum_{\substack{T'\in q(\discel') + \halfint{\genus-1} \\ T'\geq 0}}
%	a^k_\genus\big(\big(\begin{smallmatrix}
%	0 & 0\\
%	0 & T'
%	\end{smallmatrix}\big),(0,\discel')\big)\cdot e\Big( \tr T'\big(\begin{smallmatrix}
%	\tau_1' & 0\\
%	0 & \tau_4
%	\end{smallmatrix}\big) \Big)
%	\mathfrak{e}_{\discel'}
%	\\
%	&=
%	\sum_{\discel'\in D_{\lattice}^{r-1}}
%	\sum_{\substack{T_{r-1}\in q(\discel') + \halfint{r-1} \\ T_{r-1}\geq 0}}
%	e(\tr T_{r-1}\tau_1')
%	\mathfrak{e}_{\discel'}
%	\\
%	&\quad
%	\otimes
%	\underbrace{
%	\sum_{\disceltwo\in D_\lattice^s}
%	\sum_{\substack{T_{s}\in q(\disceltwo) + \halfint{s} \\ T_s\geq0}}
%	\sum_{\substack{T=\bigg(\begin{smallmatrix}
%	0 & 0 & 0 \\
%	0 & T_{r-1} & \tilde{T}^t \\
%	0 & \tilde{T} & T_s
%	\end{smallmatrix}\bigg)\in q(0,\discel',\disceltwo) + \halfint{\genus} \\ T\geq 0}}
%	a^k_\genus\big(T,(0,\discel',\disceltwo)\big)\cdot e\big(
%	\tr T_s\tau_4
%	\big)
%	\mathfrak{e}_\disceltwo}_{=\jaccoefvv{\genus}{r}{(0,\discel')}\big(\tau_4,0,\big(\begin{smallmatrix}
%	0 & 0\\
%	0 & T_{r-1}
%	\end{smallmatrix}\big)\big)
%	}{}
%	,
%	\ea
%	where~$\tau_1'\in\HH_{r-1}$.
%	We deduce~\eqref{eq;symmetryFjaccoefdiffgenvv} by comparing the right-hand side of~\eqref{eq;comPhiandJacvv} with the Fourier--Jacobi expansion of cogenus~$r-1$ of~$\eiszero{k}{\genus-1}{\lattice}\big(\begin{smallmatrix}
%	\tau_1' & 0\\
%	0 & \tau_4
%	\end{smallmatrix}\big)$.
%	The vanishing of~$\jaccoefvv{\genus}{r}{\discel}(\tau_4,0,T_{r})$ in the case of~$\alpha''\neq 0$ is deduced using~\eqref{eq;fromPhi0onkling} analogously.}
%	
%	\textcolor{blue}{By~\eqref{eq;symmetryFjaccoefdiffgenvv}, to deduce that~$\jaccoefvv{\genus}{r}{0}(\tau_4,0,0)
%	=
%	\eiszero{k}{\genus-r}{\lattice}(\tau_4)$ it is enough to show such an equality only for~$r=1$.
%	Under the latter assumption, the computations in~\eqref{eq;comPhiandJacvv} simplifies to
%	\begin{align*}
%	\eiszero{k}{\genus-1}{\lattice}(\tau_4)
%	&=
%	\sum_{\discel'\in D_\lattice^{\genus-1}}
%	\lim_{t\to\infty}
%	(\eiszero{k}{\genus}{\lattice})_{(0,\discel')}
%	\big(\begin{smallmatrix}
%	it & 0\\
%	0 & \tau_4
%	\end{smallmatrix}\big)
%	\mathfrak{e}_{\discel'}
%	\\
%	&=
%	\sum_{\discel'\in D_\lattice^{\genus-1}}
%	\sum_{\substack{T'\in q(\discel') + \halfint{\genus-1} \\ T'\geq 0}}
%	a^k_\genus\big(\big(\begin{smallmatrix}
%	0 & 0\\
%	0 & T'
%	\end{smallmatrix}\big), (0,\discel')\big)\cdot e( \tr T'
%	\tau_4)
%	\mathfrak{e}_{\discel'}
%	\\
%	&=\jaccoefvv{\genus}{1}{0}(\tau_4,0,0).\qedhere
%	\end{align*}}
%	\end{proof}
	The following result is the previously announced vector-valued generalization of the vanishing of Garrett's map.
%	It will be a key ingredient for the proof of the Lefschetz decomposition of the (cohomology class of the) Kudla--Millson theta function.
	\begin{cor}\label{cor;vanishingGmapvv}
	Let~$k>\genus+1$ and let~$\genus=r+s$ be such that~$1\leq r<s$.
	Then
	\be\label{eq;vanishingfromBösatz12vv}
	\bigpet{f}{\eiszero{k}{\genus}{\lattice}\big(
	\begin{smallmatrix}
	\tau_1 &   \\
	  & \ast
	\end{smallmatrix}\big)}=0
	\qquad
	\text{for every~$f\in S^k_{s,\lattice}$ and~$\tau_1\in\HH_r$}.
	\ee
	Here we denote by~$\ast$ the variable of integration in~$\HH_s$.
	\end{cor}
	\begin{proof}
%	Recall that we denote the Fourier expansion of~$E^k_\genus$ by
%	\[
%	E^k_\genus(\tau)=\sum_{\substack{T\in\halfint{\genus}\\ T\ge 0}}
%	a^k_\genus(T)\cdot e(\tr T\tau),\qquad\tau\in\HH_\genus.
%	\]
%	We denote the Fourier--Jacobi expansion of~$E^k_\genus$ of index~$r$ by
%	\[
%	E^k_\genus(\tau)
%	=
%	\sum_{\substack{T_r\in\halfint{r}\\ T_r\geq 0}}
%	\phi_{\genus,r}(\tau_4,\tau_2,T_r)\cdot e(\tr T_r\tau_1),
%	\qquad
%	\tau=\big(\begin{smallmatrix}
%	\tau_1 & \tau_2 \\
%	\tau_2^t & \tau_4
%	\end{smallmatrix}\big)\in\HH_\genus,
%	\]
%	where~$\tau_1\in\HH_r$ and~$\tau_4\in\HH_s$.
%	The Fourier--Jacobi coefficients~$\phi_{\genus,r}$ encode parts of the Fourier expansion of~$E^k_\genus$, more precisely we have
%	\[
%	\phi_{\genus,r}(\tau_4,\tau_2,T_r)
%	=
%	\sum_{\substack{T_s\in\halfint{s} \\ T_s\geq0}}
%	\sum_{\substack{T=\big(\begin{smallmatrix}
%	T_r & \tilde{T}^t \\
%	\tilde{T} & T_s
%	\end{smallmatrix}\big)\in\halfint{\genus} \\ T\geq 0}}
%	a^k_\genus(T)\cdot e\big(
%	\tr(2\tilde T\tau_2 + T_s\tau_4)
%	\big).
%	\]
	
	Since~$\jaccoefvv{\genus}{r}{\discel}(\tau_4,0,T_{r})$ is a Siegel modular form of weight~$k$ and genus~$s$ with respect to~$\weil{\lattice}{s}$ for every~$T_r\in q(\discel) + \halfint{r}$, we deduce that
	\be\label{eq;frmBöcherer12svv}
	\bigpet{f}{\eiszero{k}{\genus}{\lattice}\big(\begin{smallmatrix}
	\tau_1 &  \\
	 & \ast
	\end{smallmatrix}\big)}
	=
	\sum_{\discel\in D_\lattice^r}
	\sum_{\substack{T_r\in q(\discel) + \halfint{r}\\ T_r\geq 0}}
	\bigpet{f}{\jaccoefvv{\genus}{r}{\discel}(\ast,0,T_{r})}\cdot e(-\tr T_r\overline{\tau_1})
	\mathfrak{e}_\discel.
	\ee
%	Recall from Definition~\ref{def;Essentialpart} that the \emph{essential part} of cogenus~$r$ of~$\eiszero{k}{\genus}{\lattice}$ is the subseries of the Fourier--Jacobi expansion of~$\eiszero{k}{\genus}{\lattice}$ given by
%	\[
%	\espart{k}{\genus}{r}(\tau)
%	=
%	\sum_{\discel\in D_\lattice^r}
%	\sum_{\substack{T_r\in q(\discel) + \halfint{r} \\ T_r>0}}
%	\mathfrak{e}_\discel
%	\otimes
%	\jaccoefvv{\genus}{r}{\discel}(\tau_4,\tau_2,T_r)\cdot e(\tr T_r\tau_1).
%	\]
	By Theorem~\ref{thm;genBöcherervv} the diagonal restriction~$\espart{k}{\genus}{r}\big(\begin{smallmatrix}
	\tau_1 & \\
	 & \tau_4
	\end{smallmatrix}\big)$ is, with respect to~$\tau_4\in\HH_s$, a combination of Klingen Eisenstein series of genus~$s$ arising from cusp forms of genus at most~$r$.
	Since~$r<s$, we deduce from Proposition~\ref{prop;ortcusp&kling} that
	\[
	\bigpet{f}{\espart{k}{\genus}{r}\big(\begin{smallmatrix}
	\tau_1 & \\
	 & \ast
	\end{smallmatrix}\big)}=0\qquad\text{for every~$\tau_1\in\HH_r$ and~$f\in S^k_{s,\lattice}$.}
	\]
	This implies that
	\[
	\bigpet{f}{\eiszero{k}{\genus}{\lattice}\big(\begin{smallmatrix}
	\tau_1 & \\
	 & \ast
	\end{smallmatrix}\big)}
	=
	\bigpet{f}{\underbrace{\eiszero{k}{\genus}{\lattice} - \espart{k}{\genus}{r}\big(\begin{smallmatrix}
	\tau_1 & \\
	 & \ast
	\end{smallmatrix}\big)}_{\eqqcolon \tespart{k}{\genus}{r}\big(\begin{smallmatrix}
	\tau_1 & \\
	 & \ast
	\end{smallmatrix}\big)}},
	\]
	where
	\[
	\tespart{k}{\genus}{r}\big(\begin{smallmatrix}
	\tau_1 & \\
	 & \tau_4
	\end{smallmatrix}\big)
	=
	\sum_{\discel\in D_\lattice^r}
	\sum_{\substack{T_r\in q(\discel) + \halfint{r} \\
	T_r\geq 0 \\ T_r\not> 0}}
	\mathfrak{e}_\discel\otimes
	\jaccoefvv{\genus}{r}{\discel}(\tau_4,0,T_r) \cdot e(\tr T_r\tau_1).
	\]
	Hence we have
	\be\label{eq;frombös12vv}
	\bigpet{f}{\eiszero{k}{\genus}{\lattice}\big(\begin{smallmatrix}
	\tau_1 &  \\
	 & \ast
	\end{smallmatrix}\big)}
	=
	\sum_{\discel\in D_\lattice^r}
	\sum_{\substack{T_r\in q(\discel) + \halfint{r}\\ T_r\geq 0 \\ T_r \not>0}}
	\bigpet{f}{\jaccoefvv{\genus}{r}{\discel}(\ast,0,T_{r})}\cdot e(-\tr T_r\overline{\tau_1})
	\mathfrak{e}_\discel.
	\ee
%	This means that to show~\eqref{eq;vanishingfromBösatz12}, it is enough to prove that
%	\be
%	\bigpet{f}{\phi_{g,r}^k(\ast,0,T_r)}=0\qquad\text{for every~$T_r\in\halfint{r}$ such that~$T\geq0$ but~$T\not>0$.}
%	\ee
	
	We now prove~\eqref{eq;vanishingfromBösatz12vv} by induction on~$r$ and~$\genus$.
	If~$r=1$, then~\eqref{eq;frombös12vv} simplifies to
	\[
	\bigpet{f}{\eiszero{k}{\genus}{\lattice}\big(\begin{smallmatrix}
	\tau_1 & \\
	 & \ast
	\end{smallmatrix}\big)}
	=
	\sum_{\discel\in D_\lattice}
	\bigpet{f}{\jaccoefvv{\genus}{1}{\discel}(\ast,0,0)}
	\mathfrak{e}_\discel,
	\]
	which vanishes by Lemma~\ref{lemma;propFJcoefSEvv}.
	
	Let now~$r>1$.
	Suppose that~\eqref{eq;vanishingfromBösatz12vv} is true for~$r$ and every genus~$\genus$ such that~$r < \genus - r$.
	By~\eqref{eq;frmBöcherer12svv} we deduce that $\bigpet{f}{\jaccoefvv{\genus}{r}{\discel}(\ast,0,T_r)}=0$ for every~$\discel\in D_\lattice^r$, every positive semidefinite~$T_r\in q(\discel) + \halfint{r}$, and every~$f\in S^k_{\genus-r,\lattice}$.
	We show that~\eqref{eq;vanishingfromBösatz12vv} is true also for~$r+1$ and every genus~$\genus$ such that~$r+1<\genus-r-1$.
	By~\eqref{eq;frombös12vv}, it is enough to show that if~$f\in S^k_{g-r-1,\lattice}$, then
	\be\label{eq;fsvvdsvv}
	\bigpet{f}{\jaccoefvv{\genus}{r+1}{\discel}(\ast,0,T_{r+1})}=0
	\ee
	for every~$\discel\in D_\lattice^{r+1}$ and every positive semidefinite~$T_{r+1}\in q(\discel) + \halfint{r+1}$ that is \emph{not positive definite}.
	With respect to the action of~$\GL(r+1,\ZZ)$ on~$\halfint{r+1}$, every such a matrix~$T_{r+1}$ is equivalent to~$\big(\begin{smallmatrix}
	0 & 0\\
	0 & T_r
	\end{smallmatrix}\big)$, for some positive semidefinite~$T_r\in\halfint{r}$.
%	\\\textcolor{blue}{This is true, and it follows directly as for the scalar valued case.
%	In fact, if~$N$ is the level of the lattice, then~$NT\in\halfint{r+1}$.
%	One can then use the usual argument with the action of~$\GL(r+1,\ZZ)$ on~$\halfint{r+1}$.}\\
	Let~$U\in \GL(r+1,\ZZ)$ be such that~$U^t T_{r+1} U = \big(\begin{smallmatrix}
	0 & 0\\
	0 & T_r
	\end{smallmatrix}\big)$, and write~$\discel U = (\discel'',\discel')$ for some isotropic~$\discel''\in D_\lattice$.
	Since~${\jaccoefvv{\genus}{r+1}{\discel}}(\tau_4,0,T_{r+1})$ equals either zero or~$\jaccoefvv{\genus-1}{r}{\discel'}(\tau_4,0,T_r)$ by Lemma~\ref{lemma;propFJcoefSEvv}, we may deduce~\eqref{eq;fsvvdsvv} by induction.
	\end{proof}
	
	
	
	
	
	

	\end{document}
	
	
	
	
	
	
	
	
	
	
	
	
	
	
	
	
	
	
	
	
	
	
	
\documentclass[../Lefschetz_decomposition.tex]{subfiles}

	\begin{document}
	Smooth orthogonal Shimura varieties are (not necessarily compact) complete Kähler manifolds.
	In this section we recall some facts on compact Kähler manifolds and explain to what extent they are still satisfied by orthogonal Shimura varieties.
	We also introduce the Kudla--Millson theta function and special cycles, and describe the irreducible components of the latter in terms of primitive lattice vectors.
	
	\subsection{Compact Kähler manifolds and the Lefschetz decomposition}
	\label{sec;compactkähler}
	Let~$(X,\omega)$ be a compact Kähler manifold of dimension~$n$.
	We choose~$\omega^n$ to be the volume form of~$X$.
	We quickly recall some well-known properties of~$X$.
	The reader may refer to~\cite{huybrechts} for details.
	
	The standard Hermitian inner product~$(\cdot{,}\cdot)$ on the space~$A^r(X)$ of complex-valued~$r$-differential forms on~$X$ is defined as
\be\label{eq;definprodkdifforms}
	(\alpha,\beta)\coloneqq \int_X \alpha\wedge\ast\overline{\beta}\qquad\text{for every~$\alpha,\beta\in A^r(X)$,}
	\ee
	where~$\ast\colon A^r(X)\to A^{2n-r}(X)$ is the Hodge~$\ast$-operator associated to the volume form~$\omega^n$.
	
%	The standard Hermitian inner product~$(\cdot{,}\cdot)$ on the space~$A^r(X)$ of complex-valued~$r$-differential forms on~$X$ is defined as
%\be\label{eq;definprodkdifforms}
%	(\alpha,\beta)\coloneqq \int_X \langle \alpha,\beta\rangle \omega^n\qquad\text{for every~$\alpha,\beta\in A^r(X)$.}
%	\ee
%	Here we denoted by~$\langle\cdot{,}\cdot\rangle$ the Hermitian extension of the metric induced by the Kähler form of~$X$ on the fibers of the bundle~$\bigwedge^r T^*X$.
%	
%	The Hodge~$\ast$-operator is defined as the linear map~$\ast\colon A^r(X)\to A^{2n-r}(X)$ such that
%	\[
%	\alpha\wedge\ast\overline{\beta}=\langle\alpha,\beta\rangle\omega^n
%	\qquad\text{for all~$\alpha\in A^{2n-r}(X)$ and~$\beta\in A^r(X)$}.
%	\]
%%	\textcolor{blue}{\textbf{Note: }
%%	We do not make the standard choice of volume form as for Huybrechts, but we follow Kudla.
%%	In particular, for us the volume form is~$\omega^n$ and not~$\omega^n/n!$.
%%	This clarifies why we define the Hodge~$\ast$-operator as
%%	\be\label{eq;explnactorn!inblue}
%%	\alpha\wedge\ast\overline{\beta}=\langle\alpha,\beta\rangle\omega^n
%%	\qquad\text{and not as}\qquad
%%	\alpha\wedge\ast\overline{\beta}=\langle\alpha,\beta\rangle\omega^n/n!.
%%	\ee
%%	This clarifies why I corrected the formula of Lemma~\ref{lemma;formmatchonHgg} below by a factor of~$n!$. An explanation is given also inside its proof.}\\
%	We may then rewrite the inner product~\eqref{eq;definprodkdifforms} as~$(\alpha,\beta) = \int_X \alpha\wedge\ast\overline{\beta}$.
	
	Let~$H^r(X)$ be the de Rham cohomology group of index~$r$ of~$X$.
	We denote by~$\drcoho{\alpha}$ the cohomology class in~$H^r(X)$ associated to any closed form~$\alpha\in A^r(X)$.
	By Stokes Theorem, the extension to cohomology~$(\drcoho{\alpha},\drcoho{\beta})\coloneqq(\alpha,\beta)$ of the inner product~\eqref{eq;definprodkdifforms} is well-defined.
	It is well-known that the \emph{Hodge decomposition}
	\[
	H^r(X)=\bigoplus_{p+q=r} H^{p,q}(X)
	\]
	is orthogonal with respect to the inner product~$(\cdot{,}\cdot)$.
	
	We now recall a second classical decomposition of~$H^r(X)$.
	The \emph{Lefschetz operator} is the linear map~$\lefschetz\colon A^{r}(X)\to A^{r+2}(X)$, $\alpha\mapsto\alpha\wedge \omega$.
	The \emph{dual Lefschetz operator} is the operator of degree~$-2$ defined as~$\duallef\coloneqq\ast^{-1}\circ\lefschetz\circ\ast$.
	We denote with the symbols~$\lefschetz$ and~$\duallef$ also the maps induced in cohomology.
	These operators preserve the bidegree of the Hodge decomposition.
	\begin{defi}
	The \emph{primitive cohomology groups} of~$X$ are
	\bas
	\Hprim^r(X)
	&\coloneqq
	\ker\big(
	\duallef\colon H^r(X)\to H^{r-2}(X)
	\big),
	\\
	\Hprim^{p,q}(X)
	&\coloneqq
	\ker\big(
	\duallef\colon H^{p,q}(X)\to H^{p-1,q-1}(X)
	\big).
	\eas
	\end{defi}
	The following result provides a decomposition of~$H^r(X)$ known as the \emph{Lefschetz decomposition}; see~\cite[Propositions~$3.2.2$ and~$3.3.13$]{huybrechts} for details.
	\begin{thm}[Hard Lefschetz Theorem]\label{thm;Hardlefsch}
	Let~$(X,\omega)$ be a compact Kähler manifold of dimension~$n$.
	If~$r\le n$, then~$\lefschetz^{n-r}\colon H^r(X)\to H^{2n-r}(X)$ is an isomorphism.
	Moreover, for every~$r\le 2n$ we have a decomposition
	\[
	H^r(X)=\bigoplus_{j\ge 0} \lefschetz^j \Hprim^{r-2j}(X),
	\]
	which is orthogonal with respect to the inner product~$(\cdot{,}\cdot)$.
	Analogous results hold on each term of the Hodge decomposition of~$H^r(X)$.
	\end{thm}
	\begin{rem}
	It follows from Theorem~\ref{thm;Hardlefsch} that if~$r<n$, then the Lefschetz operator~$\lefschetz\colon H^r(X)\to H^{r+2}(X)$ is injective.
	The kernel of~$\lefschetz$ may be non-zero if~$r\ge n$.
	In fact, if~$m\ge0$ then~$\lefschetz^{m+1}H_0^{n-m}(X)=0$.
	\end{rem}
	
	The cohomology groups of the form~$H^{r,r}(X)$ play a key role in the present paper.
	On such groups, the inner product~$(\cdot{,}\cdot)$ may be computed as in the following result, which follows easily from~\cite[Proposition~$1.2.31$]{huybrechts}.
	\begin{lemma}\label{lemma;formmatchonHgg}
	Let~$(X,\omega)$ be a compact Kähler manifold of dimension~$n$.
	If~$\beta$ is a class in~$\lefschetz^j \Hprim^{r-j,r-j}(X)$, then
	\be\label{eq;formmatchonHgg}
	(\alpha,\beta)
	=
	\frac{(-1)^{(r-j)(2(r-j)+1)}\cdot n!\cdot j!}{(n-2r+j)!}
	\int_X\alpha\wedge\overline{\beta}\wedge\omega^{n-2r}
	\qquad
	\text{for all~$\alpha\in H^{r,r}(X)$.}
	\ee
	\end{lemma}
%	\textcolor{blue}{\textbf{Note:} There is a difference with Huybrechts formula, given by a factor of~$n!$.
%	This is due to the choice~\eqref{eq;explnactorn!inblue} of the volume form.}
%	
%	\begin{proof}
%	\textcolor{blue}{\textbf{I suggest we erase this proof, since it is not so interesting.}}
%	
%	\textcolor{blue}{We recall from~\cite[Proposition~$1.2.31$]{huybrechts} that if~$\xi\in\Hprim^{k,k}(X)$, then
%	\bes
%	\ast (\lefschetz^j\xi) =
%	\frac{(-1)^{k(2k+1)}\cdot n!\cdot j!}{(n-2k-j)!} \lefschetz^{n-2k-j}\xi.
%	\ees
%	The factor~$n!$ above does not appear in~\cite[Proposition~$1.2.31$]{huybrechts}.
%	This is due to the fact that the~$\ast$-operator in the present article differs from the one in~\cite{huybrechts} by a factor of~$n!$.}
%	
%	\textcolor{blue}{We write~$\beta=\lefschetz^j \beta'$ for some~$\beta'\in\Hprim^{r-j,r-j}(X)$ and compute
%	\begin{align*}
%	(\alpha,\beta)
%	&=
%	\int_X \alpha\wedge\ast\overline{\beta}
%	=
%	\int_X \alpha\wedge\ast(\lefschetz^j\overline{\beta'})
%	=
%	\frac{(-1)^{(r-j)(2(r-j)+1)}\cdot n! \cdot j!}{(n-2r+j)!}
%	\int_X \alpha\wedge\overline{\beta'}\wedge \omega^{n-2r+j}
%	\\
%	&=
%	\frac{(-1)^{(r-j)(2(r-j)+1)}\cdot n! \cdot j!}{(n-2r+j)!}
%	\int_X \alpha\wedge\overline{\beta}\wedge \omega^{n-2r}.
%	\qedhere
%	\end{align*}}		
%	\end{proof}
	
	We may construct~$(\cdot{,}\cdot)$ as in~\eqref{eq;definprodkdifforms} also for non-compact Kähler manifolds.
	However, it is not an inner product, since the defining integral may diverge.
	To avoid this problem, we will restrict all previous definitions to square-integrable differential forms.
	
%	---------------------------
%	
%	If~$g$ is a Riemannian metric on~$M$, we denote by~$g_p$ the inner product induced by~$g$ along the fiber at~$p$ of the tangent bundle~$TM$, i.e.~the map
%	\[
%	g_p\colon T_pM\times T_pM\longrightarrow\RR,\qquad (v,w)\longmapsto g_p(v,w).
%	\]
%	Such a metric induces a metric also along the fibers of the dual tangent bundle~$T^*M$, as follows.
%	Recall that every non-degenerate inner product on a vector space~$V$ enables us to identify~$V$ with its dual~$V^*$.
%	We denote by~$\iota\colon TM\to T^*M$ the natural isomorphism induced by~$g$ along the fibers as
%	\bes
%	\iota(v)\coloneqq g_p(\cdot{,}v)\in T^*_pM,\qquad\text{for every~$v\in T_p M$}.
%	\ees
%	The standard metric~$\langle\cdot{,}\cdot\rangle$ induced by~$g$ on the fibers of~$T^*M$ is the one defined as
%	\bes
%	\langle \alpha,\beta\rangle = g_p\big(\iota^{-1}(\alpha),\iota^{-1}(\beta)\big),\qquad\text{for every~$\alpha,\beta\in T^*_pM$.}
%	\ees
%	We denote by~$\langle\cdot{,}\cdot\rangle$ also the standard inner product induced on the fibers of the bundle~$\bigwedge^r T^*M$, defined as
%%	\be\label{eq;ginduceprodinwedgek}
%%	\langle v_1\wedge\cdots\wedge v_r,w_1\wedge\cdots\wedge w_r\rangle\coloneqq \det (\langle v_i,w_j \rangle),\qquad\text{for every~$\wedge_i v_i,\wedge_j w_j\in \bigwedge^r T^*_pM$.}
%%	\ee
%	We remark that the standard inner product induced on~$\bigwedge^r T^*M$ is usually defined with a factor of~$r!$ multiplying the right-hand side of~\eqref{eq;ginduceprodinwedgek}.
%	However, since also for the volume of subvarieties we do not consider such kind of factors, we prefer to not insert it in our definition.
%	
%	The previous inner products enable us to construct a inner product on the space~$A^r(M,\RR)$ of real differential forms on~$M$, whenever~$M$ is compact and oriented.
%%	\begin{lemma}\label{lem;inprodkdifforms}
%%	Let~$(M,g)$ be a compact oriented Riemannian manifold with volume form~$\nu_g$, and let~$\langle\cdot{,}\cdot\rangle$ be the inner product on the fibers of~$\bigwedge^r T^*M$ induced by the metric~$g$.
%%	We may define an inner product~$(\cdot{,}\cdot)_g$ on~$A^r(M)$ as
%%	\be\label{eq;definprodkdifforms}
%%	(\alpha,\beta)_g\coloneqq \int_M \langle \alpha,\beta\rangle \nu_g,\qquad\text{for every~$\alpha,\beta\in A^r(M)$.}
%%	\ee
%%	\end{lemma}
%%	The standard extension of the inner product~$(\cdot{,}\cdot)_g$ over~$\CC$ gives a Hermitian inner product on the space~$A^r(M)$ of complex-valued differential forms.\\
%	
%	From now on, we restrict our attention to Kähler manifolds.
%	Let~$(X,\omega)$ be a compact Kähler manifold.
%	The constructions above hold also if~$g$ is chosen as the Riemannian metric induced by the Kähler form~$\omega$.
%	
%	We are going to show that the Hermitian product obtained by extension over~$\CC$ of the inner product defined in~\eqref{eq;definprodkdifforms} is essentially the same as the one defined in~\eqref{eq;hermformforus}.	
%	\begin{lemma}
%	The Hodge decomposition and the Lefschetz decomposition of~$H^r(X)$ are orthogonal with respect to the Hermitian products~\eqref{eq;hermformforus} and~\eqref{eq;definprodkdifforms}.
%	\end{lemma}
%	\begin{proof}
%	The orthogonality with respect to~$(\cdot{,}\cdot)_g$ as defined in~\eqref{eq;definprodkdifforms} follows directly from~\cite[Proposition~$3.2.2$]{huybrechts}.
%	
%	We prove the orthogonality of the Hodge decomposition w.r.t.~$(\cdot{,}\cdot)$ as defined in~\eqref{eq;hermformforus}.
%	Let~${[\alpha]\in H^{p,q}(X)}$ and let~${[\beta]\in H^{p',q'}(X)}$, for some~$(p,q)\neq(p',q')$ such that~$p+q=p'+q'$.
%	Then~$\alpha\wedge\overline{\beta}\wedge\omega^{n-r}$ is a~$(p+q'+n-r,q+p'+n-r)$-form on~$X$, where~$r=p+q$.
%	This differential form is a top-degree form on~$X$.
%	The only non-zero differential forms of top-degree on~$X$ are of type~$(n,n)$.
%	We may assume w.l.o.g.\ that~$q\neq q'$.
%	Hence~$p+q'\neq r$, so~$p+q'+n-r\neq n$, implying that~$\alpha\wedge\overline{\beta}\wedge\omega^{n-r}=0$.
%	This means that~\eqref{eq;hermformforus} makes the Hodge decomposition an orthogonal decomposition.
%	
%%	We now prove the orthogonality of the Hodge decomposition w.r.t.~$(\cdot{,}\cdot)_g$ as defined in~\eqref{eq;definprodkdifforms}.
%%	Let~$\alpha$ and~$\beta$ be as above.
%%	We may choose local coordinates~$z_j=x_j+iy_j$ such that the associated differentials~$dzj$ and~$d\overline{z_j}$ give an orthonormal basis along the fibers of~$T^*X$ with respect to~$(\cdot{,}\cdot)_g$.
%%	We write locally~$\alpha$ around~$x\in X$ as
%%	\bes
%%	\alpha=\sum_{\substack{\#I=p\\ \#J=q}}\alpha_{I,J} dz_I\wedge d\overline{z}_J,
%%	\ees
%%	and the same for~$\beta$.
%%	We may then compute locally
%%	\bas
%%	\langle\alpha,\beta\rangle_g = \sum_{\substack{\#I=p\\ \#J=q}} \sum_{\substack{\#I'=p\\ \#J'=q}}\alpha_{I,J}\beta_{I',J'}
%%	\langle
%%	dz_I\wedge d\overline{z}_J,dz_{I'}\wedge d\overline{z}_{J'}
%%	\rangle_g
%%	\eas
%%	Assume w.l.o.g.\ that~$p>p'$.
%%	This implies that for every~$I,I'$ there exists~$\ell\in I$ not in~$I'$, hence the line of the matrix~$\langle
%%	dz_I\wedge d\overline{z}_J,dz_{I'}\wedge d\overline{z}_{J'}
%%	\rangle_g$ given by
%%	\bes
%%	\big(
%%	\langle dz_\ell,dz_{\ell'}\rangle_g\quad\text{with $\ell'\in I'$}
%%	\quad | \quad
%%	\langle dz_\ell,dz_{\ell'}\rangle_g\quad\text{with $\ell'\in J'$}
%%	\big)
%%	\ees
%%	is zero.
%%	We deduce then that~$\langle\alpha,\beta\rangle_g=0$.
%	
%	We now prove the orthogonality of the Lefschetz decomposition w.r.t.~$(\cdot{,}\cdot)$ as defined in~\eqref{eq;hermformforus}.
%	Let~$[\alpha],[\beta]\in H^r(X)$.
%	Suppose that~$\alpha=\lefschetz^k\alpha'$ and~$\beta=\lefschetz^\ell\beta'$ for some primitive forms~$\alpha',\beta'$ of degree respectively~$r-2k$ and~$r-2\ell$.
%	Assume that~$k<\ell$ and that both~$\alpha'$ and~$\beta'$ are non-zero.
%	By~\cite[Proposition~$1.2.30$]{huybrechts} we have that~$r-2k\le n$ and~$\lefschetz^{n-r + 2k + 1}\alpha'=0$.
%	We may then compute
%	\bas
%	\alpha\wedge\overline{\beta}\wedge \omega^{n-r} = \lefschetz^{n-r}\alpha\wedge\overline{\beta}
%	= \lefschetz^{n-r+k+\ell} \alpha'\wedge\overline{\beta'},
%	\eas
%	which is zero.
%	In fact
%	\bes
%	\lefschetz^{n-r+k+\ell} \alpha' = \lefschetz^{\ell-k} \underbrace{\lefschetz^{n-r+2k+1} \alpha'}_{=0}.
%	\ees
%	In conclusion, we deduce that~$(\alpha,\beta)=0$ since~$\alpha\wedge\overline{\beta}\wedge \omega^{n-r}$ is so.
%%
%%	Since the orthogonality of the Lefschetz decomposition w.r.t.~$(\cdot{,}\cdot)_g$ as defined in~\eqref{eq;definprodkdifforms} can be proved following the same idea, we only sketch the proof.
%%	We may rewrite~$(\cdot{,}\cdot)_g$ in terms of the Hodge~$\ast$-operator as
%%	\bes
%%	(\alpha{,}\beta)_g =
%%	\int_X \langle \alpha,\beta\rangle \omega^n = \int_X\alpha\wedge\ast\overline{\beta}.
%%	\ees
%%	The orthogonality of the Lefschetz decomposition follows from the same idea as before together with the formula~\cite[Proposition~$1.2.31$]{huybrechts} to compute~$\ast L^{\ell}\overline{\beta}$.
%	\end{proof}
%	\begin{cor}\label{cor;checkwk}
%	The Hermitian products~\eqref{eq;hermformforus} and~\eqref{eq;definprodkdifforms} on~$H^r(X)$ are essentially the same.
%	In fact, they differ at most by a non-zero factor over every term of the Lefschetz decomposition of~$H^{p,q}(X)$, where~$p+q=r$.
%	\end{cor}
%	\begin{proof}
%	We check Corollary~\ref{cor;checkwk} on every term of the Lefschetz decomposition
%	\be\label{eq;proofldcheckwk}
%	H^{p,q}(X)=H_0^{p,q}(X)\oplus \lefschetz H^{p-1,q-1}_0(X)\oplus\cdots,
%	\ee
%	for every~$p,q$ such that~$p+q=r$.
%	We begin with the primitive part of~$H^{p,q}(X)$.
%	
%	Let~$\alpha,\beta\in H^{p,q}_0(X)$.
%	Recall that we denote by~$n$ the dimension of~$X$ and by~$\ast$ the Hodge operator.
%	By~\cite[Proposition~$1.2.31$]{huybrechts} we may compute the image of~$\beta$ under the composition of the~$\ast$-operator and some power of~$\lefschetz$ as
%	\bes
%	\ast (\lefschetz^j\beta) =
%	\frac{(-1)^{r(r+1)/2}j!}{(n-r-j)!} \lefschetz^{n-r-j}\Big(\sum_{p',q'} i^{p'-q'} \pi_{p',q'} (\beta)\Big),
%	\ees
%	where~$\pi_{p',q'}$ is the projection to~$H^{p',q'}(X)$.
%	
%	We may then deduce that
%	\bas
%	(\alpha,\beta)_g=\int_X \langle \alpha,\beta\rangle \omega^n
%	=
%	\int_X \alpha\wedge \ast\overline{\beta}
%	=
%	\frac{(-1)^{r(r+1)/2}i^{q-p}}{(n-r)!} \underbrace{\int_X\alpha\wedge\overline{\beta}\wedge\omega^{n-r}}_{=(\alpha,\beta)_\omega},
%	\eas
%	where the left-hand and right-hand sides of the previous formula are respectively the Hermitian products~\eqref{eq;definprodkdifforms} and~\eqref{eq;hermformforus}.
%	
%	The check on the remaining terms of the decomposition~\eqref{eq;proofldcheckwk} is analogous, once the problem has been rewritten in terms of primitive differential forms.
%	For instance, one can see immediately that~$(\lefschetz\alpha',\lefschetz\beta')_\omega=(\alpha',\beta')_\omega$, and that
%	\bes
%	(\lefschetz\alpha',\lefschetz\beta')_g=\int_X \lefschetz\alpha'\wedge\ast \lefschetz\overline{\beta'}=\frac{(-1)^{r(r+1)/2} i^{q-p}}{(n-r-1)!}\underbrace{\int_X\alpha'\wedge\overline{\beta'}\wedge\omega^{n-r}}_{=(\alpha',\beta')_\omega}.
%	\qedhere
%	\ees
%	\end{proof}
%	
%	
%	We may construct~$(\cdot{,}\cdot)_g$ as in Lemma~\ref{lem;inprodkdifforms} also for non-compact Kähler manifolds.
%	This works also for the Hermitian extension over~$\CC$ of~$(\cdot{,}\cdot)_g$.
%	However, it is not an inner product, since the integrals~\eqref{eq;definprodkdifforms} may diverge.
%	To avoid this problem, we will restrict all previous definitions to square-integrable differential forms and drop the condition that the Kähler manifold~$(X,\omega)$ is compact.
	
	\subsection{$L^2$-cohomology groups}\label{sec;isoofseveralcoho}
	Let~$(X,\omega)$ be a (not necessarily compact) complete Kähler manifold of dimension~$n$ and volume form~$\omega^n$.
	We introduce~$L^2$-forms and~$L^2$-cohomology following the wording of~\cite[Chapter~XIV, Section~$3$]{borelwallach}, \cite{dai} and~\cite{saperzucker}.
	
	We define~$(\cdot{,}\cdot)$ on~$A^r(X)$ as in~\eqref{eq;definprodkdifforms}.
	This is not an inner product, since it may diverge.
	The \emph{$L^2$-norm} of a differential form~$\eta\in A^r(X)$ is
	\bes
	(\eta,\eta)\coloneqq
	\int_X \eta\wedge\ast\overline{\eta}.
%	\int_X \langle \eta,\eta\rangle \omega^n.
	\ees
	We denote by~$\sqforms^r(X)$ the subspace of~$A^r(X)$ containing all differential forms~$\eta$ such that
	\bes
	(\eta,\eta)<\infty\qquad\text{and}\qquad(d\eta,d\eta)<\infty,
	\ees
	i.e.\ both~$\eta$ and its exterior derivative are \emph{square-integrable}.
	We refer to the elements of~$\sqforms^r(X)$ as~\emph{$L^2$-forms}.
	The~\emph{$L^2$-cohomology}~$\sqH^*(X)$ is the cohomology arising from the complex given by~$\sqforms^*(X)$.
	
	Note that if~${\alpha,\beta\in \sqforms^r(X)}$, then~$(\alpha,\beta)$ converges.
	This is a trivial consequence of the Cauchy--Schwarz inequality, since~$\lvert(\alpha,\beta)\rvert
	\le (\alpha,\alpha)^{1/2} (\beta,\beta)^{1/2}$.
	Therefore,~$(\cdot{,}\cdot)$ is a well-defined inner product on~$\sqforms^r(X)$.
	
	\begin{rem}\label{rem;completehencestokes}
	Since~$X$ is \emph{complete}, the Stokes Theorem holds by~\cite{gaffney}, in particular the integral on~$X$ of a top-degree exact~$L^2$-form vanishes.
	This implies that~$(\cdot{,}\cdot)$ extends to an inner product on the~$L^2$-cohomology.
%	\[
%	A^{\dim X}_{(2)}(X) \longrightarrow \CC,\qquad \eta\longmapsto\int_X\eta
%	\]
%	induces a well-defined map to~$L^2$ cohomology.
	\end{rem}
%	In~$L^2$-cohomology there is an analogue of Poincaré duality, recalled in the following result.
%	\begin{lemma}[$L^2$-Poincaré duality]\label{lemma;L2poincduality}
%	The map
%	\[
%	H^r_{(2)}(M)\longrightarrow H^{2n-r}_{(2)}(M)^*,
%	\qquad
%	\alpha \longmapsto \Big(
%	\eta_\alpha\colon \beta\mapsto\int_X \alpha\wedge\overline{\beta}
%	\Big)
%	\]
%	is an isomorphism.
%	\end{lemma}
%	If the cohomology group~$H^r_{(2)}(X)$ is finite dimensional, then it is isomorphic to the space~\emph{$L^2$ harmonic forms}~$\mathcal{H}_{(2)}^r(X)$ of degree~$r$.
	
	The proof of the Hodge Theorem in the compact case of Section~\ref{sec;compactkähler} is based on certain isomorphisms from cohomology groups to spaces of \emph{harmonic forms}.
	The latter forms can be defined also in the non-compact case and may be used to prove an~$L^2$-version of the Hodge decomposition, as we are going to recall.
	
	An \emph{$L^2$-harmonic form} on~$X$ is an~$L^2$-form annihilated by the usual Laplacian~$\Delta=dd^*+d^*d$ such that also its (adjoint) differential is square-integrable.
	We denote by~$\sqharm^r(X)$ the space of~$L^2$-harmonic forms of degree~$r$.
	For details regarding the following result we refer to~\cite[pp.~$5$-$6$]{dai} and~\cite[Sections~$2.4$ and~$2.5$]{noetherlefconj}.
	\begin{thm}[$L^2$-Hodge Theorem]
	Let~$X$ be a complete Kähler manifold.
	If the cohomology group~$\sqH^r(X)$ is finite dimensional, then it is isomorphic to~$\sqharm^r(X)$ and admits a Hodge decomposition
	\[
	\sqH^r(X)=\bigoplus_{p+q=r}\sqH^{p,q}(X).
	\]
	\end{thm}
	
	The inclusion~$\sqforms^r(X)\hookrightarrow A^r(X)$ induces a homomorphism
	\be\label{eq;homonotalwiso}
	\sqH^r(X)\to H^r(X)
	\ee
	which is in general neither injective nor surjective.
	As we will see later, if~$r$ is sufficiently small and~$X$ is an orthogonal Shimura variety, then~\eqref{eq;homonotalwiso} is an isomorphism.
	
	
%	It is also known that the Hodge star operator for~$L^2$-harmonic forms
%	\bes
%	\ast\colon\mathcal{H}_{(2)}^r(M)\longrightarrow\mathcal{H}_{(2)}^{2n-r}(M)
%	\ees
%	is an isomorphism; see e.g.~\cite[p.~$4$]{dai}.
%	
%	The following result provides a~$L^2$-version of Poincaré duality.
%	\begin{lemma}\label{lemma;poincaréduality}
%	The pairing
%	\bes
%	\mathcal{H}_{(2)}^r(M)\times\mathcal{H}_{(2)}^{2n-r}(M)\longrightarrow\CC,\qquad (\alpha,\beta)\longmapsto\int_X\alpha\wedge\overline{\beta}
%	\ees
%	is non-degenerate.
%	In particular, the map
%	\bes
%	\mathcal{H}_{(2)}^r(M)\longrightarrow\mathcal{H}_{(2)}^{2n-r}(M)^*,
%	\qquad
%	\alpha\longmapsto
%	\Big(
%	\eta_\alpha\colon\beta\mapsto\int_X\alpha\wedge\overline{\beta}
%	\Big)
%	\ees
%	is an isomorphism.
%	\end{lemma}
%	\begin{proof}
%	To prove that the pairing is non-degenerate, we show that for every~${\alpha\in\mathcal{H}_{(2)}^r(M)\setminus\{0\}}$ there exists~$\beta\in\mathcal{H}_{(2)}^{2n-r}(M)$ such that~$\int_X\alpha\wedge\overline{\beta}\neq 0$.
%	If we choose~$\beta=\ast\alpha$, then
%	\bes
%	\int_X\alpha\wedge\ast\overline{\alpha} = \int_X \langle\alpha,\alpha\rangle \nu_g\neq 0,
%	\ees
%	since~$\langle\alpha,\alpha\rangle$ is a non-zero positive-valued~$L^2$-function on~$M$.
%	\end{proof}

	
%	\subsection{The~$L^2$-cohomology on orthogonal Shimura varieties}
	
	\subsection{Orthogonal Shimura varieties}
	Let~$L$ be an even indefinite lattice of signature~$(n,2)$, and let~$G=\OO(L\otimes\RR)$.
	Let~$\domain$ be one of the two connected components of the complex manifold
	\[
	\mathcal{D}=\{ z\in\lattice\otimes\CC : \text{$q(z)=0$ and~$(z,\bar{z})<0$}\}/\CC^*.
	\]
	Then~$\domain$ is (a model of) the Hermitian symmetric domain of type~IV arising from~$G$.
	
	Let~$\disk(\lattice)$ be the discriminant kernel of~$\OO(\lattice)$, namely the kernel of the natural homomorphism~$\OO(\lattice)\to\Aut(\lattice'/\lattice)$, and let~$\disk{}^+(\lattice)=\disk(\lattice)\cap \OO^+(\lattice)$, where~$\OO^+(\lattice)$ is the group of isometries of~$\lattice$ preserving the component~$\domain$.
	\begin{defi}
	An \emph{orthogonal Shimura variety} is a quotient of the form~$X=\Gamma\backslash\domain$ for some subgroup~$\Gamma\subseteq\disk{}^+(\lattice)$ of finite index.
	\end{defi}
	By the Theorem of Baily and Borel, the analytic space~$X$ admits a unique algebraic structure that makes it a quasi-projective variety.
	This may be compact only if~$n\le 2$.
	
	Assume that~$X$ is \emph{smooth}, e.g.\ by choosing~$\Gamma$ to be torsion-free.
	It is well-known that there is a unique~$G$-invariant Kähler form~$\omega$ on~$\domain$ up to a factor.
	Throughout this article, we choose~$\omega$ to be the Kähler form arising as the first Chern form of the (metrized) tautological line bundle of~$\domain$.
	The Hermitian metric on that bundle is the one induced by the bilinear form of~$L$; see~\cite[Section~$1$]{kudla-integrals} for details.
	Since~$X$ is smooth, the invariant form~$\omega$ induces a Kähler form on~$X$, which we still denote by~$\omega$.
	This makes~$X$ a complete Kähler manifold.
	
%	Throughout this article we choose~$\omega$ to be the Kähler form arising from the metric on~$\domain$ induced by the Killing form of the Lie algebra of~$G$.
%	Since~$X$ is smooth, the invariant form~$\omega$ induces a Kähler form on~$X$, which we still denote by~$\omega$.
%	This equals the first Chern form of the Hodge bundle of~$X$, and makes~$X$ a complete Kähler manifold.
	
	Let~$\BBcomp{X}$ be the Baily--Borel compactification of~$X$, and let~$IH^*(\BBcomp{X})$ be its intersection cohomology.
	Zucker's conjectures, proved by Looijenga~\cite{looijenga} and Saper--Stern~\cite{saperstern}, asserts that there is an isomorphism~$\sqH^r(X)\cong IH^r(\BBcomp{X})$ for all~$r$.
	Moreover, since the dimension of the singular locus of~$\BBcomp{X}$ is at most~$1$, by~\cite[Theorem~$5.4$ and Remark~$5.5$~(ii)]{harriszucker} we also have an isomorphism~$IH^r(\BBcomp{X})\cong H^r(X)$ for all~$r<\dim(X)-1$.
	In fact, as illustrated in~\cite[Example~$3.4$]{noetherlefconj}, the map~\eqref{eq;homonotalwiso} is an isomorphism, and so is the composition
	\be\label{eq;3isoofcoho}
	\sqH^r(X)\cong IH^r(\BBcomp{X})\cong H^r(X)
	\qquad\text{for every~$r<\dim(X)-1$.}
	\ee
	From this we deduce that every de Rham cohomology class in~$H^r(X)$ admits a~$L^2$-representative.
	By~\cite[Theorem~$5.4$]{harriszucker} the composition of isomorphisms in~\eqref{eq;3isoofcoho} is a Hodge structure morphism.
	
	The Lefschetz operator~$\lefschetz$ on~$H^r(X)$ is injective for all~$r<\dim(X)-2$, although~$X$ is only quasi-projective.
	This can be deduced in terms of the analogue of the Hard Lefschetz Theorem~\cite[Corollary~$9.2.3$]{maxim} for the intersection cohomology~$IH^r(\BBcomp{X})$.
	In fact, the Kähler class~$\omega$ of~$X$ is identified with the Chern class of an ample line bundle on~$\BBcomp{X}$; see~\cite[Sections~$2.4$ and~$2.5$]{noetherlefconj}.
	Moreover, the Lefschetz decomposition is compatible with the Hodge structure; see e.g.~\cite[Remark~$2.3$]{noetherlefconj}.
	
	We summarize the results above as follows, denoting the inner product on~$H^r(X)$ induced by~$(\cdot{,}\cdot)$ under~\eqref{eq;3isoofcoho} with the same symbol.
	\begin{thm}\label{thm;lefdecL2}
	Let~$X$ be a smooth orthogonal Shimura variety of dimension~$n$.
	Let~$r<n-1$, so that the de Rham cohomology group~$H^{r}(X)$ is identified with~$\sqH^r(X)$ as in~\eqref{eq;3isoofcoho}.
	We have a decomposition
	\be\label{eq;lefdecL2inthm}
	H^r(X)=\bigoplus_{j\ge 0} \lefschetz^j \Hprim^{r-2j}(X),
	\ee
	which is orthogonal with respect to the inner product~$(\cdot{,}\cdot)$.
	Moreover, if~$r< n-2$, then the Lefschetz operator~$\lefschetz\colon H^r(X)\to H^{r+2}(X)$ is injective.
	Analogous results hold on each term of the Hodge decomposition of~$H^r(X)$.
	\end{thm}
	The fact that the Lefschetz decomposition of Theorem~\ref{thm;lefdecL2} is orthogonal with respect to~$(\cdot{,}\cdot)$ may be proved in the same way as for the compact case.
	In fact,~$L^2$-representatives of cohomology classes in different terms of the decomposition~\eqref{eq;lefdecL2inthm} are fiberwise orthogonal with respect to~$\langle\cdot{,}\cdot\rangle$ by e.g.~\cite[Proposition~$1.2.30$]{huybrechts}.
	
	\begin{rem}
%	Let~$X$ be a smooth orthogonal Shimura variety of dimension~$n$.
	If~$\beta\in \lefschetz^j \sqHprim^{r-j,r-j}(X)$, then
	\be\label{eq;formmatchonHggsq}
	(\alpha,\beta)
	=
	\frac{(-1)^{(r-j)(2(r-j)+1)}\cdot n! \cdot j!}{(n-2r+j)!}
	\int_X\alpha\wedge\overline{\beta}\wedge\omega^{n-2r}
	\qquad
	\text{for every~$\alpha\in \sqH^{r,r}(X)$}.
	\ee
	This formula is the analog of the one in Lemma~\ref{lemma;formmatchonHgg}, and holds even if~$X$ is non-compact, with the same proof.
	In fact, the formula given by~\cite[Proposition~$1.2.31$]{huybrechts} used to prove Lemma~\ref{lemma;formmatchonHgg} is \emph{local} and does not depend on the compactness of~$X$.
	The convergence of~$(\alpha,\beta)$ is ensured by the square-integrability.
	\end{rem}
	
	
%	\begin{lemma}\label{lemma;extfromL2poincdual}
%	\textcolor{red}{This result is what we want. However, its proof is wrong.
%	Since we do not see any easy way to correct it, we decided to look at other cohomology theories; see Section~\ref{sec;digonMumfordgf}.}
%	Let~$Z$ be a subvariety of codimension~$r$ in~$X$, and let~$[Z]$ be its de Rham fundamental class.
%	Suppose that~$r<\dim(X)-1$, and let~$\zeta$ be a~$L^2$-representative of the class~$[Z]$ under the isomorphism~\eqref{eq;homonotalwiso}.
%	Then
%	\bes
%	\int_X\zeta\wedge\omega^{n-r} = \int_Z\omega^{n-r}.
%	\ees
%	\end{lemma}
%	\textcolor{black}{\textbf{Note:} It is necessary to choose a~$L^2$-representative of~$[Z]$.
%	In fact, a priori there might be some differential form in the class~$[Z]$ that is not square-integrable.}
%	\begin{proof}
%	We recall that the fundamental class~$[Z]$ is defined to be the cohomology class in~$H^{2r}(X)$ corresponding to the functional
%	\be\label{eq;deffunctdualdr}
%	\eta_Z\colon H^{2(n-r)}_c(X)\longrightarrow\CC,\qquad\xi\longmapsto\int_Z\xi
%	\ee
%	under the isomorphism~$H^{2r}(X)\to H^{2(n-r)}_c(X)^*$ induced by Poincaré duality.
%	The Poincaré duality comes from the fact that the pairing
%	\bes
%	H^{2r}(X)\times H^{2(n-r)}_c(X)\longrightarrow\CC,\qquad ([\alpha],[\beta])\longmapsto\int_X\alpha\wedge\beta
%	\ees
%	is non-degenerate.
%	In particular, we have
%	\be\label{eq;wantedintfccompsup}
%	\int_X [Z]\wedge\xi = \int_Z\xi,\qquad\text{for every~$\xi$ with compact support.}
%	\ee
%	We are going to show that~\eqref{eq;wantedintfccompsup} is true also if~$[Z]$ is replaced with its harmonic~$L^2$-representative and~$\xi$ is square-integrable.
%	This concludes the proof since we may choose~$\xi$ to be the power~$\omega^{n-r}$ of the Kähler form.
%	
%	We denote by~$[Z]_{(2)}$ and~$\zeta$ the images of~$[Z]$ under the isomorphisms
%	\[
%	\begin{alignedat}{2}
%	H^{2r}(X)&\stackrel{\cong}{\longrightarrow} H^{2r}_{(2)}(X) &&\stackrel{\cong}{\longrightarrow} \mathcal{H}^{2r}_{(2)}(X),\\
%	[Z] & \longmapsto \,\,\,\, [Z]_{(2)} && \longmapsto \,\,\,\,\zeta.
%	\end{alignedat}
%	\]
%	Note that the de Rham cohomology class of the harmonic~$L^2$-form~$\zeta$ is~$[Z]$.
%	By Lemma~\ref{lemma;poincaréduality}, the form~$\zeta$ correspond to the functional
%	\be\label{eq;auxfinproof}
%	\mathcal{H}^{2(n-r)}_{(2)}(X)\longrightarrow \CC,\qquad \xi\longmapsto\int_X\zeta\wedge\xi,
%	\ee
%	which extends the functional~$\eta_Z$ in~\eqref{eq;deffunctdualdr} to differential forms~$\xi$ which are~$L^2$-integrable.
%	
%	We conclude by showing that~$\int_X\zeta\wedge\xi=\int_Z\xi$ for every~$L^2$-form~$\xi$.
%	We follow the wording of~\cite[Proof of Proposition~$4.19$]{bruinierbook}.
%	Since~$X$ is complete, there exists a proper smooth function~${\rho\colon X\to[0,\infty)}$.
%	The value~$\rho(x)$ may be defined as the geodesic distance of~$x$ from a fixed point of~$X$.
%	For every~$R>0$, we define
%	\bes
%	\mathcal{B}_R\coloneqq\{x\in X : \rho(x)\le R\}.
%	\ees
%	Note that~$(\mathcal{B}_R)_R$ is a collection of compact subsets that cover~$X$.
%	Let~$\mu\colon[0,\infty)\to[0,1]$ be a smooth function such that~$\mu|_{[0,1]}=1$ and~$\mu|_{[2,\infty)}=0$.
%	We define the compactly supported smooth function~$\sigma_R\colon X\to[0,1]$ as~$\sigma_R(x)\coloneqq\mu\big(\rho(x)/R\big)$.
%	It is clear that
%	\be\label{eq;finstepwwj}
%	\int_X \sigma_R \cdot \zeta\wedge \xi\xrightarrow[R\to\infty]{} \int_X\zeta\wedge\xi,
%	\ee
%	for every~$L^2$-form~$\xi$ of degree~$2(n-r)$.
%	Since~$\sigma_R\cdot\xi$ is a compactly supported form, we may compute the left-hand side of~\eqref{eq;finstepwwj} as~$\int_Z\sigma_R\cdot\xi$.
%	\textcolor{red}{This is actually not so easy. In fact, this can be done for diff forms~$\xi$ which are CLOSED. But~$\sigma_R\cdot\xi$ is not so...}
%	The limit of the latter integral when~$R\to\infty$ is~$\int_Z\xi$, and equals the right-hand side of~\eqref{eq;finstepwwj}.
%	\end{proof}
%	\begin{cor}\label{cor;nonvancohoclasses}
%	Let~$r<\dim(X)-1$.
%	If~$Z$ is a subvariety of codimension~$r$ in~$X$, then the fundamental class~$[Z]\in H^{2r}(X)$ is non-zero.
%	\end{cor}
%	\begin{proof}
%	If~$[Z]=0$, then also its~$L^2$-Poincaré dual vanishes.
%	By Lemma~\ref{lemma;extfromL2poincdual}, the latter is the functional
%	\bes
%	\eta_Z\colon H^{2(n-r)}(X)\longrightarrow\CC,\qquad \xi\longmapsto\int_Z\xi,
%	\ees
%	where~$\zeta$ is the~$L^2$-form as in Lemma~\ref{lemma;extfromL2poincdual}.
%	This implies that~$\eta_Z(\omega^{n-r})=0$, but this is not possible since~$\eta_Z(\omega^{n-r})=\Vol(Z)$, and it is well-known that the volume of any subvariety of dimension at least one in a Kähler manifold is positive.
%	\end{proof}
%	Funke and Millson~\cite[Theorem~$1.3$]{funkemillson} proved that infinitely many of the Fourier coefficients of the cohomology class of the Kudla--Millson theta function are non-zero, in the setting of locally symmetric space.
%	The results of this section may be applied to deduce that in the Hermitian case, i.e.\ when~$X$ admits a complex structure, the non-vanishing can be extended to \emph{all} Fourier coefficients of~$[\KMtheta{r}]$, as stated in the following result.
%	\begin{cor}
%	If~$r<\dim(X)-1$, then all Fourier coefficients of the cohomology class~$[\KMtheta{r}]$ are non-zero.
%	\end{cor}
%	\textcolor{blue}{Caveat: when generalizing this corollary to the non-unimodular case, we should better say something like ``all Fourier coefficients for which~$Z(T,\phi)\neq\emptyset$'' instead of only ``all Fourier coefficients''.}
%	\begin{proof}
%	The Fourier coefficients of~$[\KMtheta{r}]$ are the cohomology classes of special cycles of codimension~$r$.
%	These kind of cycles are~\emph{effective}, i.e.\ they are positive sums of codimension~$r$ subvarieties.
%	Hence, by Corollary~\ref{cor;nonvancohoclasses} the cohomology class of every special cycles is non-zero.
%	\end{proof}
	
	\end{document}
\documentclass[../Lefschetz_decomposition.tex]{subfiles}

	\begin{document}
	
	\subsection{Special cycles and the Kudla--Millson theta function}\label{sec;KMthetabasics}
	
	In the eighties Kudla and Millson constructed a theta function~$\KMtheta{\gengenus}=\KMtheta{\gengenus}(\tau,z)$ in the two variables~$\tau\in\HH_\gengenus$ and~$z\in\domain$, with several remarkable properties.
	We provide here a quick overview of its construction together with the main results needed for the purposes of this article.
	We refer to~\cite{kudlamillson-harmonic1}, \cite{kudlamillson-harmonic2} and~\cite{kudlamillson-intersection} for further information.
	
	Let~$V_\RR=L\otimes\RR$.
	Kudla and Millson defined a~$G$-invariant Schwartz function~$\KMschwartz{\gengenus}$ on~$V_\RR^\gengenus$ with values in the space of smooth closed~$(\gengenus,\gengenus)$-forms on~$\domain$.
	This is first constructed on a base-point of~$\domain$ by applying a differential operator to the standard Gaussian of~$V_\RR^\gengenus$, and then spread to the whole~$\domain$ by pulling-back; see~\cite[Section~$3$]{kudlamillson-harmonic1} for further information.
	
	Let~$\schrod{\gengenus}$ be the Schrödinger model of (the restriction of) the Weil representation of~$\Mp_{2\gengenus}(\RR)$ acting on the space of Schwartz functions on~$V_\RR^\gengenus$; see e.g.~\cite[Section~$4$]{funkemillsonhyp} for a detailed definition.
	Let~$k=1+n/2$.
	The Kudla--Millson theta function of genus~$\gengenus$ associated to the lattice~$\lattice$ is defined as
	\[
	\KMtheta{\gengenus}(\tau,z)=\det y ^{-k/2}
	\sum_{\discel\in D_\lattice^\gengenus}
	\sum_{\lambda \in \discel + \lattice^\gengenus}
	\big(\schrod{\gengenus}(g_\tau)\KMschwartz{\gengenus}\big)(\lambda,z)
	\mathfrak{e}_\discel,
	\]
	for every~$\tau=x+iy\in\HH_\gengenus$ and~$z\in\domain$, where~$g_\tau$ is the standard element of~$\Sp_{2\gengenus}(\RR)$ mapping~$i\in\HH_\gengenus$ to~$\tau$.
	If~$\genus=0$, we define~$\KMtheta{0}$ to be the constant function equal to~$1$.
	
	With respect to the symplectic variable, the theta function~$\KMtheta{\gengenus}$ behaves as a (non-holomorphic) Siegel modular form of weight~$k$ and genus~$\gengenus$ with respect to the Weil representation~$\weil{\lattice}{\gengenus}$.
	With respect to~$z\in\domain$, it is a~$\disk{}^+(\lattice)$-invariant closed~$(\gengenus,\gengenus)$-form on~$\domain$.
	Hence, the theta function descends to a function on the smooth orthogonal Shimura variety~$X=\Gamma\backslash\domain$, which we still denote by~$\KMtheta{\gengenus}$.
	
	An interesting feature of that theta function is that its cohomology class behaves as a \emph{holomorphic} Siegel modular form, since the image of~$\KMtheta{\gengenus}$ under~$\bar\partial$ in the symplectic variable~$\tau$ is an exact differential form on~$X$. 
	This means that the de Rham cohomology class~$\drcoho{\KMtheta{\gengenus}}$ may be regarded as an element of~$M^k_{\gengenus,\lattice}\otimes H^{\gengenus,\gengenus}(X)$.
	
	To describe the Fourier expansion of~$\drcoho{\KMtheta{\gengenus}}$ we need to introduce the so-called \emph{special cycles} of~$X$.
	For any vector~$\lambda\in (L')^\gengenus$ we denote by~$\lambda^\perp$ the set of points in~$\domain$ orthogonal to every component of~$\lambda$.
	The special cycle of~$X$ associated with~$\discel\in D_\lattice^\gengenus$ and~$T\in q(\discel)+\halfint{\gengenus}$, where~$T\ge 0$, is
	\be\label{eq;defspcy}
	Z(T,\discel) = \Gamma\big\backslash\bigg(
	\sum_{\substack{\lambda\in\discel + \lattice^\gengenus\\ q(\lambda)=T}}
	\lambda^\perp
	\bigg).
	\ee
	If non-empty, this is a codimension~$\rk (T)$ effective cycle on~$X$.
	We denote the cohomology class induced by~$Z(T,\discel)$ in~$H^{2\rk(T)}(X,\CC)$ by~$\drcoho{Z(T,\discel)}$.
	Note that~$Z(0,0)=X$ and that if~$\alpha\neq 0$, then~$Z(0,\alpha)=\emptyset$.
	Special cycles of codimension~$1$ are also known as \emph{Heegner divisors}.
	
	The group~$\GL_\gengenus(\ZZ)$ acts on~$\halfint{\gengenus}$ under the action~$T\mapsto M^t T M$, for~$T\in\halfint{\gengenus}$ and~$M\in\GL_\gengenus(\ZZ)$.
	This induces the symmetry
	\be\label{eq;symmetryspcy}
	Z(T,\discel) = Z(M^tTM,\discel M).
	%\qquad\text{for every~$T\in\halfint{\gengenus}$ and~$M\in \GL_\gengenus(\ZZ)$.}
	\ee
	% Note: $q(\discel M)=M^t q(\discel) M$.
	
	The cohomology class~$\drcoho{\KMtheta{\gengenus}}$ is the generating series of (the cohomology classes of) the codimension~$\gengenus$ special cycles, in fact the Fourier expansion of~$\drcoho{\KMtheta{\gengenus}}$ is
	\[
	\drcoho{\KMtheta{\genus}(\tau)}=\sum_{\discel\in D_\lattice^\gengenus}\sum_{\substack{T\in q(\discel) + \halfint{\gengenus}\\ T\ge0}}
	\drcoho{Z(T,\discel)}
	\wedge
	\drcoho{-\omega}^{\gengenus-\rk(T)}
	q^T \mathfrak{e}_\discel.
	\]
	
	\begin{rem}\label{rem;resandpullbKMthetafunct}
	Let~$M\subset\lattice$ be a sublattice of~$L$.
	Let~$\Gamma_L=\disk{}^+(\lattice)$ and let~$\Gamma_M=\disk{}^+(M)$.
	We denote by~$\KMtheta{\genus}^\lattice$, resp.~$\KMtheta{\genus}^M$, the Kudla--Millson theta function on~$X_L=\Gamma_L\backslash\domain$ arising from~$L$, resp.\ on~$X_M=\Gamma_M\backslash\domain$ arising from~$M$.
	It is easy to see that if~$\pi\colon X_M\to X_\lattice$ is the quotient map
%	\\\textcolor{blue}{\textbf{This map is well-defined:}\\
%	It is enough to prove that~$\Gamma_M\subseteq\Gamma_L$.
%	This can be showed using the inclusions
%	\[
%	M \subseteq L \subseteq L' \subseteq M'
%	\qquad\implies\qquad
%	0 \subseteq L/M \subseteq L'/M \subseteq M'/M.
%	\]
%	Since~$\Gamma_M$ acts trivially on~$M'/M$, then acts trivially on~$L/M$, in particular~$\Gamma_M\subseteq O(L)$.
%	Since~$L'/L \cong  (L'/M)/ (L/M)$, we deduce that~$\Gamma_M$ acts trivially on~$L'/L$, hence~$\Gamma_M\subseteq\Gamma_L$.}\\
	and if~$\tr_{\genus,\lattice/M}\colon M^k_{\genus,M}\to M^k_{\genus,\lattice}$ is as in Lemma~\ref{lemma;genofBYandS}, then
	\be\label{eq;resandpullbKMthetafunct}
	\pi^*\big(\drcoho{\KMtheta{\genus}^\lattice}\big)
	=
	\tr_{\genus,\lattice/M}\big(\drcoho{\KMtheta{\genus}^M}\big).
	\ee
	
%	\textcolor{blue}{\textbf{Check of the previous assertion:}
%	\\
%	Let~$A^k_{\genus,\lattice}$ be the space of~$\CC[D_\lattice^\genus]$-valued smooth (not necessarily holomorphic) Siegel modular forms of weight~$k$ with respect to the representation~$\weil{\lattice}{\genus}$.
%	What is written in Lemma~\ref{lemma;genofBYandS} works also if~$M^k_{\genus,L}$ is replaced with~$A^k_{\genus,L}$.
%	To prove~\eqref{eq;resandpullbKMthetafunct} is then enough to prove that
%	\be\label{eq;resandpullbKMthetafunctblue}
%	\KMtheta{\genus}^L
%	=
%	\tr_{\genus,\lattice/M}\big(\KMtheta{\genus}^M)
%	\ee
%	as differential forms on the Hermitian symmetric domain~$\domain$.
%	\\
%	Recall that~$\KMtheta{\genus}^M =\sum_{\alpha\in D_M^\genus}\big(\KMtheta{\genus}^M\big)_\alpha\mathfrak{e}_\alpha$ with
%	\bas
%	\big(\KMtheta{\genus}^M\big)_\alpha(\tau)
%	=
%	\det y^{-k/2}
%	\sum_{\lambda\in\alpha + M^\genus}
%	\big(\schrod{\gengenus}(g_\tau)\KMschwartz{\genus}\big)(\lambda)
%	,\qquad
%	\text{$\tau\in\HH_\genus$.}
%	\eas
%	Let~$H\coloneqq (\lattice/M)^\genus$.
%	Recall that~$H\subset D_M^\genus$ and~$H\subset H^\perp$ in~$D_M^\genus$.
%	Moreover
%	\bes
%	H^\perp = (\lattice'/M)^\genus
%	\qquad
%	\text{and}
%	\qquad
%	H^\perp/H\cong D_\lattice^\genus.
%	\ees
%	We may then compute
%	\bas
%	\tr_{\genus,\lattice/M}\big(\KMtheta{\genus}^M)
%	&=
%%	\det y^{-k/2}
%	\sum_{\overline{\beta}\in D_L^\genus}
%	\sum_{\delta\in(\lattice/M)^\genus}
%	\big(\KMtheta{\genus}^M\big)_{\delta+\beta}
%%	\sum_{\lambda\in\delta + \beta + M^\genus}
%%	\big(\schrod{\gengenus}(g_\tau)\KMschwartz{\genus}\big)(\lambda)
%	\mathfrak{e}_{\overline{\beta}}
%	\\
%	&=
%	\det y^{-k/2}
%	\sum_{\overline{\beta}\in H^\perp/H}
%	\sum_{\delta\in H}
%	\sum_{\lambda\in\delta + \beta + M^\genus}
%	\big(\schrod{\gengenus}(g_\tau)\KMschwartz{\genus}\big)(\lambda)
%	\mathfrak{e}_{\overline{\beta}}
%	\\
%	&=
%	\det y^{-k/2}
%	\sum_{\beta\in H^\perp}
%	\sum_{\lambda\in \beta + M^\genus}
%	\big(\schrod{\gengenus}(g_\tau)\KMschwartz{\genus}\big)(\lambda)
%	\mathfrak{e}_{\overline{\beta}}.
%	\eas
%	We then deduce that to prove~\eqref{eq;resandpullbKMthetafunctblue} it is enough to show that
%	\be\label{eq;unioninblTYJBGHG}
%	\bigcup_{\substack{\beta\in(L'/M)^\genus \\ \overline{\beta}=\alpha}}\beta + M^\genus
%	=
%	\alpha + L^\genus
%	\qquad
%	\text{for all~$\alpha\in D_\lattice^\genus$.}
%	\ee
%	Surely if~$\beta\in(L'/M)^\genus$ is such that~$\overline{\beta}=\alpha$, then~$\beta + M^\genus \subseteq \alpha + L^\genus$, hence
%	\[
%	 \bigcup_{\substack{\beta\in(L'/M)^\genus \\ \overline{\beta}=\alpha}}\beta + M^\genus
%	\subseteq
%	\alpha + L^\genus
%	\qquad
%	\text{for all~$\alpha\in D_\lattice^\genus$.}
%	\]
%	We want to prove that the previous inclusion is an equality.
%	Suppose it is not, namely that there exists~$\lambda\in\alpha+\lattice^\genus$ such that~$\lambda\not\in\beta + M^\genus$ for all~$\beta\in(L'/M)^\genus$ such that~$\overline{\beta}=\alpha$.
%	This means that the map~$(L'/M)^\genus\to (L'/L)^\genus$, $\beta\mapsto\overline{\beta}$, is not surjective.
%	This is clearly false, hence~\eqref{eq;unioninblTYJBGHG} is true.}
	\end{rem}
	
	We conclude this section by recalling a product formula for the above theta functions; see e.g.~\cite[$(9.2)$]{kudla-algcycles}.
	As usual, we identify products of group algebras as in~\eqref{eq;identgroupalg}.
	Let~$r_1,r_2\in\ZZ_{\ge 0}$.
	If~$\tau_1\in\HH_{r_1}$ and~$\tau_2\in\HH_{r_2}$, then
	\be\label{eq;wedgeprodformforKMtheta}
	\KMtheta{r_1}(\tau_1)
	\wedge
	\KMtheta{r_2}(\tau_2)
	=
	\KMtheta{r_1+r_2}\big(
	\begin{smallmatrix}
	\tau_1 & 0\\
	0 & \tau_2
	 \end{smallmatrix}\big).
	\ee
	
%	\begin{thm}[Kudla]\label{thm;formulaintspeccycl}
%	If~$\tau_1\in\HH_{r_1}$ and~$\tau_2\in\HH_{r_2}$, then
%	\[
%	\drcoho{\KMtheta{r_1}(\tau_1)} \wedge \drcoho{\KMtheta{r_2}(\tau_2)}
%	= \drcohowide{\KMtheta{r_1+r_2}\big(
%	\begin{smallmatrix}
%	\tau_1 & \\
%	 & \tau_2
%	 \end{smallmatrix}\big)}.
%	\]
%	\end{thm}
%	\textcolor{blue}{\begin{proof}
%	We check the formula above in terms of Fourier expansions.
%	We may compute
%	\ba\label{eq;formulaintspeccycl}
%	&\drcohowide{\KMtheta{r_1+r_2}\big(
%	\begin{smallmatrix}
%	\tau_1 & \\
%	 & \tau_2
%	 \end{smallmatrix}\big)}
%	 \\
%	&=
%	\sum_{\discel\in D_\lattice^{r_1+r_2}}
%	\sum_{\substack{T\in q(\discel) + \halfint{r_1+r_2}\\ T\ge 0}}
%	\drcoho{Z(T,\discel)}\wedge \drcoho{-\omega}^{r_1+r_2-\rk T}
%	e\big(
%	\tr  T\big(\begin{smallmatrix}
%	\tau_1 & \\
%	 & \tau_2
%	 \end{smallmatrix}\big)\big)
%	 \mathfrak{e}_\discel
%	\\
%	&=
%	\sum_{\discel\in D_\lattice^{r_1+r_2}}
%	\sum_{\substack{
%	T\in q(\discel)+\halfint{r_1+r_2}
%	\\
%	T
%	= \big(\begin{smallmatrix}
%	T_1 & \ast \\
%	\ast & T_2
%	\end{smallmatrix}\big)}
%	}
%	\drcoho{Z(T,\discel)}
%	\wedge
%	\drcoho{-\omega}^{r_1+r_2-\rk T}
%	e\Big(
%	\tr \big(\begin{smallmatrix}
%	T_1\tau_1 & \ast \\
%	\ast & T_2\tau_2
%	\end{smallmatrix}\big)\Big)
%	\mathfrak{e}_\discel
%	\\
%	&=
%	\sum_{\discel_1\in D_\lattice^{r_1}}
%	\sum_{\substack{T_1\in q(\discel_1) + \halfint{r_1} \\ T_1\ge0}}
%	\sum_{\discel_2\in D_\lattice^{r_2}}
%	\sum_{\substack{T_2\in q(\discel_2) + \halfint{r_2} \\ T_2\ge 0}}
%	\Big(
%	\sum_{T
%	= \big(\begin{smallmatrix}
%	T_1 & \ast \\
%	\ast & T_2
%	\end{smallmatrix}\big)
%	}\drcoho{Z\big(T,(\discel_1,\discel_2)\big)}\wedge\drcoho{-\omega}^{r_1+r_2-\rk T}
%	\Big)
%	\\&\quad
%	\times e(\tr T_1\tau_1)\mathfrak{e}_{\discel_1}
%	\otimes
%	e(\tr T_2 \tau_2)\mathfrak{e}_{\discel_2},
%	\ea
%	where we denote by~$\ast$ any matrix making~$T$ symmetric and positive semidefinite.
%%	Remember: Principal submatrices of a positive semidefinite matrix are always positive semidefinite. (Principal submatrices of a matrix M are those matrices resulting from removing the ith row and column from M, where i runs from 1 to the order of the matrix.)
%	By~\cite[Theorem~$6.2$]{kudla-algcycles} we have
%	\begin{align*}
%	&\drcoho{Z(T_1,\discel_1)} \wedge \drcoho{Z(T_2,\discel_2)} \wedge \drcoho{-\omega}^{r_1+r_2-\rk T_1 - \rk T_2}
%	\\
%	&=
%	\sum_{T={\big(\begin{smallmatrix}
%	T_1 & \ast \\
%	\ast & T_2
%	\end{smallmatrix}\big)}}
%	\drcoho{Z\big(T,(\discel_1,\discel_2)\big)}\wedge \drcoho{-\omega}^{r_1+r_2-\rk T}.
%	\end{align*}
%	\textcolor{blue}{\textbf{Keep in mind:} For us~$Z(T)$ is a special cycle of codimension~$\rk(T)$.
%	In particular, we use the same notation of~\cite{kudla-speccyc} and not of~\cite{kudla-algcycles};
%	see for instance~\cite[(3.9)]{kudla-speccyc}.}\\
%	We may then immediately see that the right-hand side of~\eqref{eq;formulaintspeccycl} equals~$\drcoho{\KMtheta{r_1}(\tau_1)}
%	\wedge
%	\drcoho{\KMtheta{r_2}(\tau_2)}$.
%	\end{proof}}
	
	
	\subsection{Adelic and primitive special cycles}\label{sec;adelicspcy}
	
	Let~$V=\lattice\otimes\QQ$, and let~$\mathcal{S}\big(V^\genus(\Ab_f)\big)$ be the space of Schwartz functions on~$V^\genus(\Ab_f)$.
	For each prime~$p$, we let~$\lattice_p=\lattice\otimes\ZZ_p$ and let~$K_p$ be the subgroup of~$G(\QQ_p)$ which leaves~$\lattice_p$ stable and acts as the identity on~$\lattice_p'/\lattice_p$.
	Then~$K=\prod_p K_p$ is an open compact subgroup of~$G(\Ab_f)$.
	As above, we denote by~$\Gamma$ any finite index subgroup of~$K\cap G(\QQ)^+=\disk{}^+(L)$.
	
	\begin{defi}
	Given a~$K$-invariant Schwartz function~$\varphi\in\mathcal{S}(V^\genus(\Ab_f)\big){}^K$ and~$T\in\Sym_{\genus}(\QQ)$, we define the \emph{adelic special cycle}~$Z(T,\varphi,K)$ in~$X=\Gamma\backslash\domain$ as
	\[
	Z(T,\varphi,K)
	=
	\Gamma\big\backslash\sum_{\substack{x\in V^\genus \\ q(x)=T}}
	\varphi(x)\cdot x^\perp.
	\]
	\end{defi}
	
	Let~$\mathcal{S}_{L^\genus}\subset\mathcal{S}\big(V^\genus(\Ab_f)\big)$ be the subspace spanned by the characteristic functions~$\varphi_\lambda$ of the cosets~$\lambda + \widehat{\lattice^\genus}$ with~$\lambda\in (\lattice')^\genus$, where~$\widehat{\lattice^\genus}=\lattice^\genus\otimes\prod_p\ZZ_p$.
	The special cycles defined in Section~\ref{sec;KMthetabasics} are particular adelic special cycles, more precisely
	\[
	Z(T,\lambda + L^\genus)
	=
	Z(T,\varphi_\lambda,K)\qquad
	\text{for all~$\lambda\in (L')^\genus$.}
	\]
	
	We recall below the main result of~\cite{noetherlefconj} on the span of cohomology classes of adelic special cycles.
	
	\begin{thm}[Bergeron--Li--Millson--Moeglin]\label{thm;noetherlef}
	If~$\Gamma=\disk{}^+(\lattice)$ and~$\genus <(n+1)/3$, then~$H^{\genus,\genus}(X)$ is spanned by the classes of adelic special cycles of the form
	\[
	\bigdrcoho{Z(T,\varphi,K)}\wedge\bigdrcoho{\omega^{\genus-\rk T}},
	\qquad
	\text{$T\in\Sym_{\genus}(\QQ)$, $\varphi\in\mathcal{S}(V^\genus(\Ab_f)\big){}^K$.}
	\]
	\end{thm}
	It is natural to ask whether a result analogous to Theorem~\ref{thm;noetherlef} holds also for (non-adelic) special cycles.
	
	\begin{rem}\label{rem;BLMMhasimprecision}
	It is stated in~\cite[Remark~$2.4$]{noetherlefconj} that if~$\Gamma=\disk{}^+(L)$, then the adelic special divisors on~$X=\Gamma\backslash\domain$ are linear combinations of Heegner divisors.
	The argument provided in~\cite{noetherlefconj}, i.e.\ that every~$K$-invariant Schwartz function~$\varphi\in\mathcal{S}(V(\Ab_f))$ corresponds to a linear combination of characteristic functions on~$D_L$, is incorrect.
	Hence, we can not deduce from~\cite{noetherlefconj} that~$H^{1,1}(X)$ is generated by (non-adelic) Heegner divisors.
%	\textcolor{blue}{Nevertheless, this is known to be true if~$L$ splits off \emph{two hyperbolic planes}.
%	This may be proved by rewriting the irreducible components of adelic special divisors in terms of \emph{primitive Heegner divisors}; see~\cite[Lemma~$4.2$ and Lemma~$4.3$]{bruiniermöller}.}
	\end{rem}
	
	We illustrate below a non-adelic version of Theorem~\ref{thm;noetherlef}, proved in terms of primitive special cycles.
	For the sake of simplicity, we illustrate it for the cases of codimension~$1$ and~$2$.
	Together with the injectivity of the Kudla--Millson lift (known up to genus~$2$) and the main results of Section~\ref{sec;gen2}, the non-adelic version of Theorem~\ref{thm;noetherlef} will imply dimension formulas for cohomology groups of orthogonal Shimura varieties; see Corollaries~\ref{cor;dimcoho1} and~\ref{cor;dimcoho2}.
	
	Recall that a vector~$v\in L\otimes\QQ$ is primitive in~$L'$ if~$\QQ v \cap L' = \ZZ v$.
	Examples are the zero vector of~$L'$, and vectors in~$L$ of square-free norm for~$L$ unimodular.
	\begin{defi}
	Let~$\discel\in D_\lattice^\genus$.
	Let~$\primvec{\lattice}{\genus}{\discel}$ be the set of tuples~$\lambda\in\alpha +\lattice^\genus$ with primitive entries in~$L'$.
	For all~$T\in q(\alpha)+\halfint{\genus}$, the \emph{primitive special cycle}~$\Zprim(T,\discel)$ is defined as
	\[
	\Zprim(T,\discel)
	=
	\Gamma\big\backslash\bigg(
	\sum_{\substack{\lambda\in\primvec{\lattice}{\genus}{\discel}\\ q(\lambda)=T}}
	\lambda^\perp
	\bigg).
	\]
	\end{defi}
	Let~$\alpha\in D_\lattice^\genus$, and let~$\diag(\genus,\ZZ_{>0})$ be the set of diagonal~$\genus\times\genus$ matrices with coefficients in~$\ZZ_{>0}$.
	We say that~$R\in\diag(\genus,\ZZ_{>0})$ divides~$T\in q(\alpha)+\halfint{\genus}$, in short~$R|T$, if~$R^{-1}TR^{-1}\in q(\beta)+\halfint{\genus}$ for some~$\beta\in D_\lattice^\genus$.
	Note that we consider tuples in~$L^\genus\otimes\QQ$ as row vectors.
	
	Classical and primitive special cycles are related as follows.
	\begin{lemma}\label{lemma;spcyclesandprim}
	Let~$\alpha\in D_\lattice^\genus$ and~$T\in q(\alpha)+\halfint{\genus}$.
	Then
	\be\label{eq;spcyclesandprim}
	Z(T,\discel)
	=
	\sum_{\substack{R\in\diag(\genus,\ZZ_{>0})
	\\
	R|T}}
	\sum_{\substack{\beta\in D_\lattice^\genus \\ \beta R=\alpha}}
	\Zprim(R^{-1} T R^{-1},\beta)
	\ee
	and
	\be\label{eq;spcyclesandprim2}
	\Zprim(T,\discel)
	=
	\sum_{\substack{R\in\diag(\genus,\ZZ_{>0}) \\ R|T}}
	\mu(R)
	\sum_{\substack{\beta\in D_\lattice^\genus \\ \beta R = \discel}}
	Z(R^{-1}TR^{-1} , \beta),
	\ee
	where~$\mu(R)=\prod_{i=1}^\genus \mu(R_{i,i})$ for all~$R\in\diag(\genus,\ZZ_{>0})$.
	\end{lemma}
	\begin{proof}
	It is analogous to~\cite[Lemma~$4.2$]{bruiniermöller}.
%	\textcolor{blue}{The proof looks like the one of~\cite[Lemma~$4.2$]{bruiniermöller}.
%	I suggest we drop it.}
%	\\
%	Since every element of~$\lattice'$ can be written uniquely as a positive multiple of a primitive element of~$\lattice'$, we deduce that every tuple~$\lambda\in(\lattice')^\genus$ can be written as~$\mu R$ for some primitive~$\mu\in (\lattice')^\genus$ and~$R\in\diag(\genus,\ZZ_{>0})$.
%	This immediately implies~\eqref{eq;spcyclesandprim}.
%	
%	We prove~\eqref{eq;spcyclesandprim2} by induction on the number of matrices~$R\in\diag(\genus,\ZZ_{>0})$ dividing~$T$.
%	If~$T$ is not divisible by any~$R\neq I_\genus$, then~\eqref{eq;spcyclesandprim2} is obvious.
%	Otherwise, from~\eqref{eq;spcyclesandprim} and inductive hypothesis we have
%	\bas
%	\Zprim(T,\discel)
%	&=
%	Z(T,\discel)
%	-
%	\sum_{\substack{R\neq I_\genus \\ R|T}}
%	\sum_{\substack{\beta\in D_\lattice^\genus \\ \beta R = \discel}}
%	\Big(
%	\sum_{S|R^{-1}TR^{-1}}
%	\mu(S)\sum_{\substack{\sigma\in D_L^{\genus}\\ \sigma S=\beta}}
%	Z\big((SR)^{-1}T(SR)^{-1},\sigma\big)
%	\Big)
%	\\
%	&=
%	Z(T,\discel)
%	-
%	\sum_{\substack{R\neq I_\genus \\ R|T}}
%	\Big(
%	\sum_{S|R^{-1}TR^{-1}}
%	\mu(S)
%	\sum_{\substack{\beta\in D_\lattice^\genus \\ \beta R = \discel}}
%	\sum_{\substack{\sigma\in D_L^{\genus}\\ \sigma S=\beta}}
%	Z\big((SR)^{-1}T(SR)^{-1},\sigma\big)
%	\Big)
%	\\
%	&=
%	Z(T,\discel)
%	-
%	\sum_{\substack{R\neq I_\genus \\ R|T}}
%	\Big(
%	\sum_{S|R^{-1}TR^{-1}}
%	\mu(S)
%	\sum_{\substack{\sigma\in D_L^{\genus}\\ \sigma SR=\discel}}
%	Z\big((SR)^{-1}T(SR)^{-1},\sigma\big)
%	\Big).
%	\eas
%	We let~$Q=SR$ and rewrite the sum above by grouping the summands contributing to~$Z\big(Q^{-1}TQ^{-1},\sigma\big)$.
%	We obtain
%	\ba\label{eq;fromindasinbrmö} 
%	\Zprim(T,\discel)
%	&=
%	Z(T,\discel)
%	-
%	\sum_{\substack{Q\in\diag(\genus,\ZZ_{>0}) \\ Q|T}}
%	\sum_{\substack{\sigma\in D_L^\genus \\ \sigma Q=\alpha}}
%	\Big(
%	\sum_{\substack{R\neq I_\genus,\, R|T \\ QR^{-1}\in\diag(\genus,\ZZ_{>0})}}
%	\mu(QR^{-1})
%	\Big)
%	Z(Q^{-1}TQ^{-1},\sigma)
%	\\
%	&=
%	Z(T,\discel)
%	-
%	\sum_{\substack{Q\in\diag(\genus,\ZZ_{>0}) \\ Q|T}}
%	\sum_{\substack{\sigma\in D_L^\genus \\ \sigma Q=\alpha}}
%	\Big(
%	\sum_{\substack{R\in\diag(\genus,\ZZ_{>0}),\, R\neq I_\genus \\ QR^{-1}\in\diag(\genus,\ZZ_{>0})}}
%	\mu(QR^{-1})
%	\Big)
%	Z(Q^{-1}TQ^{-1},\sigma).
%	\ea
%	Since
%	\[
%	\sum_{\substack{R\in\diag(\genus,\ZZ_{>0}),\, R\neq I_\genus \\ QR^{-1}\in\diag(\genus,\ZZ_{>0})}}
%	\mu(QR^{-1})
%	=
%	\sum_{\substack{R\in\diag(\genus,\ZZ_{>0}) \\ QR^{-1}\in\diag(\genus,\ZZ_{>0})}}
%	\mu(QR^{-1})
%	- \mu(Q)
%	=
%	\delta_{Q,I_\genus} - \mu(Q),
%	\]
%%	\textcolor{blue}{\textbf{Check of the last assertion:}
%%	We should check that~$\sum_{\substack{R\in\diag(\genus,\ZZ_{>0}) \\ QR^{-1}\in\diag(\genus,\ZZ_{>0})}}
%%	\mu(QR^{-1})=\delta_{Q,I_\genus}$.
%%	This follows immediately from the definition.
%%	In fact, since
%%	\bas
%%	\mu(QR^{-1})
%%	=
%%	\prod_{j=1}^\genus\mu\big((QR^{-1})_{j,j}\big)
%%	=
%%	\prod_{j=1}^\genus\mu\big(Q_{j,j}(R^{-1})_{j,j}\big)
%%	=
%%	\prod_{j=1}^\genus\mu\big(Q_{j,j}(R_{j,j})^{-1}\big)
%%	\eas
%%	we may compute
%%	\bas
%%	\sum_{\substack{R\in\diag(\genus,\ZZ_{>0}) \\ QR^{-1}\in\diag(\genus,\ZZ_{>0})}}
%%	\mu(QR^{-1})
%%	&=
%%	\sum_{\substack{R\in\diag(\genus,\ZZ_{>0}) \\ QR^{-1}\in\diag(\genus,\ZZ_{>0})}}
%%	\prod_{j=1}^\genus\mu\big(Q_{j,j}(R_{j,j})^{-1}\big)
%%	=
%%	\prod_{j=1}^\genus
%%	\sum_{\substack{r_j\in\ZZ_{>0} \\ r_j|Q_{j,j}}}
%%	\mu\big(Q_{j,j}(r_j)^{-1}\big)
%%	\\
%%	&=
%%	\prod_{j=1}^\genus
%%	\sum_{\substack{r_j\in\ZZ_{>0} \\ r_j|Q_{j,j}}}
%%	\mu\big(r_j\big)
%%	=
%%	\prod_{j=1}^\genus
%%	\delta_{Q_{j,j},1}
%%	=
%%	\delta_{Q,I_\genus}.
%%	\eas
%%	}
%	we deduce from~\eqref{eq;fromindasinbrmö} that
%	\begin{align*}
%	\Zprim(T,\discel)
%	&=
%	Z(T,\discel)
%	-
%	\sum_{\substack{Q\in\diag(\genus,\ZZ_{>0}) \\ Q|T}}
%	\Big(\delta_{Q,I_\genus} - \mu(Q)\Big)
%	\sum_{\substack{\sigma\in D_L^\genus \\ \sigma Q=\alpha}}
%	Z(Q^{-1}TQ^{-1},\sigma)
%	\\
%	&=
%	\sum_{\substack{Q\in\diag(\genus,\ZZ_{>0}) \\ Q|T}}
%	\mu(Q)
%	\sum_{\substack{\sigma\in D_L^\genus \\ \sigma Q=\alpha}}
%	Z(Q^{-1}TQ^{-1},\sigma)
%	.
%	\qedhere
%	\end{align*}
	\end{proof}
	We denote by~$E_8$ the~$8$th root lattice, and by~$U$ the hyperbolic plane.
	They are even unimodular lattices of signature~$(8,0)$ and~$(1,1)$ respectively.
	\begin{lemma}\label{lemma;orbitsofdisck}
	Let~$L=U^{\oplus 2}\oplus E_8\oplus E$ for some positive definite even lattice~$E$.
	There exists a constant~$C=C(E)\in\QQ_{>0}$ such that the following property is satisfied for all~$\alpha\in D_L^2$.
	If~$a,b\in \alpha + (L')^{2}$ have primitive entries in~$L'$,~$q(a)=q(b)>0$ and~$q(a_1)>C$, then~$a$ and~$b$ lie in the same~$\disk{}^+(L)$-orbit.
	\end{lemma}
%	\textcolor{red}{Possible generalization:
%	\begin{lemma}
%	Let~$L=U^{\oplus 2}\oplus E_8^{\oplus(\genus-1)}\oplus E$ for some positive definite even lattice~$E$.
%	There exists a positive definite matrix~$C=C(E)\in\Sym_{\genus-1}(\QQ)$ such that the following property is satisfied for all~$\alpha\in D_L^2$.
%	If~$a,b\in \alpha + (L')^{2}$ have primitive entries in~$L'$,~$q(a)=q(b)>0$ and~$q(a_1,\dots,a_{\genus-1})-C>0$, then~$a$ and~$b$ lie in the same~$\disk{}^+(L)$-orbit.
%	\end{lemma}
%	\begin{proof}
%	We follow the wording of~\cite[Proof of Lemma~$4.4$]{freitaghermann}.
%	Let~$E$ be an even sublattice of~$L$ such that~$L$ splits in orthogonal direct sums as
%	\[
%	L=U_1 \oplus U_2\oplus E_8^1\oplus\cdots\oplus E_8^{\genus-1}\oplus E,
%	\]
%	where~$U_1$ and~$U_2$, resp.~$E_8^i$, are copies of the hyperbolic plane, resp.\ of~$E_8$.
%	\\
%	There exists an isomorphism~$U_1\oplus U_2\to \Mat_2(\ZZ)$ such that the quadratic form of~$U_1\oplus U_2$ corresponds to the determinant on~$\Mat_2(\ZZ)$.
%	\\
%	The action of~$\SL_2(\ZZ)$ on~$\Mat_2(\ZZ)$ by multiplication from both sides induces isometries in~$\OO^+(U_1\oplus U_2)$.
%	The theorem of elementary divisors for~$\SL_2(\ZZ)$ implies that every element of~$L'$ maps to~$L_1'$, where~$L_1\coloneqq \{0\}\oplus U_2\oplus E_8^1\oplus\cdots\oplus E_8^{\genus-1}\oplus E$, under an isometry in~$\OO^+(U_1\oplus U_2)$.
%	\\
%	Let
%	\[
%	C_\alpha\coloneqq \min\{
%	q(\lambda) : \text{$\lambda\in \alpha + E$}
%	\},
%	\qquad
%	\alpha\in D_{E}\cong D_{L}.
%	\]
%	Let~$C=\max_{\alpha} C_\alpha$, and let~$a,b\in(\lattice')^2$ be as in the statement of Lemma~\ref{lemma;orbitsofdisck}.
%	\end{proof}
%	}
	\begin{proof}
	We adapt the argument of~\cite[Proof of Lemma~$4.4$]{freitaghermann}.
	Let
	\[
	L=U_1 \oplus U_2\oplus E_8\oplus E,
	\]
	where~$U_1$ and~$U_2$ are copies of the hyperbolic plane.
	
	There exists an isomorphism~$U_1\oplus U_2\to \Mat_2(\ZZ)$ such that the quadratic form of~$U_1\oplus U_2$ corresponds to the determinant on~$\Mat_2(\ZZ)$.
%	\\\textcolor{blue}{
%	\textbf{Check:} The isomorphism is given by
%	\be\label{eq;stisoblue}
%	U_1\oplus U_2 \to \Mat_2(\ZZ),\qquad (a f_1 + d f_2, b f_3 + c f_4)\mapsto \left(\begin{smallmatrix}
%	a & b\\
%	-c & d
%	\end{smallmatrix}\right),
%	\ee
%	where~$f_1,f_2$, resp.~$f_3,f_4$, is a standard basis of~$U_1$, resp.~$U_2$.
%	}\\
	The action of~$\SL_2(\ZZ)$ on~$\Mat_2(\ZZ)$ by multiplication from both sides induces isometries in~$\OO^+(U_1\oplus U_2)$.
	The theorem of elementary divisors for~$\SL_2(\ZZ)$ implies that every element of~$L'$ maps to~$L_1'$ under an isometry in~${\OO^+(U_1\oplus U_2)}$, where~$L_1\coloneqq \{0\}\oplus U_2\oplus E_8\oplus E$.
%	\\\textcolor{blue}{
%	\textbf{Check:}
%	The theorem of elementary divisors for~$\SL_2(\ZZ)$ says that for every~$A\in \Mat_2(\ZZ)$, there are~$U,V\in \SL_2(\ZZ)$ and~$\alpha,\beta\in\NN$, such that
%	\[
%	UAV=\left(\begin{smallmatrix}
%	\alpha & 0\\
%	0 & \beta
%	\end{smallmatrix}\right)\qquad
%	\text{and}
%	\qquad
%	\alpha|\beta.
%	\]
%	Notes that~$\left(\begin{smallmatrix}
%	\alpha & 0\\
%	0 & \beta
%	\end{smallmatrix}\right)$ corresponds to~$(0,\alpha f_3 - \beta f_4)\in \{0\}\oplus U_2$ under the isomorphism~\eqref{eq;stisoblue}.
%	}
	
	Let
	\[
	C_\alpha\coloneqq \min\{
	q(\lambda) : \text{$\lambda\in \alpha + E$}
	\},
	\qquad
	\alpha\in D_{E}\cong D_{L}.
	\]
	Let~$C=\max_{\alpha} C_\alpha$, and let~$a,b\in(\lattice')^2$ be as in the statement of Lemma~\ref{lemma;orbitsofdisck}.
	Since~$q(a_1)>C$, there exists~$e\in\alpha+E$ such that~$q(e)<q(a_1)$.
	Since~$E_8$ represents primitively every positive integer, there exists a primitive~$\tilde e\in E_8$ such that~$q(\tilde e)=q(a_1)-q(e)$.
%	\\\textcolor{blue}{\textbf{Proof that~$E_8$ represents primitively every integer:}
%	We already know that~$E_8$ represents every positive integer.
%	In fact, the generating series~$\theta_{E_8}$ of the representation numbers of~$E_8$ is a modular form of weight~$4$ w.r.t.~$\SL_2(\ZZ)$, and is well-known that the space of such functions is~$1$-dimensional, generated by the Eisenstein series~$E^4_1$, hence all Fourier coefficients are non-zero.
%	\\
%	We now prove that~$E_8$ represents primitively every integer.
%	Let~$r(m)$, resp.~$r_\prim(m)$, be the number of lattice vectors, resp.\ primitive lattice vectors, of~$E_8$ whose norm is~$m$.
%	It is easy to see that
%	\[
%	r(m)=\sum_{n^2|m}r_\prim(m/n^2)
%	\qquad
%	\text{and}
%	\qquad
%	r_\prim(m)=\sum_{n^2|m}\mu(n)r(m/n^2).
%	\]
%	(The second formula above follows directly by substituting~$r(m/n^2)$ with the sum given by the first formula.)\\
%	We want to prove that~$r_\prim(m)>0$ for all~$m>0$.
%	Since~$\theta_{E_8}$ is a multiple of the (non-normalized) Eisenstein series
%	\[
%	\zeta(3)/2 + \sum_{m>0} \sigma_3(m)q^m,
%	\]
%	the positivity of~$r_\prim(m)$ is equivalent to
%	\[
%	\sum_{n^2|m}\mu(n) \sigma_3(m/n^2)>0.
%	\]
%	Since both~$\mu$ and~$\sigma_3$ are multiplicative, to prove the inequality above is enough to show that if~$p$ is a prime, then
%	\[
%	\sum_{n^2|p^\ell}\mu(n)\sigma_3(p^\ell/n^2)>0\qquad\text{for all~$\ell$}.
%	\]
%	This is trivial if~$\ell=1$, and if~$\ell>1$, then it follows from the fact that
%	\[
%	\sum_{n^2|p^\ell}\mu(n)\sigma_3(p^\ell/n^2)
%	=
%	\sigma_3(p^\ell) - \sigma_3(p^{\ell-2}).
%	\]}
	This implies that~$\tilde e + e$ is primitive in~$E_8\oplus E'$ (hence in~$L'$), and has the same norm and same~$D_L$-coset as~$a_1$.
	By~\cite[Lemma~$4.4$]{freitaghermann} there exists an isometry in~$\disk{}^+(L)$ mapping~$a_1$ to~$\tilde e + e$.
	The same argument works for~$b_1$ in place of~$a_1$.
	We may therefore assume that~$a_1=b_1\in E_8\oplus E'$.
%		----------------------
%		
%	We now show that we may assume that~$a_1=b_1\in E_8\oplus E'$.
%	By the elementary divisor theorem for~$\SL_2(\ZZ)$, we may assume that~$a_1\in L_1'$.
%	Let
%	\[
%	a_1=(a_1)_{U_2} + \underbrace{(a_1)_{E_8} + (a_1)_E}_{(a_1)_{E_8\oplus E}}
%	\]
%	be the decomposition of~$a_1$ w.r.t.~$L_1'=U_2\oplus E_8\oplus E'$.
%	Let~$\pi\colon L_1'\to D_{L_1}$ be the projection.
%	Clearly~$\pi(a_1)=\pi \big((a_1)_{E}\big)$.
%	
%	Suppose that~$q\big((a_1)_{U_2}\big)\ge 0$.
%	Since~$E_8$ represents every integer primitively, there exists a primitive~$\tilde a_1\in E_8$ such that
%	\[
%	q(\tilde a_1)= q\big((a_1)_{U_2}\big) + q\big((a_1)_{E_8}\big).
%	\]
%	We deduce that~$\tilde a_1 + (a_1)_{E}$ is primitive, of the same norm as~$a_1$, with the same~$D_{L_1}$-coset, which implies by~\cite[Lemma~$4.4$]{freitaghermann} that there exists an isometry in~$\Gamma$ mapping~$a_1$ to~$\tilde a_1 + (a_1)_{E}\in E_8\oplus E$.
%	
%	Suppose now that~$q\big((a_1)_{U_2}\big)< 0$.
%	Since
%	\[
%	q(a_1)-q((a_1)_{U_1})> q(a_1) > C,
%	\]
%	there exists
%	
%	Since~$E_8$ represents every integer primitively, by~\cite[Lemma~$4.4$]{freitaghermann}
%	\\
%	\\\textcolor{blue}{If~$q\big((a_1)_{U_2}\big)\ge 0$, then it is enough to show that there exists a primitive~$\tilde a_1\in E_8$ such that
%	\[
%	q(\tilde a_1)= q\big((a_1)_{U_2}\big) + q\big((a_1)_{E_8}\big).
%	\]
%	In fact, this implies that~$\tilde a_1 + (a_1)_{E}$ is primitive, of the same norm as~$a_1$, with the same~$D_{L_1}$-coset, which implies by~\cite[Lemma~$4.4$]{freitaghermann} that there exists an isometry~$\tilde\gamma\in \Gamma$ such that~$\tilde\gamma\colon a_1\mapsto\tilde a_1 + (a_1)_{E}$.
%	\\
%	We now show that there exists an isometry~$\gamma_1\in\Gamma$ mapping~$a_1$ to a~$c_1\in L_1'$ with~$q\big((c_1)_{U_2}\big)>0$.
%	W.l.o.g.\ we may assume that~$q\big((a_1)_{U_2}\big)<0$ and hence~$(a_1)_{E_8\oplus E'}$ is non-zero.
%	Let~$z$ be the primitive element of~$E_8\oplus E'$ such that~$(a_1)_{E_8\oplus E'}=mz$ for some~$m\in\ZZ_{>0}$.
%	There exists~$a'\in E_8\oplus E$ such that
%	\[
%	(z,a')=1.
%	\]
%	Furthermore, up to adding to~$a'$ some lattice vectors in~$z^\perp\subset E_8\oplus E$, we may assume that~$q(a')>2$.
%	\\
%	\textcolor{red}{Is it~$z^\perp\not=\emptyset$ for all~$z\in E_8\oplus E'$? Probably yes: Take dual basis...}\\
%	Write~$(a_1)_{U_2} = A f_3 + B f_4$ with~$A,B\in\ZZ$.
%	Since~$q\big((a_1)_{U_2}\big) = AB$ is negative, we may assume that~$A<0$ and~$B>0$.
%	Consider now the Eichler transformation~$E(-f_3,m a')$.
%	It has the property that
%	\bas
%	E(-f_3,ma')(a_1)
%	&=
%	a_1 
%	+
%	m(a_1,f_3)a'
%	-
%	\underbrace{(a_1,ma')}_{=m^2}f_3
%	-
%	q(ma')(a_1,f_3) f_3
%	\\
%	&=
%	a_1 + mBa' \underbrace{- m^2\big(1+q(a')B\big)}_{(\star)}f_3.
%	\\
%	E(f_3,ma')(a_1)
%	&=
%	a_1 - mBa' + \underbrace{m^2\big(1-q(a')B\big)}_{(\star)}f_3.
%	\eas
%	Since~$q(a')\in \ZZ_{>0}$ and~$B>0$, we deduce that~$(\star)$ is negative...
%%%%%%%%%%%%%%%%%% another approach:
%	\textbf{Check:}
%	Let
%	\[
%	a_1=(a_1)_{U_1\oplus U_2} + (a_1)_{E_8} + (a_1)_E
%	\]
%	be the decomposition of~$a_1$ w.r.t.~$L'=U_1\oplus U_2\oplus E_8\oplus E'$.
%	Let~$\pi\colon L'\to D_{L}$ be the projection.
%	Clearly~$\pi(a_1)=\pi \big((a_1)_E\big)$.
%	It is enough to show that there exists a primitive~$\tilde a_1\in E_8$ such that
%	\[
%	q(\tilde a_1)= q\big((a_1)_{U_1\oplus U_2}\big) + q\big((a_1)_{E_8}\big).
%	\]
%	In fact, this implies that~$\tilde a_1 + (a_1)_{E}$ is primitive, of the same norm as~$a_1$, with the same~$D_{L}$-coset, which implies that there exists an isometry~$\tilde\gamma\in \Gamma$ such that~$\tilde\gamma\colon a_1\mapsto\tilde a_1 + (a_1)_{E}$.
%	The same works for~$b_1$.
%	\\\textcolor{red}{Problem in this check:~$E_8$ represents only positive integers... it might be that~$q(\tilde a_1)<0$...}
%	}
	
	Since~$E_8\oplus E'$ is preserved under the action of~$\OO^+(U_1\oplus U_2)$, by the theorem of elementary divisors for~$\SL_2(\ZZ)$ we may further assume that~$a_2,b_2\in L_1'$.
	To conclude the proof it is then enough to show that there exists~$\gamma\in\Stab_\Gamma(a_1)$ such that~$\gamma(a_2)=b_2$.
	We construct~$\gamma$ in terms of Eichler transformations.
	Recall that if~$u,v\in L$ are such that~$v\perp u$ and~$u$ is isotropic, then the Eichler transformation~$E(u,v)$ defined as
	\[
	E(u,v)(x)= x - (x,u) v + (x, v) u - q(v) (x,u) u,\qquad x\in L\otimes\QQ,
	\]
	lies in the discriminant kernel of~$L$.
	
	Let~$f_1,f_2$, resp.~$f_3,f_4$, be the standard generators of~$U_1$, resp.~$U_2$.
	Since~$a_1$ and~$a_2$ are primitive in~$L_1'$ and linearly independent, there exists a basis~$\mathcal{B}$ of~$L_1'$ whose first two vectors are~$a_1$ and~$a_2$.
	Let~$a'\in L_1$ be the negative of the second vector of the basis of~$L_1$ dual to~$\mathcal{B}$.
	Then
	\[
	(a_1,a')=0
	\qquad
	\text{and}
	\qquad
	(a_2,a')=-1
	.
	\]
	Similarly, one can find~$b'\in L_1$ such that~$(b_1,b')=0$ and~$(b_2,b')=-1$.
	It is easy to check that
	\[
	E(f_1,a')(a_i)=
	\begin{cases}
	a_i, & \text{if~$i=1$,}\\
	a_i-f_1, & \text{if~$i=2$,}
	\end{cases}
	\qquad\text{and}
	\qquad
	E(f_1,b')(b_i)=
	\begin{cases}
	b_i, & \text{if~$i=1$,}\\
	b_i-f_1, & \text{if~$i=2$.}
	\end{cases}
	\]
%	\textcolor{blue}{\textbf{Check of the previous assertion regarding~$E(f_1,a')(a_i)$:}
%	\bas
%	E(f_1,a')(a_1)
%	&=
%	a_1
%	- \underbrace{(a_1,f_1)}_{=0} a'
%	+ \underbrace{(a_1,a')}_{=0}f_1
%	- q(a')\underbrace{(a_1,f_1)}_{=0}f_1
%	=
%	a_1,
%	\\
%	E(f_1,a')(a_2)
%	&=
%	a_2
%	- \underbrace{(a_2,f_1)}_{=0} a'
%	+ \underbrace{(a_2,a')}_{-1}f_1
%	- q(a')\underbrace{(a_2,f_1)}_{=0}f_1
%	=
%	a_2 - f_1.
%	\eas}
	We may further compute
	\bas
	E(f_2,b_2-a_2)(a_2-f_1)&=b_2-f_1,
	\\
	E(f_2,b_2-a_2)(a_1)&=a_1 + (a_1,b_2-a_2)f_2 = a_1,
	\eas
	where in the last equality we used that since~$q(a)=q(b)$, then
	\[
	(a_1,b_2-a_2)
	=
	(a_1,b_2)-(a_1,a_2)
	=
	(b_1,b_2)-(a_1,a_2)
	=0.
	\qedhere
	\]
	\end{proof}
	By~\cite[Lemma~$4.3$]{bruiniermöller}, if~$L$ splits off two hyperbolic planes, then the induced primitive special divisors on~$\disk{}^+(L)\backslash\domain$ are irreducible.
	The following result is a generalization of this in codimension~$2$.
	\begin{prop}\label{prop;onirredofprimcy}
	If~$X=\disk{}^+(\lattice)\backslash\domain$ and~$L$ splits off~$U^{\oplus 2}\oplus E_8$, then the primitive special cycles of codimension~$2$ arising from~$L$ are (not necessarily reduced) irreducible cycles.
	\end{prop}
	\begin{proof}
	Let~$\discel\in D_L^{2}$, and let~$T\in q(\alpha)+\halfint{2}$ be positive definite.
	We may choose~$M\in\GL_2(\ZZ)$ such that the top-left entry of~$M^t T M$ is greater than~$C$, where~$C$ is as in Lemma~\ref{lemma;orbitsofdisck}.
%	\\\textcolor{blue}{
%	\textbf{Check that such an~$M$ exists:}
%	Write~$T=\big(\begin{smallmatrix}
%	t_1 & t_{1,2}
%	\\
%	t_{1,2} & t_2
%	\end{smallmatrix}\big)$.
%	Since~$T$ is positive definite, then~$t_1,t_2>0$.
%	If~$M=\big(\begin{smallmatrix}
%	a & b
%	\\
%	c & d
%	\end{smallmatrix}\big)\in\GL_2(\ZZ)$, then
%	\[
%	M^t T M
%	=
%	\big(\begin{smallmatrix}
%	a & c
%	\\
%	b & d
%	\end{smallmatrix}\big)
%	\cdot
%	\big(\begin{smallmatrix}
%	t_1 & t_{1,2}
%	\\
%	t_{1,2} & t_2
%	\end{smallmatrix}\big)
%	\cdot M
%	=
%	\big(
%	\begin{smallmatrix}
%	a t_1 + c t_{1,2} & a t_{1,2} + c t_2
%	\\
%	\ast & \ast
%	\end{smallmatrix}
%	\big)
%	\cdot
%	\big(\begin{smallmatrix}
%	a & b
%	\\
%	c & d
%	\end{smallmatrix}\big)
%	=
%	\big(
%	\begin{smallmatrix}
%	a^2 t_1 + 2ac t_{1,2} + c^2 t_2 & \ast
%	\\
%	\ast & \ast
%	\end{smallmatrix}\big).
%	\]
%	Choose for instance~$M$ with~$a=b=1$, $c$ arbitrarily large, and~$d=1+c$.
%	Then the top-left entry of~$M^t T M$ is
%	\[
%	t_1 + 2c t_{1,2} + c^2 t_2.
%	\]
%	}\\
	By the latter lemma, the irreducible components of
	\[
	\sum_{\substack{\lambda\in\primvec{\lattice}{2}{\discel M}\\ q(\lambda)=M^tTM}}
	\lambda^\perp
	\]
	on~$\domain$ are identified under~$\disk{}^+(L)$.
	This means that the induced cycle~$\Zprim(M^tTM,\discel M)$ on~$X$ is irreducible.
	Since~$\Zprim(M^tTM,\discel M)=\Zprim(T,\alpha)$ by~\eqref{eq;symmetryspcy} and Lemma~\ref{lemma;spcyclesandprim}, also~$\Zprim(T,\alpha)$ is irreducible.
%	By Lemma~\ref{lemma;orbitsofdisck} the irreducible components of
%	\[
%	\sum_{\substack{\lambda\in\primvec{\lattice}{\genus+1}{\discel}\\ q(\lambda)=T}}
%	\lambda^\perp
%	\]
%	on~$\domain$, with~$T$ positive definite, are identified under~$\Gamma$ for every~$\discel\in D_L^{\genus+1}$.
%	This means that the induced primitive special cycle on~$X=\Gamma\backslash\domain$ is irreducible.
	\end{proof}
	
%	\begin{proof}[Old incomplete proof using split of~$U\oplus U$ of sublattice of~$L$]
%	Let~$E$ be an even sublattice of~$L$ such that~$L$ splits in orthogonal direct sums as
%	\[
%	L= E \oplus E_8^{\oplus\genus}\oplus U^{\oplus 2}.
%	\]
%	Let~$\discel\in D_\lattice^{\genus+1}$, and let~$\lambda,\mu\in \discel+\lattice^{\genus+1}$ be with primitive entries in~$L'$ and such that~$q(\lambda)=q(\mu)\eqqcolon T$.
%	We want to show that there exists an isometry~$\gamma\in\Gamma$ such that~$\gamma(\mu)=\lambda$.
%	In fact, this will imply that the irreducible components of
%	\[
%	\sum_{\substack{\lambda\in\primvec{\lattice}{\genus+1}{\discel}\\ q(\lambda)=T}}
%	\lambda^\perp
%	\]
%	on~$\domain$ are identified under~$\Gamma$.
%	Without loss of generality, we may assume that~$T$ is positive definite.
%	
%	If~$\beta$ is a tuple, we denote by~$\beta_j$ its~$j$th entry.
%	Since~$L$ splits off~$U^{\oplus 2}$, by~\cite[Lemma~$4.4$]{freitaghermann} all primitive vectors in~$L'$ with the same norm and same~$D_\lattice$-coset lie in the same~$\Gamma$-orbit.
%	Hence, since~$\lambda_1,\mu_1\in\alpha_1+\lattice$ are primitive with~$q(\lambda_1)=q(\mu_1)$, there exists~$\gamma_1\in\Gamma$ such that~$\gamma_1(\mu_1)=\lambda_1$.
%	We now show that there exist~$\gamma_j\in \Stab_{\Gamma}\big(\langle\lambda_1,\dots,\lambda_{j-1}\rangle\big)$ such that~$\gamma_j\circ\cdots\circ \gamma_1(\mu_j)=\lambda_j$ for every~$j=2,\dots,\genus+1$, so that we may choose~$\gamma=\gamma_{\genus+1}\circ\cdots\circ\gamma_1$ since then
%	\[
%	\gamma(\mu)
%	=
%	\big(
%	\gamma(\mu_1),
%	\dots,
%	\gamma(\mu_{\genus+1})
%	\big)
%	=
%	\big(
%	\gamma_1(\mu_1),
%	\gamma_2\circ\gamma_1(\mu_2),
%	\dots,
%	\gamma(\lambda_{\genus+1})
%	\big)
%	=
%	\lambda.
%	\]
%	
%	Let~$L_j$ be the orthogonal complement of~$\langle\lambda_1,\dots,\lambda_{j-1}\rangle$ in~$L$.
%	Since~$\disk{}^+(L_j)$ may be regarded as a subgroup of~$\Stab_{\Gamma}(\langle\lambda_1,\dots,\lambda_{j-1}\rangle)$,
%	\\\textcolor{blue}{
%	\textbf{Check of the last assertion:}
%	We may consider~$\OO\big(L_j\big)$ as a subgroup of~$\OO(M)$, where~$M\coloneqq L_j\oplus \langle\lambda_1,\dots,\lambda_{j-1}\rangle_\ZZ$, simply by extending every isometry of~$\OO(L_j)$ to the identity on $\langle\lambda_1,\dots,\lambda_{j-1}\rangle_\ZZ$.
%	This induces an inclusion~$\iota\colon\disk(L_j)\hookrightarrow\disk(M)$, because since every~$\xi\in\disk(L_j)$ is extended under~$\iota$ to the identity on~$\langle\lambda_1,\dots,\lambda_{j-1}\rangle_\ZZ$, the map~$\iota(\xi)$ acts on~$M'/M\cong D_{L_j}\oplus D_{\langle\lambda_1,\dots,\lambda_{j-1}\rangle_\ZZ}$ as~$\xi|_{D_{L_j}}\oplus\id|_{D_{\langle\lambda_1,\dots,\lambda_{j-1}\rangle_\ZZ}}$.
%	\\
%	The fact that~$\disk{}^+(M)$ may be regarded as a subgroup of~$\Stab_{\Gamma}\big(\langle\lambda_1,\dots,\lambda_{j-1}\rangle_\ZZ\big)$ follows from the usual argument using the inclusions
%	\[
%	M\subset L \subset L'\subset M'
%	\qquad\text{and}\qquad
%	L/M \subset L'/M \subset M'/M.
%	\]}\noindent
%	by~\cite[Lemma~$4.4$]{freitaghermann} it is enough to show that~$L_j$ splits off~$U^{\oplus 2}$ for all~$j$.
%	\\
%	\textcolor{blue}{
%	\textbf{Check of the previous assertion:}
%	Let's work for simplicity with the case of~$j=2$.
%	We have that~$L_2$ splits off~$U^{\oplus 2}$,~$\lambda_2$ and~$\gamma_{1}(\mu_2)$ are primitive in~$L'$, have the same norm, and lie in the same~$D_L$-coset.
%	We want to show that there exists an isometry~$\gamma_2\in\disk{^+}(L_2)\subset\Stab_\Gamma\big(\langle\lambda_1\rangle\big)$ that maps~$\gamma_{1}(\mu_2)$ to~$\lambda_2$.
%	\\
%	\underline{\emph{Suppose that~$\lambda_2$ and~$\gamma_{1}(\mu_2)$ are orthogonal to~$\lambda_1$.}}
%	Note that~$\lambda_2,\gamma_{1}(\mu_2)\in L_2'$.
%	Since~$\lambda_2-\gamma_{1}(\mu_2)$ lies in~$L$ and is orthogonal to~$\lambda_1$, it lies in~$L_2$.
%	Therefore, saying that~$\lambda_2$ and~$\gamma_{1}(\mu_2)$ lie in the same~$D_L$-coset implies that they lie also in the same~$D_{L_2}$-coset.
%	\\
%	Summarizing: $\lambda_2$ and~$\gamma_{1}(\mu_2)$ are in~$L_2'$, have the same norm, are in the same~$D_{L_2}$-coset.
%	If they are also primitive in~$L_2'$ or they are the same multiple of some primitive vectors in~$L_2'$, then there is an isometry~$\gamma_2\in\disk{^+}(L_2)$ mapping one to the other.
%	\\
%	\textcolor{red}{\textbf{Question for Jan:} If~$\lambda_2$ and~$\gamma_{1}(\mu_2)$ are primitive in~$L'$, are they primitive also in~$L_2'$? In an email I sent to you on 10.04.24, we agreed that they could be non-primitive in~$L_2'$, but that \textbf{they are the same multiple of some primitive vectors} in~$L_2$.
%	\\
%	\textbf{To do:} Prove that assertion!}
%	\\
%	\underline{\emph{Suppose now that~$\lambda_2$ and~$\gamma_{1}(\mu_2)$ are NOT orthogonal to~$\lambda_1$.}}
%	Then the bilinear form~$(\cdot{,}\cdot)$ of~$L$ is such that
%	\[
%	(\lambda_2,\lambda_1)=\big(\gamma_1(\mu_2),\lambda_1\big).
%	\]
%	The orthogonal projection of~$\lambda_2$ and~$\gamma_1(\mu_2)$ to~$\QQ\lambda_1$ are the same, since such projections are respectively
%	\be\label{erthpr;eq}
%	\frac{(\lambda_2,\lambda_1)}{(\lambda_1,\lambda_1)}\lambda_1
%	\qquad
%	\text{and}
%	\qquad
%	\frac{\big(\gamma_1(\mu_2),\lambda_1\big)}{(\lambda_1,\lambda_1)}\lambda_1.
%	\ee
%	Let~$\pi\colon L'\to L_2\otimes\QQ$ be the orthogonal projection.
%	It is such that
%	\[
%	\pi(\lambda_2) =
%	\lambda_2 - \frac{(\lambda_2,\lambda_1)}{(\lambda_1,\lambda_1)}\lambda_1
%	\qquad
%	\text{and}
%	\qquad
%	\pi\big(\gamma_1(\mu_2)\big) = 
%	\gamma_1(\mu_2) - \frac{\big(\gamma_1(\mu_2),\lambda_1\big)}{(\lambda_1,\lambda_1)}\lambda_1.
%	\]
%	\textcolor{red}{\textbf{To do:} Conclude the blue check by reducing this case to the case of~$\lambda_2$ and~$\gamma_1(\mu_2)$ orthogonal to~$\lambda_1$.
%	The idea is to show that it is enough to find some isometry in~$\disk{}^+(L_2)$ mapping~$\pi(\lambda_2)$ to~$\pi\big(\gamma_1(\mu_2)\big)$, since due to~\eqref{erthpr;eq} it will also map~$\lambda_2$ to~$\gamma_1(\mu_2)$.}
%	}
%	
%	We are going to show that actually~$L_j$ splits off~$E_8^{\oplus\genus + 1 - j}\oplus U^{\oplus 2}$.
%	It is enough to prove that there exists an isometry~$\xi$ of~$L$ such that~$\xi(\lambda_1)$ lies in a piece~$E'\oplus E_8^{\genus+1-j}$ of the orthogonal decomposition
%	\[
%	L'=E'\oplus E_8^{\oplus\genus}\oplus U^{\oplus 2}.
%	\]
%	\textcolor{blue}{
%	\textbf{Check that if such a~$\xi$ exists, then~$L_2$ splits~$U^{\oplus 2}$:}
%	Let~$\xi(\lambda_1)^\perp$ be the orthogonal complement of~$\xi(\lambda_1)$ in~$L$.
%	Since~$\xi(\lambda_1)\in E'\oplus E_8^{\genus}$, then~$\xi(\lambda_1)^\perp$ splits off~$U^{\oplus 2}$.
%	By applying~$\xi^{-1}$ we get that~$\xi^{-1}\big(\xi(\lambda_1)^\perp\big)=L_2$ splits off~$\xi^{-1}(U^{\oplus 2})\cong U^{\oplus 2}$.
%	}\\
%	By~\cite[Lemma~$4.4$]{freitaghermann} it is enough to show that there is~$\chi\in E'\oplus E_8$ primitive in the same~$D_L$-coset of~$\lambda_1$ such that~$q(\chi)=q(\lambda_1)$.
%	
%	The orthogonal projection~$L'\to E'$ induces an isomorphism~$\pi\colon D_\lattice\to D_E$.
%	Note that~$\lambda_1 + L = \pi(\lambda_1) + L \supseteq \pi(\lambda_1)+ E\oplus E_8$.
%	We want to find a primitive~$\lambda_1'\in\pi(\lambda_1) + E\oplus E_8$ with~$q(\lambda_1')=q(\lambda_1)$.
%	We look for a candidate of the form~$\lambda_1'=\pi(\lambda_1) + A$ for some~$A\in E_8$.
%	such that~$\pi(\lambda_1)+A$ is primitive.
%	Since~$\pi(\lambda_1)\perp A$, the primitivity of~$\lambda_1'$ is assured by the primitivity of~$A$.
%	Furthermore, since we want~$q(\lambda_1')=q(\lambda_1)$, we seek for some~$A$ such that
%	\[
%	q\big(\pi(\lambda_1)\big) + q(A) = q(\lambda_1),
%	\qquad\text{or equivalently}\qquad
%	q(A)=q(\lambda_1)-q\big(\pi(\lambda_1)\big).
%	\]
%	Note that since~$E'\oplus E_8$ is positive definite and~$\pi(\lambda_1)$ lies in the same~$D_L$-coset of~$\lambda_1$, the value~$q(A)$ is a positive integer.
%	We now show that for every positive integer~$r$ there exists~$A\in E_8$ primitive with~$q(A)=r$.
%	\\
%	\textcolor{red}{\textbf{To do:} Check the previous assertion by using that the Eisenstein series is the generating series of representation numbers of~$E_8$.
%	By working with primitive representation numbers, it should turn out that~$E_8$ represents \emph{primitively} every positive integer.}
%	\end{proof}
%	\begin{proof}[Different idea to prove (and maybe generalize) Proposition~\ref{prop;onirredofprimcy}]
%	\textcolor{blue}{\textbf{This proof is an old incomplete version. It should be discarded as soon as the proof above is completed.}}\\
%	\textbf{Case of codimension~1:}
%	It is~\cite[Lemma~$4.3$]{bruiniermöller}.
%	
%	\textbf{Case of codimension~2:}
%	Without loss of generality, we may work with special cycles associated to positive definite matrices
%	
%	Let~$\discel=(\discel_1,\discel_2)\in D_\lattice^2$, and let~$\lambda=(\lambda_1,\lambda_2),\mu=(\mu_1,\mu_2)\in \discel+\lattice^2$ be primitive in~$L'$ and such that~$q(\lambda)=q(\mu)\eqqcolon T$.
%	We want to show that there exists an isometry~$\gamma\in\Gamma$ such that~$\gamma(\mu)=\lambda$.
%	In fact, this would imply that the irreducible components of
%	\[
%	\sum_{\substack{\lambda\in\primvec{\lattice}{2}{\discel}\\ q(\lambda)=T}}
%	\lambda^\perp
%	\]
%	on~$\domain$ are identified under~$\disk{}^+(\lattice)$.
%	Without loss of generality, we may assume that~$T$ is positive definite.
%	
%	By~\cite[Lemma~$4.4$]{freitaghermann} all primitive vectors in~$L'$ with the same norm and same~$D_\lattice$-coset lie in the same~$\Gamma$-orbit.
%	Hence, since~$\lambda_1,\mu_1\in\alpha_1+\lattice$ are primitive with~$q(\lambda_1)=q(\mu_1)$, there exists also~$\gamma_1\in\Gamma$ such that~$\gamma_1(\mu_1)=\lambda_1$.
%	We want to show that there exists~$\gamma_2\in\Stab_{\Gamma}\big(\lambda_1\big)$ such that~$\gamma_2\circ \gamma_1(\mu_2)=\lambda_2$, so that we may choose~$\gamma=\gamma_2\circ\gamma_1$ since then
%	\[
%	\gamma(\mu)
%	=
%	\big(
%	\gamma_2\circ\gamma_1(\mu_1),
%	\gamma_2\circ\gamma_1(\mu_2)
%	\big)
%	=
%	\big(
%	\gamma_2(\lambda_1),
%	\lambda_2
%	\big)
%	=
%	(\lambda_1,\lambda_2)
%	=
%	\lambda.
%	\]
%	Let~$L(\lambda_1)$ be the orthogonal complement of~$\lambda_1$ in~$L$.
%	Since~$\disk{}^+(L(\lambda_1))$ may be regarded as a subgroup of~$\Stab_{\Gamma}(\lambda_1)$,
%	\\\textcolor{blue}{
%	\textbf{Check of the last assertion:}
%	Certainly we may consider~$\OO\big(L(\lambda_1)\big)$ as a subgroup of~$\OO(M)$, where~$M\coloneqq L(\lambda_1)\oplus \ZZ\lambda_1$, simply by extending every isometry of~$\OO\big(L(\lambda_1)\big)$ to the identity on~$\ZZ\lambda_1$.
%	This induces an inclusion~$\iota\colon\disk\big(L(\lambda_1)\big)\hookrightarrow\disk(M)$, because since every~$\xi\in\disk\big(L(\lambda_1)\big)$ is extended under~$\iota$ to the identity on~$\ZZ\lambda$, the map~$\iota(\xi)$ acts on~$M'/M\cong D_L(\lambda_1)\oplus D_{\ZZ\lambda}$ as~$\xi|_{D_{L(\lambda_1)}}\oplus\id|_{D_{\ZZ\lambda}}$.
%	\\
%	The fact that~$\disk{}^+(M)$ may be regarded as a subgroup of~$\Stab_{\Gamma}(\lambda_1)$ follows from the usual argument using the inclusions
%	\[
%	M\subset L \subset L'\subset M'
%	\qquad\text{and}\qquad
%	L/M \subset L'/M \subset M'/M.
%	\]}
%	by~\cite[Lemma~$4.4$]{freitaghermann} it is enough to show that~$L(\lambda_1)$ splits off two hyperbolic planes.
%	To do so, we first relate~$\mingen(D_{L(\lambda_1)})$ with~$\mingen(D_L)$, where we denote by~$\mingen(H)$ the minimum number of generators of any group~$H$.
%	
%	Up to multiply~$\lambda_1$ by a suitable integer, we may assume that~$\lambda_1$ is a primitive element of~$L$.
%	The lattice~$L$ is an overlattice of~$L(\lambda_1)\oplus\ZZ\lambda_1$, hence
%	\[
%	L(\lambda_1)\oplus\ZZ\lambda_1 \subseteq L \subseteq L'\subseteq L(\lambda_1)'\oplus(\ZZ\lambda_1)'.
%	\]
%	These inclusions induce two maps
%	\[
%	\begin{tikzcd}
%	L'/L(\lambda_1)\oplus\ZZ\lambda_1 \arrow[d, "\psi_2"] \arrow[r, "\psi_1"] & L(\lambda_1)'\oplus(\ZZ\lambda_1)'/ L(\lambda_1)\oplus\ZZ\lambda_1 = D_{L(\lambda_1)}\oplus D_{\ZZ\lambda_1}\\
%	L'/L = D_L & 
%	\end{tikzcd}
%	\]
%	The map~$\psi_1$ is injective, while~$\psi_2$ is surjective with kernel
%	\[
%	\ker(\psi_2)\cong L/L(\lambda_1)\oplus\ZZ\lambda_1.
%	\]
%	
%	Recall that if~$H$ is a finite abelian group and~$I$ is a subgroup of~$H$, then
%	\be\label{eq;mingen}
%	\mingen(H)\leq\mingen(I) + \mingen(H/I).
%	\ee
%	\textcolor{blue}{\textbf{Check of~\eqref{eq;mingen}:}
%	Let~$i_1,\dots,i_{\mingen(I)}$ be a minimal set of generators for~$I$, and let~$[h_1],\dots,[h_{\ell(H/I)}]$ be a minimal set of generators of~$H/I$.
%	Here the~$h_j$ are representatives in~$H$ of the corresponding cosets~$[h_j]\in H/I$.
%	Let~$H'$ be the subgroup of~$H$ generated by~$i_1,\dots,i_{\mingen(I)},h_1,\dots,h_{\mingen(H/I)}$.
%	Clearly~$\mingen(H')\leq\mingen(I) + \mingen(H/I)$.
%	The inequality~\eqref{eq;mingen} follows from the fact that~$H=H'$.
%	To show this, assume that it is not true, hence there is~$\xi\in H$ that does not lie in~$H'$.
%	Up to subtracting a suitable combination of the elements~$h_1,\dots,h_{\mingen(H/I)}$, we may assume that the image of~$\xi$ under the quotient map~$H\to H/I$ is trivial.
%	But this means that~$\xi\in I\subseteq H$.
%	Hence~$H=H'$.
%	}\\
%	We apply~\eqref{eq;mingen} with~$H=D_{L(\lambda_1)}\oplus D_{\ZZ\lambda_1}$ and~$I=\Imm(\psi_1)$ to deduce that
%	\ba\label{eq;ineqmingen1}
%	\mingen(D_{L(\lambda_1)}\oplus D_{\ZZ\lambda_1})
%	&\leq
%	\mingen\big(
%	\Imm(\psi_1)
%	\big)
%	+
%	\mingen\big(
%	D_{L(\lambda_1)}\oplus D_{\ZZ\lambda_1} / \Imm(\psi_1)
%	\big)
%	\\
%	&
%	\leq\mingen\big(
%	L'/L(\lambda_1)\oplus\ZZ\lambda_1
%	\big)
%	+
%	\mingen\big(
%	L(\lambda_1)'\oplus(\ZZ\lambda_1)'/L'
%	\big)
%	,
%	\ea
%	where in the last passage we used that~$\psi_1$ is injective and
%	\[
%	D_{L(\lambda_1)}\oplus D_{\ZZ\lambda_1} / \Imm(\psi_1)
%	\cong
%	L(\lambda_1)'\oplus(\ZZ\lambda_1)'/L'.
%	\]
%	Since~$\ker(\psi_2)=L/L(\lambda_1)\oplus\ZZ\lambda_1$, and since the quotient of~$L'/L(\lambda_1)\oplus \ZZ\lambda_1$ by~$\ker(\psi_2)$ is isomorphic to~$D_L$, we may apply~\eqref{eq;mingen} with~$H=L'/L(\lambda_1)\oplus\ZZ\lambda_1$ and~$I=\ker(\psi_2)$ to deduce that
%	\[
%	\mingen(L'/L(\lambda_1)\oplus\ZZ\lambda_1)
%	\leq
%	\mingen\big(L/L(\lambda_1)\oplus\ZZ\lambda_1\big) + \mingen(D_L).
%	\]
%	This and~\eqref{eq;ineqmingen1} imply that
%	\be\label{ineq;mingen2}
%	\mingen(D_{L(\lambda_1)}\oplus D_{\ZZ\lambda_1})
%	\leq
%	\mingen(D_L)
%	+
%	\mingen\big(L/L(\lambda_1)\oplus\ZZ\lambda_1\big)
%	+
%	\mingen\big(
%	L(\lambda_1)'\oplus(\ZZ\lambda_1)'/L'
%	\big).
%	\ee
%	Since~$D_{\ZZ\lambda_1}$ is cyclic, we deduce that~$\mingen(D_{L(\lambda_1)}\oplus D_{\ZZ\lambda_1}) = \mingen(D_{L(\lambda_1)}) + 1$.
%	It is easy to see by the Theorem of Elementary Divisors that if~$M$ is a sublattice of finite index in~$L$, then~$L/M\cong M'/L'$.
%	\textcolor{blue}{
%	\textbf{Check of the last assertion:}
%	By the Theorem of Elementary Divisors there exist a basis~$e_1,\dots,e_{n+2}$ of~$L$ and positive integers~$a_i$ with the property that~$a_i|a_{i+1}$, such that
%	\[
%	f_1\coloneqq a_1 e_1,\quad\dots,\quad f_{n+2}\coloneqq a_{n+2}e_{n+2}
%	\]
%	is a basis of~$M$.
%	This implies that~$L/M\cong\oplus_{i=1}^{n+2}\ZZ/a_i\ZZ$.
%	Given the basis~$e_i$ of~$L$ above, the dual~$\ZZ$-basis of~$L'=\Hom_{\ZZ}(L,\ZZ)$ is given by the homomorphisms
%	\[
%	\varepsilon_i\colon L\longrightarrow\ZZ, \qquad \text{$\varepsilon_i(e_j)\coloneqq\delta_{i,j}$ for all~$j$}.
%	\]
%	We denote by~$\varepsilon_i'$ the dual basis of~$M'$ constructed from~$f_1,\dots,f_{n+2}$ in the same way.
%	Since
%	\[
%	\varepsilon_i(f_j)
%	=
%	\varepsilon_i(a_j e_j)
%	=
%	a_j \varepsilon_i(e_j)
%	=
%	a_j \delta_{i,j}
%	=
%	a_j\varepsilon_i'(f_j),
%	\]
%	we deduce that~$\varepsilon_i=a_i\varepsilon_i'$, hence~$M'/L'\cong \oplus_{i=1}^{n+2}\ZZ/a_i\ZZ\cong L/M$.
%	}\\
%	This implies that~$\mingen\big(L/L(\lambda_1)\oplus\ZZ\lambda_1\big)
%	=
%	\mingen\big(
%	L(\lambda_1)'\oplus(\ZZ\lambda_1)'/L'
%	\big)$.
%	We may then refine~\eqref{ineq;mingen2} to
%	\be\label{ineq;mingen3}
%	\mingen(D_{L(\lambda_1)})
%	\leq
%	\mingen(D_L)
%	+
%	2\mingen\big(L/L(\lambda_1)\oplus\ZZ\lambda_1\big)
%	-
%	1.
%	\ee
%	
%	\textcolor{red}{\textbf{Question:} Is it true that~$\mingen(L/L(\lambda_1)\oplus\ZZ\lambda_1)=1$?
%	\\
%	\textcolor{black}{Probably yes. Consider the projection~$L\to L(\lambda_1)$.
%	This should induce a map~$L'\to L(\lambda_1)'$, which should induce a map~$L'/L(\lambda_1)\to D_{L(\lambda_1)}$.}
%	\\
%	Here an implementation of such an idea:\\
%	}
%	Let~$M$ be a sublattice of~$L$.
%	The following results are easy to prove:
%	\begin{enumerate}
%	\item \label{enum;3lat}
%	The orthogonal complement~$M^\perp\subseteq L$ is a primitive sublattice.
%	\item \label{enum;1lat}
%	If~$M$ is primitive, then every basis of~$M$ can be extended to a basis of~$L$.
%	\item \label{enum;2lat}
%	If~$M$ is primitive, then the restriction map~$L'\to M'$ is surjective.
%	\item \label{enum;4lat}
%	\end{enumerate}
%	\textcolor{blue}{
%	\textbf{Check of the previous list of facts.}
%	\\
%	\eqref{enum;3lat}:
%	\\
%	Suppose that~$M^\perp$ is not primitive in~$L$, i.e.~$L/M^\perp$ is not free.
%	This means that there exist~$\lambda\in L$ and~$N\in\ZZ_{>0}$ such that~$\lambda\not\in M^\perp$ and~$N\lambda\in M^\perp$.
%	Since
%	\[
%	N\lambda\in M^\perp
%	\iff
%	(N\lambda,\mu)=0\quad\text{for all~$\mu\in M$}
%	\iff
%	\lambda\in M^\perp,
%	\]
%	it is not possible that~$\lambda\not\in M^\perp$.
%	\\
%	\eqref{enum;1lat}:
%	\\
%	The fact that there is a basis of~$M$ that does not extend to a basis of~$L$ implies that there exist~$\lambda\in L$ and~$N\in\ZZ_{>0}$ such that~$\lambda\not\in M$ and~$N\lambda\in M$.
%	Hence the non-zero class~$\lambda+M$ in~$L/M$ is torsion, hence~$L/M$ is not free.
%	This is not possible, because~$M$ is primitive.
%	\\
%	\eqref{enum;2lat}:
%	\\
%	Every~$\varepsilon\in M'=\Hom_\ZZ(M,\ZZ)$ may be written as a combination~$\varepsilon=\sum_{i=1}^{\rk (M)}a_i\varepsilon_i$, where~$a_i\in\ZZ$ and~$(\varepsilon_i)_i$ is a basis of~$M'$ dual to some basis~$(m_i)_i$ of~$M$.
%	Since~$M$ is primitive, by~\eqref{enum;1lat} the basis~$(m_i)_i$ extends to a basis
%	\[
%	m_1,\dots,m_{\rk (M)},\ell_{\rk( M )+ 1},\dots,\ell_{\rk (L)}
%	\]
%	of~$L$.
%	Let~$(\varepsilon_i')_{i=1}^{\rk(L)}$ be the associated dual basis of~$L'$.
%	Every functional~$\varepsilon'\in L'$ may be written as
%	\[
%	\varepsilon'=\sum_{i=1}^{\rk(M)}
%	a_i\varepsilon_i'
%	+
%	\sum_{\rk(M)+1}^{\rk(L)}
%	b_j\varepsilon_j'
%	\]
%	for some~$a_i,b_j\in\ZZ$.
%	Since~$\varepsilon_i'=\varepsilon_i$ for all~$i=1,\dots,N$, the restriction map~$L'\to M'$ is then surjective.
%	}
%	
%	By~\eqref{enum;3lat} the sublattice~$L(\lambda_1)$ is primitive.
%	By~\eqref{enum;2lat} the restriction map~$L'\to L(\lambda_1)'$ is surjective.
%	\end{proof}
	
	\begin{cor}\label{cor;blmminnonadcase}
	Let~$X=\disk{}^+(\lattice)\backslash\domain$ be of dimension~$n>2$.
	\begin{enumerate}
	\item If~$L$ splits off~$U^{\oplus 2}$, then~$H^{1,1}(X)$ is spanned by the classes of special divisors
	\[
	\bigdrcoho{Z(m,\alpha)}\wedge\bigdrcoho{\omega^{1-\delta_{0,m}}},
	\qquad
	\text{$\alpha\in D_\lattice$, $m\in q(\alpha) + \ZZ$, $m\ge0$.}
	\]
	\item If~$L$ splits off~$U^{\oplus 2}\oplus E_8$, then~$H^{2,2}(X)$ is spanned by the classes of special cycles\label{item;spcy2cod}
	\bes%\label{eq;gencorspanHgg}
	\bigdrcoho{Z(T,\alpha)}\wedge\bigdrcoho{\omega^{2-\rk T}},
	\qquad
	\text{$\alpha\in D_\lattice^2$, $T\in q(\alpha) + \halfint{2}$, $T\ge0$.}
	\ees
	\end{enumerate}
	\end{cor}
	\begin{proof}
	We give the details for the case of codimension~$2$, since a similar proof works also in codimension~$1$.
	By Proposition~\ref{prop;onirredofprimcy} and~\cite[Lemma~$4.3$]{bruiniermöller} the primitive special cycles of codimension at most~$2$ are irreducible.
	By Lemma~\ref{lemma;spcyclesandprim} the cohomology classes
	\[
	\bigdrcoho{\Zprim(T,\alpha)}\wedge\bigdrcoho{\omega^{2-\rk T}}
	\qquad
	\text{with $\alpha\in D_\lattice^2$ and $T\in q(\alpha) + \halfint{2}$, $T\ge 0$,}
	\]
	span the same subspace of~$H^{4}(X)$ spanned by the classes of special cycles with the same indices.
	To prove Corollary~\ref{cor;blmminnonadcase}~\eqref{item;spcy2cod}, it is enough to prove that the irreducible components of the adelic special cycles of codimension~$2$ are primitive special cycles.
	
	Let~$\varphi\in\mathcal{S}\big(V^2(\Ab_f)\big){}^K$.
	Then there exists an even lattice~$M\subset V$ such that~$\varphi\in\mathcal{S}_{M^2}$; see e.g.~\cite[Fifth line of the proof of Theorem~$3.2$]{kudla-speccyc}.
	The irreducible components of the adelic special cycle~$Z(T,\varphi,K)$ are then of the form
	\be\label{eq;adspircomp}
	\Gamma\big\backslash\sum_{x\in\Gamma \tilde{x}} x^\perp
	\ee
	for some~$\tilde{x}\in V^2(\QQ)\cap\supp(\varphi)$ with~$q(\tilde{x})=T$.
	Let~$m\in\ZZ_{>0}$ be such that~$\tilde\lambda\coloneqq m\tilde{x}\in(\lattice')^2$.
	Since~$\Gamma$ acts trivially on~$D_L^2$, there is some~$\alpha\in D_\lattice^2$ such that~$\Gamma\tilde\lambda\subseteq\alpha + \lattice^2$.
	Without loss of generality, we may assume that~$\tilde\lambda$ has primitive entries in~$L'$, hence the subvariety~\eqref{eq;adspircomp} is (up to multiplicity) a primitive special cycle of coset~$\alpha$.
%	-------------------------------
%	
%	Without loss of generality we may assume that~$M\subseteq\lattice$,
%	\\\textcolor{blue}{
%	\textbf{Check of the last assertion:}
%	Suppose that~$M\not\subseteq L$, and consider the even lattice~$\tilde{M}\coloneqq M\cap L$.
%	We want to show that if~$\varphi\in\mathcal{S}_M$, then~$\varphi\in\mathcal{S}_{\tilde{M}}$.
%	We may assume w.l.o.g.\ that~$\varphi=\cha(\discel + \widehat{M})$ for some~$\discel\in D_M$.
%	Since
%	\[
%	\tilde M \subseteq M \subseteq M' \subseteq \tilde M',
%	\]
%	and since~$\rk M=\rk \tilde M$, we deduce that the quotient~$M/\tilde M$ is a finite module.
%	We may then rewrite~$M$, resp~$\widehat{M}$, as the finite union
%	\[
%	M=
%	\bigcup_{\ell\in M/\tilde M}
%	\ell + \tilde{M},\qquad
%	\text{resp.}
%	\qquad
%	\widehat{M}=
%	\bigcup_{\ell\in M/\tilde M}
%	\ell + \widehat{\tilde{M}}.
%	\]
%	These imply that
%	\[
%	\cha(\alpha + \widehat{M})
%	=
%	\sum_{\ell\in M/\tilde{M}}
%	\cha\big((\alpha + \ell) + \widehat{\tilde M}\big).
%	\]
%	}\\
%	and in fact also that~$M=mL$ for some~$m\in\ZZ_{>0}$.
%	We may then write
%	\[
%	\varphi=\sum_{\alpha\in D_M} A_\alpha\cdot\cha(\alpha + \widehat{M})
%	\]
%	for some~$A_\alpha\in\CC$.
%	
%	\textcolor{red}{To be added.}
	\end{proof}
	
	
	
	
	
	
	
	
	
	
	
	
	
	
	
	
	
	\end{document}
\documentclass[../Lefschetz_decomposition.tex]{subfiles}

	\begin{document}
	In this section we prove that~$\KMtheta{\genus}$ is square-integrable, and recall some integral formulas to compute volumes of special cycles.
	The latter are well-known for compact orthogonal Shimura varieties~\cite{kudlamillson-harmonic1} \cite{kudlamillson-harmonic2}, and have been extended to the non-compact case in~\cite{kudlamillson-tubes} and~\cite{kudlamillson-intersection}.
	\\
	
	We denote by~$X$ a smooth orthogonal Shimura variety arising from an indefinite even lattice~$L$ of signature~$(n,2)$.
%	We assume that the arithmetic group~$\Gamma$ giving rise to~$X$ is torsion-free, so that~$X$ is smooth.
%	The Kudla--Millson theta function~$\KMtheta{\gengenus}$ arising from~$L$ is a vector-valued non-holomorphic Siegel modular form of weight~$k=1+n/2$ and genus~$\gengenus$ with respect to the Weil representation~$\weil{\lattice}{\gengenus}$, as well as a closed differential~$(\gengenus,\gengenus)$-form on~$X$; see Section~\ref{sec;KMthetabasics} for details.
%	For every \emph{compactly supported}~$2(n-\gengenus)$-form~$\eta$ on~$X$, we define
%	\[
%	\theta_\eta(\tau)\coloneqq\int_X\eta\wedge\KMtheta{\gengenus}(\tau),
%	\qquad
%	\tau\in\HH_\gengenus.
%	\]
%	
%	By~\cite{kudlamillson-tubes} and~\cite{kudlamillson-intersection},~$\theta_\eta$ is a \emph{holomorphic} modular form in~$M^k_{\gengenus,\lattice}$, with Fourier coefficients given by
%	\be\label{eq;FcoefnoncohoKMtheta}
%	c_T(\theta_\eta,\discel)=\int_{Z(T,\discel)}\eta\wedge (-\omega)^{\gengenus - \rk(T)}
%	\ee
%	for every~$\discel\in D_\lattice^\gengenus$ and~$T\in q(\discel)+\halfint{\gengenus}$ with~$T\ge0$.
%%	where~$Z(T,\discel)$ is the special cycle of index~$(T,\discel)$; see~\eqref{eq;defspcy} for an explicit construction of such cycles.
%%	Here~$\halfint{\gengenus}$ is the set of symmetric half-integral~$\gengenus\times\gengenus$-matrices introduced in Section~\ref{sec;backgrvvmodforms}.
%%	A similar formula is available also under the weaker assumption that~$T$ is positive-semidefinite~\cite[Theorem~$9.3$]{kudlamillson-intersection}, up to integrating~$\eta$ against a suitable power of the dual~$\omega^*$ of the Kähler form.
%	The Fourier coefficients~\eqref{eq;FcoefnoncohoKMtheta} are computed using explicit Poincaré dual forms of the cycles~$Z(T,\discel)$, which we recall below, arising from the Kudla--Millson Schwartz function~$\KMschwartz{\gengenus}$.
%%	The latter is a~$G$-invariant element of~$\mathcal{S}(V^g)\otimes\mathcal{A}^{2\genus}(X)$.
%%	We do not recall here the definition of~$\KMschwartz{\gengenus}$, and instead refer to~\cite{kudlamillson-harmonic1} for details.
%%	We recall them below together with some extensions of their Poincaré integral formulas to powers of the Kähler form~$\omega$ of~$X$.
%	
%	The special cycle~$Z(T,\discel)$, with~$T>0$, is an effective cycle whose irreducible components are immersions in~$X=\Gamma\backslash \domain$ of orthogonal Shimura varieties of the form~$\Gamma_\lambda\backslash \domain_\lambda$, for some~$\lambda\in q(\discel) + L^\gengenus$ such that~$q(\lambda)=T$.
%	Here we denote by~$\domain_\lambda$ the Hermitian symmetric subdomain of~$\domain$ defined as the orthogonal complement~$\lambda^\perp\subset\domain$ of codimension~$\gengenus$, and by~$\Gamma_\lambda$ the subgroup of~$\Gamma$ stabilizing~$\domain_\lambda$.
%	
%	The average of the Kudla--Millson Schwartz function with respect to~$\Gamma_\lambda\backslash\Gamma$ gives (up to a factor) a Poincaré dual form~$\zeta_\lambda$ of~$\Gamma_\lambda\backslash \domain_\lambda$ in~$X$.
%	As illustrated in~\cite[Section~$5$]{kudlamillson-tubes} and~\cite[Section~$9$]{kudlamillson-intersection}, this is
%	\be\label{eq;poincdualircomp}
%	\zeta_\lambda
%	\coloneqq
%	\exp\big(2\pi\trace q(\lambda)\big)\sum_{\gamma\in\Gamma_\lambda\backslash\Gamma}\gamma^* \KMschwartz{\gengenus}(\lambda).
%	\ee
%	The de Rham Poincaré duality reads as
%	\be\label{eq;poincdualstdwpd}
%	\int_X\eta\wedge\zeta_\lambda
%	=
%	\int_{\Gamma_\lambda\backslash\domain_\lambda}\eta,\qquad\text{for every \emph{compactly supported} closed~$2(n-\gengenus)$-form~$\eta$.}
%	\ee
%	The following result extends~\eqref{eq;poincdualstdwpd} to the case where~$\eta$ is a power of the Kähler form~$\omega$ of~$X$.
%	\begin{lemma}\label{lemma;extfromL2poincdual}
%	The differential form~$\zeta_\lambda$ is such that
%	\[
%	\int_X \omega^{n-\gengenus}\wedge \zeta_\lambda = \int_{\Gamma_\lambda\backslash\domain_\lambda}\omega^{n-\gengenus}.
%	\]
%	\end{lemma}
%	\begin{proof}
%	Recall that a differential form on~$X$ is said to be \emph{bounded} if its pointwise norm is bounded on~$X$.
%	Clearly, every compactly supported form on~$X$ is bounded.
%	Since the Schwartz function~$\KMschwartz{\gengenus}(\lambda)$ is rapidly decreasing along the fibers of the quotient map~$\Gamma_\lambda\backslash\domain\to\Gamma_\lambda\backslash\domain_\lambda$, the integral formula~\cite[Proposition~$5.2$]{kudlamillson-tubes} extends~\eqref{eq;poincdualstdwpd} to bounded differential forms~$\eta$.
%%	Note that we may even assume~$\eta$ to be square-integrable on~$X$ whenever~$g<(1+n)/2$.
%%	This follows from Lemma~\ref{lemma;L2poincduality}, namely the~$L^2$-Poincaré duality, and the isomorphisms of~$L^2$ and de Rahm cohomology groups~\eqref{eq;3isoofcoho}.
%	Since the Kähler form~$\omega$ is bounded~\cite[Remark~$4.21$]{kudla-integrals},
%%	 and hence also square-integrable,
%	we then deduce the integral formula of Lemma~\ref{lemma;extfromL2poincdual}.
%	\end{proof}
%	
%	The standard Poincaré integral formula reads as
%	\be\label{eq;poincdualstdwpd}
%	\int_X\eta\wedge\zeta_\lambda
%	=
%	\int_{\Gamma_\lambda\backslash\domain_\lambda}\eta,\qquad\text{for every compactly supported closed~$2(n-\gengenus)$-form~$\eta$.}
%	\ee
%	
%	Recall that a differential form on~$X$ is said to be \emph{bounded} if its pointwise norm is bounded on~$X$.
%	Clearly, every compactly supported form on~$X$ is bounded.
%	Since the Schwartz function~$\KMschwartz{\gengenus}(\lambda,z)$ is rapidly decreasing along the fibers of the quotient map~$\Gamma_\lambda\backslash\domain\to\Gamma_\lambda\backslash\domain_\lambda$, the integral formula~\cite[Proposition~$5.2$]{kudlamillson-tubes} extends~\eqref{eq;poincdualstdwpd} to bounded differential forms~$\eta$.
%%	Note that we may even assume~$\eta$ to be square-integrable on~$X$ whenever~$g<(1+n)/2$.
%%	This follows from Lemma~\ref{lemma;L2poincduality}, namely the~$L^2$-Poincaré duality, and the isomorphisms of~$L^2$ and de Rahm cohomology groups~\eqref{eq;3isoofcoho}.
%	Since the Kähler form~$\omega$ is bounded~\cite[Remark~$4.21$]{kudla-integrals},
%%	 and hence also square-integrable,
%	we then deduce the following result.
%	\begin{lemma}\label{lemma;extfromL2poincdual}
%	We have
%	Let~$Z$ be an irreducible component of a codimension~$\gengenus$ special cycle of a~$n$-dimensional smooth orthogonal Shimura variety~$X$.
%	If~$\eta_Z$ is a closed~$L^2$-form Poincaré dual to~$Z$, then
%	\[
%	\int_X \omega^{n-\gengenus}\wedge \eta_Z = \int_Z\omega^{n-\gengenus}.
%	\]
%	\end{lemma}
%	
	Let~$\discel\in D_\lattice^\gengenus$ and let~$T\in q(\discel)+\halfint{\gengenus}$ be positive definite.
	The differential form
	\[
	\zeta(T,\discel) \coloneqq 
	\exp\big(2\pi\trace T\big)
	\sum_{\substack{\lambda\in \discel+L^\gengenus \\ q(\lambda)=T}}
	\KMschwartz{\gengenus}(\lambda)
	\]
	is a Poincaré dual form of the special cycle~$Z(T,\discel)$; see~\cite[Section~$5$]{kudlamillson-tubes}.
%	It is a finite combination of the forms~$\zeta_\lambda$.	
%	In the reminder of the current section, we explain how to generalize Lemma~\ref{lemma;extfromL2poincdual} by replacing~$\zeta_\lambda$ with~$\zeta(T,\discel)$ and the Kudla--Millson theta function~$\KMtheta{\gengenus}$.
	For any quadratic space~$V$, we denote by~$\Witt(V)$ its Witt index.
	\begin{lemma}\label{lemma;L2cohoKMtheta&fexp}
%	Let~$X$ be an orthogonal Shimura variety arising from an even indefinite lattice~$L$ of signature~$(n,2)$, and let~$k=1+n/2$.
	If~$2\gengenus < n + 1 - \Witt(L\otimes\QQ)$, then the Poincaré dual form~$\zeta(T,\discel)$ of the special cycle~$Z(T,\discel)$ and the Kudla--Millson theta function~$\KMtheta{\gengenus}$ are~$L^2$-forms.
	If moreover~$2\genus < n-1$, then the~$L^2$-cohomology class~$[\KMtheta{\gengenus}]$ is a \emph{holomorphic} Siegel modular form of genus~$\gengenus$ and weight~$k=1+n/2$ with respect to~$\weil{L}{\genus}$ with Fourier expansion
	\be\label{eq;FexpKMthetaL2}
	[\KMtheta{\gengenus}]
	=
	\sum_{\discel\in D_\lattice^\gengenus}\sum_{\substack{T\in q(\discel) + \halfint{\gengenus}\\ T\ge 0}}
	[\zeta(T,\discel)]\wedge[-\omega]^{\gengenus-\rk(T)} q^T \mathfrak{e}_\discel.
	\ee
	\end{lemma}
	\begin{proof}
	To prove that~$\KMtheta{\gengenus}$ is square-integrable, we follow the argument of~\cite[Proof of Proposition~$4.1$]{bruinierfunke}. 
	It is enough to show that the~$\discel$-component of~$\KMtheta{\gengenus}$ is square-integrable for every~$\discel\in D_\lattice^\gengenus$.
	
	Let~$V=\lattice\otimes\RR$.
	The~$\discel$-component of~$\KMtheta{\gengenus}$ equals the theta function defined as
	\[
	\theta_\discel(\tau,\KMschwartz{\gengenus})
	\coloneqq
	\det y^{-k/2}\sum_{\lambda\in\discel + \lattice^{\gengenus}} \omega_\infty(g_\tau)\KMschwartz{\gengenus}(\lambda),
	\]
	where~$\omega_\infty$ denotes the Schrödinger model of (the restriction of) the Weil representation of~$\Mp_{2\gengenus}(\RR)$
%	$\Mp_{2\gengenus}(\RR)\otimes\OO(V^\gengenus)$
acting on the Schwartz functions on~$V^\gengenus$.
	The square of the~$L^2$-norm of that differential form is
	\[
	\lVert\theta_\discel(\tau,\KMschwartz{\gengenus})\rVert^2_2
	=
	\int_X \theta_\discel(\tau,\KMschwartz{\gengenus}) \wedge \ast\overline{\theta_\discel(\tau,\KMschwartz{\gengenus})},
	\]
	where~$\ast$ is the Hodge star operator.
	Since~$\KMschwartz{\gengenus}$ is a \emph{real} differential form,
%	\\\textcolor{blue}{\textbf{For Jan:} The idea of this proof is taken from~\cite{bruinierfunke}, where you worked with locally symmetric spaces.
%	They do not have a complex structure in general. In this article they have. In particular, conjugating a~$(g,g)$-form does not mean to only conjugate the coefficients of the differential form, but also conjugate the differential~$dz_I\wedge d \bar z_J$ appearing when the form is written on local charts.
%	\\
%	Here I want to make sure that the Kudla--Millson function~$\KMschwartz{\gengenus}$ is actually a \emph{real} differential form, namely
%	\[
%	\KMschwartz{\gengenus}=\overline{\KMschwartz{\gengenus}}.
%	\]
%	Take~$g=1$ for simplicity.
%	That function is defined explicitly as a differential form with \emph{real} coefficients given by some Gaussian multiplied by some Hermit polynomials.
%	The only thing one should check is that the vectors~$\omega_{\alpha,1}\wedge\omega_{\beta,2}\in \bigwedge^2 T^*_{z_0}\mathcal{D}$ appearing in the explicit definition are, with respect to the given complex structure, as
%	\[
%	dx_\alpha\wedge dy_\beta
%	\qquad\text{and not as}\qquad
%	dz_\alpha\wedge d\bar z_\beta,
%	\]
%	since in general the former is a real form while the latter is not, namely
%	\[
%	\overline{dx_\alpha\wedge dy_\beta} = dx_\alpha\wedge dy_\beta
%	\qquad\text{but}\qquad
%	\overline{dz_\alpha\wedge d\bar z_\beta} \neq dz_\alpha\wedge d\bar z_\beta.
%	\]
%	The point is to understand the complex structure of~$T^*_{z_0}\mathcal{D}$, where the latter is identified to the space of matrices~$\RR^{n\times 2}$.
%	As said by e.g.~\cite[p.~$66$, second line]{kudla-algcycles}, in terms of matrices the complex structure is given by right multiplication by~$(\begin{smallmatrix}
%	0 & 1\\
%	-1 & 0\end{smallmatrix})$.
%	I convinced myself that~$\omega_{\alpha,1}\wedge\omega_{\beta,2}$ is a real form, and that this works also if~$\gengenus>1$.
%	If you want we could discuss about this.
%	}\\
	it is easy to check that~$\overline{\omega_\infty(g_\tau)\KMschwartz{\gengenus}}=\omega_\infty(g_{-\overline{\tau}})\KMschwartz{\gengenus}$.
	We may then compute
	\[
	\ast\overline{\theta_\discel(\tau,\KMschwartz{\gengenus})}
	=
	\overline{\theta_\discel(\tau,\ast\KMschwartz{\gengenus})}
	=
	\theta_\discel(-\overline{\tau},\ast\KMschwartz{\gengenus}).
	\]
	Let~$\underline{\discel}=(\discel,\discel)\in D_\lattice^{2\gengenus}$ and let~$\KMschwartznew{\gengenus}(\underline{v})\nu=\KMschwartz{\gengenus}(v_1)\wedge\ast\KMschwartz{\gengenus}(v_2)$ for every~$\underline{v}=(v_1,v_2)\in V^2$, where~$\nu=\omega^n$ is the volume form of~$X$.
	Then
	\[
	\theta_\discel(\tau,\KMschwartz{\gengenus}) \wedge \ast\overline{\theta_\discel(\tau,\KMschwartz{\gengenus})}
	=
	\theta_{\underline{\discel}}(\tau,-\overline{\tau},\KMschwartznew{\gengenus})\nu
	\]
%	\textcolor{red}{\textbf{Remark for Jan:} I think that in your paper with Jens you used here (in case~$g=1$) that
%	\[
%	\omega_\infty\big(
%	g_{\left(\begin{smallmatrix}
%	\tau & 0\\
%	0 & -\overline{\tau}
%	\end{smallmatrix}\right)}
%	\big)\KMschwartznew{\gengenus}(\underline{\lambda})\nu
%	=
%	\omega_\infty(g_\tau)\KMschwartz{\gengenus}(\lambda_1)
%	\wedge\ast
%	\omega_\infty(g_{-\overline{\tau}})\KMschwartz{\gengenus}(\lambda_2),
%	\]
%	In particular that the Weil representation of~$\Sp_{4\gengenus}(\RR)$ has a good behavior w.r.t.\ the inclusion~$\Sp_{2\gengenus}(\RR)\times\Sp_{2\gengenus}(\RR)\hookrightarrow\Sp_{4\gengenus}(\RR)$.
%	This reflects the behavior of the Weil representation on~$\CC[D_\lattice^{2\gengenus}]$ given by Lemma~\ref{lemma;proofweilreponsubmeta}.
%	Do you think we should say anything about this?}\\
	may be regarded as the restriction to the diagonal of a theta function of genus~$2\gengenus$, where~$(\tau,-\overline{\tau})$ is identified with the element~$\left(\begin{smallmatrix}
	\tau & 0\\
	0 & -\overline{\tau}
	\end{smallmatrix}\right)\in\HH_{2\gengenus}$, and hence
	\be\label{eq;doublingsqint}
	\lVert\theta_\discel(\tau,\KMschwartz{\gengenus})\rVert^2_2
	=
	\int_X \theta_{\underline{\discel}}(\tau,-\overline{\tau},\KMschwartznew{\gengenus})\nu.
	\ee
	This integral converges by the Weil convergence criterion~\cite{weil}, since~$2\gengenus < n + 1 - \Witt(V)$.
	
	The proof that~$\zeta(T,\discel)$ is square-integrable is similar.
	We remark that
	\[
	\zeta(T,\discel)\wedge\ast\zeta(T,\discel)
	=
	\exp(4\pi\tr T)
	\sum_{\substack{\underline{\lambda}\in \underline{\discel} + \lattice^{2\gengenus}
	\\
	q(\underline{\lambda})=\big(\begin{smallmatrix}
	T & \ast \\
	\ast & T
	\end{smallmatrix}\big)
	}}
	\KMschwartznew{\gengenus}(\underline{\lambda})\nu,
	\]
	and it is the Fourier coefficient of~$\theta_{\underline{\discel}}(\tau_1,\tau_2,\KMschwartznew{\gengenus})\nu$ of index~$(T,T)$ at~$\tau_1=\tau_2=iI_\gengenus$.
%	\\\textcolor{red}{\textbf{Question for Jan: }I do not see why the last sentence is true.
%	What is the Fourier expansion of~$\theta_{\underline{\discel}}(\tau_1,\tau_2,\KMschwartznew{\gengenus})$ and how to compute it? (This theta function is not a holomorphic modular form.) 
%	}\\
%	\textcolor{blue}{
%	\textbf{Answer: }
%	The differential form~$\theta_{\underline{\discel}}(\tau_1,\tau_2,\KMschwartznew{\gengenus})\nu$ is a periodic smooth function in~$\tau_1,\tau_2\in\HH_\gengenus$.
%	Here the periodicity is with respect to certain translations on~$\HH_\gengenus$. (These are not necessarily half-integral.)
%	Hence, the theta function has a Fourier expansion of the form
%	\[
%	\theta_{\underline{\discel}}(\tau_1,\tau_2,\KMschwartznew{\gengenus})\nu
%	=
%	\sum_{T_1,T_2}c_{T_1,T_2}(y_1,y_2)\cdot e\big(\tr (T_1\tau_1) + \tr (T_2\tau_2)\big),
%	\]
%	where the sum is over symmetric~$\gengenus\times\gengenus$ matrices, and where~$\tau_j=x_j + i y_j$.
%	The Fourier coefficients~$c_{T_1,T_2}(y_1,y_2)$ depend on~$y_j$ since the theta function is non-holomorphic.
%	In fact, the Fourier coefficient of index~$(T_1,T_2)$ may be computed explicitly as
%	\be\label{eq;inquest}
%	c_{T_1,T_2}(y_1,y_2)
%	=
%	\sum_{\substack{\lambda_1,\lambda_2\in \discel + \lattice^{\gengenus}
%	\\
%	q(\lambda_j)=T_j
%	}}
%	\KMschwartznew{\gengenus}(\lambda_1 \sqrt{y_1}, \lambda_2 \sqrt{y_2}).
%	\ee
%	}\\
%	\textcolor{red}{\textbf{Important:} I think that~\eqref{eq;inquest} is wrong, and that it should be corrected by a factor.
%	More precisely, I think that
%	\[
%	c_{T_1,T_2}(y_1,y_2)
%	=
%	\exp\big(2\pi\tr( T_1 y_1 + T_2y_2)\big)
%	\sum_{\substack{\lambda_1,\lambda_2\in \discel + \lattice^{\gengenus}
%	\\
%	q(\lambda_j)=T_j
%	}}
%	\KMschwartznew{\gengenus}(\lambda_1 \sqrt{y_1}, \lambda_2 \sqrt{y_2}).
%	\]
%	This makes the situation more easier, because then~$\zeta(T,\discel)\wedge\ast\zeta(T,\discel)$ is really the Fourier coefficient of~$\theta_{\underline{\discel}}(\tau_1,\tau_2,\KMschwartznew{\gengenus})\nu$ of index~$(T,T)$ evaluated at~$\tau_1=\tau_2=i I_\gengenus$. (Before, it was true only up to a factor.)}
%	\\
%	\textcolor{blue}{\textbf{Here the proof:}
%	From the invariance of the function~$\theta_\discel(\tau,\KMschwartz{\gengenus})$ we deduce that it admits a Fourier expansion of the form
%	\[
%	\theta_\discel(\tau,\KMschwartz{\gengenus})
%	=
%	\sum_T c_T(y)\cdot e( \tr T\tau ),
%	\]
%	where the sum runs over the set of symmetric~$g\times g$ matrices.
%	The Fourier coefficient~$c_T$ gives, for every~$y$, a differential form~$c_T(y)$.
%	We may then compute
%	\bas
%	\theta_{\underline{\discel}}(\tau_1,\tau_2,\KMschwartznew{\gengenus})\nu
%	&\coloneqq
%	\theta_\discel(\tau_1,\KMschwartz{\gengenus}) \wedge \theta_\discel(\tau_2,\ast\KMschwartz{\gengenus})
%	\\
%	&=
%	\Big(
%	\sum_T c_T(y_1)\cdot e( \tr T\tau_1 )
%	\Big)
%	\wedge
%	\Big(
%	\sum_T \ast c_T(y_2)\cdot e( \tr T\tau_2 )
%	\Big)
%	\\
%	&=
%	\sum_{T_1,T_2} c_{T_1}(y_1)\wedge\ast c_{T_2}(y_2) \cdot e(\tr T_1 \tau_1 + \tr T_2\tau_2).
%	\eas
%	This means that the Fourier coefficient of index~$(T_1,T_2)$ of~$\theta_{\underline{\discel}}(\tau_1,\tau_2,\KMschwartznew{\gengenus})\nu$ is given by~$c_{T_1}(y_1)\wedge\ast c_{T_2}(y_2)$.
%	As we will prove below, the Fourier coefficients of~$\theta_\discel(\tau,\KMschwartz{\gengenus})$ may be computed as
%	\be\label{eq;Fcoefquest}
%	c_T(y)
%	=
%	\exp(2\pi\tr Ty)\sum_{\substack{\lambda\in \discel + \lattice^{\gengenus} \\ q(\lambda)=T}}\KMschwartz{\gengenus}(\lambda\sqrt{y}).
%	\ee
%	This implies that
%	\bas
%	c_{T_1}(y_1)\wedge\ast c_{T_2}(y_2)
%	&=
%	\exp\big(2\pi\tr( T_1 y_1 + T_2y_2)\big)\Big(
%	\sum_{\substack{\lambda_1\in \discel + \lattice^{\gengenus} \\ q(\lambda_1)=T_1}}\KMschwartz{\gengenus}(\lambda_1\sqrt{y_1})
%	\Big)\wedge
%	\Big(
%	\sum_{\substack{\lambda_2\in \discel + \lattice^{\gengenus} \\ q(\lambda_2)=T_2}}\ast\KMschwartz{\gengenus}(\lambda_2\sqrt{y_2})
%	\Big)
%	\\
%	&=
%	\exp\big(2\pi\tr( T_1 y_1 + T_2y_2)\big)\sum_{\substack{\lambda_1,\lambda_2 \in \discel + \lattice^{\gengenus} \\ q(\lambda_j)=T_j}}\KMschwartznew{\gengenus}(\lambda_1\sqrt{y_1},\lambda_2\sqrt{y_2}),
%	\eas
%	which is exactly~\eqref{eq;inquest} up to a factor of~$\exp\big(2\pi\tr( T_1 y_1 + T_2y_2)\big)$.\\
%	We conclude the proof by showing~\eqref{eq;Fcoefquest}.
%	It is enough to apply the standard formula to compute Fourier coefficients.
%	The function~$\theta_\discel(\tau,\KMschwartz{\gengenus})$ is a (non-holomorphic) modular form of lever~$\Gamma'$, for some subgroup~$\Gamma'$ of~$\Sp_{2\gengenus}(\ZZ)$.
%	For simplicity, we work here under the assumption that~$\Gamma'$ is the whole modular group, in particular~$T$ and~$q(\lambda)$ are half-integral matrices for every~$\lambda\in \discel+ L^\gengenus$.
%	Standard similar formulas hold otherwise.
%	We may compute
%	\bas
%	c_T(y)
%	&=
%	\int_{\text{$x$ mod~$1$}}\theta_\discel(x+iy,\KMschwartz{\gengenus})\cdot e^{-2\pi i \tr T\tau} \, dx
%	\\
%	&=
%	e^{2\pi\tr Ty}\det y^{-k/2}\int_{\text{$x$ mod~$1$}}\sum_{\lambda\in \discel+L^\gengenus}\omega_\infty\big(\begin{smallmatrix}
%	\sqrt{y} & \\
%	 & \sqrt{y}^{-t}
%	\end{smallmatrix}\big)
%	\omega_\infty
%	\big(\begin{smallmatrix}
%	I_g & x\\
%	 & I_g
%	\end{smallmatrix}\big)
%	\KMschwartz{\gengenus}(\lambda)
%	e^{-2\pi i \tr T x}\, dx
%	\\
%	&=
%	e^{2\pi\tr Ty}\det y^{-k/2}
%	\sum_{\lambda\in \discel+L^\gengenus}\omega_\infty\big(\begin{smallmatrix}
%	\sqrt{y} & \\
%	 & \sqrt{y}^{-t}
%	\end{smallmatrix}\big)\KMschwartz{\gengenus}(\lambda)
%	\int_{\text{$x$ mod~$1$}}
%	e^{2\pi i \tr (q(\lambda)-T) x}\, dx.
%	\eas
%	Since~$\int_{\text{$x$ mod~$1$}}
%	e^{2\pi i \tr (q(\lambda)-T) x}\, dx=0$ whenever~$q(\lambda)\neq T$ and is~$1$ otherwise, we deduce that
%	\bas
%	c_T(y)
%	&=
%	e^{2\pi\tr Ty}\det y^{-k/2}
%	\sum_{\substack{\lambda\in \discel+L^\gengenus \\ q(\lambda)=T}}\omega_\infty\big(\begin{smallmatrix}
%	\sqrt{y} & \\
%	 & \sqrt{y}^{-t}
%	\end{smallmatrix}\big)\KMschwartz{\gengenus}(\lambda)
%	\\
%	&=
%	e^{2\pi\tr Ty}
%	\sum_{\substack{\lambda\in \discel+L^\gengenus \\ q(\lambda)=T}}
%	\KMschwartz{\gengenus}(\lambda \sqrt{y}).
%	\eas
%	\textbf{WHAT I WOULD DO NOW:} I would write just after the definition of~$\zeta(T,\discel)$ that such a form is the value at~$\tau=i I_\gengenus$ of the Fourier coefficient of the~$\discel$-component of~$\KMtheta{\gengenus}(\tau)$. Then I would correct the formula for~$\zeta(T,\discel)\wedge\ast\zeta(T,\discel)$ and delete all red and blue text above. I do not think that it is relevant to keep the computations I wrote in blue above, since they are "well-known".
%	}\\
	The convergence of~$\lVert\zeta(T,\discel)\rVert_2$ follows from the termwise convergence of the integral over~$X$ of the Fourier coefficients of~$\theta_{\underline{\discel}}(\tau_1,\tau_2,\KMschwartznew{\gengenus})\nu$.
	This is again guaranteed by the Weil convergence criterion.
%	\\
%	\textcolor{red}{\textbf{Question for Jan:} You said that this is ensured by what Weil proved. Not only the holw integral converge, but also the integrals giving the Fourier coefficients. May I ask for the precise reference? (Paper and page where this appears)}\\
	
	We now consider the Fourier expansion of~$[\KMtheta{\gengenus}]$ under the assumption that~$\gengenus<(n-1)/2$.
	We know from~\eqref{eq;3isoofcoho} that~$\sqH^{2\gengenus}(X)\cong H^{2\gengenus}(X)$.
	This implies that all square-integrable representatives of~$\drcoho{Z(T,\discel)}$ yield the same~$L^2$-cohomology class.
	The fact that~$[\KMtheta{\gengenus}]$ is a holomorphic Siegel modular form with Fourier expansion~\eqref{eq;FexpKMthetaL2} follows from the analogous properties of~$\drcoho{\KMtheta{\gengenus}}$ proved in~\cite{kudlamillson-intersection}; see Section~\ref{sec;KMthetabasics}.
	\end{proof}
	We define the \emph{volume} of a special cycle~$Z(T,\discel)$ to be
	\[
	\Vol(Z(T,\discel))
	\coloneqq
	\int_{Z(T,\discel)}\omega^{n-\rk(T)}.
	\]
%	Lemma~\ref{lemma;extfromL2poincdual} and Lemma~\ref{lemma;L2cohoKMtheta&fexp} provide a formula for the computation of~$\Vol(Z(T,\discel))$ in terms of any square-integrable representative of the de Rham cohomology class~$\drcoho{Z(T,\discel)}$ of~$Z(T,\discel)$.
%	In fact, if~$T\in q(\discel)+\halfint{\gengenus}$ is positive definite and~$2\gengenus < n + 1 - \Witt(L\otimes\RR)$, then
%	\be\label{eq;formulaforvolspcy}
%	\Vol(Z(T,\discel))
%	=
%	\int_X
%	\omega^{n-\gengenus}\wedge\zeta(T,\discel)
%	=
%	\int_X
%	\omega^{n-\gengenus}\wedge [Z(T,\discel)],
%	\ee
%	where in the last equality we used the~$L^2$ version of Stokes Theorem.
	The integral~$\int_X \omega^{n-\gengenus}\wedge\KMtheta{\gengenus}$ is the \emph{generating series of volumes} of special cycles of codimension~$\genus$ in~$X$.
%	By Lemma~\ref{lemma;L2cohoKMtheta&fexp} we then deduce that if~$\genus<(n-1)/2$, then
%	\be\label{eq;genseriesofvolumes}
%	\int_X \omega^{n-\gengenus}\wedge\KMtheta{\gengenus}
%	=
%	\sum_{\discel\in D_\lattice^\gengenus}\sum_{\substack{T\in q(\discel) + \halfint{\gengenus}\\ T\ge 0}}
%	(-1)^{\gengenus-\rk(T)}\Vol\big(Z(T,\discel)\big) q^T \mathfrak{e}_\discel,
%	\ee
%	i.e.~$\int_X \omega^{n-\gengenus}\wedge\KMtheta{\gengenus}$ is the \emph{generating series of volumes} of special cycles in codimension~$\genus$.
%	
%	The following result of Kudla extends~\eqref{eq;genseriesofvolumes} to higher genera, and proves that the generating series of volumes is a Siegel Eisenstein series; see~\cite[Theorem~$4.20$ and Corollary~$4.24$]{kudla-speccyc}.
%	Recall that~$E^k_{\gengenus,L}$ is the normalized vector-valued Siegel Eisenstein series associated to the zero element of~$D_\lattice^\gengenus$; see Definition~\ref{def;KlingEisseries}.
	The following result of Kudla shows that this generating series of volumes is a Siegel Eisenstein series; see~\cite[Theorem~$4.20$ and Corollary~$4.24$]{kudla-speccyc}.
	Recall the normalized Siegel Eisenstein series~$E^k_{\gengenus,L}$ from Definition~\ref{def;KlingEisseries}.
	\begin{thm}[Kudla]\label{thm;fromKudlagenservol}
	Let~$X$ be an orthogonal Shimura variety arising from an even indefinite lattice~$L$ of signature~$(n,2)$, and let~$k=1+n/2$.
	If~$\gengenus< n+1 - \Witt(\lattice\otimes\QQ)$, then%~\eqref{eq;genseriesofvolumes} holds and
	\be\label{eq;gengen2inprodkmtheta}
	\int_X \KMtheta{\gengenus}
	\wedge
	\omega^{n-\gengenus}
	=
	(-1)^\gengenus \Vol(X) E^k_{\gengenus,L}.
	\ee
	\end{thm}
	\begin{proof}	
	Let~$\Ab_f$, resp.~$\Ab$, be the ring of finite adeles, resp.~adeles, over~$\QQ$ and let~$V=L\otimes\QQ$.
	For every~$\discel\in D_\lattice^\gengenus$, we may regard the basis vector~$\mathfrak{e}_\discel$ of~$\CC[D_\lattice^\gengenus]$ as the characteristic function~$\varphi_\discel$ of~$\discel + \widehat{L}^\gengenus$ in~$V(\Ab_f)^\gengenus$, where~$\widehat{L}^\gengenus=L^\gengenus\otimes\widehat{\ZZ}$.
	The space of Schwartz--Bruhat functions on~$V(\Ab)^\gengenus$ is denoted by~$\mathcal{S}\big(V(\Ab)^\gengenus\big)$.
	
	We now follow the argument of~\cite[p.~$640$]{bruinier-yang09} and~\cite[pp.~$324$-$325$]{kudla-integrals}, given in genus~$\gengenus=1$, generalizing it to~$\gengenus\ge1$ by means of~\cite{kudla-extensions}.
	For~$s\in\CC$, let~$I_\gengenus(s,\chi_V)$ be the principal series representation of~$\Sp_{2\gengenus}(\Ab)$ induced by~$\chi_V |\cdot |^s$, where~$\chi_V$ is the quadratic character of~$\Ab^*/\QQ^*$ associated to~$V$.
	For every standard section~$\Phi(s)\in I_\gengenus(s,\chi_V)$ we define the Eisenstein series
	\[
	E(g,s;\Phi) \coloneqq \sum_{\gamma\in P(\QQ)\backslash \Sp_{2\gengenus}(\QQ)} \Phi(\gamma g,s),\qquad g\in\Sp_{2\gengenus}(\Ab),
	\]
	where~$P(\QQ)$ is the Siegel parabolic subgroup of~$\Sp_{2\gengenus}(\QQ)$.
	The series~$E(g,s;\Phi)$ converges absolutely to an automorphic form on~$\Sp_{2\gengenus}(\Ab)$ if~$\Re(s)>(\gengenus+1)/2$.
	
	Consider the~$\CC[D_\lattice^\gengenus]$-valued Eisenstein series
	\ba\label{eq;vveisseriesasinkudlabt}
	E(\tau,s;\ell)
	\coloneqq
	\sum_{\discel\in D_\lattice^\genus}\det y^{-\ell/2} E\big(g_\tau,s;\Phi_\infty^\ell \otimes\lambda_f(\varphi_\discel)\big)\mathfrak{e}_\discel,
	\ea
	where $\Phi_\infty^\ell$ is the standard section of the principal series representation of~$\Sp_{2\gengenus}(\RR)$ arising from the character~$\theta\mapsto e^{i\ell\theta}$, and~$\lambda_f\colon\mathcal{S}\big(V(\Ab)^\gengenus\big)\to I_\gengenus(s_0,\chi_V)$ is the intertwining map such that~$\lambda_f(\varphi_\discel)(g_f)=\big(\omega(g_f)\varphi_\discel\big)(0)$ for every~$g_f\in \Sp_{2\gengenus}(\Ab_f)$; see~\cite[Section~$2.1$]{bruinier-yang09} for details.
	
	Let~$s_0=(n-\gengenus+1)/2$.
	As a consequence of the Siegel--Weil formula, Kudla proved in~\cite[Theorem~$4.23$]{kudla-integrals} and~\cite[Theorem~$4.1$]{kudla-speccyc} that the~$\discel$-component of~$E(\tau,s_0;k)$ is a \emph{holomorphic} Eisenstein series and equals the \emph{generating series of volumes}
	\[
	\frac{(-1)^\gengenus}{\Vol(X)}
	\int_X\omega^{n-\gengenus}\wedge(\KMtheta{\gengenus})_\discel,
%	\sum_{m>0}\Vol\big(Z(n,\discel)\big) q^n,
	\]
	where~$(\KMtheta{\gengenus})_\discel$ is the~$\discel$-component of~$\KMtheta{\gengenus}$.
%	Note that~${\varphi_\discel(0)=\delta_{\discel,0}}$.
	
	Following the argument of~\cite[p.~$641$]{bruinier-yang09} and~\cite[Section~IV.$2$]{kudla-extensions}, the vector-valued Eisenstein series~$E(\tau,s_0;k)$ may be rewritten as a Poincaré series arising from the basis vector~$\mathfrak{e}_0$ of~$\CC[D_\lattice^\gengenus]$, namely the Eisenstein series~$\eiszero{k}{\gengenus}{\lattice}$; cf.~\cite[$(2.17)$]{bruinier-yang09}.
	\end{proof}
	
%	\subsection{Integration along the fibers}\label{sec;integalongfib}
%	We provide here some background on integration along the fibers of vector bundles, following the wording of~\cite[Chapter~$1$, Section~$6$, p.~$61$]{botttu}.\\
%	
%	Let~$\pi\colon E\to N$ be a rank~$n$ vector bundle of orientable Riemannian manifolds.
%	We recall here the map~$\pi_*\colon\mathcal{A}^*(E)\to\mathcal{A}^{*-n}(N)$, usually referred to as the \emph{integration along the fiber of~$\pi$}.
%	We may assume that~$E$ is the trivial bundle~$E=N\times\RR^n$, since the general construction is locally given in the same way.
%	
%	We denote by~$t_1,\dots,t_n$ the standard coordinates of the fiber~$\RR^n$.
%	Every differential form on~$E$ is a combination of differential forms which either do not contain the~$n$-form~$dt_1\wedge\dots\wedge dt_n$ or do.
%	The former kind of differential form may be written as
%	\[
%	f(x,t_1,\dots,t_n)\cdot(\pi^*\phi)\wedge dt_{j_1}\wedge\dots\wedge dt_{j_r},\qquad r<n,
%	\]
%	for some smooth function~$f\colon E\to N$ and some differential form~$\phi$ on~$N$,	while the latter kind can be written as
%	\[
%	f(x,t_1,\dots,t_n)\cdot(\pi^*\phi)\wedge dt_1\wedge\dots\wedge dt_n.
%	\]
%	The function~$\pi_*$ is defined by mapping the first kind of differential forms to zero, and acting on the second kind as
%	\[
%	f(x,t_1,\dots,t_n)\cdot(\pi^*\phi)\wedge dt_1\wedge\dots\wedge dt_n
%	\longmapsto
%	\phi\int_{\RR^n} f(x,t_1,\dots,t_n)\,dt_1\dots dt_n.
%	\]
%	The map~$\pi_*$ may be also defined as the adjoint of the pullback~$\pi^*$ with respect to the pairing
%	\[
%	(\alpha,\beta)\mapsto\int_E\alpha\wedge\beta.
%	\]
%	
%	\subsection{The Kudla--Millson integral formula}
%	We recall in this section the integral formula provided in~\cite[Theorem~$2.1$]{kudlamillson-tubes}.
%	This is useful to relate the volume of special cycles with their Poincaré dual.\\
%	
%	Let~$E$ and~$N$ be orientable connected complete Riemannian manifolds of dimension respectively~$m$ and~$n$.
%	We assume that~$E$ has sectional curvature bounded by~$-\rho^2$, for some~$\rho\in\RR_{>0}$, and that~$N$ is a submanifold of~$E$.
%	We denote by~$j\colon N\to E$ the inclusion map.
%	A prototypical example relevant for what follows are the Hermitian symmetric space of an orthogonal Shimura variety and the embedded symmetric domain of an irreducible component of a special cycle.
%	
%	The Riemannian structures of~$N$ and~$E$ induce a natural vector bundle~${\pi\colon E\to N}$ as follows.
%	The exponential map from the normal bundle~$\nu(N)$ of~$N$ to~$E$ defined as
%	\[
%	\nu(N)\longrightarrow E,\qquad (x,v)\longmapsto \exp_x(v)\qquad\text{for every~$x\in N$ and~$v\in\nu(N)_x$}
%	\]
%	induces an isomorphism~$\iota\colon\nu(N)\to E$.
%	The map~$\pi$ is the vector bundle structure induced by~$\nu(N)\to N$ under~$\iota$.
%	Concretely, we may define~$\pi(e)$ by moving~$e\in E$ along a geodesic to a point~$\pi(e)\in N$, where the trajectory of the geodesic is induced by the exponential map from a unique vector on the fiber of the normal bundle~$\nu(N)$.
%	That fiber of~$\nu(N)$ is by definition the one at the point~$\pi(e)\in N$.
%	
%%	We may replace~$N$ by the inclusion of its universal cover into the universal cover of~$E$.
%%	\\
%%	\textcolor{red}{\textbf{Question:} They explicitly request this.
%%	Which kind of properties of the universal cover of~$N$ are needed in what follows that~$N$ might not have?
%%	}
%%	\\	
%	We are now ready to state the integral formula obtained by Kudla and Millson~\cite[Theorem~$2.1$]{kudlamillson-tubes}, which may be considered as a generalization of~\cite[Proposition~$6.15$]{botttu}.
%	\begin{thm}[Kudla--Millson]\label{thm;KMintegralform}
%	Let~$r=r(x)$ be the geodesic distance from~$x$ to~$N$, and let~$s\ge m-n$.
%	Let~$\phi$ be a closed smooth~$s$-form on~$E$ that is \emph{rapidly decreasing along the fibers of~$\pi$}, i.e.\ it is such that
%	\[
%	||\phi(x)||\le \exp(-m\rho r)\cdot \mathcal{P}(r)\qquad\text{for some polynomial~$\mathcal{P}$,}
%	\]
%	where the norm~$||\phi(x)||$ is the one pointwise induced by the Riemannian metric.
%	If~$\eta$ is a closed bounded~${(m-s)}$-form on~$E$, then
%	\[
%	\int_E \eta\wedge\phi = \int_N j^*(\eta)\wedge \pi_*(\phi),
%	\]
%	where~$\pi_*$ is the integration along the fiber introduced in Section~\ref{sec;integalongfib}.
%	\end{thm}
%	\noindent\textcolor{red}{
%	\textbf{Question:}
%	What does it mean for a differential form to be "bounded"? \textcolor{blue}{I think that we say~$\eta$ to be bounded if~$||\eta||$, as a real function on~$E$, is bounded.}}\\
%	
%%	\textcolor{blue}{\textbf{Note:} Kudla--Millson says that~$s\ge k$, where~$k$ is the number of elements needed to create an orthonormal frame field of the normal bundle (of the universal cover) of~$N$.
%%	Anyway,~$k$ is surely less or equal than~$m-n$, so the theorem stated above should work as it is, although it is not stated in full generality.}
%%	
%	We now fix our attention to our case of interest.
%	Let~$X=\Gamma\backslash D$ be a smooth orthogonal Shimura variety of dimension~$m$.
%	(In fact, we may work with more general locally symmetric spaces.)
%	The Hermitian symmetric space~$D$ is the Grassmannian of negative definite planes in~$V=L\otimes\RR$, for some lattice~$L$ of signature~$(m,2)$.
%	
%	The irreducible components of special cycles arise from symmetric subspaces~$D_x\subseteq D$ defined as the orthogonal complement in~$D$ of the span of the entries of~$x\in V^n$.
%	For simplicity, we assume that the moment matrix of~$x$ is positive definite.
%	The irreducible components of special cycles are the immersions in~$X$ of quotient spaces of the form~$N\coloneqq\Gamma_x\backslash D_x$, where~$\Gamma_x$ is the stabilizer of some~$x$ in~$\Gamma$.
%	
%	Let~$E\coloneqq\Gamma_x\backslash D$.
%	The exponential map induces a (surjective) diffeomorphism from the normal bundle of~$D_x$ to~$D$, giving rise to a vector bundle~${\pi\colon D\to D_x}$.
%	Passing to the quotient by~$\Gamma_x$, we obtain a vector bundle structure~$E\to N$ which we still denote by~$\pi$.
%	The following result is an application of Theorem~\ref{thm;KMintegralform}; see~\cite[Proposition~$5.2$ and Corollary on p.~$30$]{kudlamillson-tubes}.
%	\begin{prop}[Kudla--Millson]\label{prop;applofKMintegform}
%	Let~$\phi$ be a closed rapidly decreasing~$2n$-form on~$E$, and let~$\eta$ be a bounded~$2(m-n)$-form on~$E$.
%	Denote
%	\[
%	\kappa(x)=\int_{E_x}\phi,\qquad\text{for every $x\in N$}.
%	\]
%	Then~$\kappa(x)=\kappa$ is constant and
%	\[
%	\int_E\eta\wedge\phi = \kappa\int_{N}\eta.
%	\]
%	Moreover, if~$\eta$ is a bounded~$2(m-n)$-form on~$X$, then for~$\tilde{\phi}\coloneqq\sum_{\Gamma_x\backslash\Gamma}\gamma^*\phi$ we have
%	\[
%	\int_X\eta\wedge\tilde{\phi} = \kappa\int_{N}\eta.
%	\]
%	\end{prop}
%	
%	In~\cite{kudlamillson-tubes}, Kudla and Millson applied the previous result replacing~$\phi$ by the Kudla--Millson theta function~$\KMtheta{n}$.
%	In this way, they were able to compute the Fourier coefficients of the Siegel modular form
%	\[
%	\HH_n\longrightarrow\CC,\qquad\tau\longmapsto\int_X\eta\wedge\KMtheta{n},
%	\]
%	for every compactly supported~$\eta$.
%	In fact, since~$\eta$ is compactly supported, then it is bounded, hence Proposition~\ref{prop;applofKMintegform} may be applied.
%%	In~\cite[Theorem~$4.20$]{kudla-integrals}, Kudla replaced~$\eta$ by~$\omega^{m-n}$, since~$\omega$ is bounded as well.
%%	
%%	I think that we could prove a similar statement replacing~$[\KMtheta{g}]$ by the cohomology class of an irreducible component of a codimension~$g$ special cycle, by means of the Localization Principle as follows; see~\cite[Proposition~$6.25$, p.~67]{botttu}.
%	
%	In what follows, we replace~$\phi$ by a Poincaré dual of an irreducible component of a special cycle, and deduce an integral formula for the computation of volumes of subvarieties.
%	To show this, we firstly recall the Localization Principle for Poincaré duals of subvarieties; see~\cite[Proposition~$6.25$, p.~67]{botttu}.
%	\begin{prop}[Localization Principle]
%	Let~$M$ be an oriented manifold, and let~$S$ be an oriented submanifold of~$M$.
%	The support of
%%	(a representative of the cohomology class given by)
%	the Poincaré dual of~$S$ can be shrunk into any given tubular neighborhood of~$S$.
%	\end{prop}
%	\begin{lemma}\label{lemma;extfromL2poincdual}
%	Let~$\gengenus<(\dim X-1)/2$.
%	Let~$Z$ be an irreducible component of a codimension~$g$ special cycle of an~$m$-dimensional smooth orthogonal Shimura variety~$X$.
%	If~$\eta_Z$ is a~$L^2$-form Poincaré dual to~$Z$, then
%	\[
%	\int_X \omega^{m-g}\wedge \eta_Z = \int_Z\omega^{m-g}.
%	\]
%	\end{lemma}
%	\begin{proof}
%	Let~$E=\Gamma_x\backslash D$ and~$N=\Gamma_x\backslash D_x$ be as in Proposition~\ref{prop;applofKMintegform}, chosen such that the immersion of the submanifold~$N$ of~$E$ in~$X=\Gamma\backslash D$ is the reduced irreducible subvariety~$Z$.
%	Since~$\gengenus<(\dim X-1)/2$, every cohomology class in~$H^{2\gengenus}(X)$ has a square-integrable representative; see Section~\ref{sec;isoofseveralcoho}.
%	
%	Let~$\phi$ be a square-integrable closed~$2g$-form on~$E$ Poincaré dual to~$N$.
%	By the Localization Principle~\cite[Proposition~$6.25$, p.~67]{botttu}, we may choose~$\phi$ such that its support is contained in a tubular neighborhood~$T$ of~$N$ whose diameter, computed with respect to the geodesic distance to~$N$, is bounded.
%	\textcolor{red}{Problem: The Localization Principle does not deal with~$L^2$-forms.
%	One should prove that we can always find that~$\phi$ in a tubular neighborhood to be square-integrable...}
%	
%	Let~$\eta$ be a compactly supported~$2(m-\gengenus)$-form on~$E$.
%	The closed differential form~$\phi$ is of compact support along the fibers.
%	We may then apply Proposition~\ref{prop;applofKMintegform} and deduce that
%	\begin{equation}\label{eq;forThclstpre}
%	\int_E \eta\wedge\phi = \kappa\int_N\eta,\qquad\text{for some constant~$\kappa$.}
%	\end{equation}
%	By Poincaré duality we also know that~$\int_E\eta\wedge\phi = \int_N\eta$, hence~$\kappa=1$.
%	Therefore Proposition~\ref{prop;applofKMintegform} also implies that~$\tilde{\phi}=\sum_{\gamma\in\Gamma_x\backslash\Gamma}\gamma^*\phi$, as a differential form on~$X$, is Poincaré dual to the subvariety~$Z$.
%	
%	We now apply Proposition~\ref{prop;applofKMintegform} choosing~$\eta$ to be the non-compactly supported form~$\omega^{m-g}$, and deduce that
%	\begin{equation}\label{eq;forThclst}
%	\int_X \omega^{m-g}\wedge[Z] = \int_X \omega^{m-g}\wedge\tilde{\phi} = \int_Z\omega^{m-g}.
%	\qedhere
%	\end{equation}
%	\end{proof}
%	
%%	Let~$\Phi$ be the Kudla--Millson Schwartz function on the quotient~$E=\Gamma_x\backslash D$.
%%	As a differential form, the function~$\Phi$ satisfies the properties of its homologous~$\phi$ in Theorem~\ref{thm;KMintegralform}.
%%	In fact, it is well-known to be \emph{closed} and \emph{rapidly decreasing along the fibers of~$\pi$}; see~\cite[Corollary on p.~$29$]{kudlamillson-tubes}.
%	
%	Together with the Siegel--Weil formula, Proposition~\ref{prop;applofKMintegform} was applied by Kudla~\cite[Theorem~$4.1$]{kudla-speccyc} to show that the generating series of volumes of special cycles is an Eisenstein series.
%	\begin{thm}[Kudla]\label{thm;fromKudlagenservol}
%	Let~$X$ be a~$n$ dimensional orthogonal Shimura variety arising from a unimodular lattice.
%	Let~$(\cdot{,}\cdot)$ be the Hermitian inner product on~$H^{\gengenus,\gengenus}(X)$ defined in~\eqref{eq;hermformforus} and let~$k=1+n/2$.
%	If~$\gengenus<n+1-\Witt(V)$, then
%	\be\label{eq;gengen2inprodkmtheta}
%	\big( [\KMtheta{\gengenus}] , [\omega]^\gengenus \big)
%	= (-1)^\gengenus \Vol(X)\cdot E^k_\gengenus.
%	\ee
%%	\bes
%%	\big([\KMtheta{r}(\tau)],[\omega]^{r}\big) = \Vol(X)\cdot E^k_r(\tau, s_0),
%%	\ees
%%	where~$s_0=(n+1-r)/2$.
%	\end{thm}
	\begin{rem}\label{rem;projthetafcohoaftervol}
	Let~$0<\gengenus<(n-1)/2$.
	Recall that~$\sqH^{\gengenus,\gengenus}(X)$ is endowed with a Hermitian inner product~$(\cdot{,}\cdot)$ induced by the Kähler form~$\omega$.
	We may use~\eqref{eq;formmatchonHggsq} and Theorem~\ref{thm;fromKudlagenservol} to compute the orthogonal projection of~$[\KMtheta{\gengenus}]$ to~$\lefschetz^\gengenus \sqH^0(X)=\CC[\omega]^\gengenus$ as
	\[
	\frac{\big([\KMtheta{\gengenus}] , [\omega]^\gengenus\big)}{\big([\omega]^\gengenus,[\omega]^\gengenus\big)}[\omega^\gengenus]
	=
	(-1)^\gengenus E^k_{\gengenus,\lattice}\cdot [\omega]^\gengenus.
	\]
%	\textcolor{blue}{
%	\textbf{Quick check of the previous formula:}
%	We use~\eqref{eq;formmatchonHggsq} applied with~$r=j=\gengenus$ to deduce
%	\bas
%	\frac{\big([\KMtheta{\gengenus}] , [\omega]^\gengenus\big)}{\big([\omega]^\gengenus,[\omega]^\gengenus\big)}[\omega^\gengenus]
%	&=
%	\frac{(-1)^0\cdot(n-g)! \cdot n! \cdot \gengenus! \int_X \KMtheta{\gengenus} \wedge \omega^{n-g}}{(-1)^0 \cdot (n-\gengenus)!\cdot n!\cdot \gengenus! \int_X\omega^n}[\omega]^\gengenus\\
%	&=
%	\frac{\int_X \KMtheta{\gengenus} \wedge \omega^{n-g}}{\Vol(X)}[\omega]^\gengenus.
%	\eas
%	Then we replace~$\int_X \KMtheta{\gengenus} \wedge \omega^{n-g}$ by what is provided by Theorem~\ref{thm;fromKudlagenservol}, and we get what we want.}\\
	This means that
	\be\label{eq;gengen2inprodkmtheta2}
	[\KMtheta{\gengenus}] + (-1)^{\gengenus+1} E^k_{\gengenus,\lattice}\cdot [\omega]^\gengenus \in M^k_{\gengenus,\lattice}\otimes\big(\CC[\omega]^\genus\big)^\perp.
	\ee
	\end{rem}
	
	
%	\textcolor{red}{\textbf{Question:} To obtain a Lefschetz decomposition of~$[\KMtheta{\genus}]$ that is \emph{maximal}, it would be great to have a result as the following one.
%	\begin{lemma}[Is this true?]\label{lemma;inrediit2911}
%	Let~$r_1$ and~$r_2$ be positive integers such that~$r_1,r_2 < (n-1)/2$ and~$r_1+r_2 < n+1 - \Witt(\lattice\otimes\QQ)$.
%	Then
%	\bas
%	\int_X \KMtheta{r_1}(\tau_1) \wedge \KMtheta{r_2}(\tau_2) \wedge \omega^{n-r_1-r_2}
%	=
%	\int_X \KMtheta{r_1+r_2}\big(\begin{smallmatrix}
%	\tau_1 & \\
%	 & \tau_2
%	\end{smallmatrix}\big)
%	\wedge \omega^{n-r_1-r_2}
%	\eas
%	for every~$\tau_1\in\HH_{r_1}$ and~$\tau_2\in\HH_{r_2}$.
%	\end{lemma}
%	The problem is that the equality
%	\[
%	\drcoho{\KMtheta{r_1}(\tau_1)} \wedge \drcoho{\KMtheta{r_2}(\tau_2)}
%	= \drcohowide{\KMtheta{r_1+r_2}\big(
%	\begin{smallmatrix}
%	\tau_1 & \\
%	 & \tau_2
%	 \end{smallmatrix}\big)}
%	 \]
%	 is available only in de Rham cohomology.
%	 And we can not use that~$\sqH^{r_1+r_2,r_1+r_2}(X)\not\cong H^{r_1+r_2,r_1+r_2}(X)$, since in general~$r_1+r_2 \not < (n-1)/2$.
%	}
	\end{document}
	
	
	
	
	
	
	
	
	
	
	
	
	
	
	
	
\documentclass[../Lefschetz_decomposition.tex]{subfiles}

	\begin{document}
	In this section we show that the genus~$\gengenus$ Kudla--Millson theta function~$[\KMtheta{\gengenus}]$ is an almost cusp form, and compute its modular decomposition.
	We also show that this modular decomposition is of geometric interest, since it simultaneously provides the Lefschetz decomposition of~$[\KMtheta{\gengenus}]$ as a cohomology class.
	
	Recall from Section~\ref{sec;mod_dec} that we denote by~$\almcusp{k}{\genus}{\lattice}$ the space of almost cusp forms of weight~$k$ and genus~$\genus$ with respect to the Weil representation~$\weil{\lattice}{\genus}$.
	Its modular decomposition for~$k>2\genus$ is
	\[
	\almcusp{k}{\genus}{\lattice}
	=
	\CC\eiszero{k}{\genus}{L}
	\oplus
	\vklisub{k}{\gengenus}{1}{\lattice}
	\oplus
	\cdots
	\oplus
	\vklisub{k}{\gengenus}{\genus-1}{\lattice}
	\oplus
	S^k_{\genus,\lattice}.
	\]
	where~$\vklisub{k}{\gengenus}{r}{\lattice}$ is the space of vector-valued Klingen Eisenstein series arising from cusp forms in~$S^k_r$.
	Recall also that the Lefschetz decomposition of~$\sqH^{\gengenus,\gengenus}(X)$ is
	\[
	\sqH^{\gengenus,\gengenus}(X)
	=
	\CC[\omega]^\gengenus \oplus \lefschetz^{\gengenus-1} \sqHprim^{1,1}(X) \oplus \cdots \oplus \lefschetz \sqHprim^{\gengenus-1,\gengenus-1}(X)\oplus \sqHprim^{\gengenus,\gengenus}(X).
	\]
%	\subsection{Projection to~$\boldsymbol{L^\gengenus H^0(X)}$}
%	The following result of Kudla~\cite[Theorem~$4.1$]{kudla-speccyc} generalizes the modularity behavior of the generating series of volumes of special cycles in general genus, and follows from an application of the Siegel--Weil formula.
%	\begin{thm}[Kudla]\label{thm;fromKudlagenservol}
%	If~$r<n+1-\Witt(V)$, then
%	\be\label{eq;gengen2inprodkmtheta}
%	\big( [\KMtheta{\gengenus}] , [\omega]^\gengenus \big)
%	= (-1)^\gengenus \Vol(X)\cdot E^k_\gengenus.
%	\ee
%%	\bes
%%	\big([\KMtheta{r}(\tau)],[\omega]^{r}\big) = \Vol(X)\cdot E^k_r(\tau, s_0),
%%	\ees
%%	where~$s_0=(n+1-r)/2$.
%	\\ \textcolor{red}{In Kudla's paper there is no factor~$(-1)^r$.
%	In fact, I think Kudla uses~$\omega$ to denote~$\omega^*$, explaining why I needed to correct that formula with a $-1$.}
%	\end{thm}
%	\textcolor{red}{\textbf{Question:} Kudla does not state the theorem as I wrote in Theorem~\ref{thm;fromKudlagenservol}.
%	In fact, he says that the generating series of volumes is the value at some special point~$s_0=n/2$ of some non-holomorphic Eisenstein series.
%	However, in~\cite[p.~$310$]{kudla-integrals} he explicitly says: \emph{``... the vanishing of the right-hand side of (2.17) at~$s=s_0$ just shows the holomorphy of the special value~$E(\tau,s_0,\varphi,\ell+2)$''}, where the latter is exactly the generating series of volumes.
%	Now, I think that he means that such an Eisenstein series is a \emph{holomorphic} Eisenstein series with respect to the variable~$\tau$.
%	Since there is (up to factor) only one scalar valued Siegel Eisenstein series, I think that the special value a~$s=s_0$ he computed is really the Siegel Eisenstein series~$E^k_\gengenus$.
%	(Up to choosing the Kähler form in the correct way, we may assume that~$E^k_\gengenus$ is normalized wih first Fourier coefficient equal to~$1$.)\\
%	\textbf{Is this correct?
%	Should I say/check anything else about Theorem~\ref{thm;fromKudlagenservol} as it is now stated?}
%	}
%	
%	This implies that the projection of~$[\KMtheta{\gengenus}]$ to~$\lefschetz^\gengenus H^0(X)$ is~$E^k_\gengenus\cdot [\omega]^\gengenus$, hence
%	\be\label{eq;gengen2inprodkmtheta2}
%	[\KMtheta{\gengenus}] - E^k_\gengenus\cdot [\omega]^\gengenus \in \big(\lefschetz^\gengenus H^0(X)\big)^\perp.
%%	= L^{\gengenus-1,\gengenus-1}H^{1,1}_0(X)\oplus\cdots\oplus L H_0^{\gengenus-1,\gengenus-1}(X) \oplus H_0^{2,2}(X).
%	\ee
%	\begin{rem}
%	By Theorem~\ref{thm;fromKudlagenservol}, since the Fourier coefficient of~${[\KMtheta{\gengenus}] - E^k_\gengenus\cdot [\omega]^\gengenus}$ associated to the zero matrix is trivial, the projection of the Kudla--Millson theta function~$[\KMtheta{\gengenus}]$ to~$\big(\lefschetz^\gengenus \sqH^0(X)\big)^\perp$ lies in~$(\kling^k_\gengenus\oplus S^k_\gengenus) \otimes \big(\lefschetz^\gengenus \sqH^0(X)\big)^\perp$.
%	\end{rem}
	
	\subsection{Comparison of Siegel and Lefschetz operators}
	Let~$\Phi_\disceltwo\colon M^k_{\gengenus,\lattice}\to M^k_{\gengenus-1,\lattice}$ be the Siegel operator of index~$\disceltwo\in D_\lattice$ introduced in Section~\ref{sec;Siegelphi}.
	We denote the natural extension of that operator to~$M^k_{\gengenus,\lattice}\otimes \sqH^{\gengenus,\gengenus}(X)$ with the same symbol.
	\begin{lemma}\label{lemma;Siegopkmtheta2}
	If~$0<g<(n -1)/2$, then the image of~$[\KMtheta{\gengenus}]$ under the operator~$\Phi_\disceltwo$ is
%	, up to a sign, the image of~$[\KMtheta{\gengenus-1}]$ under the Lefschetz operator.
%	More precisely
	\[
	\Phi_\disceltwo\big([\KMtheta{\gengenus}]\big)
	=
	\begin{cases}
	- \lefschetz\big(
	[\KMtheta{\gengenus-1}]
	\big), & \text{if~$\disceltwo=0$,}
	\\
	0, & \text{otherwise.}
	\end{cases}
	\]
	\end{lemma}
	\begin{proof}
%	Under the assumption that~$g<\frac{\dim X -1}{2}$, the~$L^2$-cohomology group~$\sqH^{2\gengenus}(X)$ and the standard de Rham cohomology group~$H^{2\gengenus}(X)$ are isomorphic; see~\eqref{eq;3isoofcoho}.
%	In particular,
	By Lemma~\ref{lemma;L2cohoKMtheta&fexp} the~$L^2$-cohomology class~$[\KMtheta{\gengenus}]$ is a holomorphic Siegel modular form with respect to~$\weil{\lattice}{\genus}$ with Fourier expansion~\eqref{eq;FexpKMthetaL2}.
	By Lemma~\ref{lemma;FexpofimPhiop} we may compute
	\begin{align*}
	\Phi_\disceltwo\big([\KMtheta{\gengenus}]\big)
	&=
	\sum_{\discel\in D_\lattice^{\genus-1}}
	\sum_{\substack{T\in q(\discel)+\halfint{\genus-1} \\ T\ge 0}} \left[ Z\left(\left(\begin{smallmatrix}
	T & 0 \\ 0 & 0
	\end{smallmatrix}\right), (\discel,\disceltwo)\right)\right]\wedge[-\omega]^{\gengenus - \rk(T)} q^{T}\mathfrak{e}_\discel.
	\end{align*}
	It is easy to see from~\eqref{eq;defspcy} that
	\[
	Z\left(\left(\begin{smallmatrix}
	T & 0 \\ 0 & 0
	\end{smallmatrix}\right),(\discel,\disceltwo)\right)
	=
	\begin{cases}
	Z(T,\discel), & \text{if~$\beta=0$,}
	\\
	\emptyset, & \text{otherwise.}
	\end{cases}
	\]
	This implies that if~$\beta\neq 0$, then~$\Phi_\disceltwo\big([\KMtheta{\gengenus}]\big)=0$, and that
	\begin{align*}
	\Phi_0\big([\KMtheta{\gengenus}]\big)
	&=
	- \Big(
	\sum_{\discel\in D_\lattice^{\genus-1}}
	\sum_{\substack{T\in q(\discel)+\halfint{\genus-1} \\ T\ge 0}}
	[Z(T,\discel)]\wedge[-\omega]^{\gengenus -1 - \rk(T)} q^{T}\mathfrak{e}_\discel
	\Big) \wedge [\omega]
	\\
	&= - [\KMtheta{\gengenus-1}]\wedge[\omega]
	= - \lefschetz\big(
	[\KMtheta{\gengenus-1}]
	\big). \qedhere
	\end{align*}
	\end{proof}
%	We recall the following intersection formula from~\cite{kudla-algcycles}.
%	\begin{thm}[Kudla]\label{thm;formulaintspeccycl}
%	Let~$T_1,T_2$ be half-integral positive semi-definite matrices.
%	The cohomology class~$[Z(T_1)]\wedge[Z(T_2)]$ of special cycles may be computed as
%	\bes
%	[Z(T_1)] \wedge [Z(T_2)] 
%	= \sum_{T={\tiny \begin{pmatrix}
%	\begin{array}{c|c}
%	T_1 & A \\
%	\hline
%	A^t & T_2
%	\end{array}
%	\end{pmatrix}}}
%	[Z(T)],
%	\ees
%	where~$A$ runs over all matrices making~$T$ positive semi-definite.
%	\end{thm}
%	\textcolor{red}{\textbf{Question:} Kudla assumes the orthogonal Shimura varieties to be \emph{compact}.
%	In our case, they are not compact.
%	I there a generalization of the following result in the non-compact case? Which reference?}
%	\begin{cor}\label{cor;formulaintspeccycl}
%	If~$\tau_1\in\HH_{r_1}$ and~$\tau_2\in\HH_{r_2}$, then
%	\[
%	[\KMtheta{r_1}(\tau_1)] \wedge [\KMtheta{r_2}(\tau_2)]
%	= \left[\KMtheta{r_1+r_2}\Big({\tiny{\begin{pmatrix}
%	\begin{array}{c|c}
%	\tau_1 & 0\\
%	\hline
%	0 & \tau_2
%	\end{array}
%	\end{pmatrix}}}\Big)\right].
%	\]
%	\end{cor}
%	\begin{proof}
%	We check the formula above in terms of Fourier expansions.
%	We may compute
%	\bas
%	\left[\KMtheta{r_1+r_2}\Big({\tiny{\begin{pmatrix}
%	\begin{array}{c|c}
%	\tau_1 & 0\\
%	\hline
%	0 & \tau_2
%	\end{array}
%	\end{pmatrix}}}\Big)\right]
%	=&
%	\sum_{T}[Z(T)] \cdot e\Big(
%	\tr \, T\cdot{\tiny{\begin{pmatrix}
%	\begin{array}{c|c}
%	\tau_1 & 0\\
%	\hline
%	0 & \tau_2
%	\end{array}
%	\end{pmatrix}}}\Big)
%	\\
%	=&
%	\sum_{T={\tiny \begin{pmatrix}
%	\begin{array}{c|c}
%	T_1 & A \\
%	\hline
%	A^t & T_2
%	\end{array}
%	\end{pmatrix}}}[Z(T)] \cdot e\Big(
%	\tr {\tiny{\begin{pmatrix}
%	\begin{array}{c|c}
%	T_1 \tau_1 & *\\
%	\hline
%	* & T_2 \tau_2
%	\end{array}
%	\end{pmatrix}}}\Big)
%	\\
%	=&
%	\sum_{T_1}\sum_{T_2}\Big(
%	\sum_{T={\tiny \begin{pmatrix}
%	\begin{array}{c|c}
%	T_1 & A \\
%	\hline
%	A^t & T_2
%	\end{array}
%	\end{pmatrix}}}[Z(T)]
%	\Big)
%	\cdot e(\tr\, T_1\tau_1)\cdot e(\tr\, T_2 \tau_2).
%	\eas
%	By Theorem~\ref{thm;formulaintspeccycl} we may continue the computation as
%	\begin{align*}
%	\left[\KMtheta{r_1+r_2}\Big({\tiny{\begin{pmatrix}
%	\begin{array}{c|c}
%	\tau_1 & 0\\
%	\hline
%	0 & \tau_2
%	\end{array}
%	\end{pmatrix}}}\Big)\right]
%	=&
%	\sum_{T_1}\sum_{T_2}
%	[Z(T_1)]\wedge[Z(T_2)]
%	\cdot e(\tr\, T_1\tau_1)\cdot e(\tr\, T_2 \tau_2)
%	\\
%	=&
%	[\KMtheta{r_1}(\tau_1)] \wedge [\KMtheta{r_2}(\tau_2)]. \qedhere
%	\end{align*}
%	\end{proof}
%	\begin{cor}
%	If~$\tau_1\in\HH_{r_1}$ and~$\tau_2\in\HH_{r_2}$, then
%	\bes
%	\big([\KMtheta{r_1}(\tau_1)],[\KMtheta{r_2}(\tau_2)]\big)
%	=
%	\Vol(X) \cdot E^k_{r_1 + r_2}	\tiny{\begin{pmatrix}
%	\begin{array}{c|c}
%	\tau_1 & 0\\
%	\hline
%	0 & \tau_2
%	\end{array}
%	\end{pmatrix}}.
%	\ees
%	\end{cor}
%	\begin{proof}
%	By Corollary~\ref{cor;formulaintspeccycl} we may compute
%	\bas
%	\big([\KMtheta{r_1}(\tau_1)],[\KMtheta{r_2}(\tau_2)]\big)
%	=&
%	\int_X \left[\KMtheta{r_1+r_2}\Big({\tiny{\begin{pmatrix}
%	\begin{array}{c|c}
%	\tau_1 & 0\\
%	\hline
%	0 & \tau_2
%	\end{array}
%	\end{pmatrix}}}\Big)\right]\wedge\omega^{n - r_1 - r_2}
%	\\
%	=&
%	\Big(\left[\KMtheta{r_1+r_2}\Big({\tiny{\begin{pmatrix}
%	\begin{array}{c|c}
%	\tau_1 & 0\\
%	\hline
%	0 & \tau_2
%	\end{array}
%	\end{pmatrix}}}\Big)\right],\omega^{n - r_1 - r_2}\Big).
%	\eas
%	We deduce the final result by Theorem~\ref{thm;fromKudlagenservol}.
%	\end{proof}
%%%%%%%%%%%%%%%%%%%%%%%%%%%%%%%%%%%%%%%%%%%%%%%%%%%%%%%%%%%%%%%%%%%%%%%%%%%%%%%%%%%%%%%%%
%	\subsection{The Lefschetz decomposition in genus~2}
%	In this section we illustrate the Lefschetz decomposition of~$[\KMtheta{2}]$.
%	This will be used in Section~\ref{sec;cogengenus} as the first inductive step to prove the Lefschetz decomposition in general genus.
%	
%	\begin{thm}\label{thm;lefschetzdecingenus2}
%	The Lefschetz decomposition of~$[\KMtheta{2}]$ is
%	\begin{align}\label{eq;Lefdecgen2coho}
%	\begin{split}
%	[\KMtheta{2}]=
%	E^k_2\cdot[\omega]^2
%	-
%	E^k_{2,1}\big([\KMtheta{1}]_0\wedge[\omega]\big)
%	+
%	[\KMtheta{2}]_0,
%	\end{split}
%	\end{align}
%	where the primitive part of the cohomology class~$[\KMtheta{2}]$ is
%	\[
%	[\KMtheta{2}]_0=[\KMtheta{2}]	 + \klin^k_{2,1}\big([\KMtheta{1}]\wedge[\omega]\big).
%	\]
%	\end{thm}
%	\begin{proof}
%	The summands on the right-hand side of~\eqref{eq;Lefdecgen2coho} sum up to~$[\KMtheta{2}]$.
%	In fact, we know that the Lefschetz decomposition~$[\KMtheta{1}]=-E^k_1\cdot[\omega]+[\KMtheta{1}]_0$ given in~\eqref{eq;lefdecgen1coho} is such that~$[\KMtheta{1}]_0\in S^k_1\otimes H_0^{1,1}(X)$, hence
%	\bes
%	\klin^k_{2,1}\big([\KMtheta{1}]\wedge[\omega]\big)
%	=
%	- E^k_2\cdot[\omega]^2 + E^k_{2,1}([\KMtheta{1}]_0\wedge[\omega]).
%	\ees
%	
%	To prove that~\eqref{eq;Lefdecgen2coho} is the Lefschetz decomposition of~$[\KMtheta{2}]$, it is enough to show that the last term of the right-hand side of~\eqref{eq;Lefdecgen2coho}, namely
%	\bes
%	[\KMtheta{2}] + \klin^k_{2,1}\big([\KMtheta{1}]\wedge[\omega]\big),
%	\ees
%	lies in~$H_0^{2,2}(X)$, which is the orthogonal complement of~$\lefschetz H^{1,1}(X)$ in~$H^{2,2}(X)$ with respect to the Hermitian inner product~$(\cdot{,}\cdot)$ defined in~\eqref{eq;hermformforus}.
%	Since~$H^{1,1}(X)$ is generated by~$[\omega]$ and the fundamental classes of Heegner divisors, we deduce that
%	\bes
%	\lefschetz H^{1,1}(X)=\big\langle
%	[\omega]^2\quad\text{and}\quad [Z(m)]\wedge[\omega]\quad\text{for all $m>0$}
%	\big\rangle_\CC.
%	\ees
%	Hence, it is enough to prove that the last summand of the right-hand side of~\eqref{eq;Lefdecgen2coho} is orthogonal to~$[Z(m)]\wedge[\omega]$ for all~$m\in\ZZ_{\ge0}$.
%%	We begin with the orthogonality with respect to~$[\omega]^2$.
%%	Recall that~${[\KMtheta{2}] - E^k_2\cdot[\omega]^2}$ is orthogonal to~$[\omega]^2$.
%%	We may then compute
%%	\bas
%%	\Big(
%%	[\KMtheta{2}] + \klin^k_{2,1}([\KMtheta{1}\wedge[\omega]),[\omega]^2
%%	\Big)\hspace{8cm}
%%	\\
%%	=
%%	\underbrace{\Big(
%%	[\KMtheta{2}] - E^k_2\cdot[\omega]^2 , [\omega]^2
%%	\Big)}_{=0}
%%	+
%%	\Big(
%%	E^k_2\cdot[\omega]^2 + \klin^k_{2,1}([\KMtheta{1}]\wedge[\omega]),[\omega]^2
%%	\Big).
%%	\eas
%%	Replacing~$[\KMtheta{1}]$ above by its Lefschetz decomposition~$-E^k_1\cdot[\omega] + [\KMtheta{1}]_0$, we may use the linearity of the operator~$\klin^k_{2,1}$ and the fact that~$\klin^k_{2,1}(E^k_1)=E^k_2$ to deduce that
%%	\bas
%%	\Big(
%%	[\KMtheta{2}] + \klin^k_{2,1} & ([\KMtheta{1}\wedge[\omega]),[\omega]^2
%%	\Big)\\
%%	=&
%%	\Big(
%%	\underbrace{E^k_2\cdot[\omega]^2 - \klin^k_{2,1}(E^k_1)\cdot[\omega]^2}_{=0},[\omega]^2
%%	\Big)
%%	+
%%	\Big(
%%	E^k_{2,1}([\KMtheta{1}]_0\wedge[\omega]),[\omega]^2
%%	\Big)\\
%%	=&
%%	\Big(
%%	E^k_{2,1}([\KMtheta{1}]_0),[\omega]
%%	\Big)
%%	=0,
%%	\eas
%%	where the last vanishing follows from the fact that the Fourier coefficients of~$E^k_{2,1}([\KMtheta{1}]_0)$ are combinations of the ones of~$[\KMtheta{1}]_0$, hence they are orthogonal to~$[\omega]$.
%%	
%%	We conclude the proof by showing that~$[\KMtheta{2}] + \klin^k_{2,1}([\KMtheta{1}]\wedge[\omega])$ is orthogonal to~$[Z(m)]\wedge[\omega]$, for every~$m\in\ZZ_{>0}$.
%	This is equivalent of proving that
%	\be\label{eq;wwwtpinauxldt2}
%	\Big(
%	[\KMtheta{2}]\wedge[Z(m)] + \klin^k_{2,1}\big([\KMtheta{1}\wedge[Z(m)]\wedge[\omega]\big) , [\omega]^3
%	\Big)
%	=
%	0,\qquad\text{for all~$m\ge0$}.
%	\ee
%	
%	Let~$\Pi_j\colon H^{j,j}(X)\to(\CC[\omega]^j)^\perp$ be the projection to the orthogonal complement of~$\CC[\omega]^j$ in~$H^{j,j}(X)$.
%	Recall that
%	\bes
%	[\KMtheta{2}] = E^k_2\cdot[\omega]^2 + \Pi_2([\KMtheta{2}])
%	\qquad\text{and}\qquad
%%	[Z(m)]=-\frac{2\sigma_{k-1}(m)}{\zeta(1-k)}[\omega] + \Pi_1\big([Z(m)]\big),
%	[Z(m)]=-c_m(E^k_1)\cdot[\omega] + \Pi_1\big([Z(m)]\big),
%	\ees
%	for every~$m\in\ZZ_{\ge0}$; see~\eqref{eq;lefdecgen1coho}.
%	These imply that
%	\begin{align}\label{eq;auxdeckmthetatwo}
%	\begin{split}
%	[\KMtheta{2}]\wedge[Z(m)]
%	=&
%	-c_m(E^k_1)\cdot E^k_2\cdot[\omega]^3
%	+
%	E^k_2\cdot[\omega]^2\wedge\Pi_1\big([Z(m)]\big)\\
%	&-
%	c_m(E^k_1)\cdot[\omega]\wedge\Pi_2([\KMtheta{2}])
%	+
%	\Pi_2([\KMtheta{2}])\wedge\Pi_1\big([Z(m)]\big).
%	\end{split}
%	\end{align}
%	Since the last three summands appearing on the right-hand side of~\eqref{eq;auxdeckmthetatwo} are orthogonal to~$\CC[\omega]^3$, we deduce that
%	\be\label{eq;somcomphipad}
%	\Big([\KMtheta{2}]\wedge[Z(m)],[\omega]^3\Big)
%	=
%	-
%	c_m(E^k_1)\Vol(X) E^k_2.
%	\ee
%	With an analogous procedure, we may show that
%	\begin{align}\label{eq;auxdeckmthetatwoagain}
%	\begin{split}
%	[\KMtheta{1}]\wedge[Z(m)]\wedge[\omega]
%	=
%	c_m(E^k_1)\cdot E^k_1\cdot[\omega]^3
%	-
%	E^k_1\cdot[\omega]^2\wedge\Pi_1\big([Z(m)]\big)
%	\hspace{2cm}\\
%	-
%	c_m(E^k_1)\cdot[\omega]^2\wedge\Pi_1([\KMtheta{1}])
%	+
%	[\omega]\wedge\Pi_1([\KMtheta{1}])\wedge\Pi_1\big([Z(m)]\big).
%	\end{split}
%	\end{align}
%	Since the last three summands appearing on the right-hand side of~\eqref{eq;auxdeckmthetatwoagain} are orthogonal to~$\CC[\omega]^3$, and since~$\klin^k_{2,1}(E^k_1)=E^k_2$, we deduce that
%	\bes
%	\Big(
%	\klin^k_{2,1}\big(
%	[\KMtheta{1}]\wedge[Z(m)]\wedge[\omega]
%	\big),[\omega]^3
%	\Big)
%	=
%	c_m(E^k_1)\Vol(X) E^k_2.
%	\ees
%	This and~\eqref{eq;somcomphipad} imply~\eqref{eq;wwwtpinauxldt2}.
%	\end{proof}
%	
%	As in the case of genus~$1$, the Lefschetz decomposition of the Kudla--Millson theta function~$[\KMtheta{2}]$ provides automatically the decomposition of such a Siegel modular form in Eisenstein and cuspidal parts.
%	\begin{cor}\label{cor;lefdecrefldecmodforms}
%	The Lefschetz decomposition of~$[\KMtheta{2}]$ is also a decomposition of~$[\KMtheta{2}]$ as a Siegel modular form of genus~$2$ with respect to
%	\bes
%	M^k_2=\CC E^k_2 \oplus M^k_{2,1} \oplus S^k_2.
%	\ees
%	More precisely, we have
%	\begin{align}
%	E^k_2\cdot[\omega]^2 & \in \CC E^k_2 \otimes \CC[\omega]^2,\label{eqcor1;lefdecrefldecmodforms}
%	\\
%	-E^k_{2,1}\big([\KMtheta{1}]_0\wedge[\omega]\big) & \in M^k_{2,1}\otimes\lefschetz H_0^{1,1}(X),\label{eqcor2;lefdecrefldecmodforms}
%	\\
%	[\KMtheta{2}]_0 & \in S^k_2\otimes H_0^{2,2}(X).\label{eqcor3;lefdecrefldecmodforms}
%	\end{align}
%	\end{cor}
%	\begin{proof}
%	The inclusion~\eqref{eqcor1;lefdecrefldecmodforms} is trivial.
%	We showed in Section~\ref{sec;gen1} that the primitive part~$[\KMtheta{1}]_0$ is an elliptic cusp form with respect to the symplectic variable, i.e.~$[\KMtheta{1}]_0\in S^k_1\otimes H^{1,1}_0(X)$.
%	This implies~\eqref{eqcor2;lefdecrefldecmodforms}.
%	To prove~\eqref{eqcor3;lefdecrefldecmodforms}, it is enough to show that~$[\KMtheta{2}] + \klin^k_{2,1}\big([\KMtheta{1}]\wedge[\omega]\big)$ lies in the kernel of the Siegel operator~$\Phi$.
%	This follows from Lemma~\ref{lemma;Siegopkmtheta2}.
%	\end{proof}
	
	\subsection{The Kudla--Millson lift and its geometric interpretation}
	The \emph{Kudla--Millson lift of genus~$g$} is the linear map defined as
	\be\label{eq;KMliftdef}
	\KMlift{\gengenus}\colon S^k_{g,\lattice}\longrightarrow \sqH^{2\gengenus}(X,\CC),\qquad f\longmapsto\pet{f}{[\KMtheta{\gengenus}]}.
	\ee
	It has been introduced by Kudla and Millson in~\cite{kudlamillson-tubes}, and may be considered as a bridge to study the geometric properties of~$X$ in terms of the arithmetic properties of Siegel modular forms.
	
	The following result, essential for the proof of the Lefschetz decomposition of~$[\KMtheta{\genus}]$, provides a geometric interpretation of the image of~$\KMlift{\genus}$.
	\begin{thm}\label{thm;imageofKMlift}
%	Let~$X=\disk{}^+(L)\backslash\domain$.
	If~$\genus < (n+2)/4$, then	the image of the Kudla--Millson lift~$\KMlift{\gengenus}$ is contained in the primitive cohomology group~$\sqHprim^{\gengenus,\gengenus}(X)$.
%	If the codimension~$\genus$ special cycles on~$X$ arising from~$L$ span the whole~$\sqH^{\genus,\genus}(X)$, then the image of~$\KMlift{\gengenus}$ is~$\sqHprim^{\gengenus,\gengenus}(X)$.
	\end{thm}
	\begin{proof}
	We begin with the case of~$\Gamma=\disk{}^+(L)$.
	Let~$K$ be the open compact subgroup of~$G(\Ab_f)$ as in Section~\ref{sec;adelicspcy}.
	By Theorem~\ref{thm;noetherlef}, the space~$\sqH^{r,r}(X)$ is generated by the cohomology classes of adelic special cycles for every~$r\leq\genus-1$.
%	, since~$g-1<(n+1)/3$ under our assumptions.
	Recall that for every~$\varphi\in\mathcal{S}\big(V^{r}(\Ab_f)\big)$, there exists an even lattice~$M\subset V$ such that~$\varphi\in\mathcal{S}_{M^{r}}$; see e.g.~\cite{kudla-speccyc}.
	Hence, for every~$r\leq\genus-1$ there exists a sufficiently small even lattice~$M_r\subset V$ such that~$\sqH^{r,r}(X)$ is generated by the adelic special cycles of the form
	\[
	[Z(T,\varphi,K)]\wedge[\omega^{r-\rk T}],
	\qquad
	\text{where~$T\in\Sym_r(\QQ)$ and~$\varphi\in(\mathcal{S}_{(M_r)^r})^K$.}
	\]
	Without loss of generality, we may assume that the~$M_r$ are all equal to the same even lattice~$M\subset V$, and also that~$M\subset L$.
	
	Saying that the image of~$\KMlift{\genus}$ is contained in~$\sqHprim^{\genus,\genus}(X)=\big(\lefschetz \sqH^{\genus-1,\genus-1}(X)\big)^\perp$ is equivalent to saying that such an image is orthogonal to~$\lefschetz^{\genus-\rk T}\big([Z(T,\varphi,K)]\big)$ for every~$T\in \Sym_{\genus-1}(\QQ)$ and~$\varphi\in (\mathcal{S}_{M^{\genus-1}})^K$.
	
	Let~$X_M=\disk{}^+(M)\backslash\domain$, and let~$\pi\colon X_M\to X$ be the quotient map.
	We denote by~$\KMtheta{\genus-1}^M$ the Kudla--Millson theta function arising from the lattice~$M$.
	Let~$f\in S^k_{\genus,\lattice}$.
	It is easy to check by using~\eqref{eq;formmatchonHggsq} and induction on~$\genus$ that if
	\be\label{eq;wwwtpingengenuspldu}
	\int_{X_M} \pi^*\big(\KMlift{\genus}(f)\big)\wedge \KMtheta{\genus-1}^M \wedge \omega^{n-(2\genus-1)}
	=0,
	\ee
	then
	\be\label{eq;wwwtpingengenusplduimpl}
	\big(
	\KMlift{\genus}(f)
	,
	\lefschetz^{\genus-\rk T}[Z(T,\varphi,K)]
	\big)=0
	\qquad
	\text{for every~$T$ and~$\varphi$ as above}.
	\ee
%	\textcolor{blue}{\textbf{Remark on integrals of pullbacks:}
%	If~$[\eta]\in\sqH^{\genus,\genus}(X)$ with~$X=\Gamma\backslash\domain$ as above, then we may regard~$\eta$ as a~$\Gamma$-invariant~$(\genus,\genus)$-form on~$\domain$.
%	Note that~$\eta$ is also~$\Gamma_M\coloneqq\disk{}^+(M)$-invariant, since~$\Gamma_M\subseteq\Gamma$.
%	The same applies for the Kähler form~$\omega$, since it is~$G(\RR)$-invariant.
%	We deduce that
%	\be\label{eq;pullbackofintrc}
%	\int_{X_M}\pi^*(\eta)
%	\wedge
%	\omega^{n-\genus}
%	=
%	\#\big(\Gamma_M\backslash\Gamma\big)
%	\int_{X}\eta\wedge\omega^{n-\genus}.
%	\ee
%	This can be proved by unfolding as follows.
%	We denote by~$\mathcal{F}_\Gamma$ and~$\mathcal{F}_{\Gamma_M}$ respectively a fundamental domain for the action of~$\Gamma$ and of~$\Gamma_M$ on~$\domain$.
%	Write~$\mathcal{F}_{\Gamma_M}=\bigcup_\gamma \gamma(\mathcal{F}_\Gamma)$ for finitely many~$\gamma\in\Gamma$.
%	The number of such~$\gamma$'s is~$\#(\Gamma_M\backslash\Gamma)$.
%	Then
%	\bas
%	\int_{X_M}\pi^*(\eta)\wedge\omega^{n-\genus}
%	&=
%	\int_{\mathcal{F}_{\Gamma_M}}
%	\eta\wedge\omega^{n-\genus}
%	\\
%	&=
%	\int_{\bigcup_\gamma \gamma(\mathcal{F}_\Gamma)}
%	\eta\wedge\omega^{n-\genus}
%	\\
%	&=
%	\sum_\gamma
%	\int_{\gamma(\mathcal{F}_\Gamma)}
%	\eta\wedge\omega^{n-\genus}
%	\\
%	&=
%	\sum_\gamma
%	\int_{\mathcal{F}_\Gamma}
%	\gamma^*(\eta\wedge\omega^{n-\genus})
%	\\
%	&=
%	\sum_\gamma
%	\int_{\mathcal{F}_\Gamma}
%	\eta\wedge\omega^{n-\genus}
%	\\
%	&=
%	\#\big(\Gamma_M\backslash\Gamma\big)
%	\int_{\mathcal{F}_\Gamma}
%	\eta\wedge\omega^{n-\genus},
%	\eas
%	where in the first to the last passage we used that~$\gamma^*(\eta\wedge\omega^{n-\genus})=\eta\wedge\omega^{n-\genus}$ for all~$\gamma\in\Gamma$, i.e.~$\eta\wedge\omega^{n-\genus}$ is~$\Gamma$-invariant. 
%	}
%	\\
%	
%	\textcolor{blue}{\textbf{Here the check by induction on~$g$ that~\eqref{eq;wwwtpingengenuspldu} implies~\eqref{eq;wwwtpingengenusplduimpl}.}\\
%	The vanishing~\eqref{eq;wwwtpingengenuspldu} is the vanishing of a modular form with Fourier expansion
%	\be\label{eq;Fexpinproofblu}
%	\sum_{\discel\in D_M^{\genus-1}}
%	\sum_{\substack{T\in q(\discel)+\halfint{\genus-1}\\ T\ge 0}}
%	\Big(
%	\int_{X_M} \pi^*\big(\KMlift{\genus}(f)\big)
%	\wedge [Z(T,\discel)]
%	\wedge
%	[-\omega]^{\genus-1-\rk T}
%	\wedge [\omega]^{n-(2\genus-1)}
%	\Big) q^T \mathfrak{e}_\discel,
%	\ee
%	where here the cohomology classes are taken in~$\sqH^\bullet(X_M)$.
%	We prove~\eqref{eq;wwwtpingengenusplduimpl} by induction on~$\genus$.
%	\\
%	\textbf{Case $\boldsymbol{g=1}$:}\\
%	The theta function~$\KMtheta{0}^M=1$ is a non-zero constant function.
%	By~\eqref{eq;formmatchonHggsq} and~\eqref{eq;pullbackofintrc} (applied with~$\eta=\KMlift{1}(f)$) the left-hand side of~\eqref{eq;wwwtpingengenuspldu} equals the left-hand side of~\eqref{eq;wwwtpingengenusplduimpl} up to a non-zero constant, hence the assertion follows immediately.
%	\\\textbf{Case $\boldsymbol{\genus > 1}$:}\\
%	We work by induction on~$\rk T$.
%	If~$\rk T =0$, namely~$T=0$, then by~\eqref{eq;formmatchonHggsq} and~\eqref{eq;pullbackofintrc} (applied with~$\eta=\KMlift{\genus}(f)$) we may rewrite the left-hand side of~\eqref{eq;wwwtpingengenusplduimpl} as
%	\bas
%	\big(\#(\Gamma_M\backslash\Gamma)\big)^{-1}\big(\pi^*\big(\KMlift{\genus}(f)\big),[\omega]^\genus\big)
%	=
%	\frac{(-1)^0\cdot n!\cdot \genus!}{(n-\genus)! \cdot \#(\Gamma_M\backslash\Gamma)}
%	\int_{X_M} 
%	\pi^*\big(\KMlift{\genus}(f)\big)\wedge\omega^{n-\genus}.
%	\eas
%	This is, up to a non-zero factor, the constant term of the Fourier expansion~\eqref{eq;Fexpinproofblu}, which vanishes by~\eqref{eq;wwwtpingengenuspldu}.}
%	
%	\textcolor{blue}{Suppose now by induction that~$\Big(\KMlift{\genus}(f),\lefschetz^{\genus-\rk T}\big([Z(T,\varphi,K)]\big)\Big)=0$ for every~$\varphi\in\big(\mathcal{S}_{M^{\genus-1}}\big)^K$ and~$T\in \Sym_{\genus-1}(\QQ)$ of rank less than some positive~$r$.
%	This implies that~$\KMlift{\genus}(f)$ is orthogonal to~$\lefschetz^{\genus-r+1}\sqH^{r-1,r-1}(X)$, since~$\sqH^{r-1,r-1}(X)$ is generated by classes of adelic special cycles arising from~$(\mathcal{S}_{M^{r-1}})^K$.
%	\\
%	Let~$T_r\in \Sym_{\genus-1}(\QQ)$ be a matrix of rank~$r$, and let~$\Pi_j\colon \sqH^{\genus,\genus}(X)\to\lefschetz^{j}\sqHprim^{\genus-j,\genus-j}(X)$ be the orthogonal projection to the~$j$-th term of the Lefschetz decomposition of~$\sqH^{\genus,\genus}(X)$.
%	We may compute
%	\begin{align*}
%	\Big(\KMlift{\genus}(f),\lefschetz^{\genus-r}\big([Z(T_r,\varphi,K)]\big)\Big)
%	&=
%	\sum_{j=\genus-r}^\genus
%	\Big(\KMlift{\genus}(f),\Pi_j\circ
%	\lefschetz^{\genus-r}\big([Z(T_r,\varphi,K)]\big)\Big)
%	\\
%	&=
%	\Big(\KMlift{\genus}(f),\Pi_{\genus-r}\circ
%	\lefschetz^{\genus-r}\big([Z(T_r,\varphi,K)]\big)\Big)
%	\\
%	&=
%	\Big(\KMlift{\genus}(f),\lefschetz^{\genus-r}\big([Z(T_r,\varphi,K)]_0\big)
%	\Big)
%	\\
%	&=C\cdot\int_X \KMlift{\genus}(f)\wedge [Z(T_r,\varphi,K)]_0 \wedge[\omega]^{n-\genus-r},
%	\end{align*}
%	where in the second passage above we used the inductive hypothesis, and where~$C$ is a non-zero constant.
%	By~\eqref{eq;pullbackofintrc} we then deduce that
%	\ba\label{eq;inbluewpofKMZ}
%	&\Big(\KMlift{\genus}(f),\lefschetz^{\genus-r}\big([Z(T_r,\varphi,K)]\big)\Big)
%	\\
%	&\quad=
%	\frac{C}{\#(\Gamma_M\backslash\Gamma)}
%	\int_{X_M}
%	\pi^*\big(\KMlift{\genus}(f)\big)\wedge \pi^*\big([Z(T_r,\varphi,K)]_0\big) \wedge[\omega]^{n-\genus-r}.
%	\ea
%	With a similar procedure, we now prove that the vanishing of~\eqref{eq;Fexpinproofblu} implies the vanishing of the integral on the right-hand side of~\eqref{eq;inbluewpofKMZ}, and hence of~$\big(\KMlift{\genus}(f),\lefschetz^{\genus-r}\big([Z(T_r,\varphi,K)]\big)\big)$.
%	The vanishing of~\eqref{eq;Fexpinproofblu} implies that
%	\[
%	\int_{X_M}\pi^*\big(\KMlift{\genus}(f)\big)
%	\wedge
%	[Z(T_r,\varphi,K)]
%	\wedge
%	\omega^{n-\genus-r}
%	=
%	0
%	\]
%	for every~$\varphi\in (\mathcal{S}_{M^{\genus-1}})^K$.
%	Since~$\sqH^{r-1,r-1}(X_M)^{\Gamma_M\backslash\Gamma}=\sqH^{r-1,r-1}(X)$, we deduce that the former cohomology group is generated by the adelic special cycles~$[Z(T,\varphi,K)]\wedge[\omega^{r-1-\rk T}]$ with~$\varphi\in(\mathcal{S}_{M^{\genus-1}})^K$ and~$T$ of rank less than~$r$.
%	With the same idea as above, using the inductive hypothesis and the projections~$\Pi_j$, one can deduce the vanishing of the right-hand side of~\eqref{eq;inbluewpofKMZ}.
%	}
	To show that the image of~$\KMlift{\genus}$ lies in~$\sqHprim^{\genus,\genus}(X)$, it is then enough to prove~\eqref{eq;wwwtpingengenuspldu}.
	
	By Lemma~\ref{lemma;genofBYandS} and Remark~\ref{rem;resandpullbKMthetafunct} we deduce that
	\[
	\pi^*\big(\KMlift{\genus}(f)\big)
	=
	\bigpet{f}{\pi^*([\KMtheta{\genus}])}
	=
	\bigpet{f}{\tr_{\genus,\lattice/M}([\KMtheta{\genus}^M])}
	=
	\bigpet{f_M}{[\KMtheta{\genus}^M]}.
	\]
	By the Fubini--Tonelli Theorem, we may then rewrite the integral in~\eqref{eq;wwwtpingengenuspldu}, keeping track of the variable~$\tau_1\in\HH_{\genus-1}$, as
	\begin{align*}
	&\int_{X_M}
	\pi^*\big(\KMlift{\genus}(f)\big)
	\wedge
	\KMtheta{\genus-1}^M(\tau_1)
	\wedge
	\omega^{n-(2\genus-1)}
	\\
	&\quad
	=
	\Bigpet{f_M}{\int_{X_M} \KMtheta{\genus-1}^M(-\overline{\tau_1}) \wedge \KMtheta{\genus}^M \wedge \omega^{n-(2\genus-1)}},
	\end{align*}
	where to deduce the equality above we used that~$\KMtheta{\genus-1}$ is a modular form with \emph{real} Fourier coefficients, and hence~$\overline{\KMtheta{\genus-1}(\tau_1)} = \KMtheta{\genus-1}(-\overline{\tau_1})$.
	By~\eqref{eq;wedgeprodformforKMtheta} and Theorem~\ref{thm;fromKudlagenservol} we may continue the computation above as
	\bas
	\,&\int_{X_M} \pi^*\big(\KMlift{\genus}(f)\big)\wedge
	\KMtheta{\genus-1}^M(\tau_1)\wedge
	\omega^{n-(2\genus-1)}
	\\
	&
	\quad
	=
	\Bigpet{f_M}{\int_{X_M} \KMtheta{2\genus-1}^M\big(\begin{smallmatrix}
	-\overline{\tau_1} & \\
	 & \ast
	\end{smallmatrix}\big) \wedge \omega^{n-(2\genus-1)}}
	\\
	&\quad
	=
	(-1)^{2\genus-1}
	\Vol(X_M)
	\Bigpet{f_M}{E^k_{2\genus-1,M}\big(\begin{smallmatrix}
	-\overline{\tau_1} & \\
	 & \ast
	\end{smallmatrix}\big)},
	\eas
	where we denoted by~$\ast$ the variable of integration in~$\HH_{\genus}$ for the Petersson inner product.
%	Note that we are allowed to use Theorem~\ref{thm;fromKudlagenservol}, since~$2\genus-1<n-1\leq n+1-\Witt(M\otimes\QQ)$.
	
	By Corollary~\ref{cor;vanishingGmapvv} we deduce that
	\[
	\Bigpet{f_M}{E^k_{2\genus-1,M}\big(\begin{smallmatrix}
	-\overline{\tau_1} & \\
	 & \ast
	\end{smallmatrix}\big)} = 0,
	\]
	which then implies that~$\KMlift{\genus}(S^k_\genus)\subseteq \sqHprim^{\gengenus,\gengenus}(X)$.	
	
	We now consider the case where~$\Gamma$ is a subgroup of finite index in~$\Gamma_L\coloneqq\disk{}^+(\lattice)$.
%	Here~$\Gamma_L$ may have torsion, hence~$X_L$ can have finite-quotient singularity.
	Let
	\[
	\pi\colon X_\Gamma\coloneqq \Gamma\backslash\domain
	\longrightarrow
	X_L\coloneqq\Gamma_\lattice\backslash\domain
	\]
	be the projection map.
	The pullback~$\pi^*\colon H^{2\genus}(X_\lattice)\to H^{2\genus}(X_\Gamma)\cong \sqH^{2\genus}(X_\Gamma)$ is injective and identifies~$H^{2\genus}(X_\lattice)$ with the subspace~$\sqH^{2\genus}(X_\Gamma)^{\Gamma\backslash\Gamma_\lattice}$ of~$\Gamma\backslash\Gamma_\lattice$-invariant classes in~$\sqH^{2\genus}(X_\Gamma)$.
	The cohomology class of the Kudla--Millson theta function attached to~$L$ induced on~$X_\Gamma$ is the pullback~$\pi^*\big([\KMtheta{\genus}]\big)\in \sqH^{2\genus}(X_\Gamma)^{\Gamma\backslash\Gamma_\lattice}$ of the one defined on~$X_\lattice$.
	Since the Kähler form~$\omega$, and hence the Hermitian inner product on~$\sqH^{2\genus}(X_\Gamma)$, is~$\Gamma_\lattice$-invariant, we deduce that
%	the Hodge and Lefschetz decompositions of~$\sqH^{2\genus}(X_\Gamma)$ preserve such invariance, in the sense that
	\[
	\sqH^{2\genus}(X_\Gamma)^{\Gamma\backslash\Gamma_\lattice}
	=
	\bigoplus_{p+q=2\genus} \sqH^{p,q}(X_\Gamma)^{\Gamma\backslash\Gamma_\lattice}
	\qquad
	\text{and}
	\qquad
	\sqH^{\genus,\genus}(X_\Gamma)^{\Gamma\backslash\Gamma_\lattice}
	=
	\bigoplus_{i\geq 0}
	\lefschetzî\sqHprim^{r-2i}(X_\Gamma)^{\Gamma\backslash\Gamma_\lattice}.
	\]
%	\textcolor{blue}{\textbf{Regarding the Hodge decomposition.}
%	Let~$\gamma\in\Gamma_L$. If~$\alpha$ is a~$(p,q)$ form on~$X_\Gamma$, then also its pullback~$\gamma^*\alpha$ is a~$(p,q)$-form on~$X_\Gamma$.
%	The check of this can be done on a local charts of~$X_\Gamma$, using the fact that~$\gamma$ acts on~$X_L$ as a holomorphic map.
%	Hence, if a differential form on~$X_\Gamma$ is~$\Gamma_L$-invariant, each term of its Hodge decomposition is~$\Gamma_L$-invariant.
%	\\
%	\textbf{Regarding the Hodge decomposition.}
%	We prove the second equality displayed above by induction.
%	If~$\genus=1$ and~$[\alpha]\in \sqH^{1,1}(X_\Gamma)$, then
%	\[
%	[\alpha]=A\cdot[\omega] + B\cdot[\alpha]_0,
%%	\sqH^{1,1}(X_\Gamma)=\CC[\omega]\oplus\sqHprim^{1,1}(X_\Gamma),
%	\qquad\text{where}\qquad
%	[\alpha]_0\in\sqHprim^{1,1}(X_\Gamma) = \big(\CC[\omega]\big)^\perp.
%	\]
%	Since~$\omega$ is $\Gamma_L$-invariant, we deduce that if~$\gamma\in\Gamma_L$ and~$[\alpha]$ is~$\Gamma_L$-invariant, then
%	\[
%	A\cdot[\omega] + B\cdot[\alpha]_0
%	=
%	[\alpha]
%	=
%	\gamma^*[\alpha]
%	=
%	A\cdot\underbrace{\gamma^*[\omega]}_{=[\omega]} + {}B\cdot\gamma^*[\alpha]_0.
%	\]
%	This implies that also~$[\alpha]_0$ is~$\Gamma_L$-invariant.\\
%	An induction argument using that
%	\[
%	\sqH^{\genus,\genus}(X_\Gamma)=\lefschetz\sqH^{\genus-1,\genus-1}(X_\Gamma)\oplus\sqHprim^{\genus,\genus}(X_\Gamma)
%	\]
%	and that the operator~$\lefschetz$ is~$\Gamma_L$-invariant concludes the proof.
%	}
%	\\
Since the pullback~$\pi^*$ maps primitive cohomology groups of~$X_L$ into primitive cohomology groups of~$X_\Gamma$, the image of~$\KMlift{\genus}$ lies in primitive cohomology also on~$X_\Gamma$.
%	By Theorem~\ref{thm;noetherlef} the space~$\sqH^{\genus,\genus}(X_\Gamma)^{\Gamma\backslash\Gamma_\lattice}$ is generated by the pullback under~$\pi$ of the classes of adelic special cycles in~$X_\lattice$.
%%	\\\textcolor{blue}{\textbf{Here a check of the last assertion:}\\
%%	It is enough to prove it in the case where~$\Gamma=\Gamma_N\coloneqq\disk{}^+(L(N))$, for some~$N$.
%%	As above, we assume~$\Gamma_N$ to be torsion-free.
%%	\\
%%	Let~$X_N\coloneqq\Gamma_N\backslash\domain$.
%%	Denote by~$SC^\genus(X)$, resp.~$SC^\genus(X_N)$, the subspace of~$H^{2\genus}(X)$, resp.~$H^{2g}(X_N)$, generated by the adelic special cycles of~$X$, resp.\ of~$X_N$.
%%	Note that since~$X_N$ is smooth, we may consider~$SC^\genus(X_N)$ as a subspace of~$\sqH^{\genus,\genus}(X_N)$.
%%	By Theorem~\ref{thm;noetherlef} we deduce that
%%	\be\label{eq;inbluespaofspblmm}
%%	SC^\genus(X_N)=\sqH^{\genus,\genus}(X_N).
%%	\ee
%%	Note that the subspace of~$\sqH^{\genus,\genus}(X_N)$ generated by the pullbacks of the special cycles of~$X$ equals~$SC^\genus(X_N)^{\Gamma_N\backslash\Gamma_L}$; see~\cite[First displayed formula in the proof of Corollary~$3.8$]{noetherlefconj}.
%%	But from~\eqref{eq;inbluespaofspblmm} we see that
%%	\[
%%	SC^\genus(X_N)^{\Gamma_N\backslash\Gamma_L}
%%	=
%%	\sqH^{\genus,\genus}(X_N)^{\Gamma_N\backslash\Gamma_L}.
%%	\]
%%	Hence~$\sqH^{\genus,\genus}(X_N)^{\Gamma_N\backslash\Gamma_L}$ is generated by the pullback of the special cycles on~$X_L$.
%%	}\\
%	By replacing the cohomology groups of~$X_\Gamma$ by their~$\Gamma\backslash\Gamma_\lattice$-invariant subspaces, one can easily see that the same proof given above for the case of~$\Gamma=\Gamma_\lattice$ works also for finite-index subgroups of~$\Gamma_L$.
	\end{proof}	
	
	\subsection{The modular Lefschetz decomposition in genus 1}\label{sec;gen1}
	In this section we describe the modular and Lefschetz decompositions of the class~$[\KMtheta{1}]$ of the genus~$1$ Kudla--Millson theta function.
%	Recall we assume that the orthogonal Shimura variety~$X$ arises from a unimodular lattice.
	
%	Recall that the volume of a special cycle~$Z(T)$, where~$T\in \halfint{\gengenus}$ with~$T\ge 0$, is
%	\[
%	\Vol\big(Z(T)\big)=\int_{Z(T)}\omega^{n-\rk(T)}.
%	\]
%	By Wirtinger's Theorem, this differs of a factor~$1/(n-r)!$ from the volume of~$Z(T)$ computed with respect to the Kähler metric associated to~$\omega$.
%	By Lemma~\ref{lemma;extfromL2poincdual} and the~$L^2$-version of Stokes Theorem~\cite{gaffney}, we may rewrite~$\Vol\big(Z(T)\big)$ in terms of the Hermitian inner product~$(\cdot{,}\cdot)$ defined in~\eqref{eq;hermformforus} as
%	\bes
%	\Vol\big(Z(T)\big)=\int_X[Z(T)]\wedge\omega^{n-\rk(T)}=([Z(T)],\omega).
%	\ees
%	Recall that the volume of a special cycle~$Z(T,\discel)$, where~$\discel\in D_\lattice^\gengenus$ and~$T\in q(\discel)+\halfint{\gengenus}$ with~$T\ge 0$, is
%	\[
%	\Vol\big(Z(T,\discel)\big)=\int_{Z(T,\discel)}\omega^{n-\rk(T)}.
%	\]
%	By Wirtinger's Theorem, this differs of a factor~$1/(n-r)!$ from the volume of~$Z(T,\discel)$ computed with respect to the Kähler metric associated to~$\omega$.
%	As illustrated in Section~\ref{sec;volofspeccyinKM}, we may rewrite~$\Vol\big(Z(T,\discel)\big)$ in terms of the Hermitian inner product~$(\cdot{,}\cdot)$ defined in~\eqref{eq;hermformforus} as
%	\bes
%	\Vol\big(Z(T,\discel)\big)=\int_X[Z(T,\discel)]\wedge\omega^{n-\rk(T)}=([Z(T,\discel)],\omega).
%	\ees
%	
%	In this section we restrict our attention to the case of codimension~$1$.
%	We denote by~$Z(m)$ the special divisor of index~$m>0$, and define~$Z(0)=-\omega$.
	Let~$n\geq 3$.
	By~\eqref{eq;3isoofcoho}, as well as~\cite[Korollar~$4.2$]{weissauer-88} for the special case of~$n=3$, we know that~$\sqH^{1,1}(X)\cong H^{1,1}(X)$.
	By Corollary~\ref{cor;blmminnonadcase} and~\cite{HeHoffman}, the space~$\sqH^{1,1}(X)$ is generated by the cohomology classes of special divisors.
	The class~$[\KMtheta{1}]$ is the generating series of cohomology classes of such divisors, more precisely
	\bes
	[\KMtheta{1}]=-[\omega]\mathfrak{e}_0 + \sum_{\discel\in D_\lattice}\sum_{\substack{m\in q(\discel)+\ZZ \\ m>0}} [Z(m,\discel)]q^m \mathfrak{e}_\discel.
	\ees
	
%	We want to compute the Lefschetz decomposition of~$[\KMtheta{1}]$.
%	To do so, we firstly compute its projection to the line~$\CC[\omega]$ generated by the Kähler class.
%	This is the generating series of the volumes of special divisors, in fact
%	\be\label{eq;prodKMtheta1omega}
%	\big( [\KMtheta{1}] , [\omega] \big)
%	= - \Vol (X) + \sum_{m>0}\Vol \big(Z(m)\big) q^m.
%	\ee
%	Such a series is strictly related with the normalized Eisenstein series of weight~$k=1+n/2$, as described in~\cite[Proposition~4.8]{bruinierkuehn} and~\cite{kudla-integrals}.
%	\begin{prop}[Bruinier--Kühn, Kudla]\label{prop;brküku}
%	Let~$X$ be an orthogonal Shimura variety of dimension~$n$ arising from a unimodular lattice, and let~$k=1+n/2$.
%	The Fourier expansion of the normalized weight~$k$ Eisenstein series~$E^k_1$ of genus~$1$ may be rewritten as
%	\bes
%	E^k_1(\tau) = 1 - \sum_{m>0} \frac{\Vol\big( Z(m)\big)}{\Vol(X)} q^m.
%	\ees
%	\end{prop}
%	By Proposition~\ref{prop;brküku} 
%	By Theorem~\ref{thm;fromKudlagenservol} we may rewrite~\eqref{eq;prodKMtheta1omega} as
%	\be\label{eq;prodKMtheta1omegawithEk}
%	\big( [\KMtheta{1}] , [\omega] \big)
%	= - \Vol(X)\cdot E^k_1,
%	\ee
	The \emph{Lefschetz decomposition} of~$[\KMtheta{1}]$ is a trivial consequence of Remark~\ref{rem;projthetafcohoaftervol}, since~\eqref{eq;gengen2inprodkmtheta2} implies for~$g=1$ that
	\be\label{eq;lefdecgen1coho}
	[\KMtheta{1}]
	=
	\Big(\underbrace{- E^k_{1,\lattice} \cdot [\omega]}_{\text{in }M^k_{1,\lattice}\otimes{}\CC[\omega]}\Big)
	+
	\Big(\underbrace{[\KMtheta{1}] + E^k_{1,\lattice} \cdot [\omega]}_{\text{in }M^k_{1,\lattice}\otimes \sqHprim^{1,1}(X)}\Big).
	\ee
%	Note that such a decomposition holds also for Siegel threefolds, namely for~$n=3$.
%	Here~$E^k_{1,\lattice}$ is the normalized weight~$k$ elliptic Eisenstein series with respect to the Weil representation~$\weil{\lattice}{1}$.
	
	Since the Fourier coefficient of index~$(0,\discel)$ of~$[\KMtheta{1}] + E^k_{1,\lattice} \cdot [\omega]$ is trivial for every~$\discel\in D_\lattice$, the primitive part of~$[\KMtheta{1}]$ lies in~$S^k_{1,\lattice}\otimes \sqHprim^{1,1}(X)$, i.e.~it behaves as a cusp form with respect to the variable~$\tau\in\HH$.
	The part in~$\CC[\omega]$ of~$[\KMtheta{1}]$ lies in~$\CC E^k_{1,\lattice} \otimes \CC[\omega]$, namely it behaves as an Eisenstein series with respect to the variable~$\tau$.
	This implies that the Lefschetz decomposition of~$[\KMtheta{1}]$ with respect to the orthogonal variable is also the modular decomposition of~$[\KMtheta{1}]$ with respect to the symplectic variable.
%	Recall that for every positive~$m$, the~$m$-th Fourier coefficient of~$E^k_1$ is $2 \sigma_{k-1}(m)/\zeta(1-k)$.
%	Denoting by~$[Z(m)]_0$ the primitive part of~$[Z(m)]$, we deduce from~\eqref{eq;lefdecgen1coho} that
%	\be\label{eq;lefdecofspecdiv}
%	[Z(m)]_0 = [Z(m)] + \frac{2 \sigma_{k-1}(m)}{\zeta(1-k)} [\omega],\qquad\text{for every~$m>0$.}
%	\ee
%	
%	\subsection{Case of genus~$2$}
%	We know that the modular decomposition of~$[\KMtheta{2}]$ is
%	\[
%	[\KMtheta{2}]
%	=
%	E^k_2\cdot[\omega]^2
%	-
%	E^k_{2,1}\big(
%	[\KMtheta{1}]_0\wedge[\omega]
%	\big)
%	+
%	\xi,
%	\]
%	where~$\xi
%	\coloneqq
%	[\KMtheta{2}]
%	-E^k_2\cdot[\omega]^2
%	+
%	E^k_{2,1}\big(
%	[\KMtheta{1}]_0\wedge[\omega]
%	\big)$.
%	
%	We want to show that the previous decomposition is the Lefschetz decomposition of~$[\KMtheta{2}]$.
%	To do so, it is enough to show that~$\xi=[\KMtheta{2}]_0$, namely~$\xi$ is primitive, i.e.\ orthogonal to~$\lefschetz \sqH^{1,1}(X)$ with respect to~$(\cdot{,}\cdot)$.
%	
%	Since~$\xi$ is a cusp form, it is enough to show that the Kudla--Millson lift
%	\[
%	\KMlift{2}\colon S^k_2\longrightarrow \sqH^{2,2}(X)
%	\qquad
%	f\longmapsto \pet{f}{[\KMtheta{2}]}
%	\]
%	has image contained in~$\sqHprim^{2,2}(X)$.
%	In fact, since~$S^k_2$ is non-trivial and generated by Poincaré series\footnote{To remember in generalization: Take hypothesis on convergence of Poincaré series.}, we deduce that every Fourier coefficient of~$\xi$ is contained in~$\sqHprim^{2,2}(X)$.
%	
%	Saying that the image of~$\KMlift{2}$ is contained in~$\sqHprim^{2,2}(X)$ is equivalent of saying that such an image is orthogonal to~$[\omega]^2$ and~$\lefschetz[Z(m)]$ for every~$m>0$.
%	In fact, the space~$\sqH^{1,1}(X)$ is generated by the cohomology classes of special cycles\footnote{
%	Here we need to assume something on~$g$ when we generalize!
%	}
%	by~\cite{noetherlefconj}.
%	
%	Let~$f\in S^k_2$.
%	We now show that
%	\ba\label{eq;ifthenlef}
%	\text{if}
%	\quad
%	\int_X\KMlift{2}(f)\wedge[\KMtheta{1}]\wedge[\omega]^{n-3}=0,
%	\quad
%	\text{then}
%	\quad
%	\big(\KMlift{2}(f),\lefschetz^{2-\rk m}[Z(m)]\big)=0
%	\quad\text{for all~$m\geq 0$}.
%	\ea
%	The vanishing on the left-hand side of~\eqref{eq;ifthenlef} implies that the modular form with Fourier expansion
%	\be\label{eq;FcoefinprodKMth2}
%	\sum_{m\ge 0} \Big(
%	\int_X \KMlift{2}(f)\wedge [Z(m)]\wedge[\omega]^{n-3}
%	\Big) q^m
%	\ee
%	\textcolor{red}{Here there is a mistake, there should another factor of~$[-\omega]$ if~$m=0$.}\\
%	vanishes, or equivalently that~$\int_X \KMlift{2}(f) \wedge[Z(m)]\wedge [\omega]^{n-3}=0$ for every~$m\ge0$.
%	By~\eqref{eq;formmatchonHggsq} we deduce that
%	\be\label{eq;nonscscr2fkm}
%	(\KMlift{2}(f),[\omega]^2)
%	=
%	\frac{(-1)^0 \cdot n!\cdot 2!}{(n-2)!}
%	\int_X\KMlift{2}(f)\wedge[\omega]^{n-2}
%	= 0,
%	\ee
%	where the last passage follows from the vanishing of the Fourier coefficient of index~$0$ of~\eqref{eq;FcoefinprodKMth2}.
%	Since~$\sqHprim^{1,1}(X)$ is generated by the primitive classes of special divisors by~\cite{noetherlefconj}, to conclude the proof of~\eqref{eq;ifthenlef} it is enough to show that~$\big(\KMlift{2}(f),\lefschetz[Z(m)]_0\big)=0$ for every~$m>0$.
%	By~\eqref{eq;formmatchonHggsq} we may compute\footnote{Note that we do not have to conjugate anything. In fact, coho classes of special cycles as well as the Kähler form are \emph{real} forms.}
%	\bas
%	\big(\KMlift{2}(f) , \lefschetz [Z(m)]_0\big)
%	=
%	\frac{(-1)^{3} \cdot n! \cdot 1 !}{(n-3)!}
%	\int_X \KMlift{2}(f) \wedge [Z(m)]_0\wedge[\omega]^{n-3},
%	\eas
%	from which we deduce that the vanishing of~$\big(\KMlift{2}(f),\lefschetz[Z(m)]_0\big)$ is equivalent to the vanishing of
%	\be\label{eq;wwwtpvaninJni}
%	\int_X \KMlift{2}(f) \wedge [Z(m)]_0\wedge[\omega]^{n-3}.
%	\ee
%	Let~$[Z(m)]_\omega$ be the orthogonal projection of~$[Z(m)]$ to~$\CC[\omega]$.
%	As proved in~\eqref{eq;nonscscr2fkm}, the integral~$\int_X \KMlift{2}(f) \wedge [Z(m)]_\omega \wedge [\omega]^{n-3}$ vanishes.
%	We may then sum such a quantity to~\eqref{eq;wwwtpvaninJni} and deduce that
%	\bas
%	\int_X &\KMlift{2}(f)\wedge[Z(m)]_0\wedge[\omega]^{n-3}
%	\\
%	&=
%	\int_X \KMlift{2}(f)\wedge[Z(m)]_0\wedge[\omega]^{n-3}
%	+
%	\int_X\KMlift{2}(f)\wedge[Z(m)]_\omega\wedge[\omega]^{n-3}
%	\\
%	&=
%	\int_X \KMlift{2}(f)\wedge\big(\underbrace{[Z(m)]_\omega + [Z(m)]_0}_{=[Z(m)]}\big)\wedge[\omega]^{n-3}
%	\\
%	&=
%	\int_X\KMlift{2}(f)\wedge[Z(m)]\wedge[\omega]^{n-3}.
%	\eas
%	This is (up to a sign) the~$m$-th Fourier coefficient of the left-hand side of~\eqref{eq;ifthenlef}, concluding the proof of such an implication.
%	
%	To show that~$\xi$ is primitive, we show the vanishing of the first integral appearing in~\eqref{eq;ifthenlef}.
%	We may rewrite it keeping track of the variables~$\tau_1\in\HH$ as
%	\bas
%	\int_X & \KMlift{2}(f)\wedge[\KMtheta{1}](\tau_1)\wedge[\omega]^{n-3}
%	\\
%	&=
%	\int_X \bigpet{f}{[\KMtheta{2}]} \wedge [\KMtheta{1}](\tau_1) \wedge[\omega]^{n-3}
%	\\
%	&=
%	\Bigpet{f}{\int_X[\KMtheta{1}](-\overline{\tau_1}) \wedge [\KMtheta{2}]\wedge[\omega]^{n-3}},
%	\eas
%	where in the last equality we used that~$[\KMtheta{1}]$ is a modular form with \emph{real} Fourier coefficients, and hence~$\overline{[\KMtheta{1}](\tau_1)} = [\KMtheta{1}](-\overline{\tau_1})$.
%	Since
%	\be\label{eq;doublKMtheta}
%	[\KMtheta{1}(\tau_1)]\wedge[\KMtheta{2}(\tau_2)]
%	=
%	\left[\KMtheta{3}\left(\begin{smallmatrix}
%	\tau_1 & \\
%	 & \tau_2
%	\end{smallmatrix}\right)\right]\qquad\text{for every~$\tau_1\in\HH$ and~$\tau_2\in\HH_2$,}
%	\ee
%	we may continue the computation above as
%	\bas
%	\int_X & \KMlift{2}(f)\wedge[\KMtheta{1}](\tau_1)\wedge[\omega]^{n-3}
%	\\
%	&=
%	\Bigpet{f}{\int_X[\KMtheta{3}]\big(\begin{smallmatrix}
%	-\overline{\tau_1} & \\
%	 & \ast
%	\end{smallmatrix}\big)\wedge[\omega]^{n-3}}
%	\\
%	&=
%	-\Vol(X)\Bigpet{f}{E^k_3\big(\begin{smallmatrix}
%	-\overline{\tau_1} & \\
%	 & \ast
%	\end{smallmatrix}\big)},
%	\eas 
%	where we denote by~$\ast$ the variable of integration in~$\HH_2$ for the Petersson inner product, and use\footnote{
%	Here we need to assume something on~$g$ when we generalize to general genus!
%	} Theorem~\ref{thm;fromKudlagenservol} to deduce the last equality.
%	Let~$\tau_2=x_2+iy_2\in\HH_2$, and let~$d\nu=d x_2 \, d y_2 / \det y_2^3$ be standard~$\Sp_4(\ZZ)$-invariant metric on~$\HH_2$.
%	By the doubling method~\cite[Satz~$13$, $(1)$]{böcherer-doubling} we deduce that
%	\bas
%	\Bigpet{f}{E^k_3\big(\begin{smallmatrix}
%	-\overline{\tau_1} & \\
%	 & \ast
%	\end{smallmatrix}\big)}
%	=
%	\int_{\Sp_4(\ZZ)\backslash\HH_2} \det y_2^k \cdot f(\tau_2)\cdot \overline{E^k_2\big(\begin{smallmatrix}
%	-\overline{\tau_1} & \\
%	 & \tau_2
%	\end{smallmatrix}\big)}\, d\nu
%	= 0.
%	\eas
%	This concludes the proof of the primitivity of~$\xi$.
%	
%	
	
	\subsection{The case of general genus}\label{sec;cogengenus}
	We are now ready to prove the main result of this article.
%	This is done by induction on the genus~$g$.
	Recall from~\eqref{eq;defofKlinblamap} that the map~$\vklin{k}{\genus}{\genus-1}{\lattice}$ is the inverse of the restriction of the Siegel operator~$\Phi_0$ to~$\vklisub{k}{\gengenus}{0}{\lattice}\oplus\cdots\oplus \vklisub{k}{\gengenus}{\genus-1}{\lattice}$.
%	and that~$\KMtheta{0}=1$ and~$\sqHprim^0(X)= \sqH^0(X)$.
%	We refer to Section~\ref{sec;mod_dec} for the notation regarding the modular decomposition of the space~$\almcusp{k}{\genus}{\lattice}$ of almost cusp forms.
	\begin{thm}\label{thm;lefschetzdecgengenus}
	Let~$\gengenus<(n+2)/4
	%\min\big\{(\dim X +1)/2,\boundgenlef\big\}
	$ be a positive integer.
	Then~$[\KMtheta{\gengenus}]$ is an almost cusp form.
	Its primitive part is
	\be\label{eq;primpartgengenus}
	[\KMtheta{\gengenus}]_0 = [\KMtheta{\gengenus}]
	+
	\lefschetz\circ\vklin{k}{\genus}{\genus-1}{\lattice}\big([\KMtheta{\gengenus-1}]\big),
	\ee
	and lies in~$S^k_{\gengenus,\lattice}\otimes \sqHprim^{\gengenus,\gengenus}(X)$.
	In particular, it is a Siegel cusp form with respect to the symplectic variable.
	The \emph{Lefschetz decomposition} of~$[\KMtheta{\gengenus}]$ equals the \emph{modular decomposition} of~$[\KMtheta{\gengenus}]$, and is given by
	\ba\label{eq;lefschetzdecgengenus}
	\,[\KMtheta{\gengenus}]
	=
	\sum_{r=0}^{\gengenus} (-\lefschetz)^{\gengenus-r}\circ
	\klinser{k}{\gengenus}{r}{\lattice} \big([\KMtheta{r}]_0\big).
	\ea
	Here the summand of index~$r$
%	~$E^k_{\gengenus,j,\lattice}\circ\lefschetz^{\gengenus-r}\big([\KMtheta{r}]_0\big)$
	lies in~$\vklisub{k}{\gengenus}{r}{\lattice} \otimes \lefschetz^{\gengenus-r} \sqHprim^{r,r}(X)$, in particular it is a Klingen Eisenstein series arising from a genus~$r$ Siegel cusp form.
	
	All previous assertions hold also if we replace~$L^2$-classes by de Rham cohomology classes.
	\end{thm}
%	\begin{rem}
%	It is natural to ask whether the modular decomposition of~$[\KMtheta{\genus}]$ provided in Theorem~\ref{thm;lefschetzdecgengenus} may be extended to larger genera.
%	We remark that the convergence of the Klingen Eisenstein series appearing in~\eqref{eq;lefschetzdecgengenus} is ensured only under the assumption that~$k=1+n/2>2\genus$ by Proposition~\ref{prop;convpropKEseries}.
%	This is equivalent to~$\gengenus<(n+2)/4$, which is exactly the bound assumed in our main result.
%%	For this reason, we may say that the decomposition in Theorem~\ref{thm;lefschetzdecgengenus} is \emph{maximal}.
%%%	Further investigations on similar decompositions in larger genus should deal with the theory of Siegel modular forms of small weight.
%	\end{rem}
%
%	\begin{rem}
%	As recalled in Section~\ref{sec;outline}, the cohomology group~$H^{\gengenus,\gengenus}$ admits a Lefschetz decomposition if~$\gengenus<\frac{\dim X-1}{2}$.
%	In Theorem~\ref{thm;lefschetzdecgengenus} we request the stricter condition that~${\gengenus<\frac{\dim X + 1}{4}}$.
%	This comes from the fact that in the following proof of Theorem~\ref{thm;lefschetzdecgengenus} we show~\eqref{eq;primpartgengenus} by using the Lefschetz decomposition of~$H^{2\gengenus-1,2\gengenus-1}(X)$, hence we should also assume that
%	\[
%	2\gengenus-1 < \frac{\dim X-1}{2},\qquad\text{or equivalently}\qquad \gengenus<\frac{\dim X + 1}{4}.
%	\]
%	Note that if the integer~$g$ is such that~$1<g<\frac{\dim X + 1}{4}$, then~$g<\frac{\dim X - 1}{2}$.
%%	Note that~$\frac{\dim X-1}{2} < \frac{\dim X + 1}{4}$ only if~$\dim X<3$, but in this case both fractions are less than~$1$.
%%	These do not create any conflict with Theorem~\ref{thm;lefschetzdecgengenus}, since in the latter we assume~$g$ to be positive.
%	\\
%	\textcolor{blue}
%	{
%	\textbf{Note:} The trick applied in Theorem~\ref{thm;lefschetzdecgengenus} under the assumption that~$\gengenus<\frac{\dim X + 1}{4}$ makes things much easier.
%	If we want to be more general and assume only that~$\gengenus<\frac{\dim X-1}{2}$, then I think one could prove the result, say for~$\gengenus=2$, by showing that for every positive definite~$2\times 2$ matrix~$T$, we have
%	\bas
%	\big( [Z(T)],[Z(\tilde{m})]_0\wedge[\omega] \big)
%	\,\,=\,\,
%	\text{$T$-th Fourier coefficient of~$E^k_{2,1}\Big(
%	\sum_{m>0}\big([Z(m)]_0,[Z(\tilde{m})]_0\big)
%	q^m\Big)$}
%	\eas
%	for every~$\tilde{m}\in\ZZ_{>0}$.
%	But I do not see how to show that.
%	}\\
%	\textcolor{red}{\textbf{Very important remark: }
%	I think that if we state the theorem on~$L^2$-cohomology, then we can drop that stronger condition!\\
%	\textbf{Natural question: }Is the differential form~$\KMtheta{\gengenus}$ always \emph{square-integrable}? If yes, and if the main result of~\cite{noetherlefconj} works also in~$L^2$-cohomology, then it could be even possible to generalize the decomposition in higher genus...}
%	\end{rem}
	\begin{proof}[Proof of Theorem~\ref{thm;lefschetzdecgengenus}]
	We prove the result by induction on the genus~$\gengenus$.
	The case of genus~$1$ has already been illustrated in Section~\ref{sec;gen1}.
	Suppose that~$\genus>1$ and that Theorem~\ref{thm;lefschetzdecgengenus} is valid for every genus less than~$g$.
	We prove that it is valid also for~$\gengenus$.
%	Since~$1<\genus<(n+2)/4$, we may assume that~$n>6$.
	
	
%	To simplify the illustration, we define~$\KMtheta{0}$ to be the constant function~$1$ on~$X$ and~$\sqHprim^0(X)\coloneqq\sqH^0(X)$.
	To prove that the right-hand side of~\eqref{eq;primpartgengenus} is a cusp form, it is enough to show that it lies in the kernel of~$\Phi_\disceltwo\colon M^k_{g,\lattice}\otimes \sqH^{g,g}(X)\to M^k_{g-1,\lattice}\otimes \sqH^{g,g}(X)$ for every~$\disceltwo\in D_\lattice$; see Lemma~\ref{lemma;FexpofimPhiop}.
	This follows easily from Lemma~\ref{lemma;Siegopkmtheta2} and Theorem~\ref{thm;Phi0onkling}.
	
	It is easy to check that the right-hand side of~\eqref{eq;lefschetzdecgengenus} equals~$[\KMtheta{\genus}]$ under the assumption of~\eqref{eq;primpartgengenus}.
%	
%	\textcolor{blue}{\textbf{Check of the previous sentence:}
%	We prove that the right-hand side of~\eqref{eq;lefschetzdecgengenus} equals the class~$[\KMtheta{g}]$.
%	We substitute~$[\KMtheta{\gengenus}]_0$ given by~\eqref{eq;primpartgengenus} in~\eqref{eq;lefschetzdecgengenus}, obtaining
%	\ba\label{eq;inproofgenlefscdec}
%	[\KMtheta{\gengenus}]
%	=
%%	(-1)^\gengenus E^k_\gengenus\cdot [\omega]^\gengenus
%%	+
%	\sum_{r=0}^{\gengenus-1} (-1)^{\gengenus-r} 
%	\lefschetz^{\gengenus-r}\circ \klinser{k}{\gengenus}{r}{\lattice}
%	\big([\KMtheta{r}]_0\big)
%%	\\
%%	&\quad
%	+
%	[\KMtheta{g}] +
%	\lefschetz \circ
%	\vklin{k}{\genus}{\genus-1}{\lattice}
%	\big([\KMtheta{\gengenus-1}]\big).
%	\ea
%	This equality holds if and only if
%	\ba\label{eq;inproofgenlefscdec2}
%	\lefschetz \circ
%	\vklin{k}{\genus}{\genus-1}{\lattice}
%	\big([\KMtheta{\gengenus-1}]\big)
%	=
%%	(-1)^{\gengenus-1}E^k_\gengenus\cdot[\omega]^\gengenus
%%	+
%	\sum_{r=0}^{\gengenus-1} (-1)^{\gengenus-1-r}
%	\lefschetz^{\gengenus-r}\circ \klinser{k}{\gengenus}{r}{\lattice}
%	\big([\KMtheta{r}]_0\big).
%	\ea
%	We rewrite the right-hand side of~\eqref{eq;inproofgenlefscdec2} as
%	\bas
%%	(-1)^{\gengenus-1}\klin^k_{\gengenus,\gengenus-1}(E^k_{\gengenus-1})\cdot[\omega]^\gengenus
%%	+
%	\sum_{r=0}^{\gengenus-1}& (-1)^{\gengenus-1-r}
%	\vklin{k}{\genus}{\genus-1}{\lattice}
%	\big(\lefschetz^{\gengenus-r}\circ\klinser{k}{\gengenus-1}{r}{\lattice}([\KMtheta{r}]_0)\big)
%	\\
%	&=
%	\lefschetz\circ
%	\vklin{k}{\genus}{\genus-1}{\lattice}
%	\Big(
%%	(-1)^{\gengenus-1}E^k_{\gengenus-1}\cdot[\omega]^\gengenus
%%	+
%	\sum_{r=0}^{\gengenus-1}(-1)^{\gengenus-1-r}
%	\lefschetz^{\gengenus-1-r} \circ
%	\klinser{k}{\gengenus-1}{r}{\lattice}
%	\big([\KMtheta{r}]_0\big)
%	\Big).
%	\eas
%	To prove~\eqref{eq;inproofgenlefscdec2} it is then enough to show that
%	\[
%	[\KMtheta{\gengenus-1}]=
%%	(-1)^{\gengenus-1}E^k_{\gengenus-1}\cdot[\omega]^{\gengenus-1}
%%	+
%	\sum_{r=0}^{\gengenus-1}(-1)^{\gengenus-1-r}
%	\lefschetz^{\gengenus-1-r}\circ
%	\klinser{k}{\gengenus-1}{r}{\lattice}
%	\big([\KMtheta{r}]_0\big).
%	\]
%	This is assumed to hold by induction.}
	
	The inductive hypothesis that~$[\KMtheta{r}]_0$ lies in~$S^k_{r,L}\otimes \sqHprim^{r,r}(X)$ for every~$r=0,\dots,\gengenus-1$ implies that~$\lefschetz^{\genus-r}\circ\klinser{k}{\gengenus}{r}{\lattice}\big([\KMtheta{r}]_0\big)$ lies in~$\vklisub{k}{\gengenus}{r}{\lattice}\otimes \lefschetz^{\gengenus-r} \sqHprim^{r,r}(X)$.
	It remains to show~\eqref{eq;primpartgengenus}, assuming by induction that~\eqref{eq;primpartgengenus} is satisfied for every genus less than~$g$.
	Since the right-hand side of~\eqref{eq;primpartgengenus} is the cuspidal part of~$[\KMtheta{\gengenus}]$, it is enough to show that the Kudla--Millson lift
	\[
	\KMlift{\genus}\colon S^k_{\genus,\lattice}\longrightarrow \sqH^{\genus,\genus}(X),\qquad f\longmapsto\pet{f}{[\KMtheta{\genus}]}
	\]
	has image contained in~$\sqHprim^{\genus,\genus}(X)$.
	In fact, the space~$S^k_{\genus,\lattice}$ is generated by Poincaré series by Lemma~\ref{eq;spanpoincseries}, hence we may deduce from Lemma~\ref{lemma;petSiegpoinandcusp} that every Fourier coefficient of the right-hand side of~\eqref{eq;primpartgengenus} is contained in~$\sqHprim^{\genus,\genus}(X)$.
	The primitivity of the image of~$\KMlift{\genus}$ follows from Theorem~\ref{thm;imageofKMlift}.
%	
%	THE TEXT BELOW HAS BEEN TRANSLATED INTO THE PROOF OF THE THEOREM THAT STATES THE PRIMITIVITY OF THE IMAGE OF THE KM LIFT
%
%	Let~$K$ be the open compact subgroup of~$G(\Ab_f)$ as in Section~\ref{sec;adelicspcy}.
%	By Theorem~\ref{thm;noetherlef}, the space~$\sqH^{r,r}(X)$ is generated by the cohomology classes of adelic special cycles for every~$r\leq\genus-1$, since~$g-1<(n+1)/3$ under our assumptions.
%	Recall that for every~$\varphi\in\mathcal{S}\big(V^{r}(\Ab_f)\big)$, there exists an even lattice~$M\subset V$ such that~$\varphi\in\mathcal{S}_{M^{r}}$; see e.g.~\cite{kudla-speccyc}.
%	Hence, for every~$r\leq\genus-1$ there exists a sufficiently small even lattice~$M_r\subset V$ such that~$\sqH^{r,r}(X)$ is generated by the adelic special cycles of the form
%	\[
%	[Z(T,\varphi,K)]\wedge[\omega^{r-\rk T}],
%	\qquad
%	\text{where~$T\in\Sym_r(\QQ)$ and~$\varphi\in(\mathcal{S}_{(M_r)^r})^K$.}
%	\]
%	Without loss of generality, we may assume that the~$M_r$ are all equal to the same even lattice~$M\subset V$, and also that~$M\subset L$.
%	
%	Saying that the image of~$\KMlift{\genus}$ is contained in~$\sqHprim^{\genus,\genus}(X)=\big(\lefschetz \sqH^{\genus-1,\genus-1}(X)\big)^\perp$ is equivalent of saying that such an image is orthogonal to~$\lefschetz^{\genus-\rk T}\big([Z(T,\varphi,K)]\big)$ for every~$T\in \Sym_{\genus-1}(\QQ)$ and~$\varphi\in (\mathcal{S}_{M^{\genus-1}})^K$.
%	
%	Let~$X_M=\disk{}^+(M)\backslash\domain$, and let~$\pi\colon X_M\to X$ be the quotient map.
%	We denote by~$\KMtheta{\genus-1}^M$ the Kudla--Millson theta function arising from the lattice~$M$.
%	Let~$f\in S^k_{\genus,\lattice}$.
%	It is easy to check by using~\eqref{eq;formmatchonHggsq} and induction on~$\genus$ that if
%	\be\label{eq;wwwtpingengenuspldu}
%	\int_{X_M} \pi^*\big(\KMlift{\genus}(f)\big)\wedge \KMtheta{\genus-1}^M \wedge \omega^{n-(2\genus-1)}
%	=0,
%	\ee
%	then
%	\be\label{eq;wwwtpingengenusplduimpl}
%	\big(
%	\KMlift{\genus}(f)
%	,
%	\lefschetz^{\genus-\rk T}[Z(T,\varphi,K)]
%	\big)=0
%	\qquad
%	\text{for every~$T$ and~$\varphi$ as above}.
%	\ee
%%	\textcolor{blue}{\textbf{Remark on integrals of pullbacks:}
%%	If~$[\eta]\in\sqH^{\genus,\genus}(X)$ with~$X=\Gamma\backslash\domain$ as above, then we may regard~$\eta$ as a~$\Gamma$-invariant~$(\genus,\genus)$-form on~$\domain$.
%%	Note that~$\eta$ is also~$\Gamma_M\coloneqq\disk{}^+(M)$-invariant, since~$\Gamma_M\subseteq\Gamma$.
%%	The same applies for the Kähler form~$\omega$, since it is~$G(\RR)$-invariant.
%%	We deduce that
%%	\be\label{eq;pullbackofintrc}
%%	\int_{X_M}\pi^*(\eta)
%%	\wedge
%%	\omega^{n-\genus}
%%	=
%%	\#\big(\Gamma_M\backslash\Gamma\big)
%%	\int_{X}\eta\wedge\omega^{n-\genus}.
%%	\ee
%%	This can be proved by unfolding as follows.
%%	We denote by~$\mathcal{F}_\Gamma$ and~$\mathcal{F}_{\Gamma_M}$ respectively a fundamental domain for the action of~$\Gamma$ and of~$\Gamma_M$ on~$\domain$.
%%	Write~$\mathcal{F}_{\Gamma_M}=\bigcup_\gamma \gamma(\mathcal{F}_\Gamma)$ for finitely many~$\gamma\in\Gamma$.
%%	The number of such~$\gamma$'s is~$\#(\Gamma_M\backslash\Gamma)$.
%%	Then
%%	\bas
%%	\int_{X_M}\pi^*(\eta)\wedge\omega^{n-\genus}
%%	&=
%%	\int_{\mathcal{F}_{\Gamma_M}}
%%	\eta\wedge\omega^{n-\genus}
%%	\\
%%	&=
%%	\int_{\bigcup_\gamma \gamma(\mathcal{F}_\Gamma)}
%%	\eta\wedge\omega^{n-\genus}
%%	\\
%%	&=
%%	\sum_\gamma
%%	\int_{\gamma(\mathcal{F}_\Gamma)}
%%	\eta\wedge\omega^{n-\genus}
%%	\\
%%	&=
%%	\sum_\gamma
%%	\int_{\mathcal{F}_\Gamma}
%%	\gamma^*(\eta\wedge\omega^{n-\genus})
%%	\\
%%	&=
%%	\sum_\gamma
%%	\int_{\mathcal{F}_\Gamma}
%%	\eta\wedge\omega^{n-\genus}
%%	\\
%%	&=
%%	\#\big(\Gamma_M\backslash\Gamma\big)
%%	\int_{\mathcal{F}_\Gamma}
%%	\eta\wedge\omega^{n-\genus},
%%	\eas
%%	where in the first to the last passage we used that~$\gamma^*(\eta\wedge\omega^{n-\genus})=\eta\wedge\omega^{n-\genus}$ for all~$\gamma\in\Gamma$, i.e.~$\eta\wedge\omega^{n-\genus}$ is~$\Gamma$-invariant. 
%%	}
%%	\\
%%	
%%	\textcolor{blue}{\textbf{Here the check by induction on~$g$ that~\eqref{eq;wwwtpingengenuspldu} implies~\eqref{eq;wwwtpingengenusplduimpl}.}\\
%%	The vanishing~\eqref{eq;wwwtpingengenuspldu} is the vanishing of a modular form with Fourier expansion
%%	\be\label{eq;Fexpinproofblu}
%%	\sum_{\discel\in D_M^{\genus-1}}
%%	\sum_{\substack{T\in q(\discel)+\halfint{\genus-1}\\ T\ge 0}}
%%	\Big(
%%	\int_{X_M} \pi^*\big(\KMlift{\genus}(f)\big)
%%	\wedge [Z(T,\discel)]
%%	\wedge
%%	[-\omega]^{\genus-1-\rk T}
%%	\wedge [\omega]^{n-(2\genus-1)}
%%	\Big) q^T \mathfrak{e}_\discel,
%%	\ee
%%	where here the cohomology classes are taken in~$\sqH^\bullet(X_M)$.
%%	We prove~\eqref{eq;wwwtpingengenusplduimpl} by induction on~$\genus$.
%%	\\
%%	\textbf{Case $\boldsymbol{g=1}$:}\\
%%	The theta function~$\KMtheta{0}^M=1$ is a non-zero constant function.
%%	By~\eqref{eq;formmatchonHggsq} and~\eqref{eq;pullbackofintrc} (applied with~$\eta=\KMlift{1}(f)$) the left-hand side of~\eqref{eq;wwwtpingengenuspldu} equals the left-hand side of~\eqref{eq;wwwtpingengenusplduimpl} up to a non-zero constant, hence the assertion follows immediately.
%%	\\\textbf{Case $\boldsymbol{\genus > 1}$:}\\
%%	We work by induction on~$\rk T$.
%%	If~$\rk T =0$, namely~$T=0$, then by~\eqref{eq;formmatchonHggsq} and~\eqref{eq;pullbackofintrc} (applied with~$\eta=\KMlift{\genus}(f)$) we may rewrite the left-hand side of~\eqref{eq;wwwtpingengenusplduimpl} as
%%	\bas
%%	\big(\#(\Gamma_M\backslash\Gamma)\big)^{-1}\big(\pi^*\big(\KMlift{\genus}(f)\big),[\omega]^\genus\big)
%%	=
%%	\frac{(-1)^0\cdot n!\cdot \genus!}{(n-\genus)! \cdot \#(\Gamma_M\backslash\Gamma)}
%%	\int_{X_M} 
%%	\pi^*\big(\KMlift{\genus}(f)\big)\wedge\omega^{n-\genus}.
%%	\eas
%%	This is, up to a non-zero factor, the constant term of the Fourier expansion~\eqref{eq;Fexpinproofblu}, which vanishes by~\eqref{eq;wwwtpingengenuspldu}.}
%%	
%%	\textcolor{blue}{Suppose now by induction that~$\Big(\KMlift{\genus}(f),\lefschetz^{\genus-\rk T}\big([Z(T,\varphi,K)]\big)\Big)=0$ for every~$\varphi\in\big(\mathcal{S}_{M^{\genus-1}}\big)^K$ and~$T\in \Sym_{\genus-1}(\QQ)$ of rank less than some positive~$r$.
%%	This implies that~$\KMlift{\genus}(f)$ is orthogonal to~$\lefschetz^{\genus-r+1}\sqH^{r-1,r-1}(X)$, since~$\sqH^{r-1,r-1}(X)$ is generated by classes of adelic special cycles arising from~$(\mathcal{S}_{M^{r-1}})^K$.
%%	\\
%%	Let~$T_r\in \Sym_{\genus-1}(\QQ)$ be a matrix of rank~$r$, and let~$\Pi_j\colon \sqH^{\genus,\genus}(X)\to\lefschetz^{j}\sqHprim^{\genus-j,\genus-j}(X)$ be the orthogonal projection to the~$j$-th term of the Lefschetz decomposition of~$\sqH^{\genus,\genus}(X)$.
%%	We may compute
%%	\begin{align*}
%%	\Big(\KMlift{\genus}(f),\lefschetz^{\genus-r}\big([Z(T_r,\varphi,K)]\big)\Big)
%%	&=
%%	\sum_{j=\genus-r}^\genus
%%	\Big(\KMlift{\genus}(f),\Pi_j\circ
%%	\lefschetz^{\genus-r}\big([Z(T_r,\varphi,K)]\big)\Big)
%%	\\
%%	&=
%%	\Big(\KMlift{\genus}(f),\Pi_{\genus-r}\circ
%%	\lefschetz^{\genus-r}\big([Z(T_r,\varphi,K)]\big)\Big)
%%	\\
%%	&=
%%	\Big(\KMlift{\genus}(f),\lefschetz^{\genus-r}\big([Z(T_r,\varphi,K)]_0\big)
%%	\Big)
%%	\\
%%	&=C\cdot\int_X \KMlift{\genus}(f)\wedge [Z(T_r,\varphi,K)]_0 \wedge[\omega]^{n-\genus-r},
%%	\end{align*}
%%	where in the second passage above we used the inductive hypothesis, and where~$C$ is a non-zero constant.
%%	By~\eqref{eq;pullbackofintrc} we then deduce that
%%	\ba\label{eq;inbluewpofKMZ}
%%	&\Big(\KMlift{\genus}(f),\lefschetz^{\genus-r}\big([Z(T_r,\varphi,K)]\big)\Big)
%%	\\
%%	&\quad=
%%	\frac{C}{\#(\Gamma_M\backslash\Gamma)}
%%	\int_X
%%	\pi^*\big(\KMlift{\genus}(f)\big)\wedge \pi^*\big([Z(T_r,\varphi,K)]_0\big) \wedge[\omega]^{n-\genus-r}.
%%	\ea
%%	With a similar procedure, we now prove that the vanishing of~\eqref{eq;Fexpinproofblu} implies the vanishing of the integral on the right-hand side of~\eqref{eq;inbluewpofKMZ}, and hence of~$\big(\KMlift{\genus}(f),\lefschetz^{\genus-r}\big([Z(T_r,\varphi,K)]\big)\big)$.
%%	The vanishing of~\eqref{eq;Fexpinproofblu} implies that
%%	\[
%%	\int_{X_M}\pi^*\big(\KMlift{\genus}(f)\big)
%%	\wedge
%%	[Z(T_r,\varphi,K)]
%%	\wedge
%%	\omega^{n-\genus-r}
%%	=
%%	0
%%	\]
%%	for every~$\varphi\in (\mathcal{S}_{M^{\genus-1}})^K$.
%%	Since~$\sqH^{r-1,r-1}(X_M)^{\Gamma_M\backslash\Gamma}=\sqH^{r-1,r-1}(X)$, we deduce that the former cohomology group is generated by the adelic special cycles~$[Z(T,\varphi,K)]\wedge[\omega^{r-1-\rk T}]$ with~$\varphi\in(\mathcal{S}_{M^{\genus-1}})^K$ and~$T$ of rank less than~$r$.
%%	With the same idea as above, using the inductive hypothesis and the projections~$\Pi_j$, one can deduce the vanishing of the right-hand side of~\eqref{eq;inbluewpofKMZ}.
%%	}
%	
%	To show that the right-hand side of~\eqref{eq;primpartgengenus} is primitive, it is then enough to prove~\eqref{eq;wwwtpingengenuspldu}.
%	By Lemma~\ref{lemma;genofBYandS} and Remark~\ref{rem;resandpullbKMthetafunct} we deduce that
%	\[
%	\pi^*\big(\KMlift{\genus}(f)\big)
%	=
%	\bigpet{f}{\pi^*([\KMtheta{\genus}])}
%	=
%	\bigpet{f}{\tr_{\genus,\lattice/M}([\KMtheta{\genus}^M])}
%	=
%	\bigpet{f_M}{[\KMtheta{\genus}^M]}.
%	\]
%	By the Fubini--Tonelli Theorem, we may then rewrite the integral in~\eqref{eq;wwwtpingengenuspldu} keeping track of the variable~$\tau_1\in\HH_{\genus-1}$ as
%	\begin{align*}
%	&\int_{X_M}
%	\pi^*\big(\KMlift{\genus}(f)\big)
%	\wedge
%	\KMtheta{\genus-1}^M(\tau_1)
%	\wedge
%	\omega^{n-(2\genus-1)}
%	\\
%	&\quad
%	=
%	\Bigpet{f_M}{\int_{X_M} \KMtheta{\genus-1}^M(-\overline{\tau_1}) \wedge \KMtheta{\genus}^M \wedge \omega^{n-(2\genus-1)}},
%	\end{align*}
%	where to deduce the equality above we used that~$\KMtheta{\genus-1}$ is a modular form with \emph{real} Fourier coefficients, and hence~$\overline{\KMtheta{\genus-1}(\tau_1)} = \KMtheta{\genus-1}(-\overline{\tau_1})$.
%	By~\eqref{eq;wedgeprodformforKMtheta} and Theorem~\ref{thm;fromKudlagenservol} we may continue the computation above as
%	\bas
%	\,&\int_{X_M} \pi^*\big(\KMlift{\genus}(f)\big)\wedge
%	\KMtheta{\genus-1}^M(\tau_1)\wedge
%	\omega^{n-(2\genus-1)}
%	\\
%	&
%	\quad
%	=
%	\Bigpet{f_M}{\int_{X_M} \KMtheta{2\genus-1}^M\big(\begin{smallmatrix}
%	-\overline{\tau_1} & \\
%	 & \ast
%	\end{smallmatrix}\big) \wedge \omega^{n-(2\genus-1)}}
%	\\
%	&\quad
%	=
%	(-1)^{2\genus-1}
%	\Vol(X_M)
%	\Bigpet{f_M}{E^k_{2\genus-1,M}\big(\begin{smallmatrix}
%	-\overline{\tau_1} & \\
%	 & \ast
%	\end{smallmatrix}\big)},
%	\eas
%	where we denoted by~$\ast$ the variable of integration in~$\HH_{\genus}$ for the Petersson inner product.
%	Note that we are allowed to use Theorem~\ref{thm;fromKudlagenservol}, since~$2\genus-1<n-1\leq n+1-\Witt(M\otimes\QQ)$.
%	
%	By Corollary~\ref{cor;vanishingGmapvv} we deduce that
%	\[
%	\Bigpet{f_M}{E^k_{2\genus-1,M}\big(\begin{smallmatrix}
%	-\overline{\tau_1} & \\
%	 & \ast
%	\end{smallmatrix}\big)} = 0.
%	\]
%	This concludes the proof of~\eqref{eq;primpartgengenus}.	
%	
%	We now consider the case where~$\Gamma$ is a torsion-free subgroup of finite index in~$\Gamma_L\coloneqq\disk{}^+(\lattice)$.
%%	Here~$\Gamma_L$ may have torsion, hence~$X_L$ can have finite-quotient singularity.
%	Let
%	\[
%	\pi\colon X_\Gamma\coloneqq \Gamma\backslash\domain
%	\longrightarrow
%	X_L\coloneqq\Gamma_\lattice\backslash\domain
%	\]
%	be the standard projection map.
%	The pullback~$\pi^*\colon H^{2\genus}(X_\lattice)\to H^{2\genus}(X_\Gamma)\cong \sqH^{2\genus}(X_\Gamma)$ is injective, and identifies~$H^{2\genus}(X_\lattice)$ with the subspace~$\sqH^{2\genus}(X_\Gamma)^{\Gamma\backslash\Gamma_\lattice}$ of~$\Gamma\backslash\Gamma_\lattice$-invariant classes in~$\sqH^{2\genus}(X_\Gamma)$.
%	The cohomology class of the Kudla--Millson theta function attached to~$L$ induced on~$X_\Gamma$ is the pullback~$\pi^*\big([\KMtheta{\genus}]\big)\in \sqH^{2\genus}(X_\Gamma)^{\Gamma\backslash\Gamma_\lattice}$ of the one defined on~$X_\lattice$.
%	Since the Kähler form~$\omega$, and hence the Hermitian inner product on~$\sqH^{2\genus}(X_\Gamma)$, is~$\Gamma_\lattice$-invariant, we deduce that the Hodge and Lefschetz decompositions of~$\sqH^{2\genus}(X_\Gamma)$ preserve such invariance, in the sense that
%	\[
%	\sqH^{2\genus}(X_\Gamma)^{\Gamma\backslash\Gamma_\lattice}
%	=
%	\bigoplus_{p+q=2\genus} \sqH^{p,q}(X_\Gamma)^{\Gamma\backslash\Gamma_\lattice}
%	\qquad
%	\text{and}
%	\qquad
%	\sqH^{\genus,\genus}(X_\Gamma)^{\Gamma\backslash\Gamma_\lattice}
%	=
%	\bigoplus_{i\geq 0}
%	\lefschetzî\sqHprim^{r-2i}(X_\Gamma)^{\Gamma\backslash\Gamma_\lattice}.
%	\]
%%	\textcolor{blue}{\textbf{Regarding the Hodge decomposition.}
%%	Let~$\gamma\in\Gamma_L$. If~$\alpha$ is a~$(p,q)$ form on~$X_\Gamma$, then also its pullback~$\gamma^*\alpha$ is a~$(p,q)$-form on~$X_\Gamma$.
%%	The check of this can be done on a local charts of~$X_\Gamma$, using the fact that~$\gamma$ acts on~$X_L$ as a holomorphic map.
%%	Hence, if a differential form on~$X_\Gamma$ is~$\Gamma_L$-invariant, each term of its Hodge decomposition is~$\Gamma_L$-invariant.
%%	\\
%%	\textbf{Regarding the Hodge decomposition.}
%%	We prove the second equality displayed above by induction.
%%	If~$\genus=1$ and~$[\alpha]\in \sqH^{1,1}(X_\Gamma)$, then
%%	\[
%%	[\alpha]=A\cdot[\omega] + B\cdot[\alpha]_0,
%%%	\sqH^{1,1}(X_\Gamma)=\CC[\omega]\oplus\sqHprim^{1,1}(X_\Gamma),
%%	\qquad\text{where}\qquad
%%	[\alpha]_0\in\sqHprim^{1,1}(X_\Gamma) = \big(\CC[\omega]\big)^\perp.
%%	\]
%%	Since~$\omega$ is $\Gamma_L$-invariant, we deduce that if~$\gamma\in\Gamma_L$ and~$[\alpha]$ is~$\Gamma_L$-invariant, then
%%	\[
%%	A\cdot[\omega] + B\cdot[\alpha]_0
%%	=
%%	[\alpha]
%%	=
%%	\gamma^*[\alpha]
%%	=
%%	A\cdot\underbrace{\gamma^*[\omega]}_{=[\omega]} + {}B\cdot\gamma^*[\alpha]_0.
%%	\]
%%	This implies that also~$[\alpha]_0$ is~$\Gamma_L$-invariant.\\
%%	An induction argument using that
%%	\[
%%	\sqH^{\genus,\genus}(X_\Gamma)=\lefschetz\sqH^{\genus-1,\genus-1}(X_\Gamma)\oplus\sqHprim^{\genus,\genus}(X_\Gamma)
%%	\]
%%	and that the operator~$\lefschetz$ is~$\Gamma_L$-invariant concludes the proof.
%%	}
%%	\\
%	By Theorem~\ref{thm;noetherlef} the space~$\sqH^{\genus,\genus}(X_\Gamma)^{\Gamma\backslash\Gamma_\lattice}$ is generated by the pullback under~$\pi$ of the classes of adelic special cycles in~$X_\lattice$.
%%	\\\textcolor{blue}{\textbf{Here a check of the last assertion:}\\
%%	It is enough to prove it in the case where~$\Gamma=\Gamma_N\coloneqq\disk{}^+(L(N))$, for some~$N$.
%%	As above, we assume~$\Gamma_N$ to be torsion-free.
%%	\\
%%	Let~$X_N\coloneqq\Gamma_N\backslash\domain$.
%%	Denote by~$SC^\genus(X)$, resp.~$SC^\genus(X_N)$, the subspace of~$H^{2\genus}(X)$, resp.~$H^{2g}(X_N)$, generated by the adelic special cycles of~$X$, resp.\ of~$X_N$.
%%	Note that since~$X_N$ is smooth, we may consider~$SC^\genus(X_N)$ as a subspace of~$\sqH^{\genus,\genus}(X_N)$.
%%	By Theorem~\ref{thm;noetherlef} we deduce that
%%	\be\label{eq;inbluespaofspblmm}
%%	SC^\genus(X_N)=\sqH^{\genus,\genus}(X_N).
%%	\ee
%%	Note that the subspace of~$\sqH^{\genus,\genus}(X_N)$ generated by the pullbacks of the special cycles of~$X$ equals~$SC^\genus(X_N)^{\Gamma_N\backslash\Gamma_L}$; see~\cite[First displayed formula in the proof of Corollary~$3.8$]{noetherlefconj}.
%%	But from~\eqref{eq;inbluespaofspblmm} we see that
%%	\[
%%	SC^\genus(X_N)^{\Gamma_N\backslash\Gamma_L}
%%	=
%%	\sqH^{\genus,\genus}(X_N)^{\Gamma_N\backslash\Gamma_L}.
%%	\]
%%	Hence~$\sqH^{\genus,\genus}(X_N)^{\Gamma_N\backslash\Gamma_L}$ is generated by the pullback of the special cycles on~$X_L$.
%%	}\\
%	By replacing the cohomology groups of~$X_\Gamma$ by their~$\Gamma\backslash\Gamma_\lattice$-invariant subspaces, one can easily see that the same proof given above for the case of~$\Gamma=\Gamma_\lattice$ works also for torsion-free finite-index subgroups.
	\end{proof}
	
%	\subsection{Remarks on Bergeron--Li--Millson--Moeglin's Remark~2.4}\label{subsec;onrmblmm}
%	\textcolor{red}{\textbf{This subsection should disappear as soon as we agree it is correct.}}
%	
%	This section should be read in parallel with~\cite[Section~$2.6$, pp.~$513$-$514$]{noetherlefconj}.
%	We follow the notation therein, working with a rational quadratic space~$V$ of signature~$(p,q)\coloneqq(n,2)$.
%	\\
%	
%	Let~$K$ be an open compact subgroup of~$G(\Ab_f)$, where~$G=\SO(V)$, and denote by~$X_K$ the associated non-connected Shimura variety constructed from the usual (adelic) double quotient; see e.g.~\cite[p.~$509$]{noetherlefconj}.
%	One has a decomposition of~$X_K$ in connected components as
%	\[
%	X_K=\coprod_{g_j} \Gamma_{g_j}\backslash D.
%	\]
%	We denote by~$Y_K=\Gamma_K\backslash D$ the connected component arising in the case of~$g_j=1$, i.e.\ with~$\Gamma_K=\Gamma_1$.
%	Recall from~\cite[Section~$2.3$, p.~$501$]{noetherlefconj} that~$\Gamma_K$ is the image of~$\widetilde{G}(\QQ)_+\cap \widetilde{K}$ in~$\SO_0(n,2)$, where~$\widetilde{G}$ and~$\widetilde{K}$ are resp.\ the general spin group~$\GSpin(V)$ and the preimage of~$K$ in~$\widetilde{G}(\Ab_f)$, and where~$\widetilde{G}(\QQ)_+$ is the subgroup of~$\widetilde{G}(\QQ)$ of elements with totally positive spinor norm.
%	
%	
%	We now recall how to construct \emph{special cycles} on~$X_K$ and~$Y_K$, following the wording of~\cite[Section~$2.6$, p.~$513$]{noetherlefconj}.
%	Given~$\textbf{\textrm{x}}\in V^r$, we let~$U(\textbf{\textrm{x}})$ be the~$\QQ$-subspace of~$V$ spanned by the components of~$\textbf{\textrm{x}}$.
%	Let~$D_{\textbf{\textrm{x}}}$ be the subset of points in~$D$ orthogonal to~$U(\textbf{\textrm{x}})$.
%	\begin{defi}
%	The \emph{natural cycle}~$c\big(U(\textbf{\textrm{x}}),g_j,K\big)$ on the connected component~$\Gamma_{g_j}\backslash D$ of~$X_K$ is defined to be the image of
%	\[
%	\Gamma_{g_j,U(\boldsymbol{\mathrm{x}})}\backslash D_{\textbf{\textrm{x}}} \longrightarrow \Gamma_{g_j}\backslash D,
%	\]
%	where~$\Gamma_{g_j,U(\mathrm{\boldsymbol{x}})}$ is the stabilizer of~$U(\textbf{\textrm{x}})$ in~$\Gamma_{g_j}$.
%	\end{defi}
%	
%	For any~$\beta\in\Sym_{r}(\QQ)$, we set
%	\[
%	\Omega_\beta
%	=
%	\big\{
%	\textbf{\textrm{x}} \in V^r : q(\textbf{\textrm{x}})=\beta
%	\quad\text{and}\quad
%	\dim U(\textbf{\textrm{x}})=\rk\beta
%	\big\}.
%	\]
%	\begin{defi}
%	To any~$K$-invariant Schwartz function~$\varphi\in\mathcal{S}\big(V(\Ab_f)^r\big)$ and any~$\beta\in\Sym_{r}(\QQ)$ we define the \emph{special cycle}~$Z(\beta,\varphi,K)$ on~$X_K$ as the linear combination
%	\[
%	Z(\beta,\varphi,K)
%	=
%	\sum_j
%	\sum_{\substack{\textbf{\textrm{x}} \in\Omega_\beta
%	\\
%	\text{mod}\, \Gamma_{g_j}
%	}}
%	\varphi(g_j^{-1}\textbf{\textrm{x}})\cdot c\big(U(\textbf{\textrm{x}}),g_j,K\big).
%	\]
%	Similarly, the \emph{special cycle} of the same indices on the connected component~$Y_K$ is
%	\be\label{eq;defspccyconnibr}
%	Z_{Y_K}(\beta,\varphi,K)
%	=
%	\sum_{\substack{\textbf{\textrm{x}} \in\Omega_\beta
%	\\
%	\text{mod}\, \Gamma_K
%	}}
%	\varphi(\textbf{\textrm{x}})\cdot c\big(U(\textbf{\textrm{x}}),1,K\big).
%	\ee
%	\end{defi}
%	We now focus on~\cite[Remark~$2.4$]{noetherlefconj}, rewriting the special cycles above in terms of lattices.
%	
%	Let~$\Lambda$ be an even lattice such that~$V=\Lambda\otimes\QQ$.
%	For each prime~$p$, we let~$\Lambda_p=\Lambda\otimes\ZZ_p$ and let~$K_p$ be the subgroup of~$G(\QQ_p)$ which leaves~$\Lambda_p$ stable and acts as the identity on~$\Lambda_p^\lor/\Lambda_p$.
%	We then set~$K=\prod_p K_p$.
%	
%	\begin{lemma}\label{lemma;intbdBLMM}
%	We have that
%	\[
%	\Gamma_K=\{\gamma\in\SO^+(\Lambda) : \text{\normalfont$\gamma$ acts trivially on~$\Lambda^\lor/\Lambda$}\},
%	\]
%	namely~$\Gamma_K$ is the subgroup~$\widetilde{\SO}{}^+(\Lambda)$ of the discriminant kernel given by elements with positive spinor norm.
%	\end{lemma}
%	\textcolor{blue}{\textbf{NOTE:} This is not what is written in~\cite[Remark~$2.4$]{noetherlefconj}, in fact I think there is a typo.
%	They wrote that
%	\[
%	\Gamma_K=\{\gamma\in\SO(\Lambda) : \text{\normalfont$\gamma$ acts trivially on~$\Lambda^\lor/\Lambda$}\},
%	\]
%	namely it is a subgroup of the whole~$\SO(\Lambda)$ instead of~$\SO^+(\Lambda)$.
%	\\
%	I think this is a typo, because in~\cite[Section~$2.3$, second paragraph]{noetherlefconj} they defined~$\Gamma_K$ to be a subgroup of~$\widetilde{G}(\QQ)_+$, which is the subgroup of~$\GSpin(V)$ of elements with \emph{totally positive spinor norm.}
%	Hence also the elements of~$\Gamma_K$ must have positive spinor norm.
%	}
%	\begin{proof}[Proof of Lemma~\ref{lemma;intbdBLMM}]
%	Since~$\Gamma_K$ is the image of~$\widetilde{G}(\QQ)_+\cap \widetilde{K}\subset\widetilde{G}(\Ab_f)$ in~$\SO_0(n,2)$, if~$\gamma\in\Gamma_K$, then it is the diagonal embedding in~$G(\Ab_f)$ of some~$\gamma\in\SO^+(\Lambda\otimes\QQ)$.
%	Moreover we have
%	\[
%	\gamma(\Lambda_p)=\Lambda_p\qquad\text{for every prime~$p$},
%	\]
%	hence~$\gamma\in G(\ZZ_p)=\SO(\Lambda\otimes\ZZ_p)$ for every~$p$.
%	Since~$\ZZ=\{x\in\QQ : x\in\ZZ_p\,\,\text{for all~$p$}\}$, we deduce that~$\gamma\in\SO^+(\Lambda)$.
%	Since~$K_p$ is a subgroup of~$G(\QQ_p)$ which acts as the identity on~$\Lambda_p^\lor/\Lambda_p$, we deduce that every~$\gamma$ acts as the identity on~$\Lambda_p^\lor/\Lambda_p$ for every prime~$p$, hence acts as the identity also on~$\Lambda^\lor/\Lambda$.
%	\end{proof}
%	
%	\textcolor{red}{\textbf{For Jan:}}\\
%	For every even lattice~$\Lambda$ in~$V$, we denote by~$\mathcal{S}_\Lambda\subset\mathcal{S}\big(V(\Ab_f)\big)$ the subspace spanned by the characteristic functions~$\varphi_\lambda$ of the cosets~$\lambda + \widehat{\Lambda}$ with~$\lambda\in \Lambda^\lor$, where~$\widehat{\Lambda}=\Lambda\otimes\prod_p\ZZ_p$.
%%	\begin{lemma}[Wrong!!!]\label{lemma;KinvimplSlambda}
%%	Let~$K$ be the open compact in~$G(\Ab_f)$ arising from the lattice~$\Lambda$ as above.
%%	If~$\varphi\in\mathcal{S}\big(V(\Ab_f)\big)$ is~$K$-invariant, then~$\varphi\in\mathcal{S}_\Lambda$.
%%	\end{lemma}
%%	As recalled in~\cite[Fifth line of the proof of Theorem~$3.2$]{kudla-speccyc}, every~$\varphi\in\mathcal{S}\big(V(\Ab_f)\big)$ lies in~$\mathcal{S}_L$ for some even lattice~$L$.
%%	Without loss of generality\footnote{
%%	Consider the general case.
%%	Let~$L$ be an even lattice in~$V$ such that~$\varphi\in\mathcal{S}_L$.
%%	Hence~$\varphi$ as a combination of characteristic functions
%%	\[
%%	\varphi=\sum_i A_i\varphi_{\lambda_i},
%%	\]
%%	where~$A_i\in\CC$ and~$\lambda_i\in L^\lor$.
%%	We may assume that the~$A_i$ are pairwise different, and so also the~$\lambda_i$.
%%	Note that~$\varphi(\mu_i)=A_i$ if and only if~$\mu_i\in\lambda_i + \widehat{L}$.
%%	Since~$\varphi$ is~$K$-invariant, we deduce that~$\varphi(\mu_i)=\varphi\big(k_p(\mu_i)\big)$ for every~$k_p\in K_p$ and every~$\mu_i\in \lambda_i + L\otimes\ZZ_p$.
%%	Hence
%%	\[
%%	k_p(\lambda_i+L_p)=\lambda_i+L_p\qquad\text{for every~$k_p\in K_p$ and every prime~$p$.}
%%	\]
%%	This is the same as~\eqref{eq;fromKinvad}, hence the assertion follows as for the case of~$\varphi=\varphi_\lambda$.
%%	}
%%	we may assume that~$\varphi=\varphi_\lambda$, where~$\varphi_\lambda$ is the characteristic function of~$\lambda + \widehat{L}$ for some~$\lambda\in L^\lor$.
%%	Since~$\varphi_\lambda$ is~$K$-invariant by hypothesis, we deduce that
%%	\be\label{eq;fromKinvad}
%%	k_p(\lambda+L_p)=\lambda+L_p\qquad\text{for every~$k_p\in K_p$ and every prime~$p$.}
%%	\ee
%%	This is equivalent to the fact that
%%	\be\label{eq;fromnelnwbkehgvfxmöoköd}
%%	k_p(\lambda)=\lambda\quad\text{in~$L^\lor_p/L_p$}\quad\text{and}\quad
%%	k_p(L_p)=L_p\quad
%%	\text{for every~$k_p\in K_p$  and every prime~$p$.}
%%	\ee
%%	Since by definition~$K_p$ is the subgroup of~$G(\QQ_p)$ which leaves~$\Lambda_p$ stable and acts as the identity on~$\Lambda_p^\lor/\Lambda_p$, we deduce from~\eqref{eq;fromnelnwbkehgvfxmöoköd} that~$L$ is a sublattice of~$\Lambda$
%%	\\\textcolor{red}{\textbf{Question for Jan:} Is the last assertion correct? Why?}\\
%%	and hence~$\varphi_\lambda\in\mathcal{S}_L\subseteq\mathcal{S}_\Lambda$.
%%	
%	We can not expect that if~$\varphi\in\mathcal{S}\big(V(\Ab_f)\big)$ is~$K$-invariant, with~$K$ arising from~$\Lambda$ as above, then~$\varphi\in\mathcal{S}_\Lambda$.
%	However, it is possible to explicitly rewrite the special divisors arising from~$K$-invariant Schwartz functions in terms of "pieces" of the divisors arising from~$\mathcal{S}_\Lambda$, as explained in the following.
%	\\
%	
%	Let~$\varphi\in\mathcal{S}\big(V(\Ab_f)\big)^K$.
%	Then there exists an even lattice~$L\subset V$ such that~$\varphi\in\mathcal{S}_L$; see e.g.~\cite[Fifth line of the proof of Theorem~$3.2$]{kudla-speccyc}.
%	Without loss of generality we may assume that~$L\subseteq\Lambda$,
%	\\\textcolor{blue}{
%	\textbf{Check of the last assertion:}
%	Suppose that~$L\not\subseteq\Lambda$, and consider the even lattice~$\tilde{L}\coloneqq L\cap\Lambda$.
%	We want to show that if~$\varphi\in\mathcal{S}_L$, then~$\varphi\in\mathcal{S}_{\tilde{L}}$.
%	We may assume w.l.o.g.\ that~$\varphi=\cha(\discel + \widehat{L})$ for some~$\discel\in D_L$.
%	Since
%	\[
%	\tilde L \subseteq L \subseteq L' \subseteq \tilde L',
%	\]
%	and since~$\rk L=\rk \tilde L$, we deduce that the quotient~$L/\tilde L$ is a finite module.
%	We may then rewrite~$L$, resp~$\widehat{L}$, as the finite union
%	\[
%	L=
%	\bigcup_{\ell\in L/\tilde L}
%	\ell + \tilde{L},\qquad
%	\text{resp.}
%	\qquad
%	\widehat{L}=
%	\bigcup_{\ell\in L/\tilde L}
%	\ell + \widehat{\tilde{L}}.
%	\]
%	These imply that
%	\[
%	\cha(\alpha + \widehat{L})
%	=
%	\sum_{\ell\in L/\tilde{L}}
%	\cha\big((\alpha + \ell) + \widehat{\tilde L}\big).
%	\]
%	}\\
%	and in fact also that~$L=m\Lambda$ for some~$m\in\ZZ_{>0}$.
%	We may then write
%	\[
%	\varphi=\sum_{\alpha\in D_L} A_\alpha\cdot\cha(\alpha + \widehat{L})
%	\]
%	for some~$A_\alpha\in\CC$.
%	
%	The action of~$K$ on~$V$ induces an action of~$K$ on~$D_L$.
%	This induces an action on~$\mathcal{S}_L$ given by permutations of characteristic functions.
%	Since~$\varphi$ is~$K$-invariant, we deduce that its coefficients satisfy
%	\[
%	A_{\kappa(\alpha)} = A_\alpha\qquad\text{for all~$\kappa\in K$.}
%	\]
%	This implies that we may rewrite the locally finite sum on the Hermitian symmetric domain~$\domain$ giving rise to the special divisor of indices~$\varphi$ and~$\beta\in\QQ_{>0}$ as
%	\ba\label{eq;splitspcyad}
%	\sum_{\substack{x\in V \\ q(x)=\beta}}
%	\varphi(x) x^\perp
%	=
%	\sum_{\alpha\in D_L}
%	A_\alpha
%	\sum_{\substack{\lambda\in \alpha + L \\ q(\lambda) = \beta}}
%	\lambda^\perp
%	=
%	\sum_{\xi\in K\backslash D_L}
%	A_\xi
%	\underbrace{\sum_{\alpha\in\xi}
%	\sum_{\substack{\lambda\in\alpha + L\\
%	q(\lambda)=\beta}}
%	\lambda^\perp}_{\eqqcolon \bigstar_{\xi,\beta}},
%	\ea
%	where~$A_\xi$ is the value~$A_\alpha$ obtained from one (hence any)~$\alpha\in\xi$.
%	\\
%	\textcolor{red}{\textbf{IMPORTANT REMARK:} The LHS of~\eqref{eq;splitspcyad} is a combination of divisors from~$\mathcal{S}_\Lambda$ if and only if each term~$\bigstar_{\xi,\beta}$ is so.
%	In fact,~$\varphi$ can be any~$K$-invariant Schwartz function, hence we may choose~$A_\xi=0$ for all~$\xi\in K\backslash D_L$ but one.}
%	\\	
%	We are going to show that the term~$\bigstar_{\xi,\beta}$ appearing on the right-hand side of~\eqref{eq;splitspcyad} is actually a "piece" of the locally finite sum on~$\domain$ giving rise to a special divisor arising from~$\mathcal{S}_\Lambda$.
%	
%	Since~$L=m\Lambda$, we have
%	\ba
%%	\Lambda'&=\{
%%	\text{$v\in V$ : $(v,\lambda)\in \ZZ$ for all $\lambda\in\Lambda$} 
%%	\}
%%	\\
%	L'&=\{
%	\text{$v\in V$ : $(v,\lambda)\in \frac{1}{m}\ZZ$ for all $\lambda\in\Lambda$}
%	\}
%	=
%	\frac{1}{m}\Lambda'.
%	\ea
%	The multiplication by~$m$ on~$V$, i.e.\ the map~$f_m\colon V\to V$, $v\mapsto mv$, induces a bijective map~$f_m\colon L'\to\Lambda'$ such that~$f_m(L)=m^2\Lambda$.
%	This implies that the map
%	\be\label{eq;identDLquotm2}
%	D_L\longrightarrow\Lambda'/m^2\Lambda,\qquad
%	\alpha + m\Lambda\longmapsto m\alpha + m^2\Lambda
%	\ee
%	is an isomorphism.
%%	From now on we identify~$D_L$ as in~\eqref{eq;identDLquotm2}.
%	
%	The multiplication-by-$m$ map~$f_m$ passes to the quotient.
%	We denote by~$F_m$ the induced map
%	\bas
%	F_m\colon D_L\cong\Lambda'/m^2\Lambda&\longrightarrow D_\Lambda,\\
%	\gamma + m^2\Lambda&\longmapsto \gamma + \Lambda.
%	\eas
%	\textcolor{blue}{\textbf{Note:} The map~$F_m$ does NOT multiplies~$\gamma$ by~$m$, because the multiplication has been already carried over in the identification~\eqref{eq;identDLquotm2}.}
%	\\
%	Let~$\mu$ be a \emph{primitive} element of~$\Lambda$, then
%	\[
%	F_m^{-1}(\gamma+\Lambda)
%	=
%	\Big\{
%	\gamma + m^2\Lambda,
%	\gamma + \mu + m^2\Lambda,
%	\dots,
%	\gamma + (m^2-1)\mu + m^2\Lambda
%	\Big\}.
%	\]
%%	from which we deduce that the preimage of an element of~$D_\Lambda$ under~$F_m$ has always cardinality~$m^2$.
%	We remark that the union of all the co-sets in~$F_m^{-1}(\gamma+\Lambda)$ give~$\gamma + \Lambda$, namely
%	\be\label{eq;unionpreim}
%	\bigcup_{j=0}^{m^2-1}
%	\gamma + j\mu + m^2\Lambda
%	=\gamma + \Lambda.
%	\ee
%	
%	The group~$K$ acts as the identity on~$D_\Lambda$, but in general permutes the elements of~$D_L$.
%	Using that the action of~$K$ on~$L'$ is linear, it is easy to see that
%	\[
%	F_m\big(\kappa (\alpha)\big) = F_m(\alpha)
%	\qquad
%	\text{for all~$\kappa\in K$ and~$\alpha\in D_L$}.
%	\]
%	This implies that~$F_m\colon D_L\to D_\Lambda$ passes to the quotient as
%	\[
%	F_m'\colon K\backslash D_L\longrightarrow D_\Lambda,
%	\qquad
%	[\alpha]\longmapsto F_m(\alpha).
%	\]
%	\textcolor{blue}{\textbf{Note:} The quotient~$K\backslash D_L$ is not a group.}
%	\\
%	Hence, we may split~$F_m^{-1}(\gamma + \Lambda)$ into a disjoint union of~$K$-orbits.
%	This means that there exist~$\theta_1,\dots,\theta_\ell\in F_m^{-1}(\gamma + \Lambda)$ such that
%	\bes
%	F_m^{-1}(\gamma + \Lambda)
%	=
%	\coprod_{j=1}^\ell
%	\mathcal{O}_j(\gamma),\qquad
%	\mathcal{O}_j(\gamma)\coloneqq
%	\big\{
%	\kappa(\theta_j) + m^2\Lambda : \kappa\in K
%	\big\}\subseteq F_m^{-1}(\gamma + \Lambda).
%	\ees
%	We may then rewrite~\eqref{eq;unionpreim} as a
%	\ba\label{eq;unionpreim2}
%	\coprod_{j=1}^{\ell}
%	\coprod_{\alpha\in\mathcal{O}_j(\gamma)}
%	\alpha + m^2\Lambda
%	=\gamma + \Lambda.
%	\ea
%	We now rewrite~$\bigstar_{\xi,\beta}$ as
%	\bas
%	\bigstar_{\xi,\beta}
%	&=
%	\sum_{\alpha\in\xi}
%	\sum_{\substack{\lambda \in \alpha+m\Lambda \\ q(\lambda)=\beta}}
%	\lambda^\perp
%	=
%	\sum_{\alpha\in\xi}
%	\sum_{\substack{\lambda \in \frac{1}{m}(m\alpha + m^2\Lambda) \\ q(\lambda)=\beta}}
%	\lambda^\perp
%	=
%	\sum_{\alpha\in\xi}
%	\sum_{\substack{\lambda \in m\alpha + m^2\Lambda \\ q(\lambda)=m^2\beta}}
%	\lambda^\perp
%	=
%	\sum_{\substack{\lambda \in \bigcup_{\kappa\in K}\kappa(\theta_j) + m^2\Lambda \\ q(\lambda)=m^2\beta}}
%	\lambda^\perp
%	\eas
%	for some~$\theta_j$ in~$F_m^{-1}(\gamma + \Lambda)$, where~$\gamma \in F_m(\xi)+\Lambda$.
%	We then use~\eqref{eq;unionpreim2} to deduce that
%	\bas
%	\bigstar_{\xi,\beta}
%	&=
%	\sum_{\substack{\lambda \in \coprod_{\alpha\in\mathcal{O}_j(\gamma)}\alpha + m^2\Lambda \\ q(\lambda)=m^2\beta}}
%	\lambda^\perp,
%	\eas
%	namely~$\bigstar_{\xi,\beta}$ is \emph{just a piece} of the locally finite sum
%	\[
%	\sum_{\substack{\lambda \in \gamma + \Lambda \\ q(\lambda)=m^2\beta}}
%	\lambda^\perp
%	\]
%	giving rise to the Heegner divisors~$H_{\gamma,m^2\beta}=H_{F_m(\xi),m^2\beta}$.
%	\\
%	\textcolor{red}{
%	\textbf{The only possible strategy I have in mind to prove that~$\bigstar_{\xi,\beta}$ is a combination of Heegner divisors is:}
%	\begin{itemize}
%	\item
%	Assume that~$\Lambda$ splits off two orthogonal hyperbolic planes, e.g.~$\Lambda$ is unimodular.
%	Then the \emph{primitive Heegner divisors} give irreducible divisors on~$X$, hence~$\bigstar_{\xi,\beta}$ must be a combination of Heegner divisors.
%	\\
%	To do: Write~$\bigstar_{\xi,\beta}$ as a combination of primitive Heegner divisors.
%	\item
%	Take off the assumption that~$\Lambda$ splits off two hyperbolic planes.
%	Perhaps the linear combination above of~$\bigstar_{\xi,\beta}$ in terms of primitive Heegner divisors is still true.
%	Since primitive Heegner divisors are always combinations of Heegner divisors even when they are not irreducible, we deduce that~$\bigstar_{\xi,\beta}$ is a combination of Heegner divisors.
%	\item
%	Increase the codimension.
%	Is it still possible to rewrite~$\bigstar_{\xi,\beta}$ in terms of primitive special cycles?
%	\end{itemize}
%	}
%	
	
	
	
	
	\end{document}
	
	
	
	
	
	
	
	
	
	
	
	
	
	
	
	
	
	
	
	
	
	
	
	
	
	
	
	
\documentclass[../Lefschetz_decomposition.tex]{subfiles}

	\begin{document}
	
	In this section we illustrate some geometric applications of the Lefschetz decomposition of~$[\KMtheta{\genus}]$ and the geometric interpretation of the lift~$\KMlift{\genus}$.
%	We will geometrically characterize the image of the Kudla--Millson lift, and deduce dimension formulas for the cohomology groups of~$X$, as well as of the span of special cycles in Chow groups.
	
	\subsection{More on the image of the Kudla--Millson lift}\label{sec;imageKMlift}
	From the Lefschetz decomposition of~$[\KMtheta{\gengenus}]$ we may sharpen Theorem~\ref{thm;imageofKMlift} as follows.
%	\begin{cor}
%	Let~$X=\disk{}^+(L)\backslash\domain$.
%	If~$\genus < (n+2)/4$, then
%%	~$\gengenus<\min \{(\dim X+1)/4,(\dim X-1)/3\}$.
%	the image of the Kudla--Millson lift~$\KMlift{\gengenus}$ is contained in the primitive cohomology group~$\sqHprim^{\gengenus,\gengenus}(X)$.
%	Furthermore, if~$g\leq 2$ and~$L$ splits off~$E_8^{\genus-1}\oplus U^{\oplus 2}$, then the image of~$\KMlift{\gengenus}$ is the whole~$\sqHprim^{\gengenus,\gengenus}(X)$.
%	\end{cor}
%	\textcolor{red}{The claim in Corollary~\ref{cor;imageofKMlift} under the split of~$E_8^{\genus-1}\oplus U^{\oplus 2}$ depends on Proposition~\ref{prop;onirredofprimcy}.
%	It is ok as soon as the proof of the latter proposition is completed!}
	\begin{cor}\label{cor;imageofKMliftissurj}
	Let~$\genus < (n+2)/4$.
	If the codimension~$\genus$ special cycles on~$X$ arising from~$L$ span the whole~$\sqH^{\genus,\genus}(X)$, then the image of~$\KMlift{\gengenus}$ is~$\sqHprim^{\gengenus,\gengenus}(X)$.
	\end{cor}
	A criterion on~$L$ for which~$H^{\genus,\genus}(X)$ is spanned by the codimension~$\genus$ special cycles is given by Corollary~\ref{cor;blmminnonadcase} for small~$g$.
	\begin{proof}
%	The space of Siegel cusp forms~$S^k_{\gengenus,\lattice}$ is orthogonal to the space of Klingen Eisenstein series~$\vklisub{k}{\gengenus}{r}{\lattice}$ for every~$r<\gengenus$ by Proposition~\ref{prop;ortcusp&kling}.
%	By Theorem~\ref{thm;lefschetzdecgengenus} we deduce that~$\KMlift{\gengenus}(f)=\pet{f}{[\KMtheta{g}]_0}$
%%	\begin{align*}
%%	\KMlift{\gengenus}(f)
%%	&=
%%	\pet{f}{[\KMtheta{\gengenus}]}
%%%	\\
%%%	&
%%	=
%%	\sum_{r=0}^{\gengenus-1}
%%	(-1)^{\gengenus-r}
%%	\lefschetz^{\gengenus-r} \circ
%%	\underbrace{\pet{f}{\klinser{k}{\gengenus}{r}{\lattice}\big([\KMtheta{r}]_0\big)}}_{=0}
%%	+
%%	\pet{f}{[\KMtheta{g}]_0},
%%	\end{align*}
%	for every cusp form~$f\in S^k_{\gengenus,\lattice}$.
%	It follows that the image of~$\KMlift{\gengenus}$ is a subspace of~$\sqHprim^{\gengenus,\gengenus}(X)$.
%	
	By Lemma~\ref{eq;spanpoincseries} the space~$S^k_{\gengenus,\lattice}$ is generated by Poincaré series~$\poincv{\gengenus}{\discel}{T}{\lattice}$, where~$\discel\in D_\lattice^\genus$ and~$T\in q(\discel)+\halfint{\genus}$ is positive definite.
	Together with Lemma~\ref{lemma;petSiegpoinandcusp} and Theorem~\ref{thm;lefschetzdecgengenus}, this implies that the image of~$\KMlift{\genus}$ is generated by the primitive classes~$[Z(T,\discel)]_0$ of codimension~$\genus$ special cycles, with~$T>0$.
	These primitive classes span the whole~$\sqHprim^{\gengenus,\gengenus}(X)$ by assumption.
%	
%	Since~$\genus<(n+1)/3$, by~\cite[Theorem~$3.6$]{noetherlefconj} the space~$\sqH^{\gengenus,\gengenus}(X)$ is generated by~$L^2$-classes of special cycles.
%	This implies that the projection of such cohomology classes to~$\sqHprim^{\gengenus,\gengenus}(X)$ is a set of generators of the latter cohomology group.
%	To prove that the image of~$\KMlift{\gengenus}$ is the whole~$\sqHprim^{\gengenus,\gengenus}(X)$, it is then enough to show that for every~$\discel\in D_\lattice^\genus$ and every positive definite matrix~$T\in q(\discel)+\halfint{\genus}$, there exists a cusp form~$f_{\discel,T}$ such that~$\KMlift{\gengenus}(f_{\discel,T})=[Z(T,\discel)]_0$.
%	In fact, since~$[\KMtheta{\gengenus}]_0$ is the generating series of the primitive classes of codimension~$g$ special cycles, by Lemma~\ref{lemma;petSiegpoinandcusp} we may choose~$f_{\discel,T}$ to be the Poincaré series~$\poincv{\gengenus}{\discel}{T}{\lattice}$.
%%	We remark that~$P_T$ converges to a well-defined cusp form since~$g<k/2$; see e.g.~\cite[p.~$90$]{klingen}.
%	Up to a non-zero factor, the pairing~$\pet{[\KMtheta{\gengenus}]_0}{\poincv{\gengenus}{\discel}{T}{\lattice}}$ gives the Fourier coefficient of~$[\KMtheta{\gengenus}]_0$ of indices~$T$ and~$\discel$, namely~$[Z(T,\discel)]_0$.
	\end{proof}
	
	In~\cite[Theorem, p.~$4$]{kudlamillson-tubes} Kudla and Millson claim that the image of~$\KMlift{\gengenus}$ is the span in~$\sqH^{\gengenus,\gengenus}(X)$ of the classes of codimension~$\gengenus$ special cycles.
	This is in contrast with Theorem~\ref{thm;imageofKMlift} and Corollary~\ref{cor;imageofKMliftissurj}, since that span does not lie in~$\sqHprim^{\gengenus,\gengenus}(X)$.
	In fact, the proof of~\cite[Theorem, p.~$4$]{kudlamillson-tubes} contains an error.
	The goal of the remaining part of this section is to illustrate the idea behind the proof of~\cite[Theorem, p.~$4$]{kudlamillson-tubes} and to provide a counterexample to it.
	The latter does not rely on the main results of this article. 
	\\
	
%	Let~$g<\min\big\{
%	(\dim X -1)/2,\boundgenlef
%	\big\}$.
	To prove~\cite[Theorem, p.~$4$]{kudlamillson-tubes} Kudla and Millson used the following incorrect claim, which for simplicity we state only in the case of genus~$g=1$ and under the assumption that~$L$ is \emph{unimodular}:
	For every closed compactly supported~$(n-1,n-1)$-differential form~$\eta$ on~$X$, if~$\pet{\theta_\eta}{P_m}=0$, then~$\int_{Z(m)}\eta=0$.
	Here~$\theta_\eta\coloneqq\int_X \eta\wedge\KMtheta{1}\in M^k_1$,~$P_m$ is the normalized Poincaré series of positive index~$m$, and~$Z(m)$ is the~$m$-th special divisor on~$X$.
%	If a closed compactly supported~$(n-1,n-1)$-differential form~$\eta$ on~$X$ is such that
%	\be\label{eq;orthofromKM}
%	\int_X
%	\eta
%	\wedge
%	\KMlift{1}(f)=0\qquad
%	\text{for every~$f\in S^k_1$,}
%	\ee
%	then
%	\be\label{eq;KMmist}
%	\int_{Z(m)}\eta=0
%	\ee
%	for every special divisor~$Z(m)$ of positive index~$m$.
%	(and associated to~$0\in \{0\}=D_\lattice$).
	
	Their attempt of proof of the claim above is based on an imprecise analytic continuation, which we now explain.
%	For simplicity, we illustrate it only in the case of genus~$g=1$ and under the assumption that~$L$ is \emph{unimodular}.
%	The function
%	\[
%	\theta_\eta\coloneqq \int_X\eta\wedge\KMtheta{1}
%	\]
%	is a weight~$k$ elliptic modular form with respect to~$\SL_2(\ZZ)$ with Fourier expansion
%	\[
%	\theta_\eta(\tau)=-\int_X\eta\wedge\omega + \sum_{m>0}\Big(\int_{Z(m)}\eta\Big) q^m.
%	\]
%%	Let~$\petempty$ be the genus~$1$ Petersson inner product.
%	Since
%%	\Big(\pet{[\KMtheta{1}]}{f},[\eta]\Big)
%%	=
%	$\int_X
%	\eta\wedge
%	\pet{[\KMtheta{1}]}{f}
%	=
%	\bigpet{\theta_\eta}{f}$ for every~$f\in S^k_1$, then~\eqref{eq;orthofromKM} implies that
%	\be\label{eq;assumptionKMkist}
%	\bigpet{\theta_\eta}{f}=0
%	\qquad\text{for every~$f\in S^k_1$.}
%	\ee
	Define the \emph{non-holomorphic Poincaré series of index~$m$} as
	\[
	P_m(\tau,s)
	= 
	c(s)
	\sum_{\gamma\in\Gamma_\infty\backslash\SL_2(\ZZ)}
	(y^sq^m)|_k\gamma
%	=
%	c(s) \cdot\sum_{\gamma\in\Gamma_\infty\backslash\SL_2(\ZZ)}
%	\frac{e\big( m\gamma\cdot\tau  \big)}{(c\tau+d)^k} \Im(\gamma\cdot\tau)^s
	,\qquad s\in \CC,
	\]
	where~$c(s)$ is the normalization constant such that~$\pet{f}{P_m(\cdot,s)}=c_m(f)$ for every~${f\in S^k_1}$, and~$\Gamma_\infty$ is the parabolic subgroup of translations in~$\SL_2(\ZZ)$.
	Note that~$P_m(\tau,0)=P_m(\tau)$.
%	If~$s=0$, these are the well-known holomorphic Poincaré series spanning the space of cusp forms~$S^k_1$.
%	In this case the normalizing constant is
%	\be\label{eq;normconstholpsergen1}
%	c(0)=\frac{(4\pi m)^{k-1}}{(k-2)!}.
%	\ee
%	
%	By~\eqref{eq;assumptionKMkist}, we know that
%	\be\label{eq;KMmisbunf}
%	\bigpet{\theta_\eta}{P_m(\cdot,s)}\big|_{s=0}
%	=
%	\bigpet{\theta_\eta}{P_m(\cdot,0)}
%	=
%	0.
%	\ee
%	The idea in~\cite{kudlamillson-tubes} is to compute~$\bigpet{\theta_\eta}{P_m(\cdot,s)}$ using an unfolding argument, and then evaluate this in~$s=0$ to deduce~\eqref{eq;KMmist}.
%	This is the dubious part of the proof.
%	The evaluation in~$s=0$ must be computed \emph{after a suitable analytic continuation}, since the unfolding process applied to~$\bigpet{\theta_\eta}{P_m(\cdot,s)}$ is justified only for~$s$ with large real part.
%	The problem, as we will explain soon, is that the value at~$s=0$ of that analytic continuation differs from~$\bigpet{\theta_\eta}{P_m(\cdot,0)}$.
%	\\
%	
%	\textcolor{red}{This means that the function~$s\mapsto \bigpet{\theta_\eta}{P_m(\cdot,s)}$ is not analytic in the half-plane~$\Re(s)>-\varepsilon$.
%	Otherwise, the uniqueness of the analytic continuation would imply that the value at~$s=0$ of that analytic continuation equals~$\bigpet{\theta_\eta}{P_m(\cdot,0)}$.
%	\\
%	
%	With the usual unfolding trick, it is possible to prove that if~$\Re(s)>1$, then
%	We now illustrate the unfolding process of~\cite{kudlamillson-tubes}.
%	We may rewrite
%	\begin{align}\label{eq;unfoldingfstKMmist}
%	\begin{split}
%	\bigpet{\theta_\eta}{P_m(\cdot,s)}
%	=
%	c(s)\int_{\SL_2(\ZZ)\backslash\HH}y^k\theta_\eta(\tau)
%	\sum_{\gamma\in\Gamma_\infty\backslash\SL_2(\ZZ)}
%	\frac{\Im(\gamma\cdot\tau)^s e(-m\gamma\cdot\overline{\tau})}{(c\overline{\tau}+d)^k} \frac{dx\,dy}{y^2}
%	\\
%	=
%	c(s)\int_{\SL_2(\ZZ)\backslash\HH}
%	\sum_{\gamma\in\Gamma_\infty\backslash\SL_2(\ZZ)}
%	\theta_\eta(\gamma\cdot\tau)
%	\cdot e(-m\gamma\cdot\overline{\tau})
%	\cdot
%	\Im(\gamma\cdot\tau)^{k+s}
%	\frac{dx\,dy}{y^2}.
%	\end{split}
%	\end{align}
%	We want to apply the unfolding trick on the right-hand side of~\eqref{eq;unfoldingfstKMmist}, hence we want to rewrite the integral over~$\SL_2(\ZZ)\backslash\HH$ as an integral over the unfolded quotient~$\Gamma_\infty\backslash\HH$.
%	To do so, we need to switch the summation with the integration over~$\SL_2(\ZZ)\backslash\HH$.
%	This is possible only if~$\Re(s)$ is large enough.
%	In fact, if we formally unfold the integral as
%	\be\label{eq;unfoldedKMmist}
%	\bigpet{\theta_\eta}{P_m(\cdot,s)}
%	=
%	c(s)\int_{\Gamma_\infty\backslash\HH}\theta_\eta(\tau)\cdot e(-m\overline{\tau}) y^{k+s-2}\,dx\,dy.
%	\ee
%	The integrand above is not integrable on the unfolded domain~$[0,1]\times[0,\infty)$ for~$\Re(s)$ small.
%	The problem comes from the integration with respect to~$y\in[0,\infty)$, since the modular form~$\theta_\eta$, which has in general a non-trivial Eisenstein part, is~$O(y^{-k})$ when~$y\to 0$.
	
	As illustrated in~\cite[p.~$25$]{kudlamillson-tubes}, if~$\Re(s)>1$, then one can apply the unfolding trick to deduce that
%	we may substitute the Fourier expansion~$\theta_\eta=\sum_{\ell\ge 0}\int_{Z(\ell)}\eta q^\ell $ in~\eqref{eq;unfoldedKMmist} and deduce that
	\begin{align}\label{eq;unfoldingKMmistend}
	\begin{split}
	\bigpet{\theta_\eta}{P_m(\cdot,s)}
%	=
%	\int_0^\infty\int_0^1
%	\sum_{\ell\ge 0} c_\ell(\theta_\eta)\cdot e(\ell x)\cdot\exp(-2\pi\ell y)\cdot e(-mx)\cdot \exp(-2\pi m y) y^{k+s-2}\,dx\,dy
%	\\
%	&=
%	\sum_{\ell\ge 0} c(s)\cdot c_\ell(\theta_\eta)
%	\int_0^\infty
%	\exp\big(-2\pi(\ell+m) y\big) y^{k+s-2}\,dy
%	\int_0^1
%	e\big((\ell-m)x\big)\,dx
%	\\
%	&=
%	c(s)\cdot c_m(\theta_\eta)\int_0^\infty y^{k+s-2}\exp(-4\pi my)\,dy
	=
	\frac{c(s)\cdot\Gamma(k+s-1)}{(4\pi m)^{k+s-1}}\int_{Z(m)}\eta.
	\end{split}
	\end{align}
%	where in the last equality we used that~$c_m(\theta_\eta)=\int_{Z(m)}\eta$.
	Although~\eqref{eq;unfoldingKMmistend} is justified only for~$\Re(s)>1$, we may consider the analytic continuation of the right-hand side of~\eqref{eq;unfoldingKMmistend} and evaluate it at~$s=0$.
	This equals~$\int_{Z(m)}\eta$.
	
	The problem is that such value does not equal~$\bigpet{\theta_\eta}{P_m(\cdot,0)}$, and in fact the vanishing of~$\bigpet{\theta_\eta}{P_m(\cdot,0)}$ does not imply the vanishing of~$\int_{Z(m)}\eta$.
	This is shown in the following result, based on the fact that~$\bigpet{\theta_\eta}{P_m(\cdot,0)}=c_m(\theta_\eta)$ only if~$\theta_\eta$ is a cusp form, which is not the case for certain~$\eta$.
	\begin{lemma}
	The vanishing of~$\bigpet{\theta_\eta}{P_m(\cdot,0)}$ does not imply the vanishing of~$\int_{Z(m)}\eta$.
	\end{lemma}
	\begin{proof}
%	We prove the result in two different independent ways.
%	
%	\textbf{First proof:}
%	If~$\eta$ is a closed compactly supported~$(n-1,n-1)$-form on~$X$, then~$\theta_\eta$ is a weight~$k$ modular form.
	We rewrite the weight~$k$ modular form~$\theta_\eta$ with respect to the modular decomposition~${M^k_1=\CC E^k_1\oplus S^k_1}$ in Eisenstein and cuspidal parts as~$\theta_\eta=a_\eta E^k_1 + \tilde{\theta}_\eta$, for some constant~$a_\eta\in\CC$ and some cusp form~$\tilde{\theta}_\eta\in S^k_1$.
	
	Let~$\eta$ be such that~$a_\eta\neq 0$.
	An example of such a~$\eta$ is given by the Poincaré dual of a compact orthogonal Shimura curve, as follows.
	Let~$X=\Gamma\backslash \domain$, where~$\domain$ is the Hermitian symmetric domain associated to some unimodular lattice~$L$ of signature~$(n,2)$.
	The orthogonal Shimura subvarieties of codimension~$r$ in~$X$ arise from Hermitian symmetric subdomains~$\domain_W$ in~$\domain$ defined as the orthogonal complement in~$\domain$ of some positive definite~$r$-dimensional subspace~$W\subseteq L\otimes\QQ$ .
	We choose~$r=n-1$ and~$W$ such that~$W^\perp$ is anisotropic in~$L\otimes\QQ$.
	The subvariety arising from~$\domain_W$ is then a \emph{compact} curve~$C$ in~$X$.
	The Poincaré dual of~$C$ is a~$(n-1,n-1)$-differential form~$\eta$ of compact support, and the associated modular form~$\theta_\eta$ has Fourier expansion
	\[
	\theta_\eta(\tau)=-\Vol(C) + \sum_{m>0}\Big(\int_{Z(m)}\eta\Big)q^m.
	\]
	This is not a cusp form, since the constant Fourier coefficient is non-zero.
	
	The left-hand side of~\eqref{eq;unfoldingKMmistend} evaluated at~$s=0$ may be computed as
	\[
	\bigpet{\theta_\eta}{P_m(\cdot,0)}=\bigpet{\tilde{\theta}_\eta}{P_m(\cdot,0)}=c_m(\tilde{\theta}_\eta),
	\]
	since the Eisenstein part of~$\theta_\eta$ is orthogonal to the holomorphic Poincaré series of index~$m$.
	
	The right-hand side of~\eqref{eq;unfoldingKMmistend} evaluated at~$s=0$ is the~$m$-th Fourier coefficient of~$\theta_\eta$.
	This differs from the~$m$-th Fourier coefficient of the cuspidal part~$\tilde{\theta}_\eta$, since~$c_m(E^k_1)\neq 0$ for every~$m$.
%	
%	\textbf{Second proof:} The unfolding method applied to deduce~\eqref{eq;unfoldingKMmistend} does not depend on the fact that~$\eta$ is compactly supported.
%	In fact, we may choose any square-integrable differential form, e.g.\ the Kähler form~$\omega$.
%	
%	By Theorem~\ref{thm;fromKudlagenservol} we have~$\theta_\omega=-\Vol(X)\cdot E^k_1$.
%	Since cusp forms are orthogonal to the Eisenstein series with respect to the Petersson inner product, the left-hand side of~\eqref{eq;unfoldingKMmistend} with~$\eta=\omega$ and~$s=0$ vanishes.
%	The right-hand side of~\eqref{eq;unfoldingKMmistend} with~$s=0$ equals~$\Vol(Z(m))$, which is positive for all~$m>0$.
	\end{proof}
	
	
	\subsection{Dimensions of cohomology groups}\label{sec;dimcoho}
	Let~$\genus < (n+2)/4$, and let
	\be\label{eq;tocohomap}
	\modcoho{\genus}{\lattice}\colon \almcusp{k}{\genus}{\lattice}\longrightarrow \sqH^{\genus,\genus}(X)
	\ee
	be the linear map constructed on each term of the modular decomposition of~$\almcusp{k}{\genus}{\lattice}$ as
	\[
	\klinser{k}{\gengenus}{r}{\lattice}(f)
	\longmapsto
	(-\lefschetz)^{\genus-r} \circ \KMlift{r}(f)\qquad \text{for all~$f\in S^k_{r,\lattice}$}.
	\]
	As immediate consequence of Theorem~\ref{thm;imageofKMlift}, we may deduce the following result.
	\begin{cor}\label{cor;ifinjthenformcoho}
%	Let~$X=\disk{}^+(L)\backslash\domain$ and
	Let~$\genus<(n+2)/4$.
	Then
%	\min \{(\dim X+1)/4,(\dim X-1)/3\}.
%	then~$\modcoho{\genus}{\lattice}$ is \emph{surjective} and
	\[
	\dim\lefschetz^{\genus-r} H_0^{r,r}(X)\geq\dim \KMlift{r}\big(S^k_{r,L}\big)\qquad\text{for every~$0\leq r\leq\genus$}.
	\]
	If the Kudla--Millson lift~$\KMlift{r}$ is injective for every~$r\leq \genus$, then so is~$\modcoho{\genus}{\lattice}$ and
	\[
	\dim\lefschetz^{\genus-r} H_0^{r,r}(X)\geq\dim S^k_{r,\lattice}
	\qquad\text{and}\qquad
	\dim H^{\genus,\genus}(X)\geq\dim \almcusp{k}{\genus}{\lattice}.
	\]
%	If moreover~$L$ splits off~$E_8^{\oplus\genus-1}\oplus U^{\oplus 2}$, then~$\modcoho{\genus}{\lattice}$ is an isomorphism and
%	\[
%	\dim\lefschetz^{\genus-r} H_0^{r,r}(X)=\dim S^k_{r,\lattice}
%	\qquad\text{and}\qquad
%	\dim H^{\genus,\genus}(X)=\dim \almcusp{k}{\genus}{\lattice}.
%	\]
	\end{cor}
%	\textcolor{red}{The last sentence of Corollary~\ref{cor;ifinjthenformcoho} is ok only after Proposition~\ref{prop;onirredofprimcy} is proved!
%	}
%	
%	\textcolor{red}{Perhaps we could delete the last sentence of the previous corollary and just change Corollary~\ref{cor;dimcoho1}.}
	\begin{proof}%[Proof of Corollary~\ref{cor;ifinjthenformcoho}]
	By~\eqref{eq;3isoofcoho} and~\cite{weissauer-88} the group~$\sqH^{\genus,\genus}(X)$ is isomorphic to and induces a Lefschetz decomposition on~$H^{\genus,\genus}(X)$.
	The operators~$\lefschetz^{\genus-r}\colon \sqH^{r,r}(X)\to \sqH^{\genus,\genus}(X)$ are injective by Theorem~\ref{thm;lefdecL2}.
	Corollary~\ref{cor;ifinjthenformcoho} is then a consequence of Theorem~\ref{thm;imageofKMlift}.
%	\[
%	\dim\lefschetz^j H_0^{\genus-j,\genus-j}(X)
%	=
%	\dim \sqHprim^{\genus-j,\genus-j}(X)
%	=
%	\dim \KMlift{\genus-j}\big(S^k_{\genus-j,L}\big)\qquad\text{for every~$j$}.\qedhere
%	\]
	\end{proof}
	\begin{rem}
	The dual space~$(\almcusp{k}{\genus}{\lattice})^*$ is generated by the \emph{coefficient extraction functionals}
	\[
	c_{T,\discel}
	\colon \almcusp{k}{\genus}{\lattice}\longrightarrow\CC,
	\qquad
	f\longmapsto c_T(f_\discel).
	\]
	We identify~$\vklisub{k}{\gengenus}{r}{\lattice}$ with its dual space by means of the Petersson product of~$S^k_{r,\lattice}$ as
	\[
	\vklisub{k}{\gengenus}{r}{\lattice}\longrightarrow
	\big(\vklisub{k}{\gengenus}{r}{\lattice}\big)^*,
	\qquad
	\klinser{k}{\gengenus}{r}{\lattice}(f)\longmapsto
	\Big(
	\klinser{k}{\gengenus}{r}{\lattice}(f')\mapsto \pet{f'}{f}
	\Big).
	\]
	It is easy to use such identifications to show that surjectivity and injectivity results on~$\modcoho{\genus}{\lattice}$ implies analogous results on the linear map
	\[
	\dumodcoho{\genus}{\lattice} \colon
	(\almcusp{k}{\genus}{\lattice})^* \longrightarrow \sqH^{\genus,\genus}(X),
	\qquad
	c_{T,\discel}\longmapsto [Z(T,\discel)].
	\]
	Such a map has been used in~\cite{bruiniermöller}, \cite{zuffetti-equid} and~\cite{zuffetti-conescod2} to study cones generated by special divisors in~$\Pic(X)\otimes\RR$ and cones of codimension~$2$ special cycles in~$\CH^{2}(X)\otimes\RR$.
	Important for that purpose is the injectivity of~$\dumodcoho{\genus}{\lattice}$, which has been studied only for~$\genus=1$ with a different approach; see~\cite{bruinierbook} and~\cite{borcherds-grassmannians}.
	By Corollary~\ref{cor;ifinjthenformcoho} we may deduce the injectivity of~$\dumodcoho{\genus}{\lattice}$ from the injectivity of Kudla--Millson lifts.
	This solves the problem raised in~\cite[Remark~$4.13$]{zuffetti-conescod2}.
%	
%	\textcolor{blue}{
%	\textbf{Background to check that the injectivity of~$\modcoho{\genus}{\lattice}$ implies the injectivity of~$\dumodcoho{\genus}{\lattice}$:}\\
%	For simplicity we assume~$L$ to be \emph{unimodular}, and sometimes drop it from the notation.
%	Consider the commutative diagram}
%	\bes
%	\begin{tikzcd}
%	(M^k_\genus)^* \arrow[d]
%	\arrow[r, "\psi_{\genus,\lattice}"]
%	&
%	\sqH^{\genus,\genus}(X) \arrow[d, "\pi_r"]
%	\\
%	(M^k_{\genus,r})^*
%	\arrow[r, "\psi_{\genus,r}"]
%	&
%	\lefschetz^{\genus-r} \sqHprim^{r,r}(X)
%	\end{tikzcd}
%	\ees
%	\textcolor{blue}{where the left-hand arrow is the restriction map, the right-hand arrow~$\pi_r$ is the orthogonal projection to the~$r$-term of the Lefschetz decomposition of~$\sqH^{\genus,\genus}(X)$, and where~$\psi_{\genus,r}$ is the linear map defined as
%	\[
%	\psi_{\genus,r}
%	\colon
%	c_T|_{M^k_{\genus,r}}
%	\longmapsto \pi_r\big([-\omega]^{\genus - \rk T} \wedge [Z(T)]\big),
%	\qquad
%	\text{$T\in\halfint{\genus}$ positive semidefinite.}
%%	\sum_j \lambda_j c_j|_{M^k_{\genus,r}}\mapsto \sum_j \lambda_j [Z(T_j)]_r.
%	\]
%%	Here we denote by~$\pi_r\colon\sqH^{\genus,\genus}(X)\to\lefschetz^{\genus-r}\sqH^{r,r}(X)$ the orthogonal projection to the~$r$-term of the Lefschetz decomposition of~$\sqH^{\genus,\genus}(X)$.\\
%	The fact that~$\psi_{\genus,r}$ is \emph{well-defined} follows from Theorem~\ref{thm;lefschetzdecgengenus} with the same idea to prove that~$\psi_{\genus,\lattice}$ is well-defined.
%	To do so, we need to check that if a combination of restricted coefficient extraction functionals is zero in~$(M^k_{\genus,r})^*$, then such a combination must map to zero under~$\psi_{\genus,r}$.
%%	Note that
%%	\[
%%	\psi_{\genus,r}(c_T|_{M^k_{\genus,r}})
%%	=
%%	\pi_r\big([-\omega]^{\genus - \rk T} \wedge [Z(T)]\big)
%%	=
%%	c_T\big(\pi_r([\KMtheta{\genus}])\big).
%%	\]
%	Since~$\pi_r([\KMtheta{\genus}])$ is a modular form in~$M^k_{\genus,r}$ by Theorem~\ref{thm;lefschetzdecgengenus}, we deduce that if a combination~$\sum_j \lambda_j c_{T_j}|_{M^k_{\genus,r}}$ vanishes in~$(M^k_{\genus,r})^*$ for some~$\lambda_j\in\CC$, then
%	\[
%	0
%	=
%	\sum_j \lambda_j c_{T_j}|_{M^k_{\genus,r}}\big(\pi_r([\KMtheta{\genus}])\big)
%	=
%	\sum_j \lambda_j \pi_r\big([-\omega]^{\genus - \rk T_j}\wedge[Z(T_j)]\big).
%	\]
%	Hence the image of~$\sum_j \lambda_j c_{T_j}|_{M^k_{\genus,r}}$ maps to zero under~$\psi_{\genus,r}$.
%	\\
%	We now construct an identification~$\iota\colon M^k_{\genus,r}\to(M^k_{\genus,r})^*$.
%	To do so, we define an inner product~$((\cdot{,}\cdot))$ on~$M^k_{\genus,r}$ by means of the Petersson product of~$S^k_r$ as
%	\[
%	\big(\big(E^k_{\genus,r}(f_1),E^k_{\genus,r}(f_2)\big)\big)
%	\coloneqq\pet{f_1}{f_2},
%	\qquad
%	f_1,f_2\in S^k_r.
%	\]
%	We may now identify~$M^k_{\genus,r}$ with its dual in the standard way as
%	\[
%	\iota\colon M^k_{\genus,r}\longrightarrow(M^k_{\genus,r})^*,
%	\qquad
%	E^k_{\genus,r}(f)\longmapsto\big(\big(\,\cdot\,,E^k_{\genus,r}(f)\big)\big).
%	\]
%	Consider the following diagram.}
%	\be\label{eq;diagramwithtikzcdib}
%	\begin{tikzcd}
%	(M^k_{\genus,r})^*
%	\arrow[r, "\psi_{\genus,r}"]
%	&
%	\lefschetz^{\genus-r}\sqHprim^{r,r}(X) \arrow[d, "id"]
%	\\
%	M^k_{\genus,r}
%	\arrow[u, "\iota"]
%	\arrow[r, "\varphi_{\genus,r}"]
%	&
%	\lefschetz^{\genus-r} \sqHprim^{r,r}(X)
%	\end{tikzcd}
%	\ee
%	\textcolor{blue}{where~$\varphi_{\genus,r}\coloneqq (-\lefschetz)^{\genus-r}\circ\KMlift{r}\circ\Phi^{\genus-r}$, i.e.\ it is the restriction of~$\modcoho{\genus}{\lattice}$ to~$M^k_{\genus,r}$.
%	\begin{lemma}\label{lemma;riblue?}
%	The diagram~\eqref{eq;diagramwithtikzcdib} commutes.
%	\end{lemma}
%	\begin{proof}[Proof of Lemma~\ref{lemma;riblue?}]
%	It is enough to show the commutativity on a set of generators of~$M^k_{\genus,r}$.
%	Since the Poincaré series of genus~$r$ span the whole~$S^k_r$, it is enough to check the commutativity on the genus~$g$ Klingen Eisenstein series arising from Poincaré series of genus~$r$.
%	Let~$T\in\halfint{r}$ be positive definite, and let~$P_T\in S^k_r$ be the associated Poincaré series.
%	For simplicity we assume that it is \emph{normalized} so that~$\pet{f}{P_T}=c_T(f)$ for every~$f\in S^k_r$.
%	By Corollary~\ref{cor;imageofKMlift} the image of~$P_T$ under~$\varphi_{\genus,r}$ is
%	\ba\label{eq;comminbluetbd?}
%	\varphi_{\genus,r}\big(E^k_{\genus,r}(P_T)\big)
%	&=
%	(-\lefschetz)^{\genus-r}\circ\KMlift{r}\circ\Phi^{\genus-r}\big(E^k_{\genus,r}(P_T)\big)
%	=
%	(-\lefschetz)^{\genus-r}\circ\KMlift{r}(P_T)
%	\\
%	&=
%	(-\lefschetz)^{\genus-r}([Z(T)]_0).
%	\ea
%	The value of the functional~$\iota\big(E^k_{\genus,r}(P_T)\big)=((\cdot , E^k_{\genus,r}(P_T)))$ on~$E^k_{\genus,r}(f')$, with~$f'\in S^k_r$, is the Fourier coefficient~$c_T(f')$.
%	The latter is the Fourier coefficient of~$E^k_{\genus,r}(f')$ of index~$\big(\begin{smallmatrix}
%	T & 0\\
%	0 & 0
%	\end{smallmatrix}\big)\in\halfint{\genus}$.
%	This means that
%	\[
%	\iota\big(E^k_{\genus,r}(P_T)\big)
%	=
%	c_{\big(\begin{smallmatrix}
%	T & 0\\
%	0 & 0
%	\end{smallmatrix}\big)}|_{M^k_{\genus,r}},
%	\]
%	hence
%	\bas
%	\psi_{\genus,r}\Big(\iota\big(E^k_{\genus,r}(P_T)\big)\Big)
%	&=
%	\psi_{\genus,r}\Big(
%	c_{\big(\begin{smallmatrix}
%	T & 0\\
%	0 & 0
%	\end{smallmatrix}\big)}|_{M^k_{\genus,r}}
%	\Big)
%	=
%	\pi_r\Big(
%	[-\omega]^{\genus-r}\wedge
%	\underbrace{\big[Z\big(\begin{smallmatrix}
%	T & 0\\
%	0 & 0
%	\end{smallmatrix}\big)\big]}_{=Z(T)}
%	\Big)
%	\\
%	&=
%	[-\omega]^{\genus-r}\wedge [Z(T)]_0
%	=
%	(-\lefschetz)^{\genus-r}\big([Z(T)]_0\big).
%	\eas
%	This and~\eqref{eq;comminbluetbd?} imply that the diagram~\eqref{eq;diagramwithtikzcdib} commutes.
%	\end{proof}
%	\textbf{Check that injectivity of~$\modcoho{\genus}{\lattice}$ implies injectivity of~$\dumodcoho{\genus}{\lattice}$:}\\
%	Let~$\xi\in (M^k_\genus)^*$ be such that~$\xi\neq 0$.
%	There exists~$r$ such that the restriction~$\xi|_{M^k_{\genus,r}}$ is a non-zero element of~$(M^k_{\genus,r})^*$.
%	Suppose also that~$\psi_{\genus,\lattice}(\xi)=0$.
%	We then have that~$\xi|_{M^k_{\genus,r}}\neq 0$ and~$\psi_{\genus,r}(\xi|_{M^k_{\genus,r}})=0$.\\
%	Let~$\xi'\coloneqq\iota^{-1}(\xi|_{M^k_{\genus,r}})$.
%	Since~$\iota$ is an isomorphism, we deduce that~$\xi'\neq 0$.
%	By Lemma~\ref{lemma;riblue?} and since~$\psi_{\genus,r}(\xi|_{M^k_{\genus,r}})=0$, we deduce that also~$\varphi_{\genus,r}(\xi')=0$.
%	But this is not possible, since~($\varphi_{\genus,\lattice}$ and hence)~$\varphi_{\genus,r}$ is assumed to be injective.
%	}
	\end{rem}
	
	The injectivity of the Kudla--Millson lift of genus~$g=1$ has been studied in several articles.
	The first author proved in~\cite{bruinierbook} and~\cite{bruinier-converse} that~$\KMlift{1}$ is injective whenever the lattice~$L$ splits off a hyperbolic plane and a scaled hyperbolic plane.
%	It is known that~$\KMlift{\genus}$ is injective if~$\genus=1$ and if~$L$ splits off two hyperbolic planes~\cite{bruinierbook}.
%	The latter condition is always satisfied for~$L$ unimodular.
%	This injectivity criterion in genus~$1$ has been generalized to a larger families of lattices in~\cite{bruinier-converse} and~\cite{zuffetti-unfolding}, as well as to locally symmetric spaces of orthogonal type in~\cite{bruinierfunke}, \cite{stein-converse} and~\cite{metzlerzuffetti}.
%	The first author and Funke~\cite{bruinierfunke} proved that~$\KMlift{1}$ is injective whenever the lattice~$L$ is unimodular with~$\rk L>4$.
%	This has been extended in~\cite{bruinier-converse} to (not necessarily unimodular) lattices that split off a hyperbolic plane and a scaled hyperbolic plane.
	Further works of the first author and Funke~\cite{bruinierfunke}, Stein~\cite{stein-converse} and Metzler and the second author~\cite{metzlerzuffetti} provide further generalizations to locally symmetric spaces of orthogonal type.
	
	The second author illustrated in~\cite{zuffetti-unfolding} a new method to prove the injectivity in genus~$1$, with the advantage of paving the ground for a strategy to deduce injectivity in higher genus.	
	This has been implemented in genus~$2$ by Kiefer and the second author in~\cite{kieferzuffetti}.
	They proved the injectivity of~$\KMlift{2}$ under the assumption that~$\lattice$ splits off two hyperbolic planes, and that the orthogonal complement~$\lattice^+$ of such hyperbolic planes is such that~$\lattice^+\otimes\ZZ_p$ splits off two hyperbolic planes over~$\ZZ_p$ for every prime~$p$.
	An example of~$L^+$ satisfying the latter condition is the~$8$th root lattice~$E_8$; see e.g.~\cite[Satz~$15.8$]{kneser} and~\cite[Theorem~$5.2.2$]{kitaoka}.
	
	The following results follow directly from the injectivity results stated above and Corollaries~\ref{cor;blmminnonadcase}, \ref{cor;imageofKMliftissurj} and~\ref{cor;ifinjthenformcoho}.
%	We denote by~$U(N)$ the hyperbolic plane whose quadratic form is twisted by a positive factor~$N$.
	\begin{cor}\label{cor;dimcoho1}
	Let~$X=\disk{}^+(L)\backslash\domain$ with~$L$ of rank~$rk(L)>4$.
	If~$L$ splits off~$U^{\oplus 2}$, then
	\[
	\dim H^{1,1}_0(X)=\dim S^k_{1,\lattice}
	\qquad\text{and}\qquad
	\dim H^{1,1}(X)= \dim \almcusp{k}{1}{\lattice}= 1 + \dim S^k_{1,\lattice}.
	\]
	\end{cor}
	\begin{cor}\label{cor;dimcoho2}
	Let~$X=\disk{}^+(L)\backslash\domain$.
	If~$\lattice$ splits off~$E_8\oplus U^{\oplus 2}$,
%	and if the orthogonal complement~$\lattice^+$ of~$U^{\oplus 2}$ is such that~$\lattice^+\otimes\ZZ_p$ splits off two hyperbolic planes for every prime~$p$,
	then~$\dim H^{2,2}_0(X)=\dim S^k_{2,\lattice}$, $\dim \lefschetz H^{1,1}_0(X)=\dim S^k_{1,\lattice}$ and
	\[
	\dim H^{2,2}(X)= \dim \almcusp{k}{2}{\lattice}= 1 + \dim S^k_{1,\lattice} + \dim S^k_{2,\lattice}.
	\]
	\end{cor}
%	We remark that Corollary~\ref{cor;dimcoho1} was firstly proved by the first author in~\cite[Theorem~$1.5$]{bruinier-converse} with a different method.
	
	\begin{ex}
	Interesting examples of orthogonal Shimura varieties are the (coarse) moduli spaces of irreducible holomorphic symplectic manifolds.
	Several of them satisfy the hypothesis of Corollaries~\ref{cor;dimcoho1} and~\ref{cor;dimcoho2}, for instance the moduli spaces of quasi-polarized K3 surfaces, which are orthogonal Shimura varieties arising from the lattices
	\[
	L_d=\langle 2d \rangle\oplus U^{\oplus 2}\oplus E_8^{\oplus2},
	\]
	where~$d\in\ZZ_{>0}$ is the degree of the polarization.
	Note that we choose the bilinear form on the middle cohomology groups of~K3 surfaces to be the negative of the intersection paring.
	
	For K3 surfaces the degree is the only invariant of a polarization.
	This is not the case for the Hilbert schemes~$S^{[m]}$, where~$S$ is a K3 surface.
	We refer to~\cite[Section~$3$]{GHS-compositio} for a general overview, and restrict here to the \emph{split case}, where the moduli space of deformation~K3${}^{[m]}$ manifolds arises from a lattice of the form
	\[
	L_{t,d}=\langle 2t \rangle\oplus\langle 2d \rangle\oplus U^{\oplus 2}\oplus E_8^{\oplus2}
	\] 
	for some~$t,d\in\ZZ_{>0}$.
	Also in this case the hypothesis of Corollary~\ref{cor;dimcoho2} are satisfied.
	\end{ex}
	
	We remark that explicit formulas to compute~$\dim S^k_{\genus,\lattice}$ are available in the literature only if either~$\genus=1$, or~$L$ is unimodular and~$\genus$ is small; see~\cite[Section~$7$]{borcherds-reflection}, \cite{fischer}, \cite[Section~$5.4$]{taibi}, \cite{ChenevierLannes}.
	The computation of~$\dim S^k_{\genus,\lattice}$ for a general lattice~$L$ is still an interesting open problem.
	
	We remark that the dimension formulas stated in this section hold also for the space spanned by the \emph{rational classes} of special cycles in the Chow group~$\CH^\genus(X)_\CC\coloneqq\CH^\genus(X)\otimes\CC$.
	We define that subspace as
	\[
	\SC^\genus(X)
	\coloneqq
	\Big\langle
	\{Z(T,\discel)\}\cdot\{\omega\}^{\genus-\rk T}
	:
	\text{$\discel\in D_\lattice^\genus$, $T\in q(\discel) + \halfint{\genus}$ and $T\geq 0$}
	\Big\rangle_\CC
	\subseteq\CH^\genus(X)_\CC.
	\]
	As a consequence of Kudla's Modularity Conjecture, proved by the first author and Raum in~\cite{bruinierraum}, it is easy to deduce the following result with the same wording of Corollary~\ref{cor;ifinjthenformcoho}.
	\begin{cor}\label{cor;dimspanCH}
	Let~$X=\disk{}^+(L)\backslash\domain$ and let~$\genus<(n+2)/4$.
	If the Kudla--Millson lift~$\KMlift{r}$ is \emph{injective} for every~$r\leq \genus$, then~$\dim SC^\genus(X)=\dim \almcusp{k}{\genus}{\lattice}$.
%	In particular, under the assumptions of Corollary~\ref{cor;dimcoho1} and Corollary~\ref{cor;dimcoho2} we have resp.~$\dim \Pic(X)_\CC=\dim \almcusp{k}{1}{\lattice}$ and~$\dim SC^2(X)=\dim \almcusp{k}{2}{\lattice}$.
	\end{cor}
%	
%	\textcolor{blue}{\textbf{An extended version of the result above is given in the following subsection.}}
%	\subsection{Dimension of the span of special cycles in Chow groups}\label{sec;Chowdim}
%	
%	\textcolor{blue}{We may consider an analogue of the Kudla--Millson theta function in the Chow group~$\CH^\genus(X)_\CC\coloneqq\CH^\genus(X)\otimes\CC$ by replacing Poincaré duals of special cycles with \emph{rational classes} of special cycles.
%	Kudla's Modularity Conjecture, proved by Bruinier and Raum~\cite{bruinierraum}, states that the generating series of rational classes of special cycles converges to a (holomorphic) Siegel modular form.
%	Equivalently, we have that
%	\[
%	\{\KMtheta{\genus}\}(\tau)
%	\coloneqq
%	\sum_{\discel\in D_\lattice^\genus}
%	\sum_{\substack{T\in q(\discel) + \halfint{\genus} \\ T\geq 0}}
%	\{Z(T,\discel)\}\cdot\{\omega^*\}^{\genus-\rk T}
%	q^T \mathfrak{e}_\discel
%	\in \CH^\genus(X)_\CC \otimes M^k_{\genus,\lattice},
%	\]
%	where~$\{Z(T,\discel)\}\in\CH^{\rk T}(X)_\CC$ is the rational class of the special cycle~$Z(T,\discel)$, and~$\{\omega^*\}$ is the dual class of the Hodge bundle on~$X$.}
%	
%	\textcolor{blue}{The goal of this section is to generalize the dimension formulas of Corollary~\ref{cor;ifinjthenformcoho} to the span of special cycles in~$\CH^{\genus}(X)_\CC$.
%	We define such a subspace as
%	\[
%	\SC^\genus(X)
%	\coloneqq
%	\Big\langle
%	\{Z(T,\discel)\}\cdot\{\omega^*\}^{\genus-\rk T}
%	:
%	\text{$\discel\in D_\lattice^\genus$, $T\in q(\discel) + \halfint{\genus}$ and $T\geq 0$}
%	\Big\rangle_\CC
%	\subseteq\CH^\genus(X)_\CC.
%	\]
%	\begin{cor}\label{cor;dimspanCH}
%	Let~$\genus<(n+2)/4$.
%	If the Kudla--Millson lift~$\KMlift{r}$ is \emph{injective} for every~$r\leq \genus$, then
%	\[
%	\dim SC^\genus(X)=\dim \almcusp{k}{\genus}{\lattice}.
%	\]
%	In particular, under the assumptions on~$L$ of Corollary~\ref{cor;dimcoho1} and Corollary~\ref{cor;dimcoho2} we have respectively
%	\[
%	\dim \Pic(X)_\CC=\dim \almcusp{k}{1}{\lattice}
%	\qquad\text{and}\qquad
%	\dim SC^2(X)=\dim \almcusp{k}{2}{\lattice}.
%	\]	
%	\end{cor}}
%	\begin{proof}
%	\textcolor{blue}{
%	Let~$\KMliftchow{r}\colon S^k_{r,\lattice}\to \CH^r(X)_\CC$ be the linear map defined as~$f\mapsto \bigpet{f}{\{\KMtheta{r}\}}$, and let~$\lefschetz_{\{\omega\}}\colon\CH^r(X)_\CC\to\CH^{r+1}(X)_\CC$ be the map given by intersecting with the Hodge class~$\{\omega\}$.
%	From the commutative diagram}
%	\bes
%	\begin{tikzcd}
%	S^k_{r,\lattice} \arrow[dr, "\KMlift{r}"'] \arrow[r, "\KMliftchow{r}"]
%	&
%	\CH^r(X)_\CC \arrow[d]
%	\arrow[r, "\lefschetz_{\{\omega\}}^{\genus-r}"]
%	&
%	\CH^{\genus}(X)_\CC \arrow[d]
%	\\
%	&
%	H^{r,r}(X)
%	\arrow[r, "\lefschetz^{\genus-r}"]
%	&
%	H^{\genus,\genus}(X)
%	\end{tikzcd}
%	\ees
%	\textcolor{blue}{where the vertical arrows are the class maps, we deduce from the injectivity of~$\KMlift{r}$ that also~$\KMliftchow{r}$ is injective for every~$r\le\genus$.
%	This implies that~$\dim\KMliftchow{r}(S^k_{r,\lattice})=\dim S^k_{r,\lattice}$.
%	Furthermore, from the injectivity of~$\lefschetz^{\genus-r}$, and hence of~$\lefschetz^{\genus-r}\circ\KMlift{r}$, we deduce that also~$\lefschetz_{\{\omega\}}^{\genus-r}\circ\KMliftchow{r}$ is injective.	
%	It is easy to see that the images of the latter maps span~$\SC^\genus(X)$.
%	\\
%	---------------------------
%	\\
%	\textbf{Check of the last assertion:} It is enough to show that if~$T\in q(\discel) + \halfint{\genus}$, with~$\discel\in D_\lattice^\genus$, is \emph{positive definite}, then
%	\be\label{eq;ewwtpakmsjd}
%	\{Z(T,\discel)\} \in \Big\langle
%	\bigcup_{r\leq\genus} \lefschetz_{\{\omega\}}^{\genus-r}\circ\KMliftchow{r}(S^k_{r,\lattice})
%	\Big\rangle_\CC.
%	\ee
%	As in Theorem~\ref{thm;lefschetzdecgengenus}, the cuspidal part of~$\{\KMtheta{\genus}\}$ is
%	\[
%	\{\KMtheta{\genus}\}^{\operatorname{cusp}}
%	\coloneqq
%	\{\KMtheta{\genus}\} + \lefschetz_{\{\omega\}}
%	\circ\vklin{k}{\genus}{\genus-1}{\lattice}\big(\{\KMtheta{\gengenus-1}\}\big).
%	\]
%	Denote by~$c_{T,\discel}(f)$ the Fourier coefficient of index~$(T,\discel)$ of any modular form~$f$.
%	We may compute
%	\bas
%	\KMlift{\genus}(\poincv{\gengenus}{\discel}{T}{\lattice})
%	&=
%	\bigpet{\poincv{\gengenus}{\discel}{T}{\lattice}}{\{\KMtheta{\genus}\}}
%	=
%	\bigpet{\poincv{\gengenus}{\discel}{T}{\lattice}}{\{\KMtheta{\genus}\}^{\operatorname{cusp}}}
%	=
%	C\cdot
%	c_{T,\discel}\big(\{\KMtheta{\genus}\}^{\operatorname{cusp}}\big)
%	\\
%	&=
%	C\cdot\underbrace{c_{T,\discel}\big(\{\KMtheta{\genus}\}\big)}_{=\{Z(T,\discel)\}}
%	+
%	C\cdot\lefschetz_{\{\omega\}}\Big(
%	c_{T,\discel}\big(
%	\vklin{k}{\genus}{\genus-1}{\lattice}\big(\{\KMtheta{\gengenus-1}\}
%	\big)
%	\Big),
%	\eas
%	where~$C$ is a non-zero constant, from which we deduce that
%	\be\label{eq;ewwtpakmsjd2}
%	\{Z(T,\discel)\}
%	=
%	\KMlift{\genus}(\poincv{\gengenus}{\discel}{T}{\lattice})/C
%	-
%	\lefschetz_{\{\omega\}}\Big(
%	c_{T,\discel}\big(
%	\vklin{k}{\genus}{\genus-1}{\lattice}\big(\{\KMtheta{\gengenus-1}\}
%	\big)
%	\Big).
%	\ee
%	We now show~\eqref{eq;ewwtpakmsjd} by induction on~$g$.
%	\\
%	\textbf{Case~$g=1$.} The equation~\eqref{eq;ewwtpakmsjd2} reads~$\{Z(T,\discel)\}=\KMlift{1}(\poincv{1}{\discel}{T}{\lattice}) - c_{T,\discel}(E^k_{1,\lattice})\cdot\{\omega\}$, from which~\eqref{eq;ewwtpakmsjd} immediately follows.
%	\\
%	\textbf{Case~$\genus>1$.}
%	We see from~\eqref{eq;ewwtpakmsjd2} that to show~\eqref{eq;ewwtpakmsjd} it is enough to prove that the coefficient~$c_{T,\discel}\big(
%	\vklin{k}{\genus}{\genus-1}{\lattice}\big(\{\KMtheta{\gengenus-1}\}
%	\big)$ lies in the span on the right-hand side of~\eqref{eq;ewwtpakmsjd} where we replace~$\genus$ by~$\genus-1$.
%	This is assumed to hold by induction.\\
%	-----------------------\\
%%	Since linearly independent vectors in the image of the cycle map are images of linearly independent vectors,
%	We deduce that
%	\[
%	\SC^\genus(X)=\bigoplus_{r\leq\genus} \lefschetz_{\{\omega\}}^{\genus-r}\circ\KMliftchow{r} (S^k_{r,\lattice}).
%	\]
%	Corollary~\ref{cor;ifinjthenformcoho} implies that~$\dim\SC^\genus(X) = \dim H^{\genus,\genus}(X) = \dim\almcusp{k}{\genus}{\lattice}$.
%	}
%	\end{proof}
%	
%
%
%
%
%	
%	
%	
%	------------------------
%	
%	\textcolor{blue}{Old stuff:}
%	
%	In this section we illustrate an application of the Lefschetz decomposition of the Kudla--Millson theta function~$[\KMtheta{1}]$ to extend the Kudla--Millson lift from cusp forms to the whole space of modular forms.
%	To simplify the illustration, in this section we assume the lattice~$\lattice$ to be \emph{unimodular}.\\
%	
%	It is well-known that if~$V$ is a complex vector space endowed with a Hermitian inner product~$(\cdot{,}\cdot)$, then there is a canonical isomorphism of~$V$ with its dual~$V^*$ given by
%	\be\label{eq;identgenvecinprod}
%	V\longrightarrow V^*,\qquad v\longmapsto \Big(w\mapsto(w,v)\Big).
%	\ee
%	We are going to illustrate that the same is true for~$V=M^k_1$ endowed with a regularized Petersson inner product.
%	
%	The \emph{Petersson inner product} is the Hermitian product on~$S^k_1$ defined as
%	\bes
%	\pet{f}{g} = \int_{\SL_2(\ZZ)\backslash\HH}y^k f(\tau) \,\overline{g(\tau)}\,\frac{dx\,dy}{y^2},\qquad\text{for all~$f,g\in S^k_1$},
%	\ees
%	where~$\tau=x+iy\in\HH$.
%	It can be extended as a pairing between cusp forms and elements of~$M^k_1$, i.e.\ allowing Eisenstein series.
%	The defining integral does not converges whenever both~$f$ and~$g$ are Eisenstein series.
%	
%	It is possible to regularize such an integral, extending the Petersson inner product to the whole~$M^k_1$, by defining the value that the non-converging integral~$\pet{E^k_1}{E^k_1}$ should assume.
%	\begin{defi}
%	 A \emph{regularized Peterrson inner product} on~$M^k_1$ is an inner product~$\petregempty$ such that its restriction to~$S^k_1$ equals~$\petempty$ and makes~$S^k_1$ orthogonal to the Eisenstein series~$E^k_1$.
%	\end{defi}
%	Since an inner product is always positive definite, every regularization~$\petregempty$ is such that the value~$\petreg{E^k_1}{E^k_1}$ is always positive.
%	Note that to define a regularization it is enough to extend~$\petempty$ by choosing the (positive) value of~$\petreg{E^k_1}{E^k_1}$.
%	
%	Let~$k>2$.
%	For every~$n\in\ZZ_{>0}$, we denote by~$\poinc{k}{n}$ the~$n$-th Poincaré series normalized such that
%	\bes
%	\pet{f}{\poinc{k}{n}} = c_n(f),\qquad\text{for every~$f\in S^k_1$,}
%	\ees
%	where we denote by~$c_n(f)$ the~$n$-th Fourier coefficient of~$f$.
%	Recall that since~$\petempty$ is non-degenerate, the Poincaré series span the whole~$S^k_1$.
%	
%	We are now ready to illustrate how to identify~$M^k_1$ with its dual, whenever it is endowed with some regularized Petersson inner product.
%	We denote by~$c_n$ the coefficient extraction functional on~$M^k_1$ defined as~$f\mapsto c_n(f)$.
%	\begin{lemma}\label{eq;lemmaapplLDgen1}
%	Let~$\petregempty$ be a regularized Petersson inner product on~$M^k_1$.
%	The canonical identification of~$M^k_1$ and~$(M^k_1)^*$ induced by~$\petregempty$ is the bijective map~$\ident\colon (M^k_1)^*\to M^k_1$ defined by	
%	\bes
%	c_0\longmapsto\frac{E^k_1}{\petreg{E^k_1}{E^k_1}},\qquad\text{and}\qquad
%	c_n\longmapsto \poinc{k}{n} + \frac{c_n(E^k_1)}{\petreg{E^k_1}{E^k_1}} E^k_1,\quad \text{for every~$n\ge1$.}
%	\ees
%	\end{lemma}
%	\begin{proof}
%	If~$f\in M^k_1$, then~$f=a\cdot E^k_1 + g$ for some~$a\in \CC$ and~$g\in S^k_1$.
%	We may then rewrite
%	\bes
%	c_0(f)= a = \Bigpetreg{a\cdot E^k_1}{\frac{ E^k_1}{\petreg{E^k_1}{E^k_1}}} = \Bigpetreg{f}{\frac{ E^k_1}{\petreg{E^k_1}{E^k_1}}}.
%	\ees
%	Since this is true for every~$f\in M^k_1$, we deduce that~$c_0$ is dual to~$E^k_1/\petreg{E^k_1}{E^k_1}$ with respect to~$\petregempty$; cf.~\eqref{eq;identgenvecinprod}.
%	
%	The case of~$c_n$, with~$n>0$, is analogous.
%	We may compute
%	\bas
%	c_n(f) &= c_n(g) + a \cdot c_n(E^k_1)
%	=
%	\pet{g}{\poinc{k}{n}} + \Bigpetreg{a\cdot E^k_1}{\frac{ c_n(E^k_1)}{\petreg{E^k_1}{E^k_1}} E^k_1 }\\
%	&=
%	\Bigpetreg{ f }{ \poinc{k}{n} + \frac{ c_n(E^k_1)}{\petreg{E^k_1}{E^k_1}} E^k_1 }.
%	\eas
%	Since this is true for all~$f\in M^k_1$, we deduce that~$c_n$ is dual to~$\poinc{k}{n} + c_n(E^k_1)\cdot E^k_1 / \petreg{E^k_1}{E^k_1}$.
%	\end{proof}
%	In~\cite{bruiniermöller}, Bruinier and Möller studied the cone generated by the rational classes of Heegner divisors in terms of Fourier coefficients of elliptic modular forms.
%	This was possible thanks to the injectivity of the linear function
%	\bas
%	\psi\colon (M^k_1)^*\longrightarrow H^2(X,\CC),
%	\eas
%	defined as the one mapping~$c_m$ to the cohomology class~$[Z(m)]$ of the~$m$-th Heegner divisor, for every~$m\in\ZZ_{>0}$, and~$c_0\mapsto-[\omega]$.
%	
%	The identification provided by Lemma~\ref{eq;lemmaapplLDgen1} enables us to explain the relation of the map~$\psi$ with the Kudla--Millson lift, as follows.
%	Recall that the (cohomological) Kudla--Millson lift of genus~$1$ is the map defined as
%	\bas
%	\KMlift{1}\colon S^k_1\longrightarrow H^2(X,\CC),\qquad f\longmapsto \pet{f}{[\KMtheta{1}]}.
%	\eas
%	Clearly, any regularized Petersson inner product~$\petregempty$ induces an extension of the lift~$\KMlift{1}$ from~$S^k_1$ to the whole~$M^k_1$, which we call a \emph{regularized Kudla--Millson lift} and denote with the same symbol.
%	\begin{thm}
%	Let~$\petregempty$ be a regularized Petersson inner product, and let~${\ident\colon (M^k_1)^*\to M^k_1}$ be the identification of Lemma~\ref{eq;lemmaapplLDgen1}.
%	The induced regularized Kudla--Millson lift~$\KMlift{1}$ is such that the following diagram commutes.
%	\bes
%	\begin{tikzcd}
%	(M^k_1)^* \arrow[d, "\ident"] \arrow[r, "\psi"]
%	&
%	H^2(X,\CC) \arrow[d]
%	\\
%	M^k_1 \arrow[r, "\KMlift{1}"']
%	&
%	H^2(X,\CC)
%	\end{tikzcd}
%	\ees
%	Here the map~$H^2(X)\to H^2(X)$ on the right-hand side is given by \emph{conjugation}.
%	In particular, the map~$\psi$ is injective if and only if the regularized Kudla--Millson lift is injective.
%	\end{thm}
%	\begin{proof}
%	We show that~$\KMlift{1}\circ\ident (c_m) = \overline{\psi(c_m)}$ for every~$m\in\ZZ_{\ge0}$.
%	If~$m>0$, then by Lemma~\ref{eq;lemmaapplLDgen1} we deduce that
%	\bas
%	\KMlift{1}\circ\ident (c_m) =& \KMlift{1}(\poinc{k}{m}) + \frac{c_m(E^k_1)}{\petreg{E^k_1}{E^k_1}} \KMlift{1}(E^k_1)\\
%	=& \bigpet{\poinc{k}{m}}{[\KMtheta{1}]} + \frac{c_m(E^k_1)}{\petreg{E^k_1}{E^k_1}} \bigpetreg{E^k_1}{[\KMtheta{1}]}.
%	\eas
%	We proceed with the computation using the Lefschetz decomposition of~$[\KMtheta{1}]$ as illustrated in Section~\ref{sec;gen1}.
%	We deduce that
%	\bas
%	\KMlift{1}\circ\ident (c_m) &=
%	\bigpet{\poinc{k}{m}}{[\KMtheta{1}] + E^k_1\cdot[\omega]} + \frac{c_m(E^k_1)}{\petreg{E^k_1}{E^k_1}} \bigpetreg{E^k_1}{-E^k_1\cdot[\omega]}\\
%%	=& [H_n] + c_n(E^k_1)\cdot[\omega] - c_n(E^k_1)\cdot[\omega]\\
%	&= [Z(m)]_0 - c_m(E^k_1)\cdot [\omega] = [Z(m)] = \psi(c_m).
%	\eas
%	The case of~$n=0$ is analogous.
%	\\\textcolor{blue}{Here there might be a slight modification to be made, if we change the inner product on~$H^2(X)$ coming from previous sections.}
%	\end{proof}
%	\begin{cor}
%	We have that~$\KMlift{1}(E^k_1) = - \petreg{E^k_1}{E^k_1}\cdot[\omega]$.
%	In particular, we may rewrite the Bruinier--Funke relation of Borcherds' lift~$\Lambda^B$ and the Kudla--Millson lift as
%	\bes
%	\Lambda^B(f) = \KMlift{1}\Big(\xi_{1-n/2}(f) - a^+(0) E^k_1/\petreg{E^k_1}{E^k_1}\Big),
%	\ees
%	where~$a^+(0)$ is the constant Fourier coefficient of the weak Maass form~$f$.
%	\end{cor}
%	
%	
	
	
	\end{document}
% $ biblatex auxiliary file $
% $ biblatex bbl format version 3.1 $
% Do not modify the above lines!
%
% This is an auxiliary file used by the 'biblatex' package.
% This file may safely be deleted. It will be recreated as
% required.
%
\begingroup
\makeatletter
\@ifundefined{ver@biblatex.sty}
  {\@latex@error
     {Missing 'biblatex' package}
     {The bibliography requires the 'biblatex' package.}
      \aftergroup\endinput}
  {}
\endgroup

\datalist[entry]{nyt/global//global/global}
  \entry{noetherlefconj}{article}{}
    \name{author}{4}{}{%
      {{hash=BN}{%
         family={Bergeron},
         familyi={B\bibinitperiod},
         given={N.},
         giveni={N\bibinitperiod},
      }}%
      {{hash=LZ}{%
         family={Li},
         familyi={L\bibinitperiod},
         given={Z.},
         giveni={Z\bibinitperiod},
      }}%
      {{hash=MJ}{%
         family={Millson},
         familyi={M\bibinitperiod},
         given={J.},
         giveni={J\bibinitperiod},
      }}%
      {{hash=MC}{%
         family={Moeglin},
         familyi={M\bibinitperiod},
         given={C.},
         giveni={C\bibinitperiod},
      }}%
    }
    \strng{namehash}{BNLZMJMC1}
    \strng{fullhash}{BNLZMJMC1}
    \field{labelnamesource}{author}
    \field{labeltitlesource}{title}
    \field{labelalpha}{BLMM17}
    \field{sortinit}{B}
    \field{sortinithash}{B}
    \field{number}{2}
    \field{title}{The {N}oether-{L}efschetz conjecture and generalizations}
    \field{volume}{208}
    \field{journaltitle}{Invent. Math.}
    \field{year}{2017}
  \endentry

  \entry{doubling}{article}{}
    \name{author}{1}{}{%
      {{hash=BS}{%
         family={B{\"o}cherer},
         familyi={B\bibinitperiod},
         given={S.},
         giveni={S\bibinitperiod},
      }}%
    }
    \strng{namehash}{BS1}
    \strng{fullhash}{BS1}
    \field{labelnamesource}{author}
    \field{labeltitlesource}{title}
    \field{labelalpha}{B{\"o}c83}
    \field{sortinit}{B}
    \field{sortinithash}{B}
    \field{number}{1}
    \field{pages}{21\bibrangedash 46}
    \field{title}{\"{U}ber die {F}ourier-{J}acobi-{E}ntwicklung {S}iegelscher
  {E}isensteinreihen}
    \field{volume}{183}
    \field{journaltitle}{Math. Z.}
    \field{year}{1983}
  \endentry

  \entry{borcherds-grassmannians}{article}{}
    \name{author}{1}{}{%
      {{hash=BR}{%
         family={Borcherds},
         familyi={B\bibinitperiod},
         given={R.E.},
         giveni={R\bibinitperiod},
      }}%
    }
    \strng{namehash}{BR1}
    \strng{fullhash}{BR1}
    \field{labelnamesource}{author}
    \field{labeltitlesource}{title}
    \field{labelalpha}{Bor98}
    \field{sortinit}{B}
    \field{sortinithash}{B}
    \field{number}{3}
    \field{title}{Automorphic forms with singularities on {G}rassmannians}
    \field{volume}{132}
    \field{journaltitle}{Invent. Math.}
    \field{year}{1998}
  \endentry

  \entry{borcherds-reflection}{article}{}
    \name{author}{1}{}{%
      {{hash=BR}{%
         family={Borcherds},
         familyi={B\bibinitperiod},
         given={R.E.},
         giveni={R\bibinitperiod},
      }}%
    }
    \strng{namehash}{BR1}
    \strng{fullhash}{BR1}
    \field{labelnamesource}{author}
    \field{labeltitlesource}{title}
    \field{labelalpha}{Bor00}
    \field{sortinit}{B}
    \field{sortinithash}{B}
    \field{number}{2}
    \field{pages}{319\bibrangedash 366}
    \field{title}{Reflection groups of {L}orentzian lattices}
    \field{volume}{104}
    \field{journaltitle}{Duke Math. J.}
    \field{year}{2000}
  \endentry

  \entry{borelwallach}{book}{}
    \name{author}{2}{}{%
      {{hash=BA}{%
         family={Borel},
         familyi={B\bibinitperiod},
         given={A.},
         giveni={A\bibinitperiod},
      }}%
      {{hash=WN}{%
         family={Wallach},
         familyi={W\bibinitperiod},
         given={N.},
         giveni={N\bibinitperiod},
      }}%
    }
    \list{publisher}{1}{%
      {American Mathematical Society, Providence, RI}%
    }
    \strng{namehash}{BAWN1}
    \strng{fullhash}{BAWN1}
    \field{labelnamesource}{author}
    \field{labeltitlesource}{title}
    \field{labelalpha}{BW00}
    \field{sortinit}{B}
    \field{sortinithash}{B}
    \field{edition}{Second}
    \field{series}{Mathematical Surveys and Monographs}
    \field{title}{Continuous cohomology, discrete subgroups, and
  representations of reductive groups}
    \field{volume}{67}
    \field{year}{2000}
  \endentry

  \entry{bruinierbook}{book}{}
    \name{author}{1}{}{%
      {{hash=BJ}{%
         family={Bruinier},
         familyi={B\bibinitperiod},
         given={J.H.},
         giveni={J\bibinitperiod},
      }}%
    }
    \list{publisher}{1}{%
      {Springer-Verlag, Berlin}%
    }
    \strng{namehash}{BJ1}
    \strng{fullhash}{BJ1}
    \field{labelnamesource}{author}
    \field{labeltitlesource}{title}
    \field{labelalpha}{Bru02}
    \field{sortinit}{B}
    \field{sortinithash}{B}
    \field{series}{Lecture Notes in Mathematics}
    \field{title}{Borcherds products on {O}(2, {$l$}) and {C}hern classes of
  {H}eegner divisors}
    \field{volume}{1780}
    \field{year}{2002}
  \endentry

  \entry{bruinier-converse}{article}{}
    \name{author}{1}{}{%
      {{hash=BJ}{%
         family={Bruinier},
         familyi={B\bibinitperiod},
         given={J.H.},
         giveni={J\bibinitperiod},
      }}%
    }
    \strng{namehash}{BJ1}
    \strng{fullhash}{BJ1}
    \field{labelnamesource}{author}
    \field{labeltitlesource}{title}
    \field{labelalpha}{Bru14}
    \field{sortinit}{B}
    \field{sortinithash}{B}
    \field{pages}{315\bibrangedash 342}
    \field{title}{On the converse theorem for {B}orcherds products}
    \field{volume}{397}
    \field{journaltitle}{J. Algebra}
    \field{year}{2014}
  \endentry

  \entry{bruinierfunke}{incollection}{}
    \name{author}{2}{}{%
      {{hash=BJ}{%
         family={Bruinier},
         familyi={B\bibinitperiod},
         given={J.H.},
         giveni={J\bibinitperiod},
      }}%
      {{hash=FJ}{%
         family={Funke},
         familyi={F\bibinitperiod},
         given={J.},
         giveni={J\bibinitperiod},
      }}%
    }
    \list{publisher}{1}{%
      {Cambridge Univ. Press, Cambridge}%
    }
    \strng{namehash}{BJFJ1}
    \strng{fullhash}{BJFJ1}
    \field{labelnamesource}{author}
    \field{labeltitlesource}{title}
    \field{labelalpha}{BF10}
    \field{sortinit}{B}
    \field{sortinithash}{B}
    \field{booktitle}{Moonshine: the first quarter century and beyond}
    \field{pages}{12\bibrangedash 39}
    \field{series}{London Math. Soc. Lecture Note Ser.}
    \field{title}{On the injectivity of the {K}udla-{M}illson lift and
  surjectivity of the {B}orcherds lift}
    \field{volume}{372}
    \field{year}{2010}
  \endentry

  \entry{bruiniermöller}{article}{}
    \name{author}{2}{}{%
      {{hash=BJ}{%
         family={Bruinier},
         familyi={B\bibinitperiod},
         given={J.H.},
         giveni={J\bibinitperiod},
      }}%
      {{hash=MM}{%
         family={M\"{o}ller},
         familyi={M\bibinitperiod},
         given={M.},
         giveni={M\bibinitperiod},
      }}%
    }
    \strng{namehash}{BJMM1}
    \strng{fullhash}{BJMM1}
    \field{labelnamesource}{author}
    \field{labeltitlesource}{title}
    \field{labelalpha}{BM19}
    \field{sortinit}{B}
    \field{sortinithash}{B}
    \field{number}{3}
    \field{title}{Cones of {H}eegner divisors}
    \field{volume}{28}
    \field{journaltitle}{J. Algebraic Geom.}
    \field{year}{2019}
  \endentry

  \entry{bruinierraum}{article}{}
    \name{author}{2}{}{%
      {{hash=BJ}{%
         family={Bruinier},
         familyi={B\bibinitperiod},
         given={J.H.},
         giveni={J\bibinitperiod},
      }}%
      {{hash=WRM}{%
         family={Westerholt-Raum},
         familyi={W\bibinithyphendelim R\bibinitperiod},
         given={M.},
         giveni={M\bibinitperiod},
      }}%
    }
    \strng{namehash}{BJWRM1}
    \strng{fullhash}{BJWRM1}
    \field{labelnamesource}{author}
    \field{labeltitlesource}{title}
    \field{labelalpha}{BWR15}
    \field{sortinit}{B}
    \field{sortinithash}{B}
    \field{pages}{e7, 30}
    \field{title}{Kudla's modularity conjecture and formal {F}ourier-{J}acobi
  series}
    \field{volume}{3}
    \field{journaltitle}{Forum Math. Pi}
    \field{year}{2015}
  \endentry

  \entry{bruinier-yang09}{article}{}
    \name{author}{2}{}{%
      {{hash=BJ}{%
         family={Bruinier},
         familyi={B\bibinitperiod},
         given={J.H.},
         giveni={J\bibinitperiod},
      }}%
      {{hash=YT}{%
         family={Yang},
         familyi={Y\bibinitperiod},
         given={T.},
         giveni={T\bibinitperiod},
      }}%
    }
    \strng{namehash}{BJYT1}
    \strng{fullhash}{BJYT1}
    \field{labelnamesource}{author}
    \field{labeltitlesource}{title}
    \field{labelalpha}{BY09}
    \field{sortinit}{B}
    \field{sortinithash}{B}
    \field{number}{3}
    \field{pages}{631\bibrangedash 681}
    \field{title}{Faltings heights of {CM} cycles and derivatives of
  {$L$}-functions}
    \field{volume}{177}
    \field{journaltitle}{Invent. Math.}
    \field{year}{2009}
  \endentry

  \entry{ChenevierLannes}{book}{}
    \name{author}{2}{}{%
      {{hash=CG}{%
         family={Chenevier},
         familyi={C\bibinitperiod},
         given={G.},
         giveni={G\bibinitperiod},
      }}%
      {{hash=LJ}{%
         family={Lannes},
         familyi={L\bibinitperiod},
         given={J.},
         giveni={J\bibinitperiod},
      }}%
    }
    \list{publisher}{1}{%
      {Springer, Cham}%
    }
    \strng{namehash}{CGLJ1}
    \strng{fullhash}{CGLJ1}
    \field{labelnamesource}{author}
    \field{labeltitlesource}{title}
    \field{labelalpha}{CL19}
    \field{sortinit}{C}
    \field{sortinithash}{C}
    \field{note}{Kneser neighbors of Niemeier lattices, Translated from the
  French by Reinie Ern\'{e}}
    \field{pages}{xxi+417}
    \field{series}{Ergebnisse der Mathematik und ihrer Grenzgebiete. 3. Folge.
  A Series of Modern Surveys in Mathematics [Results in Mathematics and Related
  Areas. 3rd Series. A Series of Modern Surveys in Mathematics]}
    \field{title}{Automorphic forms and even unimodular lattices}
    \field{volume}{69}
    \field{year}{2019}
  \endentry

  \entry{christian}{article}{}
    \name{author}{1}{}{%
      {{hash=CU}{%
         family={Christian},
         familyi={C\bibinitperiod},
         given={U.},
         giveni={U\bibinitperiod},
      }}%
    }
    \strng{namehash}{CU1}
    \strng{fullhash}{CU1}
    \field{labelnamesource}{author}
    \field{labeltitlesource}{title}
    \field{labelalpha}{Chr62}
    \field{sortinit}{C}
    \field{sortinithash}{C}
    \field{pages}{257\bibrangedash 307}
    \field{title}{\"{U}ber {H}ilbert-{S}iegelsche {M}odulformen und
  {P}oincar\'{e}sche {R}eihen}
    \field{volume}{148}
    \field{journaltitle}{Math. Ann.}
    \field{year}{1962}
  \endentry

  \entry{dai}{incollection}{}
    \name{author}{1}{}{%
      {{hash=DX}{%
         family={Dai},
         familyi={D\bibinitperiod},
         given={X.},
         giveni={X\bibinitperiod},
      }}%
    }
    \list{publisher}{1}{%
      {Cambridge Univ. Press, Cambridge}%
    }
    \strng{namehash}{DX1}
    \strng{fullhash}{DX1}
    \field{labelnamesource}{author}
    \field{labeltitlesource}{title}
    \field{labelalpha}{Dai11}
    \field{sortinit}{D}
    \field{sortinithash}{D}
    \field{booktitle}{Topology of stratified spaces}
    \field{series}{Math. Sci. Res. Inst. Publ.}
    \field{title}{An introduction to {$L^2$} cohomology}
    \field{volume}{58}
    \field{year}{2011}
  \endentry

  \entry{fischer}{book}{}
    \name{author}{1}{}{%
      {{hash=FJ}{%
         family={Fischer},
         familyi={F\bibinitperiod},
         given={J.},
         giveni={J\bibinitperiod},
      }}%
    }
    \list{publisher}{1}{%
      {Springer-Verlag, Berlin}%
    }
    \strng{namehash}{FJ1}
    \strng{fullhash}{FJ1}
    \field{labelnamesource}{author}
    \field{labeltitlesource}{title}
    \field{labelalpha}{Fis87}
    \field{sortinit}{F}
    \field{sortinithash}{F}
    \field{pages}{iv+184}
    \field{series}{Lecture Notes in Mathematics}
    \field{title}{An approach to the {S}elberg trace formula via the {S}elberg
  zeta-function}
    \field{volume}{1253}
    \field{year}{1987}
  \endentry

  \entry{freitaghermann}{article}{}
    \name{author}{2}{}{%
      {{hash=FE}{%
         family={Freitag},
         familyi={F\bibinitperiod},
         given={E.},
         giveni={E\bibinitperiod},
      }}%
      {{hash=HC}{%
         family={Hermann},
         familyi={H\bibinitperiod},
         given={C.F.},
         giveni={C\bibinitperiod},
      }}%
    }
    \strng{namehash}{FEHC1}
    \strng{fullhash}{FEHC1}
    \field{labelnamesource}{author}
    \field{labeltitlesource}{title}
    \field{labelalpha}{FH00}
    \field{sortinit}{F}
    \field{sortinithash}{F}
    \field{number}{2}
    \field{pages}{203\bibrangedash 287}
    \field{title}{Some modular varieties of low dimension}
    \field{volume}{152}
    \field{journaltitle}{Adv. Math.}
    \field{year}{2000}
  \endentry

  \entry{funkemillsonhyp}{article}{}
    \name{author}{2}{}{%
      {{hash=FJ}{%
         family={Funke},
         familyi={F\bibinitperiod},
         given={J.},
         giveni={J\bibinitperiod},
      }}%
      {{hash=MJ}{%
         family={Millson},
         familyi={M\bibinitperiod},
         given={J.},
         giveni={J\bibinitperiod},
      }}%
    }
    \strng{namehash}{FJMJ1}
    \strng{fullhash}{FJMJ1}
    \field{labelnamesource}{author}
    \field{labeltitlesource}{title}
    \field{labelalpha}{FM02}
    \field{sortinit}{F}
    \field{sortinithash}{F}
    \field{number}{4}
    \field{pages}{409\bibrangedash 444}
    \field{title}{Cycles in hyperbolic manifolds of non-compact type and
  {F}ourier coefficients of {S}iegel modular forms}
    \field{volume}{107}
    \field{journaltitle}{Manuscripta Math.}
    \field{year}{2002}
  \endentry

  \entry{gaffney}{article}{}
    \name{author}{1}{}{%
      {{hash=GM}{%
         family={Gaffney},
         familyi={G\bibinitperiod},
         given={M.P.},
         giveni={M\bibinitperiod},
      }}%
    }
    \strng{namehash}{GM1}
    \strng{fullhash}{GM1}
    \field{labelnamesource}{author}
    \field{labeltitlesource}{title}
    \field{labelalpha}{Gaf54}
    \field{sortinit}{G}
    \field{sortinithash}{G}
    \field{title}{A special {S}tokes's theorem for complete {R}iemannian
  manifolds}
    \field{volume}{60}
    \field{journaltitle}{Ann. of Math. (2)}
    \field{year}{1954}
  \endentry

  \entry{GHS-compositio}{article}{}
    \name{author}{3}{}{%
      {{hash=GV}{%
         family={Gritsenko},
         familyi={G\bibinitperiod},
         given={V.},
         giveni={V\bibinitperiod},
      }}%
      {{hash=HK}{%
         family={Hulek},
         familyi={H\bibinitperiod},
         given={K.},
         giveni={K\bibinitperiod},
      }}%
      {{hash=SG}{%
         family={Sankaran},
         familyi={S\bibinitperiod},
         given={G.K.},
         giveni={G\bibinitperiod},
      }}%
    }
    \strng{namehash}{GVHKSG1}
    \strng{fullhash}{GVHKSG1}
    \field{labelnamesource}{author}
    \field{labeltitlesource}{title}
    \field{labelalpha}{GHS10}
    \field{sortinit}{G}
    \field{sortinithash}{G}
    \field{number}{2}
    \field{pages}{404\bibrangedash 434}
    \field{title}{Moduli spaces of irreducible symplectic manifolds}
    \field{volume}{146}
    \field{journaltitle}{Compos. Math.}
    \field{year}{2010}
  \endentry

  \entry{harriszucker}{article}{}
    \name{author}{2}{}{%
      {{hash=HM}{%
         family={Harris},
         familyi={H\bibinitperiod},
         given={M.},
         giveni={M\bibinitperiod},
      }}%
      {{hash=ZS}{%
         family={Zucker},
         familyi={Z\bibinitperiod},
         given={S.},
         giveni={S\bibinitperiod},
      }}%
    }
    \strng{namehash}{HMZS1}
    \strng{fullhash}{HMZS1}
    \field{labelnamesource}{author}
    \field{labeltitlesource}{title}
    \field{labelalpha}{HZ01}
    \field{sortinit}{H}
    \field{sortinithash}{H}
    \field{number}{85}
    \field{title}{Boundary cohomology of {S}himura varieties. {III}. {C}oherent
  cohomology on higher-rank boundary strata and applications to {H}odge theory}
    \field{journaltitle}{M\'{e}m. Soc. Math. Fr. (N.S.)}
    \field{year}{2001}
  \endentry

  \entry{HeHoffman}{article}{}
    \name{author}{2}{}{%
      {{hash=HH}{%
         family={He},
         familyi={H\bibinitperiod},
         given={H.},
         giveni={H\bibinitperiod},
      }}%
      {{hash=HJ}{%
         family={Hoffman},
         familyi={H\bibinitperiod},
         given={J.W.},
         giveni={J\bibinitperiod},
      }}%
    }
    \strng{namehash}{HHHJ1}
    \strng{fullhash}{HHHJ1}
    \field{labelnamesource}{author}
    \field{labeltitlesource}{title}
    \field{labelalpha}{HH12}
    \field{sortinit}{H}
    \field{sortinithash}{H}
    \field{number}{3}
    \field{pages}{769\bibrangedash 801}
    \field{title}{Picard groups of {S}iegel modular 3-folds and
  {$\theta$}-liftings}
    \field{volume}{22}
    \field{journaltitle}{J. Lie Theory}
    \field{year}{2012}
  \endentry

  \entry{hirzebruch-zagier}{article}{}
    \name{author}{2}{}{%
      {{hash=HF}{%
         family={Hirzebruch},
         familyi={H\bibinitperiod},
         given={F.},
         giveni={F\bibinitperiod},
      }}%
      {{hash=ZD}{%
         family={Zagier},
         familyi={Z\bibinitperiod},
         given={D.},
         giveni={D\bibinitperiod},
      }}%
    }
    \strng{namehash}{HFZD1}
    \strng{fullhash}{HFZD1}
    \field{labelnamesource}{author}
    \field{labeltitlesource}{title}
    \field{labelalpha}{HZ76}
    \field{sortinit}{H}
    \field{sortinithash}{H}
    \field{pages}{57\bibrangedash 113}
    \field{title}{Intersection numbers of curves on {H}ilbert modular surfaces
  and modular forms of {N}ebentypus}
    \field{volume}{36}
    \field{journaltitle}{Invent. Math.}
    \field{year}{1976}
  \endentry

  \entry{huybrechts}{book}{}
    \name{author}{1}{}{%
      {{hash=HD}{%
         family={Huybrechts},
         familyi={H\bibinitperiod},
         given={D.},
         giveni={D\bibinitperiod},
      }}%
    }
    \list{publisher}{1}{%
      {Springer-Verlag, Berlin}%
    }
    \strng{namehash}{HD1}
    \strng{fullhash}{HD1}
    \field{labelnamesource}{author}
    \field{labeltitlesource}{title}
    \field{labelalpha}{Huy05}
    \field{sortinit}{H}
    \field{sortinithash}{H}
    \field{note}{An introduction}
    \field{series}{Universitext}
    \field{title}{Complex geometry}
    \field{year}{2005}
  \endentry

  \entry{kieferzuffetti}{misc}{}
    \name{author}{2}{}{%
      {{hash=KP}{%
         family={Kiefer},
         familyi={K\bibinitperiod},
         given={P.},
         giveni={P\bibinitperiod},
      }}%
      {{hash=ZR}{%
         family={Zuffetti},
         familyi={Z\bibinitperiod},
         given={R.},
         giveni={R\bibinitperiod},
      }}%
    }
    \strng{namehash}{KPZR1}
    \strng{fullhash}{KPZR1}
    \field{labelnamesource}{author}
    \field{labeltitlesource}{title}
    \field{labelalpha}{KZ23}
    \field{sortinit}{K}
    \field{sortinithash}{K}
    \verb{eprint}
    \verb 2307.15809
    \endverb
    \field{title}{The Kudla-Millson lift of Siegel cusp forms}
    \field{eprinttype}{arXiv}
    \field{eprintclass}{math.NT}
    \field{year}{2023}
  \endentry

  \entry{kitaoka}{book}{}
    \name{author}{1}{}{%
      {{hash=KY}{%
         family={Kitaoka},
         familyi={K\bibinitperiod},
         given={Y.},
         giveni={Y\bibinitperiod},
      }}%
    }
    \list{publisher}{1}{%
      {Cambridge University Press, Cambridge}%
    }
    \strng{namehash}{KY1}
    \strng{fullhash}{KY1}
    \field{labelnamesource}{author}
    \field{labeltitlesource}{title}
    \field{labelalpha}{Kit93}
    \field{sortinit}{K}
    \field{sortinithash}{K}
    \field{pages}{x+268}
    \field{series}{Cambridge Tracts in Mathematics}
    \field{title}{Arithmetic of quadratic forms}
    \field{volume}{106}
    \field{year}{1993}
  \endentry

  \entry{klingen-defKES}{article}{}
    \name{author}{1}{}{%
      {{hash=KH}{%
         family={Klingen},
         familyi={K\bibinitperiod},
         given={H.},
         giveni={H\bibinitperiod},
      }}%
    }
    \strng{namehash}{KH1}
    \strng{fullhash}{KH1}
    \field{labelnamesource}{author}
    \field{labeltitlesource}{title}
    \field{labelalpha}{Kli67}
    \field{sortinit}{K}
    \field{sortinithash}{K}
    \field{pages}{30\bibrangedash 43}
    \field{title}{Zum {D}arstellungssatz f\"{u}r {S}iegelsche {M}odulformen}
    \field{volume}{102}
    \field{journaltitle}{Math. Z.}
    \field{year}{1967}
  \endentry

  \entry{klingen}{book}{}
    \name{author}{1}{}{%
      {{hash=KH}{%
         family={Klingen},
         familyi={K\bibinitperiod},
         given={H.},
         giveni={H\bibinitperiod},
      }}%
    }
    \list{publisher}{1}{%
      {Cambridge University Press, Cambridge}%
    }
    \strng{namehash}{KH1}
    \strng{fullhash}{KH1}
    \field{labelnamesource}{author}
    \field{labeltitlesource}{title}
    \field{labelalpha}{Kli90}
    \field{sortinit}{K}
    \field{sortinithash}{K}
    \field{series}{Cambridge Studies in Advanced Mathematics}
    \field{title}{Introductory lectures on {S}iegel modular forms}
    \field{volume}{20}
    \field{year}{1990}
  \endentry

  \entry{kneser}{book}{}
    \name{author}{1}{}{%
      {{hash=KM}{%
         family={Kneser},
         familyi={K\bibinitperiod},
         given={M.},
         giveni={M\bibinitperiod},
      }}%
    }
    \list{publisher}{1}{%
      {Springer-Verlag, Berlin}%
    }
    \strng{namehash}{KM1}
    \strng{fullhash}{KM1}
    \field{labelnamesource}{author}
    \field{labeltitlesource}{title}
    \field{labelalpha}{Kne02}
    \field{sortinit}{K}
    \field{sortinithash}{K}
    \field{note}{Revised and edited in collaboration with Rudolf Scharlau}
    \field{pages}{viii+164}
    \field{title}{Quadratische {F}ormen}
    \field{year}{2002}
  \endentry

  \entry{kudla-algcycles}{article}{}
    \name{author}{1}{}{%
      {{hash=KS}{%
         family={Kudla},
         familyi={K\bibinitperiod},
         given={S.S.},
         giveni={S\bibinitperiod},
      }}%
    }
    \strng{namehash}{KS1}
    \strng{fullhash}{KS1}
    \field{labelnamesource}{author}
    \field{labeltitlesource}{title}
    \field{labelalpha}{Kud97}
    \field{sortinit}{K}
    \field{sortinithash}{K}
    \field{number}{1}
    \field{title}{Algebraic cycles on {S}himura varieties of orthogonal type}
    \field{volume}{86}
    \field{journaltitle}{Duke Math. J.}
    \field{year}{1997}
  \endentry

  \entry{kudla-integrals}{article}{}
    \name{author}{1}{}{%
      {{hash=KS}{%
         family={Kudla},
         familyi={K\bibinitperiod},
         given={S.S.},
         giveni={S\bibinitperiod},
      }}%
    }
    \strng{namehash}{KS1}
    \strng{fullhash}{KS1}
    \field{labelnamesource}{author}
    \field{labeltitlesource}{title}
    \field{labelalpha}{Kud03}
    \field{sortinit}{K}
    \field{sortinithash}{K}
    \field{number}{3}
    \field{title}{Integrals of {B}orcherds forms}
    \field{volume}{137}
    \field{journaltitle}{Compositio Math.}
    \field{year}{2003}
  \endentry

  \entry{kudla-speccyc}{incollection}{}
    \name{author}{1}{}{%
      {{hash=KS}{%
         family={Kudla},
         familyi={K\bibinitperiod},
         given={S.S.},
         giveni={S\bibinitperiod},
      }}%
    }
    \list{publisher}{1}{%
      {Cambridge Univ. Press, Cambridge}%
    }
    \strng{namehash}{KS1}
    \strng{fullhash}{KS1}
    \field{labelnamesource}{author}
    \field{labeltitlesource}{title}
    \field{labelalpha}{Kud04}
    \field{sortinit}{K}
    \field{sortinithash}{K}
    \field{booktitle}{Heegner points and {R}ankin {$L$}-series}
    \field{series}{Math. Sci. Res. Inst. Publ.}
    \field{title}{Special cycles and derivatives of {E}isenstein series}
    \field{volume}{49}
    \field{year}{2004}
  \endentry

  \entry{kudla-extensions}{incollection}{}
    \name{author}{1}{}{%
      {{hash=KS}{%
         family={Kudla},
         familyi={K\bibinitperiod},
         given={S.S.},
         giveni={S\bibinitperiod},
      }}%
    }
    \list{publisher}{1}{%
      {Birkh\"{a}user Boston, Boston, MA}%
    }
    \strng{namehash}{KS1}
    \strng{fullhash}{KS1}
    \field{labelnamesource}{author}
    \field{labeltitlesource}{title}
    \field{labelalpha}{Kud08}
    \field{sortinit}{K}
    \field{sortinithash}{K}
    \field{booktitle}{Eisenstein series and applications}
    \field{pages}{205\bibrangedash 237}
    \field{series}{Progr. Math.}
    \field{title}{Some extensions of the {S}iegel-{W}eil formula}
    \field{volume}{258}
    \field{year}{2008}
  \endentry

  \entry{kudlamillson-harmonic1}{article}{}
    \name{author}{2}{}{%
      {{hash=KS}{%
         family={Kudla},
         familyi={K\bibinitperiod},
         given={S.S.},
         giveni={S\bibinitperiod},
      }}%
      {{hash=MJ}{%
         family={Millson},
         familyi={M\bibinitperiod},
         given={J.J.},
         giveni={J\bibinitperiod},
      }}%
    }
    \strng{namehash}{KSMJ1}
    \strng{fullhash}{KSMJ1}
    \field{labelnamesource}{author}
    \field{labeltitlesource}{title}
    \field{labelalpha}{KM86}
    \field{sortinit}{K}
    \field{sortinithash}{K}
    \field{number}{3}
    \field{pages}{353\bibrangedash 378}
    \field{title}{The theta correspondence and harmonic forms. {I}}
    \field{volume}{274}
    \field{journaltitle}{Math. Ann.}
    \field{year}{1986}
  \endentry

  \entry{kudlamillson-harmonic2}{article}{}
    \name{author}{2}{}{%
      {{hash=KS}{%
         family={Kudla},
         familyi={K\bibinitperiod},
         given={S.S.},
         giveni={S\bibinitperiod},
      }}%
      {{hash=MJ}{%
         family={Millson},
         familyi={M\bibinitperiod},
         given={J.J.},
         giveni={J\bibinitperiod},
      }}%
    }
    \strng{namehash}{KSMJ1}
    \strng{fullhash}{KSMJ1}
    \field{labelnamesource}{author}
    \field{labeltitlesource}{title}
    \field{labelalpha}{KM87}
    \field{sortinit}{K}
    \field{sortinithash}{K}
    \field{number}{2}
    \field{pages}{267\bibrangedash 314}
    \field{title}{The theta correspondence and harmonic forms. {II}}
    \field{volume}{277}
    \field{journaltitle}{Math. Ann.}
    \field{year}{1987}
  \endentry

  \entry{kudlamillson-tubes}{article}{}
    \name{author}{2}{}{%
      {{hash=KS}{%
         family={Kudla},
         familyi={K\bibinitperiod},
         given={S.S.},
         giveni={S\bibinitperiod},
      }}%
      {{hash=MJ}{%
         family={Millson},
         familyi={M\bibinitperiod},
         given={J.J.},
         giveni={J\bibinitperiod},
      }}%
    }
    \strng{namehash}{KSMJ1}
    \strng{fullhash}{KSMJ1}
    \field{labelnamesource}{author}
    \field{labeltitlesource}{title}
    \field{labelalpha}{KM88}
    \field{sortinit}{K}
    \field{sortinithash}{K}
    \field{number}{1}
    \field{title}{Tubes, cohomology with growth conditions and an application
  to the theta correspondence}
    \field{volume}{40}
    \field{journaltitle}{Canad. J. Math.}
    \field{year}{1988}
  \endentry

  \entry{kudlamillson-intersection}{article}{}
    \name{author}{2}{}{%
      {{hash=KS}{%
         family={Kudla},
         familyi={K\bibinitperiod},
         given={S.S.},
         giveni={S\bibinitperiod},
      }}%
      {{hash=MJ}{%
         family={Millson},
         familyi={M\bibinitperiod},
         given={J.J.},
         giveni={J\bibinitperiod},
      }}%
    }
    \strng{namehash}{KSMJ1}
    \strng{fullhash}{KSMJ1}
    \field{labelnamesource}{author}
    \field{labeltitlesource}{title}
    \field{labelalpha}{KM90}
    \field{sortinit}{K}
    \field{sortinithash}{K}
    \field{number}{71}
    \field{pages}{121\bibrangedash 172}
    \field{title}{Intersection numbers of cycles on locally symmetric spaces
  and {F}ourier coefficients of holomorphic modular forms in several complex
  variables}
    \field{journaltitle}{Inst. Hautes \'{E}tudes Sci. Publ. Math.}
    \field{year}{1990}
  \endentry

  \entry{looijenga}{article}{}
    \name{author}{1}{}{%
      {{hash=LE}{%
         family={Looijenga},
         familyi={L\bibinitperiod},
         given={E.},
         giveni={E\bibinitperiod},
      }}%
    }
    \strng{namehash}{LE1}
    \strng{fullhash}{LE1}
    \field{labelnamesource}{author}
    \field{labeltitlesource}{title}
    \field{labelalpha}{Loo88}
    \field{sortinit}{L}
    \field{sortinithash}{L}
    \field{number}{1}
    \field{title}{{$L^2$}-cohomology of locally symmetric varieties}
    \field{volume}{67}
    \field{journaltitle}{Compositio Math.}
    \field{year}{1988}
  \endentry

  \entry{maxim}{book}{}
    \name{author}{1}{}{%
      {{hash=ML}{%
         family={Maxim},
         familyi={M\bibinitperiod},
         given={L.G.},
         giveni={L\bibinitperiod},
      }}%
    }
    \list{publisher}{1}{%
      {Springer, Cham}%
    }
    \strng{namehash}{ML1}
    \strng{fullhash}{ML1}
    \field{labelnamesource}{author}
    \field{labeltitlesource}{title}
    \field{labelalpha}{Max19}
    \field{sortinit}{M}
    \field{sortinithash}{M}
    \field{series}{Graduate Texts in Mathematics}
    \field{title}{Intersection homology \& perverse sheaves---with applications
  to singularities}
    \field{volume}{281}
    \field{year}{[2019] \copyright 2019}
  \endentry

  \entry{metzlerzuffetti}{misc}{}
    \name{author}{2}{}{%
      {{hash=MI}{%
         family={Metzler},
         familyi={M\bibinitperiod},
         given={I.},
         giveni={I\bibinitperiod},
      }}%
      {{hash=ZR}{%
         family={Zuffetti},
         familyi={Z\bibinitperiod},
         given={R.},
         giveni={R\bibinitperiod},
      }}%
    }
    \strng{namehash}{MIZR1}
    \strng{fullhash}{MIZR1}
    \field{labelnamesource}{author}
    \field{labeltitlesource}{title}
    \field{labelalpha}{MZ23}
    \field{sortinit}{M}
    \field{sortinithash}{M}
    \verb{eprint}
    \verb 2312.00572
    \endverb
    \field{title}{Injectivity of the genus 1 Kudla-Millson lift on locally
  symmetric spaces}
    \field{eprinttype}{arXiv}
    \field{eprintclass}{math.NT}
    \field{year}{2023}
  \endentry

  \entry{namikawa}{book}{}
    \name{author}{1}{}{%
      {{hash=NY}{%
         family={Namikawa},
         familyi={N\bibinitperiod},
         given={Y.},
         giveni={Y\bibinitperiod},
      }}%
    }
    \list{publisher}{1}{%
      {Springer, Berlin}%
    }
    \strng{namehash}{NY1}
    \strng{fullhash}{NY1}
    \field{labelnamesource}{author}
    \field{labeltitlesource}{title}
    \field{labelalpha}{Nam80}
    \field{sortinit}{N}
    \field{sortinithash}{N}
    \field{pages}{viii+162}
    \field{title}{Toroidal compactification of {S}iegel spaces}
    \field{volume}{812}
    \field{year}{1980}
  \endentry

  \entry{peterson}{misc}{}
    \name{author}{1}{}{%
      {{hash=PA}{%
         family={Peterson},
         familyi={P\bibinitperiod},
         given={A.},
         giveni={A\bibinitperiod},
      }}%
    }
    \strng{namehash}{PA1}
    \strng{fullhash}{PA1}
    \field{labelnamesource}{author}
    \field{labeltitlesource}{title}
    \field{labelalpha}{Pet15}
    \field{sortinit}{P}
    \field{sortinithash}{P}
    \verb{eprint}
    \verb 1511.06887
    \endverb
    \field{title}{Modular forms on the moduli space of polarised K3 surfaces}
    \field{eprinttype}{arXiv}
    \field{year}{2015}
  \endentry

  \entry{saperstern}{article}{}
    \name{author}{2}{}{%
      {{hash=SL}{%
         family={Saper},
         familyi={S\bibinitperiod},
         given={L.},
         giveni={L\bibinitperiod},
      }}%
      {{hash=SM}{%
         family={Stern},
         familyi={S\bibinitperiod},
         given={M.},
         giveni={M\bibinitperiod},
      }}%
    }
    \strng{namehash}{SLSM1}
    \strng{fullhash}{SLSM1}
    \field{labelnamesource}{author}
    \field{labeltitlesource}{title}
    \field{labelalpha}{SS90}
    \field{sortinit}{S}
    \field{sortinithash}{S}
    \field{number}{1}
    \field{title}{{$L_2$}-cohomology of arithmetic varieties}
    \field{volume}{132}
    \field{journaltitle}{Ann. of Math. (2)}
    \field{year}{1990}
  \endentry

  \entry{saperzucker}{incollection}{}
    \name{author}{2}{}{%
      {{hash=SL}{%
         family={Saper},
         familyi={S\bibinitperiod},
         given={L.},
         giveni={L\bibinitperiod},
      }}%
      {{hash=ZS}{%
         family={Zucker},
         familyi={Z\bibinitperiod},
         given={S.},
         giveni={S\bibinitperiod},
      }}%
    }
    \strng{namehash}{SLZS1}
    \strng{fullhash}{SLZS1}
    \field{labelnamesource}{author}
    \field{labeltitlesource}{title}
    \field{labelalpha}{SZ91}
    \field{sortinit}{S}
    \field{sortinithash}{S}
    \field{booktitle}{Several complex variables and complex geometry, {P}art 2
  ({S}anta {C}ruz, {CA}, 1989)}
    \field{series}{Proc. Sympos. Pure Math.}
    \field{title}{An introduction to {$L^2$}-cohomology}
    \field{volume}{52}
    \field{year}{1991}
  \endentry

  \entry{scheit}{article}{}
    \name{author}{1}{}{%
      {{hash=SN}{%
         family={Scheithauer},
         familyi={S\bibinitperiod},
         given={N.R.},
         giveni={N\bibinitperiod},
      }}%
    }
    \strng{namehash}{SN1}
    \strng{fullhash}{SN1}
    \field{labelnamesource}{author}
    \field{labeltitlesource}{title}
    \field{labelalpha}{Sch15}
    \field{sortinit}{S}
    \field{sortinithash}{S}
    \field{pages}{1\bibrangedash 43}
    \field{title}{Some constructions of modular forms for the {W}eil
  representation of {${\rm SL}_2(\Bbb{Z})$}}
    \field{volume}{220}
    \field{journaltitle}{Nagoya Math. J.}
    \field{year}{2015}
  \endentry

  \entry{stein-converse}{misc}{}
    \name{author}{1}{}{%
      {{hash=SO}{%
         family={Stein},
         familyi={S\bibinitperiod},
         given={Oliver},
         giveni={O\bibinitperiod},
      }}%
    }
    \strng{namehash}{SO1}
    \strng{fullhash}{SO1}
    \field{labelnamesource}{author}
    \field{labeltitlesource}{title}
    \field{labelalpha}{Ste23}
    \field{sortinit}{S}
    \field{sortinithash}{S}
    \verb{eprint}
    \verb 2306.17660
    \endverb
    \field{title}{A converse theorem for Borcherds products and the injectivity
  of the Kudla-Millson theta lift}
    \field{eprinttype}{arXiv}
    \field{eprintclass}{math.NT}
    \field{year}{2023}
  \endentry

  \entry{taibi}{article}{}
    \name{author}{1}{}{%
      {{hash=TO}{%
         family={Taibi},
         familyi={T\bibinitperiod},
         given={O.},
         giveni={O\bibinitperiod},
      }}%
    }
    \strng{namehash}{TO1}
    \strng{fullhash}{TO1}
    \field{labelnamesource}{author}
    \field{labeltitlesource}{title}
    \field{labelalpha}{Tai17}
    \field{sortinit}{T}
    \field{sortinithash}{T}
    \field{number}{2}
    \field{pages}{269\bibrangedash 344}
    \field{title}{Dimensions of spaces of level one automorphic forms for split
  classical groups using the trace formula}
    \field{volume}{50}
    \field{journaltitle}{Ann. Sci. \'{E}c. Norm. Sup\'{e}r. (4)}
    \field{year}{2017}
  \endentry

  \entry{weil}{article}{}
    \name{author}{1}{}{%
      {{hash=WA}{%
         family={Weil},
         familyi={W\bibinitperiod},
         given={A.},
         giveni={A\bibinitperiod},
      }}%
    }
    \strng{namehash}{WA1}
    \strng{fullhash}{WA1}
    \field{labelnamesource}{author}
    \field{labeltitlesource}{title}
    \field{labelalpha}{Wei65}
    \field{sortinit}{W}
    \field{sortinithash}{W}
    \field{pages}{1\bibrangedash 87}
    \field{title}{Sur la formule de {S}iegel dans la th\'{e}orie des groupes
  classiques}
    \field{volume}{113}
    \field{journaltitle}{Acta Math.}
    \field{year}{1965}
  \endentry

  \entry{weissauer-88}{article}{}
    \name{author}{1}{}{%
      {{hash=WR}{%
         family={Weissauer},
         familyi={W\bibinitperiod},
         given={R.},
         giveni={R\bibinitperiod},
      }}%
    }
    \strng{namehash}{WR1}
    \strng{fullhash}{WR1}
    \field{labelnamesource}{author}
    \field{labeltitlesource}{title}
    \field{labelalpha}{Wei88}
    \field{sortinit}{W}
    \field{sortinithash}{W}
    \field{pages}{100\bibrangedash 156}
    \field{title}{Differentialformen zu {U}ntergruppen der {S}iegelschen
  {M}odulgruppe zweiten {G}rades}
    \field{volume}{391}
    \field{journaltitle}{J. Reine Angew. Math.}
    \field{year}{1988}
  \endentry

  \entry{zh;phd}{online}{}
    \name{author}{1}{}{%
      {{hash=ZW}{%
         family={Zhang},
         familyi={Z\bibinitperiod},
         given={W.},
         giveni={W\bibinitperiod},
      }}%
    }
    \strng{namehash}{ZW1}
    \strng{fullhash}{ZW1}
    \field{labelnamesource}{author}
    \field{labeltitlesource}{title}
    \field{labelalpha}{Zha09}
    \field{sortinit}{Z}
    \field{sortinithash}{Z}
    \field{note}{PhD Thesis}
    \field{title}{Modularity of Generating Functions of Special Cycles on
  Shimura Varieties}
    \verb{url}
    \verb http://www.math.columbia.edu/~thaddeus/theses/2009/zhang.pdf
    \endverb
    \field{year}{2009}
  \endentry

  \entry{zuffetti-conescod2}{article}{}
    \name{author}{1}{}{%
      {{hash=ZR}{%
         family={Zuffetti},
         familyi={Z\bibinitperiod},
         given={R.},
         giveni={R\bibinitperiod},
      }}%
    }
    \strng{namehash}{ZR1}
    \strng{fullhash}{ZR1}
    \field{labelnamesource}{author}
    \field{labeltitlesource}{title}
    \field{labelalpha}{Zuf22}
    \field{extraalpha}{1}
    \field{sortinit}{Z}
    \field{sortinithash}{Z}
    \field{number}{10}
    \field{pages}{7385\bibrangedash 7441}
    \field{title}{Cones of special cycles of codimension 2 on orthogonal
  {S}himura varieties}
    \field{volume}{375}
    \field{journaltitle}{Trans. Amer. Math. Soc.}
    \field{year}{2022}
  \endentry

  \entry{zuffetti-unfolding}{article}{}
    \name{author}{1}{}{%
      {{hash=ZR}{%
         family={Zuffetti},
         familyi={Z\bibinitperiod},
         given={R.},
         giveni={R\bibinitperiod},
      }}%
    }
    \strng{namehash}{ZR1}
    \strng{fullhash}{ZR1}
    \field{labelnamesource}{author}
    \field{labeltitlesource}{title}
    \field{labelalpha}{Zuf24}
    \field{sortinit}{Z}
    \field{sortinithash}{Z}
    \field{number}{1}
    \field{pages}{Paper No. 10}
    \field{title}{Unfolding and injectivity of the {K}udla--{M}illson lift of
  genus 1}
    \field{volume}{307}
    \field{journaltitle}{Math. Z.}
    \field{year}{2024}
  \endentry

  \entry{zuffetti-equid}{misc}{}
    \name{author}{1}{}{%
      {{hash=ZR}{%
         family={Zuffetti},
         familyi={Z\bibinitperiod},
         given={Riccardo},
         giveni={R\bibinitperiod},
      }}%
    }
    \strng{namehash}{ZR2}
    \strng{fullhash}{ZR2}
    \field{labelnamesource}{author}
    \field{labeltitlesource}{title}
    \field{labelalpha}{Zuf22}
    \field{extraalpha}{2}
    \field{sortinit}{Z}
    \field{sortinithash}{Z}
    \verb{eprint}
    \verb 2208.07745
    \endverb
    \field{title}{Cones of orthogonal Shimura subvarieties and
  equidistribution}
    \field{eprinttype}{arXiv}
    \field{eprintclass}{math.AG}
    \field{year}{2022}
  \endentry
\enddatalist
\endinput

% $ biblatex auxiliary file $
% $ biblatex bbl format version 3.1 $
% Do not modify the above lines!
%
% This is an auxiliary file used by the 'biblatex' package.
% This file may safely be deleted. It will be recreated as
% required.
%
\begingroup
\makeatletter
\@ifundefined{ver@biblatex.sty}
  {\@latex@error
     {Missing 'biblatex' package}
     {The bibliography requires the 'biblatex' package.}
      \aftergroup\endinput}
  {}
\endgroup

\datalist[entry]{anyt/global//global/global}\enddatalist
\endinput

% $ biblatex auxiliary file $
% $ biblatex bbl format version 3.1 $
% Do not modify the above lines!
%
% This is an auxiliary file used by the 'biblatex' package.
% This file may safely be deleted. It will be recreated as
% required.
%
\begingroup
\makeatletter
\@ifundefined{ver@biblatex.sty}
  {\@latex@error
     {Missing 'biblatex' package}
     {The bibliography requires the 'biblatex' package.}
      \aftergroup\endinput}
  {}
\endgroup

\datalist[entry]{anyt/global//global/global}\enddatalist
\endinput
